% 本文件由 LaTeX 排版 Agent 根据有效 translation.json 自动生成。
% author=Origen; work_id=1015; title=若望福音注释

\chapter{若望福音注释}

% structure: p0001 source=h1 target=omitted reason=page-title-already-rendered-by-parent

\Emph{卷I}\newline{}\Emph{卷II}\newline{}\Emph{卷IV}\newline{}\Emph{卷V}\newline{}\Emph{卷VI}\newline{}\Emph{卷X}\par

\SourceRecord{Translated by Allan Menzies.；Ante-Nicene Fathers；Vol. 9.；Edited by Allan Menzies.；Buffalo, NY: Christian Literature Publishing Co.,；1896.；<http://www.newadvent.org/fathers/1015.htm>.}

% structure: p0001 source=h1 target=section reason=page-title

\section{\Emph{若望福音注释（卷一）}}

% structure: p0002 source=h2 target=subsection reason=relative-heading-rank; base=h2

\subsection{一、基督徒如何是属神的以色列。}

那自古以来被称为天主之民的民族分为十二支派，此外还有其他支派，其中有肋未族系，这族系本身又分为各种司祭和肋未的附属等级，按不同的方式服务于天主。同样，在我看来，基督的全部子民，当我们从内心隐藏之人的角度来审视时（\BibleRef{罗 2:29}），那被称为\Quote{内心为犹太人}且受心灵割损的子民，也以更奥秘的方式具有支派的特点。这从若望在\WorkTitle{默示录}中可以更清楚地看出，尽管其他先知也丝毫没有向那些有耳能听的人隐瞒实情。\newline{}若望这样说（\BibleRef{默 7:2}）：\BibleQuote{我又看见另一位天使，从日出之地升上来，拿着永生天主的印，大声向那四位得权柄伤害大地和海洋的天使呼喊说：\InnerQuote{你们不可伤害大地、海洋和树木，等我们在我们天主的众仆额上印上印。}我听见被印的人数，以色列子孙各支派中受印的共有十四万四千；犹大支派中受印的有一万二千，勒乌本支派中受印的有一万二千。}\newline{}接着，在稍后处（\BibleRef{默 14:1}）他继续说：\BibleQuote{我又看见，看，羔羊站在熙雍山上，同他在一起的，有十四万四千人，都有他的名号和在他额上写着他父的名号。我听见有声音来自天上，好像大水的声音，又像巨雷的声音；我所听见的，好像弹琴者在弹琴时所发的声音。他们在宝座前，并在那四个活物和长老前唱新歌；除了从地上赎回来的那十四万四千人外，谁也不能学会那歌。这些人没有与妇女在一起，因为他们是童贞；这些羔羊无论往哪里去，他们都跟随他。这些是从人类中赎回来，作为初熟之果献给天主和羔羊；在他们口中没有谎言，因为他们是没有瑕疵的。}\newline{}这是若望针对那些信仰基督的人说的，因为他们即使按肉身不能追溯至圣祖的后裔，却仍然是从各支派中被聚集的。这一点我们可以从关于他们的话进一步推断：\BibleQuote{你们不可伤害}，他说，\BibleQuote{大地、海洋和树木，等我们在我们天主的众仆额上印上印。我听见被印的人数，以色列子孙各支派中受印的共有十四万四千。}\par

% structure: p0004 source=h2 target=subsection reason=relative-heading-rank; base=h2

\subsection{二、默示录中受印的十四万四千人是外邦世界皈依基督者。}

这些在额上受印（\BibleRef{默 7:3}）的人，来自以色列子孙的每一支派，共十四万四千；后来若望说这十四万四千人额上写着羔羊和他父的名号，他们是童贞，没有与妇女玷污自己。那么，他们额上的印除了羔羊和他父的名号，还能是什么呢？两处经文都说他们的额上有印；一处提到印，另一处似乎表明印内包含羔羊和他父的名号。这些从支派中被选出的人，正如我们之前所述，就是那些童贞者。但按肉身属于以色列的信徒人数很少，我们敢断言他们几乎达不到十四万四千之数。因此，显然那没有与妇女玷污自己的十四万四千人必定是由外邦世界归向圣言的人组成。这样，关于各支派初熟之果是其童贞者的说法才能成立。因为经文接着说：\BibleQuote{这些是从人类中赎回来，作为初熟之果献给天主和羔羊；在他们口中没有谎言，因为他们是没有瑕疵的。}关于十四万四千人的说法无疑允许奥秘的诠释，但此时在此处比较先知中关于从外邦人中蒙召者的类似教导，不仅不必要，而且会偏离我们的目的。\par

% structure: p0006 source=h2 target=subsection reason=relative-heading-rank; base=h2

\subsection{三、在属神的以色列中，大司祭是那些献身于圣经研究的人。}

但这一切对我们有何意义呢？你读这些话时，一定会问，哦，\Person{盎博罗削}，你真是天主的人，在基督内的人，寻求不仅作一个人，更作一个属神的人。\BibleRef{格前2:14} 其意义是这样的：支派中的人通过肋未人和司铎向天主奉献\OriginalTerm{什一和初果}{tithe and first fruits}；他们并不把自己所有的一切都视为什一或初果。相反，肋未人和司铎除了什一和初果外别无所有；然而他们又藉着大司祭向天主奉献什一，我相信也奉献初果。同样，那些接近基督信仰研究的人也是如此。我们大多数人将大部分时间用于今生的事物，只将少数特殊行为奉献给天主，因而好比那些与司铎交往甚少、不费太多时间履行宗教义务的支派成员。但那些献身于天主圣言、除服务天主外别无他事的人，考虑到两种职业的差异，我们很自然可以称他们为我们的肋未人和司铎。而那些比同胞履行更杰出职务的，或许可称为大司祭，是按亚郎的品位，而非默基瑟德的品位。有人可能会反对说，将大司祭的名号用在人身上未免太过大胆，因为许多先知篇章提到耶稣本人是唯一的大司祭，如\BibleRef{希4:14}\BibleQuote{我们有一位经过诸天的大司祭，就是天主子耶稣。} 对此我们回答说：宗徒已清楚定义了他的意思，并宣称先知论及基督说过：\BibleQuote{你按照默基瑟德的品位，永远作司祭。} 而不是按照亚郎的品位。因此我们说，人可以按照亚郎的品位成为大司祭，但按照默基瑟德的品位，只有天主的基督才可以。\par

% structure: p0008 source=h2 target=subsection reason=relative-heading-rank; base=h2

\subsection{四、福音的研究是这些基督信仰司铎所奉献的初果。}

如今我们全部活动都奉献给天主，我们整个生命都致力于在神圣事物上进步。那么，如果我们渴望拥有上述由众多初果所组成的全部初果——如果这个看法不错的话——在我们彼此经历身体的分离之后，我们的初果除了在福音的研究中，还能包含什么呢？因为我们敢说，福音是全部圣经的初果。那么，自从我们来到\Place{亚历山大里亚}的时候起，我们活动的初果若不是在圣经的初果中，还能在哪里呢？然而，不可忘记，初果与初生之物并不相同。因为初果是在所有果实成熟之后奉献的，而初生之物却是在所有果实之前。至于在众教会中流传并被相信是神圣的圣经，说初生之物是梅瑟法律，而初果是福音，并不为错。因为是在所有先知（他们直到主耶稣为止一直说预言）的果实之后，完美的圣言才显露出来。\par

% structure: p0010 source=h2 target=subsection reason=relative-heading-rank; base=h2

\subsection{五、全部圣经都是福音；但福音书卓越超群于其他圣经。}

然而，此处有人可能提出异议，援引刚才提出的\Quote{初果最后展开}的概念，并说宗徒大事录和宗徒的书信是在福音之后写成的，这就破坏了我们关于福音是全部圣经初果的论证。对此我们必须回答：在基督内明智的人，即那些从流传的书信中获益、并看出这些书信有法律和先知所保存的证言为之担保的人，他们坚信宗徒著作应被称为明智且可信赖的，并具有极大的权威，但它们不与\BibleRef{格后 6:18}\BibleQuote{万军的上主这样说。}居于同等地位。关于这一点，请考虑圣保禄的话。当他宣告\BibleRef{弟后 3:16}\BibleQuote{凡受天主默感所写的圣经，为教训都是有益的，}他是否包括他自己的著作？或者他是否不包括他的说法：\BibleRef{格前 7:12}\BibleQuote{我说，不是主说，}以及\BibleRef{格前 7:17}\BibleQuote{我在各教会中这样吩咐，}还有\BibleRef{弟后 3:11}\BibleQuote{我在安提约基雅、依科尼雍、吕斯特辣所遭遇的事，}以及他凭自己权威所写的类似内容，这些并不完全具有出自神圣默感之话语的特点？我们还必须表明，旧约圣经并非福音，因为它并未指出那将要来临的一位，只是预告祂并宣告祂在未来的来临；而新约圣经全都是福音。它不仅在福音的开头说：\BibleRef{若 1:29}\BibleQuote{请看天主的羔羊，除免世罪者；}而且还包含许多对祂的赞美，以及许多祂的教导，因着祂，福音才成为福音。此外，如果天主在教会中设立了\BibleRef{弗 4:11}宗徒、先知和传福音者（gospellers）、牧者和教师，我们必须首先查考传福音者的职务是什么，并注意：传福音者不仅叙述救主如何治愈生来失明的人，\BibleRef{若 9:1}或复活已发臭的死人，\BibleRef{若 11:39}或陈述祂行了何等奇事；如此界定传福音者的职务后，我们便不会犹豫，也在这类非叙述性的、而是劝勉性的、旨在坚固对耶稣使命之信仰的言论中，找到福音；这样我们就得出这样的结论：凡宗徒所写的，都是福音。对于这第二个定义，可能会有人反对，认为书信并未冠以\Quote{福音，}而我们错误地将福音之名应用于整个新约。但对此我们回答：圣经中常有这样的情况，即两个或更多的事物或人同名时，该名称最显著地属于其中之一。因此，救主说：\BibleRef{玛 23:8}\BibleQuote{不要在地上称人为师傅；}而宗徒却说在教会中设立了师傅。后者因此不会成为福音训令严格意义上的师傅。同样，书信中的福音，当与耶稣言行、受难和言论的叙述相比较时，便不会延伸到其中的每一个字。不：福音是全部圣经的初果，而我们将我们一切行为的初果——我们希望这些行为会如愿以偿——奉献给这圣经的初果。\par

% structure: p0012 source=h2 target=subsection reason=relative-heading-rank; base=h2

\subsection{六、四部福音。若望福音是四部中的初果。诠释它所必需的资格。}

如今福音共有四部。这四部福音可说是教会信德的要素，藉着这些要素，那在基督内与天主和好的整个世界得以组成；正如保禄所说（\BibleRef{格后 5:19}）：\BibleQuote{天主在基督内，使世界与自己和好}；耶稣背负了这世界的罪；因为经上正是论及这教会世界而写道（\BibleRef{若 1:29}）：\BibleQuote{请看天主的羔羊，除免世罪者}。福音既然有四部，我认为福音的初果就是你们托付我尽力探讨的那部，即\BibleBook{若望}{福音}；这部福音所谈论的那位的族谱早已有人记述，但它却从那位尚无任何族谱之前的起点开始谈论祂。因为\Person{玛窦}为期待那出自亚巴郎和达味后裔而来的希伯来人写作，他说（\BibleRef{玛 1:1}）：\BibleQuote{耶稣基督，达味之子，亚巴郎之子的族谱}。而\Person{马尔谷}知道他所写的是什么，叙述了福音的开端；我们或许可以在\Person{若望}中找到他的目标：在起初已有圣言，圣言就是天主。但是\Person{路加}，虽然他在\BibleBook{}{宗徒大事录}开头说：\BibleQuote{在先前，我写了所有耶稣开始所行和所教导的事}，却把关于耶稣最伟大、最完备的言论留给了那靠在耶稣胸怀的那位。因为上述几位都没有像\Person{若望}那样明确宣告祂的天主性；\Person{若望}让祂说：\BibleQuote{我是世界的光}，\BibleQuote{我是道路、真理和生命}，\BibleQuote{我是复活}，\BibleQuote{我是门}，\BibleQuote{我是善牧}；并在\BibleBook{}{默示录}中说：\BibleQuote{我是阿耳法和敖默加，起始和终结，首先的和最后的}。所以我们可以大胆说：全部圣经的初果是福音，而福音的初果是\BibleBook{若望}{福音}。除非他曾靠在耶稣的胸怀，并从耶稣那里接受了\Person{玛利亚}作为自己的母亲，没有人能领悟它的含义。那要成为另一位若望的人必须成为这样的人，且如若望一样，由耶稣亲自将祂的本来面目显示给他。因为如果\Person{玛利亚}，正如那些以健全心智赞美她的人所声明的，除了耶稣之外没有别的儿子，但耶稣却对祂的母亲说（\BibleRef{若 19:26}）：\BibleQuote{女人，看，你的儿子}，而不是说\Quote{看，你也有这个儿子}，那么祂实际上是对她说：\Quote{看，这是你所生的耶稣}。难道不是所有成全的人不再活着（\BibleRef{迦 2:20}），而是基督在他内活着吗？如果基督在他内活着，那么对\Person{玛利亚}论及他时就说：\Quote{看，你的儿子基督}。那么，我们必须有怎样的心智，才能以相称的方式解释这部著作？尽管它被托付给普通言语的尘世宝库，即任何路人皆可阅读的文字，而且当任何用肉耳倾听的人大声诵读时，也能听到。关于这部著作，我们该说什么呢？凡是准确领悟其内容的人，应当能真诚地说：\BibleQuote{我们拥有基督的心意，为能知道天主所赐给我们的事}。可以引用\Person{保禄}的一句话来支持整个新约都是福音的主张。他在某处写道（\BibleRef{罗 2:16}）：\BibleQuote{按照我的福音}。如今，我们并没有\Person{保禄}的著作通常被称为福音。但是他宣讲和所说的一切都是福音；他宣讲和所说的，他也惯常书写，因此他所写的也就是福音。但如果\Person{保禄}所写的是福音，那么\Person{伯多禄}所写的也同样是福音，总之，所有为延续基督在世间居留的知识，并为预备祂的第二次来临，或为使祂的来临成为那些愿意接受天主圣言之灵魂的现世现实而说的或写的一切，都是福音——那时祂正站在门口敲门，要进入它们内。\par

% structure: p0014 source=h2 target=subsection reason=relative-heading-rank; base=h2

\subsection{七、福音中所宣告的美好事物}

但是现在是时候探究称号\Quote{福音}的意义，以及这些书卷为何拥有这个标题。福音是一篇含有事物应许的论述，这些事物因其带来的益处，自然会在人听到并相信这应许时便使听者欢喜。即使我们根据听者的处境来定义它，这样的论述也仍然不失为福音。福音要么是暗示某种善的事物对信者实际临在的话语，要么是应许所期待的善即将到来的话语。所有这些定义都适用于那些被称为福音的书卷。因为每一部福音都是有益于相信且不错解它们之人的宣告集；它给他带来益处，并自然使他欢喜，因为它讲述了一切受造物中的首生者（见：\BibleRef{哥 1:15}）\Person{基督耶稣}，为了人类，且为他们的得救，而与人类同住。此外，每一部福音都讲述了良善的父在子内与那些有心接纳祂之人同住的寄居，这对每位信者是显而易见的；而且，这些书卷宣布了一个先前已被期待的善，这并不难看出。因为\Person{若翰洗者}几乎是代表全体人民说话，当他派人到\Person{耶稣}那里问（见：\BibleRef{玛 11:3}）\Quote{你是那要来的那位，还是我们该期待另一位？}因为对人民来说，\OriginalTerm{默西亚}{Messiah}是所期待的善，众先知早已预言；他们所有人都一样，虽在律法和先知之下，却将希望寄托于祂，正如\Person{撒玛黎雅妇人}作证时所说（见：\BibleRef{若 4:25}）\Quote{我知道那称为基督的默西亚要来；祂来时，会告诉我们一切。}\Person{西满}和\Person{克罗帕}也是如此，当他们彼此谈论所有发生在\Person{耶稣基督}自己身上的事——那时祂已复活，虽然他们不知道祂已从死者中复活——就这样说（见：\BibleRef{路 24:18}）\Quote{祢是\Place{耶路撒冷}唯一寄居的人，竟不知道这几天在那里所发生的事吗？祂说：\InnerQuote{什么事？}他们回答说：\InnerQuote{就是关于\Place{纳匝肋}人\Person{耶稣}的事。祂是先知，在天主和全体百姓面前，言行都有大能。司祭长和我们的首领竟把祂交出去，判了死罪，钉在十字架上。但我们原希望祂就是要救赎\Place{以色列}的那位。}}此外，\Person{西满伯多禄}的兄弟\Person{安德肋}找到了自己的兄弟\Person{西满}，对他说（见：\BibleRef{若 1:42}）\Quote{我们找到了\OriginalTerm{默西亚}{Messiah}，翻译出来就是基督。}稍后，\Person{斐理伯}找到了\Person{纳塔乃耳}，对他说（见：\BibleRef{若 1:46}）\Quote{我们找到了\Person{梅瑟}在律法中所记载的，先知们也所记载的那位，就是\Person{若瑟}的儿子，\Place{纳匝肋}人\Person{耶稣}。}\par

% structure: p0016 source=h2 target=subsection reason=relative-heading-rank; base=h2

\subsection{八、福音如何使圣经的其他书卷也成为福音。}

现在有人可能会对我们第一个定义提出异议，因为它会包含那些不称为福音的书卷。因为律法和先知在我们看来也是书卷，含有事物应许，这些事物因将赋予听者的益处，自然会在人领受信息时使听者喜悦。对此可以这样说：在基督的寄居之前，律法和先知，由于解释其中奥秘的那位尚未到来，并未传达属于我们福音定义的那种应许；但是救主，当祂与人同住并使福音以身体形式出现时，藉着福音使一切事物都显为福音。在此，我认为引用一个……不多……却涵盖所有的例子并非离题。因为当祂揭去了律法和先知中的幔子，并借祂的天主性向人子证明神性在运行，祂为所有渴望成为祂智慧门徒的人开启了道路，使他们理解\Person{梅瑟}律法中哪些是真实和真正的——古人所敬拜的是那些事物的预像和阴影——以及那些历史叙述中哪些是真实的，这些历史是\BibleQuote{作为预像发生在他们身上的事}（见：\BibleRef{格前 10:11}），但\BibleQuote{这些事是为了我们这些末世的人而写的。}因此，无论基督与谁同住，他既不在这座山也不在耶路撒冷朝拜天主；他知道天主是神，并用精神朝拜祂，在精神和真理中朝拜；他不再用预像朝拜万有的父和创造者。因此，在因基督寄居而产生的福音之前，旧的作品中没有一部是福音。但是福音，即新盟约，将我们从文字的旧中解放出来，藉知识之光为我们点燃精神的新，这是一种永不陈旧的事物，它在新约中安家，但也存在于所有圣经中。因此，那使我们即使在旧约中也能发现福音临在的福音，本身在特殊意义上配得福音之名。\par

% structure: p0018 source=h2 target=subsection reason=relative-heading-rank; base=h2

\subsection{九、属肉体的福音与属神的福音。}

然而，我们不可忘记，基督与人同在的寓居，在祂身体寓居之前，就已以一种理智的方式，临于那些更成全、非孩童、且不在启蒙师和监护之下的人。他们在心中看见了时期已满——族长们、仆人梅瑟、以及看见基督荣耀的众先知。正如在祂显明且身体的来临之前，祂已临于那些成全的人；同样，在祂的来临被宣报给众人之后，对于那些仍是孩童的人（因为他们仍在启蒙师和监护之下，尚未到达时期的满全），基督的前驱来寓居，即适合仍处于孩童心智的\OriginalTerm{话语}{logoi}，因此正当地被称为启蒙师。但圣子自己，荣耀的天主，圣言，尚未来临；祂等待那些将要接纳祂神性的天主之人必须进行的预备。这一点我们也要牢记：正如法律含有未来美物的影子，这些美物由那按真理宣布的法律所指示，同样，福音也教导基督奥迹的影子，即那被认为任何人都能理解的福音。若望所称的永恒福音，以及可恰当地称为属神福音的，向那些愿意理解的人清楚地呈现了关于天主圣子的一切事，包括祂话语中的奥迹以及祂事迹所隐示的事。据此我们可以得出结论：正如那外表上是犹太人、肉身受割礼的人一样，基督徒和圣洗也是如此。保禄和伯多禄早期在表面上是犹太人并受了割礼，但后来从基督领受了也要在内里成为这样；因此，外表上他们为许多人的救恩而作犹太人，并因着经济，他们不仅在言语上承认是犹太人，而且以行动表明。关于他们的基督徒身份也是如此。正如保禄若不按理智表明必要时给弟茂德行割礼，以及照常理发许愿，总之，为犹太人成为犹太人，以赢得犹太人，就无法有益于那些按肉身的犹太人——同样，那负责许多人益处的人，也不能仅凭那内在的基督徒身份来按应有的方式运作。那绝不可能使他改善那些遵循外在基督徒身份的人，或引导他们达到更高更好的事物。因此，我们必须在身体上和属神上都是基督徒，并且在有需要时运用身体的（属肉体的）福音，人借此对属血肉的人说，他只认识耶稣基督并被钉在十字架上。但若我们找到那些在圣神内成全、结出圣神的果实、并热爱天上智慧的人，必须使他们分享那圣言，这圣言在成为血肉之后，又复活到起初与天主同在时的状态。\par

% structure: p0020 source=h2 target=subsection reason=relative-heading-rank; base=h2

\subsection{十、耶稣自己如何是福音。}

以上对福音本质的探讨不可视为无用；它使我们看到感官的福音与理智和属神的福音之间有何区别。我们现在要做的是将感官的福音转化为属神的福音。\par

因为感官福音的叙述若不发展到属神的地步，又能有什么意义呢？它将微不足道或毫无价值；任何人都能阅读并确知它所讲述的事实——仅此而已。但我们现在的全部精力要致力于渗透福音深奥的涵义，并探索除去预象后的真理。现在，福音所说的一切应视为美物的应许；我们必须说，宗徒们在福音中宣告的美物就是耶稣。他们宣告的一个美物是复活；但在某种意义上，复活就是耶稣，因为耶稣说：\BibleQuote{\BibleRef{若 11:25} 我是复活。}耶稣向穷人宣讲那些为圣人所存留的事物，召唤他们领受天主的应许。圣经为宗徒们的福音宣告以及我们救主的福音宣告作证。达味论及宗徒们，也许也论及传福音者说：\Quote{主要赐给那以大力传道者话语；这是万军之主、所爱者的君王；}同时教导我们，令人信服的不是巧妙结构的言辞、不是表达方式、不是训练有素的雄辩，而是神圣能力的传递。因此保禄也说：\Quote{我要知道的不是虚浮的话语，而是能力；因为天主的国不在于话语，而在于能力。}在另一处：\BibleQuote{\BibleRef{格前 2:4} 并且我的言语和宣讲，不是用智慧动听的言辞，而是藉圣神和能力的证明。}西满和克罗帕为此作证说：\BibleQuote{\BibleRef{路 24:32} 当祂在路上给我们开解圣经时，我们的心不是火热的吗？}宗徒们，由于天主赐给讲者大量的能力，就拥有大能，按达味的话：\Quote{主将话语赐给传道者，他们有大能力。}依撒意亚也说：\Quote{那传报喜讯者的脚步是多么美丽！}他看见宗徒们的宣告是多么美丽和适时，他们行走在那说\Quote{我是道路}的祂内，并赞美那些行走在基督耶稣理智之道上、通过那道进入天主的人的脚。那些脚美丽的人传报喜讯，就是耶稣。\par

% structure: p0023 source=h2 target=subsection reason=relative-heading-rank; base=h2

\subsection{十一、耶稣是一切美物；因此福音是多方面的。}

但愿无人惊奇，若我们以为耶稣在福音中以多种美善之名称被宣告。若我们察看那称呼天主子的事物的名字，就会明白耶稣是何种美善，那些双足佳美之人所传扬的正是祂。一件美善是生命，但耶稣是生命。另一件美善是世界的光，当它是真光、是人的光时，而这一切都说是天主子。又有另一件美善，一人可设想为生命或光之外的，即真理。第四件是时间之外的，即通向真理的道路。我们的救主教导说祂就是这一切，当祂说：\BibleQuote{我是道路、真理和生命。}\BibleRef{若 14:6}啊，那岂不美善吗？摆脱尘世与可朽，复活过来，从主那里获得这恩惠，因为祂就是复活，如祂所说：\BibleQuote{我是复活。}\BibleRef{若 11:25}但门也是美善，通过它进入至高的福乐。如今基督说：\BibleQuote{我是门。}\BibleRef{若 10:9}又何必论及智慧呢？\BibleQuote{上主在祂造化的起头，在太初之初就造了我，作为祂行事的起始，}\BibleRef{箴 8:22}祂的父因她而喜乐，喜爱她多姿多彩的理智之美，唯独心智之眼才能看见，并激发那辨识她神圣与属天魅力者对她的爱。确实，天主的智慧是美善，与那先前由双足佳美之人所传扬的其他美善一同被宣告。天主的德能是第八件我们列举的美善，就是基督。也不可忽略那圣言，祂是万有之父后的天主。因为这也是一美善，不亚于其他。那么，有福的是那些接受这些美善，并从传报它们佳音的人手中领受它们的人，那些人双足佳美。事实上，甚至格林多人中的一个，保禄向他宣告自己只知道耶稣基督和被钉十字架的祂，若他学习那为我们而成了人的那位，并如此领受祂，他便会与我们所说美善的开端认同；藉着那人耶稣，他会被造就成天主的人，藉着祂的死，他会向罪而死。因为\BibleQuote{基督，祂既死了，就一次而永远地死了。}\BibleRef{罗 6:10}但从祂的生命，既然\BibleQuote{祂活着，是活于天主，}每一位与祂的复活相似的人都领受了那活于天主的生命。但谁能否认公义——本质的公义——是美善，以及本质的圣化与本质的救赎呢？那些传扬耶稣的人传扬这些事，说\BibleQuote{祂是由天主造成了我们的公义、圣化和救赎。}\BibleRef{格前 1:30}因此，我们将有无数关于祂的著作，显示耶稣是众多美善的集合；因为从那些几乎无法计数且已写下的东西中，我们可以对实际存在于祂内的事物略作推测，\BibleQuote{天主乐意使整个圆满的性体居住在那里，}而它们并不包含在著作中。我为何说\Quote{不包含在著作中}？因为若望在此论及整个世界，说：\BibleQuote{我想连世界本身也容不下所要写的书。}\BibleRef{若 21:25}如今说宗徒传扬救主，即是说他们传扬这些美善。因为这位就是那从良善之父领受了自己应为这些美善的那一位，好使每人从耶稣那里领受他所能领受的这一件或那些事物，从而享受美善。但宗徒们，他们双足佳美，以及那些效法他们而寻求传报佳音的人，若非耶稣自己首先向他们传报佳音，就不能如此行，正如依撒意亚所说：\BibleQuote{我亲自说话，我就在这里，如那在山上报佳音者的脚，如那传报和平者，传报佳善者；因为我要使我的救恩被听见，说：天主要作王统治你，熙雍！}\BibleRef{依 52:6}因为那说话者声称自己亲自临在的山是什么，岂不是那些不低于世上最高最大者吗？这些需要新盟约的能干仆役去寻求，为能遵守那诫命：\BibleQuote{你这位传报佳音给熙雍的，请上高山上；你这位传报佳音给耶路撒冷的，请用力扬声！}\BibleRef{依 40:9}如今，若耶稣自己向那些将传报佳音的人传报佳音，这并不奇怪，因为天主子向那些不能通过他人认识祂的人传报关于祂自己的佳音。祂登上高山，向他们传报佳善，而祂自己是由祂的良善之父所教导，\BibleQuote{祂使自己的太阳上升，光照恶人，也光照善人；降雨给义人，也给不义的人，}\BibleRef{玛 5:45}祂并不轻视那些心灵贫穷的人。祂向他们传报佳音，正如祂亲自给我们作证，当祂拿起依撒意亚书并诵读：\BibleQuote{上主的神临于我身，因为上主给我傅了油，派遣我向贫穷人传报佳音，向俘虏宣告释放，向盲者宣告复明。祂把书卷闭上，交给仆役，就坐下了。当众人的眼都注视祂时，祂说：今天这段圣经应验在你们耳中了。}\par

% structure: p0025 source=h2 target=subsection reason=relative-heading-rank; base=h2

\subsection{十二、福音也记载了加于耶稣的恶行。}

不可忘记，在这部福音中，包含了所有向耶稣所行的善事；例如，那个曾是罪人而后悔改的妇人的故事，她真正从邪恶状态中痊愈后，蒙恩将香膏倒在耶稣身上，使满屋都闻到馨香。因此也有话说：\Quote{无论这福音在万民中传到何处，也要述说这女人所作的，作为对她的纪念。}显然，凡向耶稣的门徒所行的，就是向祂行的。祂指着那些受到善待的门徒，对善待他们的人说：\BibleRef{玛 25:40}\Quote{你们作在这小子中一个最小的身上，就是作在我身上。}因此，我们向邻人所行的每一件善事，都记在福音里，就是那写在天上版上、被一切配得认识万有者所诵读的福音。但另一方面，福音也有为那些向耶稣行恶的人定罪的部分。犹达斯的背叛、恶毒群众的呼喊说：\BibleRef{若 19:6}\Quote{除掉这样的人，从地上除掉！}以及\Quote{钉死祂，钉死祂！}那些用荆棘给祂戴冠冕者的嘲笑，以及诸如此类的一切，都记在福音中。因此我们看到，凡出卖耶稣门徒的，被视为出卖耶稣本人。对当时仍是迫害者的扫罗，\BibleRef{宗 9:4}说：\Quote{扫禄，扫禄，你为什么迫害我？}以及\Quote{我就是你所迫害的耶稣。}还有人仍用荆棘给耶稣戴冠冕并羞辱祂，就是那些被今生的忧虑、财富和享乐所窒息的人，他们虽接受了天主的话，却未能使之结果实。\BibleRef{路 8:14}因此我们必须谨慎，免得我们也用自己的荆棘给耶稣戴冠冕，被记在福音里，并在这种角色中被那些认识那在万有之内、临在于所有理性而圣洁的生命中的耶稣的人所诵读，学习祂如何被傅油、被款待、被光荣，或另一方面，祂如何被羞辱、被嘲弄、被鞭打。这一切都必须说出来；这是我们证明的一部分：我们的善行，以及那些跌倒之人的罪，都被记载在福音里，或为永生，或为羞辱和永远的憎恶。\par

% structure: p0027 source=h2 target=subsection reason=relative-heading-rank; base=h2

\subsection{十三、天使也是福音的宣讲者。}

如果人类中有一些人被尊为宣讲福音的职务，并且耶稣亲自传报佳音，向穷人宣讲福音，那么那些被天主造为神灵、成为火焰的使者，万有之父的仆役，难道能被排除在宣讲福音之外吗？因此，一位天使站在牧羊人旁边，有大光照耀他们，说：\BibleRef{路 2:10}\BibleQuote{不要害怕；看，我给你们报告一个关乎万民的大喜讯：今天在达味城中，为你们诞生了一位救主，就是主基督。}当人们尚不知晓福音的奥秘时，那些比人更高、天上的居民，天主的军队，赞美天主说：\Quote{光荣归于至高天主，在地上平安，善意临于人。}天使说完这话，就离开牧羊人升天去了，留给我们揣摩：因耶稣诞生而传给我们的喜乐，如何是至高之处光荣归于天主；他们谦卑自下直至地上，然后返回安息之处，藉着耶稣基督在至高之处光荣天主。天使也惊叹因耶稣而要在世上带来的和平——这地上的战场，路济弗尔、早晨之星从天坠落之处，要被耶稣争战并消灭。\par

% structure: p0029 source=h2 target=subsection reason=relative-heading-rank; base=h2

\subsection{十四、旧约——由若翰预表——是福音的开端。}

除了我们已经说过的，关于福音还有这一点要考虑：首先，福音是基督耶稣的福音，祂是整个得救群体的元首；正如马尔谷所说：\BibleRef{谷 1:1}\Quote{天主之子耶稣基督福音的开始。}然后，福音也是宗徒的福音；因此保禄\BibleRef{罗 2:16}说：\Quote{按我的福音。}但福音的开端——就其范围而言，它有开始、延续、中间和结束——无非是整个旧约。就此而言，若翰是旧约的预表，或者，如果我们考虑新约与旧约的联系，若翰代表旧约的结束。因为同一马尔谷说：\BibleQuote{天主之子耶稣基督福音的开始，正如依撒意亚先知所记载的：看，我派遣我的使者在你面前，他要在你前面预备你的道路。旷野中有呼号者的声音：你们当预备主的道路，修直他的路径。}在这里，我必须惊奇，持异议者怎能将两约归于两个不同的神。这些话，即使没有其他，也足以定他们的错。因为如果他们认为若翰属于另一个神，属于造物主，且如他们所持，不认识新神性，那么若翰怎能成为福音的开端？天使不仅受托一次短暂传福音的职务，即对牧羊人的那一次。因为在末了，有一位崇高飞翔的天使，持着福音，要把它传给万民，因为良善的圣父并未完全遗弃那些离开祂的人。载伯德的儿子若望在《默示录》中说：\BibleRef{默 14:6}\BibleQuote{我又看见另一位天使飞在空中，有永远的福音要传给住在地上的人，就是各邦国、各支派、各异语、各民族，大声说：你们应当敬畏天主，将光荣归于他，因为他审判的时候到了；应当敬拜那创造天地海和众水泉源的。}\par

% structure: p0031 source=h2 target=subsection reason=relative-heading-rank; base=h2

\subsection{十五、福音在旧约中，实际上也在整个宇宙中。祈求帮助以理解当前这部著作的神秘意义。}

因此，既然我们已经表明，若按一种解释，福音的开端就是全部旧约，并且以若翰的形象为标记，那么——以免有人称这仅为无根据的断言——我们应补充\WorkTitle{宗徒大事录}中关于厄提约丕雅女王太监与斐理伯的记述。据说，斐理伯从\BibleBook{}{依撒意亚}的这段经文开始：\Quote{他被牵去如同羔羊被宰杀，又如同羔羊在剪毛者前缄默不语，}并由此向他宣讲了主耶稣。如果\BibleBook{}{依撒意亚}不是福音开端的一部分，那么他怎能以先知开端并宣讲耶稣呢？由此我们可以为开头的论断——即每部神圣圣经都是福音——提供证据。如果宣讲福音的人传报喜讯，而所有在耶稣肉身降临之前说话的人都宣讲基督（祂就是我们看为的喜讯），那么他们所有人的话都在某种意义上属于福音的一部分。而当福音被宣告传遍全世界时，我们就推断它确实在全世界被宣讲；这并非仅指在这地上区域，而是指在天地之整体秩序中，或自天地之间。我们何必进一步讨论福音是什么？我们所说的已足够。除了我们引用的段落——这些段落绝非不恰当或不适用于我们的目的——还可以从圣经中收集许多类似的内容，以便清楚看出，借着福音所倾注的、在耶稣基督内的喜讯的光荣，这福音由人和天使，而且我相信也由掌权者、异能者、\BibleRef{弗 1:21}宝座、宰制者，以及一切被称呼的名号——不仅在此世，也在来世——甚至由基督亲自传布。在此，让我们结束在阅读作品本身之前必须说的话。现在，我们祈求天主藉耶稣基督并借圣神助佑我们，使我们能阐明当前文字中所珍藏的奥秘含义。\par

% structure: p0033 source=h2 target=subsection reason=relative-heading-rank; base=h2

\subsection{十六、\Quote{开始}的含义。（1）论空间。}

\Quote{\Emph{在起初已有圣言。}}\BibleRef{若 1:1} 并非只有希腊人认为\Quote{开始}一词有多种含义。任何人若收集圣经中该词出现的段落，并为了精确解释它而留意该词在每个段落中所代表的意义，就会发现这词在圣言中也具有多种含义。我们谈到一种与转进有关的开始。此处它涉及道路和长度。这在下列经文中显明：\BibleRef{箴 16:5}\Quote{行正义是善途的开端。}因为善途是漫长的，所以首先要考虑与行动相关的问题，而这方面在\Quote{行正义}一词中呈现；静观方面则随后加以考虑。在后者中，其终点最终安息于所谓的万物复兴中，因为若经上所说真实：\BibleRef{格前 15:25}\Quote{因为祂必须为王，直到把一切仇敌放在祂脚下。最后被毁灭的仇敌是死亡。}那时，对于那些因着与祂同在的圣言而来到天主面前的人，将只剩下一种活动，即认识天主，好使他们因着认识父而被找到，成为祂的圣子，正如如今除了圣子以外无人认识父。因为若有人仔细探究，那些将要认识父的人（认识父的那位向他们启示父）将在何时认识父，并思考人如今如何只透过镜子与谜语观看，从未学会按应当的方式去认识，那么他有理由说，没有人——甚至没有宗徒，没有先知——认识父，除非他成为与父合一的人，如同子和父合一一样。若有人说，我们之所以离题至此，是为了探讨\Quote{开始}一词的某一种含义，那么我们必须表明，这离题对于我们达到目的而言是必要且有益的。因为如果我们谈到一种与转进、道路及其长度有关的开始，并且我们被告知行正义是善途的开端，那么我们就需要知道，每条善途如何以行正义为开端，以及如何在这样的开端之后达到静观，又如何以此方式达到静观。\par

% structure: p0035 source=h2 target=subsection reason=relative-heading-rank; base=h2

\subsection{十七、（2）论时间。创造的开端。}

此外，还有关于来源的\OriginalTerm{元始}{arché}，正如这句经文所示：\BibleRef{创 1:1} \BibleQuote{在起初天主创造了天地。}然而，这一含义在约伯传的经文中更加明显：\BibleRef{约 40:19} \BibleQuote{这是天主创造的元始，是为供祂的天使嘲笑而造的。}有人会以为，在世界的创造中，诸天和大地是最先被造的。但第二段经文提出了一个更好的看法，即：许多存有被赋予形体，其中最先被造的是被称为龙的受造物，但在另一段经文中\BibleRef{约 3:8}被称为巨兽（\OriginalTerm{利维坦}{Leviathan}），主驯服了它。我们必须查考这事：当圣人们过着远离物质和任何形体的真福生活时，那龙是否从纯洁的生命堕落，变得适宜被束缚在物质和形体中，以致主能藉着风暴和云彩说：\BibleQuote{这是天主创造的元始，是为供祂的天使嘲笑而造的。}然而，有可能那龙并非绝对是主创造的元始，而是有许多受造物被造为有形体供天使嘲笑，龙是其中第一个，而其他的能够在形体中存在而不受这种责难。但并非如此。因为太阳的灵魂被安置在一个形体中，整个受造界——宗徒论及它说：\Quote{整个受造界一同叹息、同受产痛直到如今}，也许下面这段话也是论及同一件事：\Quote{受造界被屈服于虚妄，并非出于自愿，而是由于那使它屈服于希望的那一位；}如此，形体可能处于虚妄中，行形体之事，如同在形体中的人必须……。在形体中的人，虽不情愿，却行形体之事。因此受造界被屈服于虚妄，并非出于自愿；但那不情愿地行形体之事的人，是为了希望而行其所行，正如我们可说，保禄希望留在肉身内，并非出于自愿，而是为了希望。因为虽然他认为\BibleRef{斐 1:23}离去与基督同在一起更好，但他希望留在肉身内，为他人谋利益，并在所希望的事上长进——不仅为他本人，也为蒙他恩惠的人——也不是不合理的。这一术语\Quote{元始}的含义也将用于智慧在箴言中发言的那段经文。我们读到：\Quote{上主，} \Quote{在祂造化的元始，在祂作之初，就造了我。}在此，这术语可解释为我们所谈的第一种应用，即道路的元始：它说：\Quote{上主，} \Quote{造了我作为祂道路的元始。}有人可以有理有据地断言，天主自身就是万物的元始，并进而说——这是显而易见的——圣父是圣子的元始，造物主（\OriginalTerm{造物主}{Demiurge}）是造物主工程的元始，总而言之，天主是所有存在者的元始。我们的这句经文支持这一观点：\BibleQuote{在元始已有圣言。}在圣言中，人可看见圣子，并且因为祂在圣父内，祂可说是在元始。\par

% structure: p0037 source=h2 target=subsection reason=relative-heading-rank; base=h2

\subsection{十八、(3) 关于实体。}

第三，元始可以是事物由之而出的东西，即构成万物的基质。然而，这是那些主张物质本身非受造者的观点，我们信徒不能认同，因为我们相信天主从无中造出了现有的万物，正如玛加伯中七位殉道者的母亲所教导的\BibleRef{加下 7:28}，也如\WorkTitle{牧者}中的忏悔天使所训示的。\par

% structure: p0039 source=h2 target=subsection reason=relative-heading-rank; base=h2

\subsection{十九、(4) 关于原型与肖像。}

除了这些含义，还有我们按形式谈论的\OriginalTerm{元始}{arché}的含义；因此，如果一切受造物的首生者\BibleRef{哥 1:15}是不可见天主的肖像，那么圣父就是祂的元始。同样，基督是按天主的肖像所造之人的元始。因为如果人是按照肖像而造，而肖像又是按照圣父，那么在第一种情况下，圣父是基督的元始；在另一种情况下，基督是人的元始，而人不是按照祂所是的肖像被造，而是按照肖像被造。这一例子与我们的经文一致：\BibleQuote{在元始已有圣言。}\par

% structure: p0041 source=h2 target=subsection reason=relative-heading-rank; base=h2

\subsection{二十、(5) 关于元素及其所构成之物。}

还有一种关于学习的\OriginalTerm{元始}{arché}，例如我们说字母是语法的元始。宗徒据此说：\BibleRef{希 5:12} \Quote{按时间说，你们本应做导师，却还需要有人教导你们天主圣言元始的元素。}与学习有关的元始是双重的：首先是按其本质，其次是相对于我们；正如我们论基督可说，按其本质，祂的元始是神性，但相对于我们——我们因元始的浩大而无法把握关于祂的全部真理——祂的元始是祂的人性，正如祂被传报给婴孩：\Quote{耶稣基督，并祂被钉在十字架上。}如此看来，基督按其本性是学习的元始，因为祂是天主的智慧和德能；但为我们，圣言成了血肉，好居住在我们中间，我们只能以此方式首先接受祂。或许这就是为何祂不但是一切受造物的首生者，也被称为人亚当。因为保禄说祂是亚当：\BibleRef{格前 15:45} \Quote{最后的亚当成了赋予生命的灵。}\par

% structure: p0043 source=h2 target=subsection reason=relative-heading-rank; base=h2

\subsection{二十一、(6) 关于设计与实行。}

此外，我们谈论行动的元始，其中有一个在开始之后才显现的设计。可以考虑，上智是否应被视为天主工程的元始，因为正是在这个意义上，上智是它们的本原。\par

% structure: p0045 source=h2 target=subsection reason=relative-heading-rank; base=h2

\subsection{二十二、圣言在元始，即在上智内，上智在万物存在之前已在理念中包含万物。基督作为上智的特性先于祂的其他特性。}

关于\\OriginalTerm{元始}{arche}这个词，我们同时想到许多含义。现在我们必须问自己，对于经文\\BibleQuote{在起初已有圣言}，我们应采纳哪一种含义。显然，我们可以立即排除与转变或道路及其长度相关的含义。也很明显，与起源相关的含义也不适用于我们的目的。然而，人们可能会想到那种指向作者、指向产生效果者的含义，例如我们读到：\\BibleQuote{天主命令，它们就被创造}。因为基督在某种意义上就是\\OriginalTerm{工匠}{demiurge}，圣父对祂说：\\BibleQuote{要有光}和\\BibleQuote{要有穹苍}。但基督作为元始而是工匠，因为祂是智慧。正是因祂是智慧，祂才被称为元始。因为智慧在\\BibleBook{}{箴言}中说：\\BibleRef{箴 8:22} \\BibleQuote{天主在祂道路的元始创造了我在先，为祂的事工}，这样圣言就存在于一个元始中，即智慧中。从对万有之沉思与思想的结构来看，这被视为智慧；但从理性受造物认识思想对象的那一面来看，这被视为圣言。这并不奇怪，因为我们先前说过，救主是许多善的事物，如果祂本身包含第一、第二和第三等级的思想的话。这正是若望论及圣言时所表达的：\\BibleRef{若 1:3} \\BibleQuote{那被造的在祂内是生命}。于是生命在圣言中产生了。一方面，圣言就是基督，圣言，祂与圣父同在，万有藉祂而造成；另一方面，生命就是天主子，祂说：\\BibleRef{若 14:6} \\BibleQuote{我是道路、真理和生命}。因此，如同生命在圣言中产生，圣言也在元始中。然而，请思考我们是否可以将\\Quote{元始}的这一含义用于经文\\BibleQuote{在起初已有圣言}，从而得出这样的含义：万有是按照智慧、按照存在于祂思想中的系统模式而产生的。因为我认为，正如房屋或船舶是根据建筑师或设计者的草图建造和成型，房屋或船舶的元始存在于他心中的草图和设计算式中，同样，万有是按照天主在智慧中清楚设定的未来设计的样式而产生的。我们还应该补充说，天主可以说是创造了有灵的智慧之后，就让智慧从她内部的理型中，将万有的实际涌现、塑造和形式赋予存在之物和物质。但我认为，如果允许这样说，真实存在的元始就是天主子，祂说：\\BibleRef{默 22:13} \\BibleQuote{我是元始和终结，\\OriginalTerm{阿尔法}{\GreekText{Α}}和\\OriginalTerm{奥米加}{\GreekText{Ω}}，最初的和最末的}。然而，我们必须记住，祂并非对于祂的每一个名称都是元始。因为当祂是生命时，祂怎么可能是元始呢？生命在圣言中产生，而圣言显然是生命的元始。颇为明显的是，祂作为死者中的首生者也不是元始。如果我们仔细审视祂所有的称号，就会发现祂只是在作为智慧时才被称为元始。即便作为圣言，祂也不是元始，因为圣言在元始中。因此，人或许敢于说，智慧先于一切表达在\\Quote{首生者}称号中的思想。现在，天主完全是唯一且单纯的；但我们的救主，由于许多原因，既然天主立祂为赎罪祭和整个受造界的初果\\BibleRef{罗 3:25}，便成为许多事物，甚至可能是所有事物；整个受造界，只要能够获得救赎，就需要祂。因此，祂成了人的光，因为人因邪恶而黑暗，需要那照耀在黑暗中而不被黑暗所胜的光；如果人不在黑暗中，祂就不会成为人的光。关于祂是死者中的首生者，也可观察到同样的事。因为如果女人没有受骗，亚当没有堕落，为不朽而造的人得到了不朽，那么祂就不会下到坟墓，也不会死去——既然没有罪——祂对人的爱也不会要求祂死；如果祂没有死，祂就不能成为死者中的首生者。我们也可以问，如果人没有被抛给无理性的野兽并变得像它们一样，祂是否会成为牧人。因为如果天主拯救人和野兽，祂就通过赐予它们牧人来挽救那些祂所救的野兽，因为它们不能有王。因此，如果我们收集耶稣的称号，问题就出现了：其中哪些是后来授予祂的，并且如果圣人们从一开始就持守于真福，这些称号就永远不会具有如此的重要性。也许惟独智慧会留下，或许圣言也会留下，或许生命，或许真理，而不是其他那些祂为我们而取的称号。那些需要天主子却已经变得不需要祂作为医治病人的医生、作为牧人、作为救赎的人，实在有福；他们只需要祂作为智慧、作为圣言和正义，或者任何其他适合于如此完美以致能以祂最美丽的身份接受祂的人的称号。关于\\BibleQuote{在起初}的讨论就到此为止。\par

% structure: p0047 source=h2 target=subsection reason=relative-heading-rank; base=h2

\subsection{二十三、\\Quote{圣言}这一称号应按与基督其他称号相同的方法来解释。天主的圣言并非天主的单纯属性，而是一个独立的位格。当祂被称为圣言时，其含义是什么。}

然而，让我们更仔细地思量那\BibleQuote{在起初已有的圣言}是什么。当我思量关于基督所说的事，甚至那些真诚信仰祂的人所讲论的，我不禁时常感到惊奇。虽然有无数名号可以应用于我们的救主，但他们忽略了大部分；即便他们记得这些名号，他们也宣称这些称号不应按字面意义理解，而是比喻性地。但一触及\Quote{\OriginalTerm{圣言}{Logos}}这个称号，并重复说唯基督是天主圣言时，他们却又不一致，不像对待其他称号那样，去探寻\Quote{圣言}一词背后的含义。我对一般基督徒在此事上的愚昧感到惊讶。我不讳言，这根本就是愚昧。天主子在一处说：\BibleQuote{我是世界的光}；在另一处说：\BibleQuote{我是复活}；又说：\BibleQuote{我是道路、真理和生命}。经上也记着：\BibleQuote{我是门}；我们也有这话：\BibleQuote{我是善牧}。当撒玛黎雅妇人说：\BibleQuote{我们知道默西亚——即称为基督的那位——要来；祂来了，必将一切事告诉我们}，耶稣回答：\BibleQuote{同你说话的，就是祂}。再者，当祂洗门徒的脚时，用这些话\BibleRef{若 13:13}宣称自己是他们的师傅和主：\BibleQuote{你们称我师傅和主，说得对，我本来就是}。祂也明确宣布自己是天主子，祂说：\BibleRef{若 10:36}\BibleQuote{父所祝圣并派到世界上来的那位，你们却对祂说：你说亵渎的话，因为我说了我是天主子？}以及\BibleRef{若 17:1}\BibleQuote{父啊，时辰到了，请光荣祢的子，使子也光荣祢}。我们也发现祂宣称自己是一位君王，如祂回答比拉多的提问\BibleRef{若 18:33}\BibleQuote{祢是犹太人的君王吗？}说：\BibleQuote{我的国不属于这世界；如果我的国属于这世界，我的臣仆必要战斗，使我不致被交给犹太人；但我的国不是从这世界来的}。我们也读过这些话：\BibleQuote{我是真葡萄树，我父是园丁}；又说：\BibleQuote{我是葡萄树，你们是枝条}。此外，还有这话：\BibleQuote{我是生命的食粮，从天上降下来，赐给世界生命}。目前这些经文已经足够，它们取自福音的宝库，其中祂都声称自己是天主子。但在\WorkTitle{若望的默示录}中，祂也说：\BibleRef{默 1:18}\BibleQuote{我是元始和终末，是永生的，我曾死过。看，我永远活着}。又说：\BibleRef{默 22:13}\BibleQuote{我是\OriginalTerm{阿耳法}{\GreekText{Α}}和\OriginalTerm{敖默加}{\GreekText{Ω}}，是元始和终末，是初和终}。此外，细心研读圣经的人可以从先知书中收集不少类似段落，比如祂称自己（\BibleRef{依 49:2}）为\Quote{被选的箭}、\Quote{天主的仆人}以及\Quote{外邦人的光}。依撒意亚也说：\BibleRef{依 49:6}\BibleQuote{祂自母腹中就呼我的名字，使我的口如利剑，将我藏在祂手的荫庇下，祂对我说：你是我的仆人以色列，我将在你内受到光荣}。稍后又说：\BibleQuote{我的天主将成为我的力量。祂对我说：你被称为我的仆人，复兴雅各伯的支派，领回以色列的离散者，这是一件大事。看，我立你作外邦人的光，使你成为救恩，直到地极}。在耶肋米亚书中，\BibleRef{耶 11:19}祂也自比羔羊：\BibleQuote{我像一只温顺的羔羊被牵去宰杀}。这些以及其他类似的话，祂都应用于自己。除此之外，还可以从福音、宗徒书信和先知书中收集无数应用于天主子称号，如同福音作者陈述他们对祂的观点，宗徒根据所学颂扬祂，或先知预言祂的来临，用各种名称预报有关祂的事。于是若翰称祂为天主的羔羊，说：\BibleRef{若 1:29}\BibleQuote{看，天主的羔羊，除免世罪者}，并用这些话宣告祂是个人：\BibleRef{若 1:30}\BibleQuote{这就是我以前所说的：在我以后来的一位，成了在我以前的，因祂原是先我而有的}。若望在他的公函中说，祂是我们灵魂在父面前的\OriginalTerm{护慰者}{Paraclete}：\BibleQuote{如果有人犯了罪，我们在父那里有一位护慰者，就是正义的耶稣基督}；又补充说祂是我们罪的\OriginalTerm{赎罪祭}{propitiation}；同样保禄也说祂是赎罪祭：\BibleQuote{天主设立了祂作赎罪祭，通过信祂的血，以显示祂的正义，因为祂用忍耐宽容了先时所犯的罪过}。同样，按保禄所言，祂被宣称为天主的智慧和天主的德能，如在致格林多人书中所记：\BibleQuote{基督是天主的德能和天主的智慧}。又补充说祂也是圣化和救赎：\BibleQuote{祂是由天主成为我们的……}他说，\BibleQuote{智慧、正义、圣化和救赎}。但他在写给希伯来人的书信中也教导我们，基督是\OriginalTerm{大司祭}{High-Priest}：\BibleRef{希 4:14}\BibleQuote{既然我们有一位伟大的大司祭，祂已穿透诸天，就是天主子耶稣，就让我们坚持所宣认的信仰}。先知们还有其它名称来称呼祂。雅各伯祝福诸子时说（\BibleRef{创 49:10}）：\BibleQuote{犹大，你的兄弟将赞美你；你的手在你仇敌的颈上；你父亲的儿子要向你俯伏。犹大是一只小狮；我儿，你由猎物上来；你屈身伏卧，如雄狮，如母狮，谁能惹起你？}我们现在不能在这些话上久留，以证明论及犹大的话适用于基督。可能有人引用反对意见，即：\BibleQuote{权杖不离犹大，权柄不离他双足之间，直到那应得者来到}。这最好留待他时澄清。但依撒意亚确知基督被以雅各伯和以色列之名提及，他说：\BibleRef{依 42:1}\BibleQuote{雅各伯是我的仆人，我必扶持祂；以色列是我的选民，我心悦纳祂。祂将向外邦人宣布正义。祂不争竞，不喧嚷，街上也无人听见祂的声音。压伤的芦苇，祂不折断；将熄的灯心，祂不吹灭，直到祂使正义得胜。外邦人将仰望祂的名}。这些预言所指的就是基督，玛窦在他的福音中指明，他凭记忆引用说：\BibleQuote{这话是要应验……祂不争竞，不喧嚷}等等。达味也被称为基督，如厄则克耳在预言中对牧者所说，加上出自天主的话：\BibleRef{则 34:23}\BibleQuote{我要兴起我的仆人达味，他必牧放他们}。因为要兴起作圣者牧者的，不是先祖达味，而是基督。依撒意亚也称基督为\Quote{苗芽}和\Quote{花朵}：\BibleRef{依 11:1}\BibleQuote{由叶瑟的树干将生出一条嫩枝，由它的根上将开出一朵花。上主的神将安息在祂身上，是智慧和聪敏的神，是超见和刚毅的神，是知识和敬畏上主的神}。在圣咏集中，我们的主被称为\Quote{石头}：\BibleQuote{匠人弃而不用的石头，却成了屋角的基石。这是上主的所作所为，在我们眼中神奇莫测}。福音证明，路加在宗徒大事录中也证实，这块石头就是基督。福音如此说：\BibleRef{玛 21:42}\BibleQuote{你们没有读过这段经文吗？匠人弃而不用的石头，却成了屋角的基石。凡是跌在这石头上的，必被摔碎；这石头落在谁身上，就把他压成粉碎}。路加在宗徒大事录中写道：\BibleRef{宗 4:11}\BibleQuote{这石头是你们建筑工人弃而不用的，却成了屋角的基石}。应用于救主的名号中，有一个是祂自己未曾说出的，而是若望所记载的——\BibleQuote{在起初与天主同在的圣言，天主圣言}。我们值得花片刻留意那些忽视我们所提大部分伟大名号、反将这个视为最重要名号的学者。对于前几个称号，他们会寻求任何人的解释；但面对这一个，他们却另辟蹊径，问道：\Quote{当天主子被称为\InnerQuote{圣言}时，祂是什么？}他们最常用的经文是圣咏中的：\BibleQuote{我的心灵涌出优美的言辞}；他们想象天主子是父的言语，好比储存在音节中的话语，因此如果我们进一步审察他们，他们并不承认祂有任何独立的\OriginalTerm{本体}{hypostasis}，对祂的本质也不清楚。我不是说他们混淆了祂的属性，而是指祂自有一位本质这个事实。因为没有人能理解，被称为\Quote{圣言}的，怎能是一位子？这样的一个有生命的\Quote{圣言}，不与父分离，因此没有自立性，就不是子；如果祂是子，就让他们说天主圣言是一个单独的存在，并有自己的本质。因此，我们坚持，正如我们在上述每个称号的情形中，从称号转向它所暗示的概念，应用并证明天主子如何恰当地被这样描述，当我们发现祂被称为\Quote{圣言}时，也必须遵循同样的方法。在所有上述情况中，不拘泥于所用的术语，却探究基督应在何种意义上被理解为\Quote{门}，如何是\Quote{葡萄树}，为何是\Quote{道路}；但唯独在祂被称为\Quote{圣言}的情形中，却走另一条路，这是何等任意！因此，为使我们关于\Quote{天主子如何是圣言}这一问题所说的话更具权威性，我们必须从我们最初谈论的那些应用于祂的名号开始。我们无法否认，这在有些人看来是多余和离题的，但有思想的读者不会认为探究这些称号所表达的概念是无用的；留意这些事，将为下文铺平道路。一旦我们进入关于救主的神学讨论，我们以所能有的勤奋去寻求和发现关于祂所教导的种种，我们必然会更了解祂，不仅了解祂作为\Quote{圣言}的特性，也了解祂的其他特性。\par

% structure: p0049 source=h2 target=subsection reason=relative-heading-rank; base=h2

\subsection{二十四、基督为光；祂与祂的门徒如何成为世界的光。}

於是，祂說祂是世界的光；我們必須連同這個稱號一起考察那些與之平行的稱號；事實上，有些人認為它們不僅是平行的，而且與它相同。祂是真光，也是外邦人的光。在我們面前的這部福音的開端，祂是人的光：\BibleQuote{那受造的，} \BibleRef{若 1:3} 它說，\BibleQuote{在祂內是生命，這生命是人的光；光在黑暗中照耀，黑暗沒有勝過它。} 稍後，在同一段落中，祂被稱為真光：\BibleRef{若 1:9} \BibleQuote{那照亮每個人的真光，正在進入這世界。} 在\WorkTitle{依撒意亞}中，祂是外邦人的光，如我們之前所說。\BibleQuote{看，\BibleRef{依 49:6} 我立祢作外邦人的光，使祢成為救恩，直到地極。} 如今，世界可感覺的光是太陽，其後非常相稱的是月亮，同一稱號也可應用於星辰；但這些世界的光在梅瑟記載中是在第四日受造的，它們照在地上的事物上，它們不是真光。但是救主照耀在擁有理智和最高理性的受造物上，使它們的心靈能看見它們適當的視覺對象，因此祂是理智世界的光，也就是說，是在可感覺世界中的理性靈魂的光；如果在這世界中有任何超出這些的存在，祂宣布自己是我們的救主。祂實在是那個世界中最決定性和最卓越的部分，我們不妨說，祂是產生主之大日的太陽。關於這日，祂對那些分享祂光的人說：\BibleQuote{趁著白日，\BibleRef{若 9:4} 做工；黑夜來到，就沒有人能做工了。我在世界上的時候，是世界的光。} 然後祂對門徒說：\Quote{你們是世界的光，} 以及 \Quote{讓你們的光在人前照耀。} 因此，我們看到教會，新娘，與月亮和星辰相似；門徒們有一道光，是他們自己的或是從真太陽借來的，所以他們能照亮那些自己沒有任何光源的人。我們可以說，保祿和伯多祿是世界的光；而他們那些自己得光照但不能光照別人的門徒，就是宗徒們作為其光的世界。但是救主作為世界的光，照亮的不是身體，而是藉祂無體的能力照亮無體的理智，使我們在如同被太陽光照後，能夠辨別其餘屬於心靈的事物。正如太陽照耀時，月亮和星辰失去發光的能力，同樣，被基督光照並接受祂光線的人，不需要服務的宗徒和先知——我們必須有勇氣宣告這個真理——也不需要天使；我還要說，當他們成為那首生之光的門徒時，甚至不需要更大的權能。對於那些不接受基督太陽光線的人，服務的聖人們提供一種遠小於前者的光照；這種光照恰是那些人能夠接受的，並且完全充滿他們。基督，世界的光，又是真光，與感覺的光相區別；凡是可感覺的，都不是真的。然而，雖然可感覺的與真的不同，但並不意味著可感覺的就是假的，因為可感覺的與理智的之間可以有類比；並非所有不是真的東西都能正確地被稱為假的。現在我要問，世界的光與人的光是否是同一回事；我認為前者所指的光的能力比後者更高，因為世界在某種意義上不僅是人。保祿在他的第一封致格林多人的書信中表明：\BibleRef{格前 4:9} \BibleQuote{我們成了一臺戲給世界、天使和人觀看。} 從某種意義上也可以認為，世界就是那個受造界，它正從敗壞奴役中被解放，進入天主子女光榮的自由；受造界以熱切期待等候著天主子女的顯揚。我們還提請注意，可以在\BibleQuote{我是世界的光}這句話和\BibleQuote{你們是世界的光}這句話之間進行比較。有些人認為耶穌的真門徒比其他受造物更大；有些人從這些門徒的自然成長中尋找原因，另一些人從他們更艱苦的奮鬥中推論出來。因為那些在血肉之中的存在，比那些在有形體的身體之中的存在，有更大的勞苦和更充滿危險的生命；天上的光體若披上地上的身體，可能無法完成我們這免於危險和錯誤的生命。傾向於這個論證的人可以引用聖經中那些論及人的最高貴經文，以及福音直接對人說的這一事實；對於受造界，或按我們的理解，對於世界，說得並不那麼多。我們讀到，\BibleRef{若 17:21} \BibleQuote{如同我和祢是一體，願他們也在我們內合而為一，} 以及 \BibleQuote{我在哪裡，我的僕人也要在那裡。} 這些說法顯然是關於人的；而對於受造界，則說它從敗壞奴役中被解放，進入天主子女光榮的自由。可以補充說，即使它被解放，也不會參與天主子女的榮耀。持此觀點的人也不會忘記，那萬有的首生者，尊崇人超越一切，成為人；而且，成為東方之星的不是天空中的任何星座，而是另一種秩序中的一位，為這目的而被指定，服務於對耶穌的認識；無論它與其他星辰相似，還是可能更好，它作為那最卓越者的標記。如果聖人們以他們的困苦誇耀，因為\BibleRef{羅 5:3} \BibleQuote{困苦生忍耐，忍耐生老練，老練生望德，望德不使人蒙羞，} 那麼，受痛苦的受造界不能有與人同樣的忍耐，同樣的老練，同樣的望德，而是有另一種程度，因為\BibleRef{羅 8:20} \BibleQuote{受造界被屈從於虛妄，不是出於自願，而是由於那使它屈從的一位，基於望德。} 現在，那些不願將如此偉大的屬性賦予人的人，將轉向另一方向，說那屈服於虛妄的受造物呻吟並遭受比在這帳幕中呻吟者更大的痛苦，因為她在服事虛妄時受苦到了時間的極限——不，比人長久許多倍？因為她做這事不是出於自願，而是因為屈從於虛妄違反她的本性，並且她沒有得到生命的最佳安排，就是當她被釋放、世界被毀滅並從身體的虛妄中得釋放時所將得到的。然而，這裡我們似乎伸展得太遠，目標超出了目前問題的需要。因此，我們可以回到出發點，問：救主為何被稱為世界的光、真光和人的光？現在，我們看到祂被稱為真光是相對於世界可感覺的光而言；而世界的光與人的光是同一回事，或者我們至少可以探究它們是否相同。這個討論並非多餘。有些學者從救主是聖言這一陳述中根本得不到任何東西；對我們來說，重要的是確保我們自己不是任性而關注這個概念，以免被指責。如果它可以按比喻意義理解，我們就不應按字面理解。當我們對\Quote{世界的光}這一表達以及上面提到的許多類似術語運用神秘和寓意方法時，對於這個表達也應當如此。\par

% structure: p0051 source=h2 target=subsection reason=relative-heading-rank; base=h2

\subsection{基督是复活。}

现在，祂被称为人的光、真光，以及圣言的光，因为祂照亮并照耀人的更高部分，或者简言之，一切理性受造物的更高部分。同样地，由于祂的能力，祂使人除去旧的死气，穿上那\Emph{卓越的}生命，以致那些真正接受了祂的人从死者中复活，因此祂被称为复活。祂这样做不仅是在一个人说\BibleQuote{我们藉着圣洗与基督同葬，也一同复活了}（\BibleRef{罗 6:4}）的时刻，更是当一个人脱去所有属于死亡的事物，在此世行走在属于祂——圣子——的新生命中时。我们常\BibleQuote{在身体上带着主耶稣的死}（\BibleRef{格后 4:10}），借此收获极大的益处，\BibleQuote{使主耶稣的生命也彰显在我们身上。}\par

% structure: p0053 source=h2 target=subsection reason=relative-heading-rank; base=h2

\subsection{基督是道路。}

然而，那种在智慧中的进步，被那些在其中寻求救恩的人发现，能为他们完成所需之事，既关乎在圣言中阐明真理，也关乎按照真义行事，这使我们理解基督如何是道路。在这条路上，我们不可携带任何东西（\BibleRef{玛 10:10}），既不带钱囊，也不带外衣；我们必须旅行，连手杖也不可带，脚上也不可穿鞋。因为这条路本身足以供应我们旅途的一切所需；凡行在这条路上的人，一无所缺。他穿着适合赴婚宴的衣裳；在这条路上，他不会遇到任何令他烦恼的事。撒罗满说：\BibleQuote{没有人能寻得蛇在岩石上的道路。}（\BibleRef{箴 30:19}）我要加上说，也没有任何其他野兽的道路。因此，在这条路上不需要手杖，因为其上没有任何敌对生物的踪迹，它的坚硬（也因此被称为\OriginalTerm{岩石}{petra}）使它无法容纳任何有害之物。\par

% structure: p0055 source=h2 target=subsection reason=relative-heading-rank; base=h2

\subsection{基督是真理。}

此外，独生子是真理，因为祂按照圣父的旨意，完全清晰地包含了万物的全部理据，并且作为真理，按各受造物的相称程度分施予它。若有人询问，圣父按其丰富的上智和知识的深渊所知道的一切，我们的救主是否也知道；又若他以为如此可以光荣圣父，而表明圣父知道的一些事是圣子所不知道的，尽管圣子本可与那无始的天主有同等的理解，我们必须提醒他，正因为祂是真理，祂才是救主，并补充说，如果祂是完全的真理，那么没有一样真理是祂所不知道的；真理不可因缺乏那些按照那些人的说法只有圣父知道的事物而跛行。否则，就请指出有些事物是知道的，而真理的名称却不适用于它们，它们是在真理之上的。\par

% structure: p0057 source=h2 target=subsection reason=relative-heading-rank; base=h2

\subsection{基督是生命。}

显然，那纯粹而不掺杂任何其他元素的生命之原则，居于那一切受造物的首生者内；从祂那里，凡有分于基督的人，就活出那真实的生命；而所有被认为与此分离而活着的人，既然没有真光，也就没有真实的生命。\par

% structure: p0059 source=h2 target=subsection reason=relative-heading-rank; base=h2

\subsection{基督是门，也是牧人。}

但正如人不能在圣父内，或与圣父同在，除非从下往上攀升，先来到圣子的天主性，藉此被引领而达至圣父本身的真福，同样地，救主拥有\Quote{门}的称号。又因祂是爱人的，并赞同人灵向更好之事的冲动，甚至连那些不急于趋向圣言，而像羊一样有软弱和温顺、缺乏一切精确和理性的人，祂也是牧人。因为主拯救人和牲畜，以色列和犹大不仅撒了人的种子，也撒了牲畜的种子。\BibleRef{耶 31:27}\par

% structure: p0061 source=h2 target=subsection reason=relative-heading-rank; base=h2

\subsection{基督是受傅者（基督），也是君王。}

除了这些称号外，我们必须在著作的开端也考虑基督的称号，也要考虑君王的称号，并比较两者，以找出它们之间的区别。在第四十四篇圣咏中记载：\BibleQuote{你喜爱正义，憎恶不义，因此天主、你的天主、以喜油傅你，使你作基督，超越你的同伴。}祂的喜爱正义和憎恶不义因此是祂身上附加的资格；祂的受傅不是与祂的存在同时，也不是祂从起初所继承的。受傅是进入王权的象征，有时也是进入司祭职的象征；那么我们是否必须得出结论，天主子（Son of God）的王权不是继承的，也不是与生俱来的？但是，万物的首生者（First-born of all creation）本身岂不就是正义，却因为祂喜爱正义而后来成为君王，这怎能想象？我们不能看不出，祂作为人乃是基督，就祂的灵魂而言，那是人的灵魂，会烦恼和痛苦；但就祂内的神性而言，祂被理解为君王。我在第七十一篇圣咏中找到了支持，该篇说：\BibleQuote{天主，请将祢的判断权赐给君王，将祢的正义赐给君王的儿子，使他以正义审判祢的百姓，以公平审判祢的穷人。}这篇圣咏虽然是对\Person{撒罗满}说的，但显然是基督的预言；值得一问的是，预言希望天主将判断权赐给哪位君王，赐给哪位君王的儿子，以及所说的是哪位君王的正义。于是我认为，所谓君王是万物的首生者的主导本性，因着其卓越而获得判断权；而那被祂所取、由这本性按正义而形成和塑造的人，便是君王的儿子。我更加倾向于这样认为，因为这两个存在在此处被放在同一句话中，被说成好像不是两个，而是一个。因为救主使二者成为一体，\BibleRef{弗 2:14} 即是说，祂使它们按照那在万有之先已在祂内成为一体的两个的原型而成为一体。我所说的两个是指人性，因为每个人的灵魂都与圣神相融合，每个得救的人也因此成为属神的。正如有些人，因着他们温良平和的性情，基督是他们的牧人，如我们先前所说，尽管他们较少受理性引导；同样，也有些人，基督是他们的君王，即那些在接近信仰时更多受理性部分引导的人。在君王之下的人们中也有区别：有些人以更奥秘、更隐藏、更神圣的方式体验祂的统治，另一些人则以不那么完善的方式。我认为，那些由理性引导、脱离一切感官作用而看见无形之物——\Person{保禄}所说的\BibleQuote{不可见的}或\BibleQuote{未见之事}——的人，是由独生子（Only-begotten）的主导本性所统治；而那些只进步到与感官事物打交道的理性，并因这些事物而荣耀其造物主的人，也是由圣言（Word），由基督所管理。我们若在救主身上区分这些概念，无需感到反感；我们在祂的本体中也作同样的区分。\par

% structure: p0063 source=h2 target=subsection reason=relative-heading-rank; base=h2

\subsection{三十一、基督为师为傅。}

所有人都清楚，我们的主如何为那些努力追求敬虔的人作教师和解释者，另一方面，为那些有奴仆之心、因恐惧而受束缚（\BibleRef{罗 8:15}）、在进步中奔向智慧、并配得拥有智慧的人作主人。\Quote{仆人不知道主人所愿的}，因为他已不再是他的主人，而成了他的朋友。主亲自教导这一点，因为祂对那些仍是仆人的听众说：\BibleRef{若 13:13}\BibleQuote{你们称我为师傅、为主，说得对，因为我本来就是}，但在另一处，\BibleRef{若 15:15}\BibleQuote{我不再称你们为仆人，因为仆人不知道主人的意愿，但我称你们为朋友}，因为\BibleRef{路 22:28}\BibleQuote{你们在我的各样诱惑中常同我在一起}。因此，那些生活在恐惧中——这恐惧是由天主向那些不是好仆人的所要求的，正如我们在玛拉基亚所读到的：\BibleQuote{我若是主人，我的敬畏在哪里呢？}——的人，是那位被称为他们救主的主人的仆人。\par

% structure: p0065 source=h2 target=subsection reason=relative-heading-rank; base=h2

\subsection{三十二、基督为子。}

但没有一个证词明确陈述救主的高贵诞生；但当这话对祂说：\BibleQuote{祢是我的儿子，我今日生祢}，这是由天主对祂说的，在天主那里，所有的时间都是今日，因为我认为，天主那里没有晚上，也没有早晨，只有时间延伸，与祂的无始及不可见的生命一起。在祂那里，儿子被生的日子就是今日，因此祂诞生的开始未被发现，正如祂的日子没有被发现一样。\par

% structure: p0067 source=h2 target=subsection reason=relative-heading-rank; base=h2

\subsection{三十三、基督为真葡萄树，并为粮。}

在我们所说之上，还必须加上圣子如何是真葡萄树。那些以与先知恩宠相称的方式理解这句格言的人，将毫不困难地领会这一点：\Quote{酒使人心愉悦}。因为如果心脏是理智的部分，而使它喜悦的是那最令人喜悦的、可饮用的圣言，它使我们脱离人性事物，让我们感到自己被灵感充满，并用一种并非非理性而是神圣的醉意使我们沉醉——我认为，这种醉意就是\Person{若瑟}使他的兄弟们欢乐的那种（\BibleRef{创 43:34}）——那么，显然那位这样带来酒使人心中喜悦的，就是\Emph{真}葡萄树。祂是真葡萄树，因为祂所结的葡萄就是真理，门徒是祂的枝条，他们也结出真理作为果实。要指出葡萄树与食粮之间的区别有些困难，因为祂不仅说祂是葡萄树，还说祂是生命之粮。或许，正如食粮滋养人、使人强壮，并被称为坚固人心，而酒则相反，使人愉悦、欢喜并消融人心；同样，伦理学问，为学习并实践它的人带来生命，是生命之粮，却不能恰当地称为葡萄树的果实；而隐秘而神秘的思辨，使人心欢喜，使接受它们的人感到被灵感充满，在主内喜乐，并渴望不仅被滋养，而且要获得幸福，这些被称为真葡萄树的汁液，因为它们是自它流出的。\par

% structure: p0069 source=h2 target=subsection reason=relative-heading-rank; base=h2

\subsection{三十四、基督作为首先的和末后的；祂也是介于二者之间的。}

此外，我们必须询问：在默示录中祂为何被称为\Quote{首先者与最后者}？而作为首先者，祂与阿耳法及元始有何不同？作为最后者，祂与敖默加及终结又有何不同？在我看来，现存的理性受造物具有多种形式，有些是首先的，有些是第二的，有些是第三的，如此直到最后。然而，准确说出何者为首先、第二者为怎样的存在、何者真正可称为第三，并推至系列末端——这并非人的工作，而是超越我们本性的。尽管如此，我们仍将勉为其难，在此稍作停留，并对该事略作观察。有些神祇，天主是他们的神，正如我们在预言中所听见的：\Quote{感谢众神之神}，以及\Quote{众神之神发言，并召唤了大地}。如今，根据福音，\BibleRef{玛 20:2} \BibleQuote{天主不是死人的天主，而是活人的天主}。因此，那些神祇是活着的，天主是他们的神。宗徒在致格林多人书信中也写道：\BibleRef{格前 8:5} \BibleQuote{因为虽有称为神的，或在天上或在地上，就如那许多的神，许多的主}，所以我们论及这些神祇时，视其为真实存在。此外，除了那些以天主为神的神祇之外，还有另一些被称为宝座者，又有被称为宰制者、主权者，以及权能者。\BibleRef{弗 1:21} \BibleQuote{远超一切称名的名号，不但在今世，而且在来世}，这句话使我们相信，除了这些之外，还有另一些对我们较为陌生的存在；其中一类，希伯来人称之为\OriginalTerm{撒巴依}{Sabai}，由此形成撒巴奥特（\OriginalTerm{撒巴奥特}{Sabaoth}），后者是他们的统治者，而祂正是天主。在所有这些之上，还有有死的理性存在者——人。万有的天主首先创造了一类在尊荣上位居首位的理性受造物；我认为这就是被称为神祇的，第二等，姑且称之为宝座，第三等，无疑就是宰制者。如此我们按次序下降到最后的理性族类，或许那正是人。救主便以超乎保禄的方式，成了对一切人成为一切，为要赢得一切人或成全他们；显然，祂对众人成了人，对天使则成了天使。关于祂成为人，任何信徒均无疑问；但关于祂成为天使，若我们仔细观察天使的显现和言辞，便会找到相信的理由：在某些情形中，天使的能力似乎属于祂。有几处经文，天使说话的方式暗示了这一点，例如\Quote{上主的天使在火焰中显现。他说：我是亚巴郎的天主，依撒格的天主，雅各伯的天主}。依撒意亚也说：\BibleRef{依 9:6} \BibleQuote{祂的名字称为伟大谋士的天使}。因此，救主是首先的与最后的，并非因为祂不包括中间者，而是藉此表明祂成了万有。然而，请思考：最后者是人，还是那些称为地底下的存在，其中有魔鬼（全部或部分）？我们也必须询问，救主藉达味先知所说的那些事：\Quote{我成了无人援助的人，在死者中自由}。祂由童贞女降生及其卓越的生活，已证明祂超越人，在死者中亦然。祂是那里唯一自由的人，祂的灵魂没有被遗弃在地狱。如此，祂便是首先的与最后的。再者，如果存在天主的文字（确有其事），圣人们藉着阅读这些文字，可说他们已读了写在诸天版上的内容；那么，这些用以解读天上事物的文字，其细分为微小部分、从\OriginalTerm{阿耳法}{\GreekText{Α}}直到\OriginalTerm{敖默加}{\GreekText{Ω}}的，就是天主之子。此外，祂是元始与终结，但并非在祂所有的方面都相同。因为照箴言所教导的，祂是元始，乃因祂是智慧；经上记载：\Quote{上主在祂道路的开始，在祂化工之先，已创建了我}。就祂是圣言而言，祂不是元始。\Quote{圣言在起初已有}。因此，在祂的诸方面中，有一个先来且是元始，第二个在元始之后，第三个，如此直到终结，就如祂说过：我是元始，因为我是智慧；第二者，或许因为我是不可见的；第三者，因为我是生命，因为\Quote{在祂内受造的是生命}。一个善于考察并辨别圣经含义的人，无疑会发现系列中的许多环节；但我不能说他能找到全部。\Quote{元始与终结}这个短语，我们通常用于一个完成统一体的事物；一所房屋的元始是其基础，终结是其护栏。我们不得不想到这预象，因为基督是被选为角石的石头，指向那得救者的庞大身体。因为基督，独生子，是万有并在万有中；祂作为元始临在于祂所取的人性中，作为终结临在于最后的一位圣人中，也临在于中间者；或者，祂作为元始临在于亚当中，作为终结临在于祂在世的生命，正如经上所说：\Quote{最后亚当成了赋予生命的灵}。这话与我们对首先与最后的解释十分吻合。\par

% structure: p0071 source=h2 target=subsection reason=relative-heading-rank; base=h2

\subsection{三十五、作为生活者与死者的基督。}

在关于首位与末位、开始与终结的论述中，我们将这些话有时指涉理性受造物的各种形态，有时指涉对天主圣子的不同理解。因此，我们区分了\Quote{首位}与\Quote{开始}，\Quote{末位}与\Quote{终结}，并区分了\OriginalTerm{阿尔法}{\GreekText{Α}}与\OriginalTerm{敖默加}{\GreekText{Ω}}的特殊含义。不难看出为何祂被称为\BibleRef{默 1:17}\BibleQuote{生活的与死去的}，并在死后成为永生者。因为祂起初的生命并未救助我们——我们深陷罪恶之中——祂便降临到我们的死亡之境，为使祂死于罪恶，我们\BibleRef{格后 4:10}\BibleQuote{身上常带着耶稣的死}，好能获得祂那永恒的生命。因为那些常在身上带着耶稣之死的人，也必获得耶稣的生命，彰显在他们的身上。\par

% structure: p0073 source=h2 target=subsection reason=relative-heading-rank; base=h2

\subsection{三十六、基督作为利剑。}

我们讨论过的\WorkTitle{新约}经文，是祂关于自己所说的话。然而，依撒意亚却说，\BibleRef{依 49:2}祂的口被圣父安置如同利剑，祂隐藏在圣父的手影之下，被造成如同精选的箭，藏在圣父的箭囊中，被万有的天主称为祂的仆人、以色列和外邦人的光明。天主圣子的口是利剑，因为\BibleRef{希 4:12}\BibleQuote{天主的话是生活的，有效能的，比各种双刃的剑还锐利，直穿入灵魂和神魂，关节与骨髓的分离，且能辨别心中的思想和意念。}确实，祂来不是为在地上，即有形可见的事物中，带来和平，而是带来利剑，并且——如果可以这样说——要斩断灵魂与肉体那不幸的友谊，使灵魂将自己交托给那与肉体对立的圣神，从而与天主缔结友谊。因此，按照先知的话，祂使自己的口成为利剑，成为锋利的剑。谁能看见那么多被神性之爱所伤的人，像\WorkTitle{雅歌}中那位抱怨自己受伤的女子：\BibleRef{歌 2:5}\BibleQuote{我因爱成伤}，并找到那为爱天主而伤了那么多灵魂的箭矢，除了那说\BibleQuote{祂使我成为精选的箭}的那一位呢？\par

% structure: p0075 source=h2 target=subsection reason=relative-heading-rank; base=h2

\subsection{三十七、基督作为仆人、天主的羔羊，以及若翰所不认识的那个人。}

再者，请任何人思量耶稣如何对待祂的门徒，并非如同坐席者，而是如同服事者；并且，天主子虽为那被罪奴役者的自由而取了仆人的形体，他就不难明白圣父为何对祂说：\Quote{祢是我的仆人}，随后又说：\Quote{祢被称为我的仆人，这是一件大事。}因为我们毫不迟疑地说，基督的善在更伟大、更神圣的光中，更肖似圣父的肖像，因为\BibleRef{斐 2:6}\Quote{祂贬抑自己，成了服从的，至死，甚至死在十字架上}，胜过祂视自己与天主同等为可紧持不放的，并且为世界的得救而退缩于成为仆人。因此，祂说：\BibleRef{依 49:5}，愿意教导我们：在接受了这种奴仆状态时，祂从圣父那里领受了伟大的恩赐：\Quote{我的天主必要作我的力量。祂对我说：祢被称为我的仆人，这是一件大事。}因为如果祂没有成为仆人，祂就不会复兴雅各伯的支派，也不会转回以色列遗民的心，也不会成为外邦人的光明，以达救恩直到地极。祂成为仆人并非大事，即使圣父称之为大事，因为这是相对于祂被称为无辜的羊和羔羊而言的。因为天主的羔羊成了如同被牵去宰杀的羔羊，为除免世罪。祂向万有供给理智（\OriginalTerm{逻各斯}{\GreekText{λογος}}），却成了如同在剪毛者前缄默的羔羊，好使祂的死——一种作为对抗敌对势力的药剂，也作为那些向真理敞开心灵者的罪的药剂——净化我们。因为基督的死削弱了那些攻击人类的权势，并以超出言词的能力从罪中释放了每位信徒的生命。既然祂不断除罪，直到一切仇敌都被消灭，最后是死亡，好使整个世界从罪中得释放，因此若翰指着祂说：\BibleRef{若 1:29}\Quote{请看，天主的羔羊，除免世罪者！}经文不是说祂将来要除去，也不是说祂现在正在除去，也不是说祂已除去而现在不再除去。祂除罪仍在进行，祂正从世界上的每个人除去罪，直到罪从整个世界除去，救主将预备并完成好的王国交给圣父——一个完全没有罪残留，因而预备好接受圣父为君的王国——并且另一方面，正等待接受天主所赐予的一切，完满的，每一部分，在那段经句\BibleRef{格前 5:28}应验的时刻：\Quote{为叫天主成为万物中的万有。}此外，我们听说一位被称为在若翰以后来的人，祂在我以前已存在，并且原先就在。这是要教导我们：天主子的那人性——与祂的神性混合的人——也比祂从玛利亚的诞生更古老。若翰说他不知道这个人，但当他还在依撒伯尔腹中，听到玛利亚的问候声传入匝加利亚妻子的耳中而欢喜跳跃时，难道他不知道祂吗？因此，请思量\Quote{我不认识祂}这句话是否可能指肉身存在之前的时期。虽然祂在取得肉身之前他并不认识祂，但在母腹中他已认识祂，也许他在这里学到了关于祂的某种新知识，超越他以前所知道的：即圣神降下并停留在谁身上，谁就是要以圣神和火施洗的那一位。他从母腹中已认识祂，但并非关于祂的一切。他也许不知道当看见圣神降下并停留在祂身上时，这就是那以圣神和火施洗的一位。然而，祂确实是一个人，而且是第一个人，若翰并不知道。\par

% structure: p0077 source=h2 target=subsection reason=relative-heading-rank; base=h2

\subsection{三十八、基督作护慰者、赎罪祭和天主的德能。}

但我们所提的名字没有一个表达祂在圣父面前为我们代表，为人的天性辩护并作补赎：护慰者、赎罪祭、和好如初。祂在若望书信中被称为护慰者：\BibleRef{若一 2:1}\Quote{如果有人犯了罪，我们在圣父那里有护慰者，耶稣基督，正义者。}在同一书信中，祂也被说成是为我们的罪的赎罪祭。同样，在《罗马书》中，祂被称为赎罪祭：\Quote{天主公开立定为赎罪祭，因信德。}这预表曾在圣殿的最内部分，至圣所，即置于两革鲁宾之上的金赎罪盖。然而，若没有天主的德能，祂怎能成为护慰者、赎罪祭和赎罪祭呢？这德能终结我们的软弱，洋溢于信徒的灵魂，并由耶稣施行——耶稣实在是先于这德能，且本身就是天主的德能——祂使人能说：\BibleRef{斐 4:13}\Quote{我藉着坚固我的耶稣基督，能应付一切。}由此我们知道，那自称\Quote{那称为\InnerQuote{伟大}的天主德能}的西满魔术士，被交于丧亡和毁灭，他和他的金钱一同毁灭。相反，我们承认基督为真正的天主德能，相信我们与祂共享，因祂是那德能，一切赋有能力的万物。\par

% structure: p0079 source=h2 target=subsection reason=relative-heading-rank; base=h2

\subsection{三十九、基督作智慧、圣化和救赎。}

然而，我们不可默然不语，祂本是天主的智慧，故以此名称之。因为万有之父、天主的智慧并非像人类思想的幻象那样，仅以幻像领会祂的实体。凡能构思一种无形体的存在，其具有多种思辨，延伸至存在之物的理性，是活的，并可说是赋有灵魂的，他必看见那超越一切受造物的天主的智慧如何论及自己，当她说：\BibleQuote{天主造了我，是祂道路的开始，为祂的工程。}（\BibleRef{箴 8:22}）藉此创造行动，整个受造物得以存在，并非不能容纳那神圣智慧——按此智慧受造物被带进存在；因为天主，按先知达味所说，以智慧创造了万物。许多事物藉智慧的帮助而存在，却不把握那创造它们的事物；确实少有事物不仅把握那关乎自身的智慧，而且把握那涉及许多其他事物的智慧，即基督——祂是全部的智慧。每位智者，按他拥抱智慧的程度，就如此分享基督，因祂就是智慧；正如每个大有能力的人，按他有能力的程度，就如此分享基督，因祂就是能力。关于圣化和救赎，也可作同样的思考；因为耶稣自己成了我们的圣化和救赎。我们各人以此圣化而被圣化，以此救赎而被救赎。此外，请思考宗徒所加的\Quote{为我们}一词是否具有特殊意义。他说：\Quote{基督成了从天主而来的我们的智慧、正义、圣化和救赎。}在其他经文中，他论及基督是智慧，没有任何此类限制，论及祂是能力，说基督是天主的德能和天主的智慧，尽管我们可能以为祂并非绝对地是天主的智慧或天主的德能，而只是为我们。现在，关于智慧和能力，我们有陈述的两种形式：相对的和绝对的；但关于圣化和救赎，情况并非如此。因此，请思考，既然\BibleQuote{那圣化的和那些被圣化的都是出于一}（\BibleRef{希 2:11}），那么圣父是否是那圣化我们之主的圣化，正如基督是我们的头，天主是祂的头。但基督是我们的救赎，因为我们成了囚徒，需要赎金。我不探究祂自己的救赎，因为祂虽在一切事上如我们一样受过试探，却没有罪，祂的仇敌从未使祂被俘。\par

% structure: p0081 source=h2 target=subsection reason=relative-heading-rank; base=h2

\subsection{四十、基督作为正义；作为德谬哥、善天主的代理人，以及大司祭。}

在厘清了\Quote{为我们}与\Quote{绝对地}——圣化和救赎是\Quote{为我们}而非绝对地，智慧和救赎既为我们又绝对地——之后，我们不可省略探究同一段经文中正义的地位。基督相对地为我们而成为正义，从这话清楚可见：\Quote{祂成了从天主而来的我们的智慧、正义、圣化和救赎。}如果我们在绝对意义上找不到祂是正义，正如祂是绝对地天主的智慧和德能，那么我们必须探究：对于基督自己，正如圣父是圣化，圣父是否也是正义。我们知道，在天主面前没有不义；\BibleQuote{祂是正义和圣洁的主}（\BibleRef{若 7:18}）；\BibleQuote{祂的审判是正义的}（\BibleRef{默 16:5}），且祂是正义的，祂以正义安排一切。\par

异端者出于自身目的，在义人与善人之间作了区分。他们并未将此事解释清楚，但他们认为造物主是义的，而基督之父是善的。我认为，若仔细考察，这一区分可应用于圣父与圣子：圣子是义，并已领受权力\BibleRef{若 5:27}执行审判，因为祂是人子，将以正义审判世界；而圣父则向那些藉圣子之义受过管教的人施行善。这是在圣子王国之后；那时圣父将在祂的工程中彰显祂的名——善，当天主成为万物中的万有。或许，救主藉祂的义在适当的时节预备一切，并藉祂的圣言、祂的安排、祂的惩戒，以及（若我可这样表达）祂属灵的医治辅助，安排万有最终领受圣父的善。正是出于对那善的感知，祂回应了那位以\Quote{善师}称呼独生子的人\BibleRef{希 2:9}，并说：\Quote{你为什么称我为善？除了独一天主、圣父以外，没有一个是善的。}我们在别处已处理过此事，尤其在讨论那位大于造物主的问题时；我们将基督视为造物主，而圣父是大于祂的那位。那么，祂是如此伟大的存在——护慰者、赎罪、赎罪祭、同情我们软弱的那位，\Quote{在一切事上像我们一样受过试探，只是没有罪}\BibleRef{希 4:15}；因此，祂是一位伟大的大司祭，将自己作为一次而永远的祭品献上，不仅为人，也为每一个理性的受造物。因为祂\Quote{免于天主}为每个人都尝了死味。在《希伯来书》的一些抄本中，经文为：\Quote{因天主的恩宠。}现在，无论祂是\Quote{免于天主}为每个人都尝了死味（祂不仅为人死，也为其他一切理性存在而死），还是\Quote{因天主的恩宠}为每个人都尝了死味，祂都是为一切\Quote{免于天主}的人而死，因为祂是因天主的恩宠为每个人尝了死味。若说祂只为人类的罪尝了死味，而不为除人以外任何堕入罪恶的其他存在（例如星辰）尝死味，那无疑是荒谬的。因为就连星辰在天主眼中也不洁净，正如我们在约伯传中所读到的：\BibleRef{约 25:5}\Quote{星辰在祂眼前也不洁净}，除非这被视为夸张。因此，祂是一位伟大的大司祭，因为祂将万有恢复至圣父的王国，并安排受造各部分的任何缺陷得以补足，充满圣父的光荣。这位大司祭，基于我们尚未论及的某种概念，被称为犹达斯，为要使那些隐秘的犹太人\BibleRef{罗 2:29}，不以雅各伯之子犹达为名，而是从祂得名，因为他们是祂的弟兄，并因所得的自由赞美祂。因为正是祂使他们自由，从敌人手中拯救他们，将手按在他们背上予以制服。当祂将敌对势力踏于脚下，独自在圣父面前时，祂便是雅各伯和以色列；因此，正如我们由祂成为光（因为祂是世界之光），同样，我们成为雅各伯，因为祂被称为雅各伯；成为以色列，因为祂被称为以色列。\par

% structure: p0084 source=h2 target=subsection reason=relative-heading-rank; base=h2

\subsection{四十一、基督作为杖、花、石。}

现在，祂从那些以色列子民所立的君王那里领受了王国，这君王并非奉天主命令，甚至没有求问天主，就开始了君主制。因此，祂作战是为上主作战，从而为祂的儿子、祂的子民预备和平；这或许就是祂被称为达味的理由。然后，祂被称为杖；\BibleRef{依 11:1}祂对于那些需要更严厉、更严格管教、且不肯顺从天主之爱与温和的人，便是如此。因此，若祂是杖，祂必须\Quote{发出}；祂并不留在自身内，而是似乎超越其先前状态。于是，祂发出并成为杖，却并不停留在杖中；在杖之后，祂成为上升的花；在成为杖之后，祂向那些因祂是杖而受到眷顾的人显明为花。因为\Quote{天主必以杖惩罚他们的过犯}，即基督。但\Quote{祂的慈爱却不离开他}，因为祂将怜悯他；凡圣子所怜悯的，圣父也怜悯。有一种解释使祂对不同的人成为杖和花：对需要惩罚的人成为杖，对正被拯救的人成为花；但我更倾向上述说法。然而，我们在此必须补充：或许，从终局看，若基督向某人成为杖，祂也向他成为花；而接受祂为花的人未必也必须认识祂是杖。然而，正如一朵花比另一朵更完美，植物即使不结完美果实也被说成开花，同样，完美的人领受基督超越花的部分。另一方面，那些认识祂为杖的人将同杖一起分享，不是祂的完美，而是花——那在果实之前的。最后，在我们谈到圣言一词之前，基督曾是一块石头，被建筑匠所弃，却成了屋角的头石；因为活石被建造在宗徒和先知这些石头之上，耶稣基督自己为主，作为屋角石，因为祂是在活人之地用活石所建建筑的一部分；因此，祂被称为石。我们所说的一切，是要表明那些人在基督有如此多名称时，仅取\Quote{圣言}这一称呼，而不像考察祂其他称号那样查问它是在何种意义上使用，其做法是何等反复无常、毫无根据；他们实在应当追问：论及天主子，说祂是圣言、是天主、在起初与圣父同在、万物由祂而造成，究竟是什么意思。\par

% structure: p0086 source=h2 target=subsection reason=relative-heading-rank; base=h2

\subsection{四十二、论基督作为圣言的多种方式。}

既然基督因祂启迪世界（祂是这世界的光）的活动而被称作世界之光，又因祂使真诚追随祂的人摆脱死亡、复活并穿上新生命而被称作复活，同样，因另一种活动，祂被称为牧者、教师、君王、\OriginalTerm{特选之箭}{Chosen Shaft}、仆人，此外还有\OriginalTerm{护慰者}{Paraclete}、\OriginalTerm{赎罪}{Atonement}和\OriginalTerm{赎罪祭}{Propitiation}。依同样方式，祂也被称为\OriginalTerm{圣言}{Logos}，因为祂从我们身上除去一切非理性，使我们真正有理性的，以致我们做所有事，甚至吃喝，都为光荣天主，并借着圣言为光荣天主而履行生命中的普通职责以及属于更高级阶段的事务。因为如果我们借着分享祂而被高举、被启迪，或许也被牧养和治理，那么显然，当我们以神圣方式成为有理性的，当祂驱走我们身上非理性和死亡的东西时，因为祂是圣言（理性）和复活。然而，要思考是否所有人都在某种程度上分享祂作为圣言的身份。关于这一点，宗徒教导我们，祂不可在被寻找者之外寻求，而凡专心寻求祂的人，就在自己内找到祂：\BibleQuote{你们心里不要说\BibleRef{罗 10:6}：\InnerQuote{谁能升到天上？}那就是把基督带下来；或者：\InnerQuote{谁能下到深渊？}那就是把基督从死者中带上来。但圣经说了什么？\InnerQuote{圣言离你很近，就在你口里，就在你心里。}}似乎基督自己就是那被寻求的圣言。但当主亲自说\BibleQuote{我若没有来向他们说过话，他们就没有罪；但如今他们的罪无可推诿。}\BibleRef{若 15:22}，我们在祂的话中能找到的唯一意义是：圣言亲自说，那些祂（理性）尚未完全临到的人，没有罪责；而那些分享了祂却行为违背祂所宣告的祂在我们内圆满临在的理念的人，若有罪，就是有罪的。只有如此解读，这句话才是真实的：\Quote{我若没有来向他们说过话，他们就没有罪。}若将此话应用在可见的基督身上（许多人认为应当如此），那么，祂没有到过的人就没有罪，这怎么可能是真的呢？那样的话，所有在救主降世之前生活的人都会无罪，因为那时耶稣尚未在肉身中显现。而且，所有从未听闻祂宣讲的人也将无罪，若无罪，则显然不受审判。但人内的圣言——我们说过全人类都有份——有两种意义：第一，指理念的充实，这发生在每个越过少年期的人身上（神迹例外）；第二，指成全，只发生在成全者身上。因此，\Quote{我若没有来向他们说过话，他们就没有罪；但如今他们的罪无可推诿}这句话应按第一种意义理解；而\BibleQuote{凡在我以前来的，都是贼和强盗，羊没有听从他们。}\BibleRef{若 10:8}这句话应按第二种意义理解。因为在理性成全之前，人内没有什么是可称赞的；一切都是不完全和有缺陷的，并且绝不能使我们内那些被寓意为羊的非理性元素服从。也许第一种意义应在\Quote{圣言成了血肉}这句话中认出来，第二种则在\Quote{圣言就是天主}中。因此，我们必须观察人间事务中在\Quote{圣言（理性）成了血肉}和\Quote{圣言就是天主}之间有什么可看到的。当圣言成了血肉时，我们能否说它在某种程度上被破碎稀释了，并且能否说它从那时起恢复，直到它再次成为起初的样子：天主圣言，与圣父同在的圣言，若望所见荣耀的圣言，确是从圣父而来的独生者。但圣子也可以被称为圣言，因为祂宣报祂父的奥秘，祂的父是理智，正如那被称为圣言的圣子一样。因为就如我们的话是心智所感知之事的传报者，同样，天主的圣言认识父，因为没有向导就没有受造物能接近父，祂启示祂所认识的父。因为除了圣子，没有人认识父\BibleRef{玛 11:27}，以及圣子愿意启示的人。而且，祂作为圣言，就是\BibleQuote{伟大谋略的使者}\BibleRef{依 9:5}，政权担在祂肩上；因为祂借着忍受十字架进入了祂的国度。此外，在\BibleBook{}{默示录}中，那\Quote{忠信真实者}（圣言）被描述为骑着一匹白马，我认为这些称号表明真理的圣言在我们中居住时对我们说话的声音清晰。此处几乎不是展示\Quote{马}一词如何在神圣学习中为鼓励我们的章节中经常使用的地方。我只引用其中两处：\Quote{马是为得救而欺诈的}和\Quote{有人依靠战车，有人依靠马，但我们却要因上主我们天主的名字而欢欣。}我们也不能忽略\WorkTitle{第四十四篇圣咏}中的一段，许多作者常引用它，仿佛他们理解它：\BibleQuote{我的心涌出美好的言语，我向君王陈述我的作品，我的笔是敏捷文士的舌头。你在世人中最美。}假设这是天主圣父所说，那么祂的心是什么，美好的言语要依照祂的心出现？如果如这些作者所假设，圣言不需要解释，那么心也应按字面理解。但认为天主的心如我们身体的器官一样是祂的一部分，这是相当荒谬的。我们必须提醒这些作者，正如我们论及天主的\Quote{手}、\Quote{臂}、\Quote{指}时，我们不按字面阅读，而是探究如何以健全的意义理解它们，使之堪配天主，同样，祂的\Quote{心}应理解为祂的理性能力，祂借此安排万有，而祂的\Quote{话}则宣告祂的内心所存。但谁向祂的受造物中那些堪当且超越自身者宣告父的谋划呢？除了救主，还有谁？\Quote{涌出}一词或许不是没有意义的；可能用了一百个其他词；可以说\Quote{我的心产生了一句美好的言语}或\Quote{我的心说了一句美好的言语}。但在涌出时，隐藏的风进入世界，同样，父可能不是持续地启示真理的观点，而是仿佛以涌出的方式，而圣言具有这样产生之物的特性，因此被称为不可见天主的肖像。因此，我们可以同意对这些话的通常理解，即认为它们是圣父所言。然而，并非当然就是天主自己宣告这些事。为什么不能是一位先知呢？充满圣神而无法自持，他关于基督的预言说：\BibleQuote{我的心涌出美好的言语，我向君王陈述我的作品，我的笔是敏捷文士的舌头。你在世人中最美。}然后对基督自己说：\BibleQuote{恩宠倾注在你嘴上。}如果说话者是圣父，在说了\BibleQuote{恩宠倾注在你嘴上}之后，祂如何能接着说\Quote{因此天主永远祝福了你}，以及后文\Quote{因此天主，你的天主，以喜乐之油膏你，胜过你的同伴}？那些希望圣父是说话者的人也许会引用这些话：\BibleQuote{女儿，请听，请看，请侧耳，忘记你的百姓和你的父家。}可以说，先知不能用\Quote{女儿，请听}这些话对教会说话。然而，不难表明，圣咏中经常改变说话者，因此这段话中的\Quote{女儿，请听}可能出自圣父，尽管整篇圣咏不是。关于圣言的讨论，我们在此可加上一段：\BibleQuote{诸天因上主的话而建立，万象因祂口中的气而成全。}有些人将此归于救主和圣神。然而，这段话不一定意味着更多，只不过诸天是由天主的理性（logos）建立的，如同我们说房屋由建筑师的计划（logos）建造，船由造船者的计划（logos）建造一样。同样，诸天是由天主的圣言建立（坚固）的，因为它们是更神圣的物质，因此被称为坚固的；它们大部分时间不流动，也不像世界的其他部分（尤其是下部）那样容易融化。由于这种区别，诸天被特别说成是由天主的圣言构成的。\par

那么，这句话首先说：\BibleQuote{在起初已有圣言}；我们应充分注视这一点；但我们在\BibleBook{}{箴言}中引用的证词引导我们首先放置智慧，并认为智慧先于宣告她的圣言。那么，我们必须注意到，圣言在起初，即在智慧中，始终存在。它在被称为起初的智慧中，并不妨碍它与天主同在并是天主，而且它不单与天主同在，而是与天主同在起初，在智慧中。因为接着他说：\BibleQuote{他在起初就与天主同在}。他本可以说\Quote{他与天主同在}，但既然他在起初，所以他就在起初与天主同在，并且\BibleQuote{万物是藉着他而造成的}，他在起初，因为天主在智慧中创造了万物，正如达味告诉我们的。并且为了让我们明白圣言有自己的确定位置和领域，作为一位在自己内具有生命（且是一个独立的位格）者，我们也必须谈论诸能力，而非能力。\Quote{万军的上主这样说（钦定本作hosts）}我们经常读到；有一些受造物，理性的和神圣的，被称为能力：其中基督是最崇高和最优秀的，他不仅被称为天主的智慧，也被称为他的能力。因此，正如天主的诸多能力各有自己的形式，救主与它们不同，同样，基督，即使我们内的圣言在形式方面并不在我们之外，也会从我们至今的讨论中被理解为圣言，他在起初、在智慧中拥有其存在。至今，对于\BibleQuote{在起初已有圣言}这句话，这些可能就足够了。\par

\SourceRecord{Translated by Allan Menzies.；Ante-Nicene Fathers；Vol. 9.；Edited by Allan Menzies.；Buffalo, NY: Christian Literature Publishing Co.,；1896.；<http://www.newadvent.org/fathers/101501.htm>.}

% structure: p0001 source=h1 target=section reason=page-title

\section{若望福音释义（卷二）}

% structure: p0002 source=h2 target=subsection reason=relative-heading-rank; base=h2

\subsection{一}

\BibleQuote{\Emph{圣言与天主同在，圣言就是天主。}}在前一节中，我可敬的兄弟\Person{安博}（按福音塑造的兄弟），我们已在目前能力所及的范围内，讨论了什么是福音，以及起初\OriginalTerm{圣言}{Logos}所在的\Quote{起初}是什么，以及那在起初的圣言是什么。现在我们来思考面前著作中的下一个点：圣言如何与天主同在。为此，记住那被称为圣言的曾降临到某些人身上，是有益的；正如\BibleQuote{上主的话\BibleRef{欧 1:1}传给贝厄黎的儿子欧瑟亚}，以及\BibleQuote{\BibleRef{依 2:1}传给阿摩兹的儿子依撒意亚的话，关于犹大和关于耶路撒冷}，以及\BibleQuote{\BibleRef{耶 14:1}关于旱灾传给耶肋米亚的话}。我们必须探究这话如何降临到欧瑟亚，以及它又如何降临到阿摩兹的儿子依撒意亚，以及再次关于旱灾降临到耶肋米亚；这样的比较也许能帮助我们找出圣言如何与天主同在。一般人只会简单地将先知们所说的话视为降临到他们身上的上主的话或圣言。然而，这难道不可以是：正如我们说这个人到那个人那里去，同样地，我们现在神学讨论的那位圣子圣言，由父派遣，降临到欧瑟亚——按历史说，即降临到贝厄黎的儿子、先知欧瑟亚，但按奥秘说，降临到那被拯救的人，因为欧瑟亚，按照语源，意为\Emph{被拯救者}；以及降临到贝厄黎的儿子，贝厄黎语源意为井泉，因为每个被拯救的人都成为那从深处涌出的泉水——天主的智慧——的儿子。圣人成为井泉之子，这丝毫不奇怪。从他的英勇事迹中，他常被称为儿子：或因他的行为在人前照耀而被称为光明之子，或因他拥有超越一切理解的天主的平安而被称为平安之子，或再一次因智慧带给他的帮助而被称为智慧之子；因为经上说：\BibleQuote{智慧\BibleRef{玛 11:19}因她的儿女而显为义}。因此，那藉着天主的圣神洞察万事，甚至洞察天主的深奥之事，以致能呼喊：\BibleQuote{啊，天主的智慧和知识的财富何其深奥！\BibleRef{罗 11:33}}的人，他就能成为井泉之子，上主的话就降临到他身上。同样地，圣言也降临到依撒意亚，教导在末日将降临于犹大和耶路撒冷的事；它也降临到因神圣昂扬而被举起的耶肋米亚。因为\OriginalTerm{雅威}{Iao}按语源意为举起、昂扬。现在，圣言降临到那些原先不能接纳那作为圣言的天主子来临的人；但祂并没有降临到天主那里，好像祂之前不与其同在一样。圣言永远与圣父同在；因此经上说：\BibleQuote{圣言与天主同在。}祂并没有成为与天主同在，而是\Quote{同在}这个词用于圣言，因为祂在起初就与天主同时同在，既不与起初分离，也不失去圣父。同样，祂并非在起初之后才进入起初，也非在与天主分离后才与天主同在。因为在万世和极遥远的年代之前，圣言就在起初，圣言与天主同在。因此，为了找出\Quote{圣言与天主同在}这句话的含义，我们引用了关于先知们的话——祂如何降临到欧瑟亚、依撒意亚和耶肋米亚——并且我们注意到\Quote{成了}与\Quote{同在}之间的区别，这绝非偶然。我们还得补充说：在祂降临到先知们时，祂以知识的真光光照先知们，使他们看见以前一直存在、但他们直到那时还未理解的事物。然而，与天主同在时，祂就是天主，仅仅因为祂与祂同在。或许正是因为\Person{若望}在\OriginalTerm{圣言}{Logos}中看到某种这样的次序，他才没有将\BibleQuote{圣言就是天主}这句放在\BibleQuote{圣言与天主同在}之前。他安排不同句子的顺序并不妨碍每个命题的力量被单独而完全地看到。一个命题是：\BibleQuote{在起初已有圣言}，第二个：\BibleQuote{圣言与天主同在}，然后：\BibleQuote{圣言就是天主。}句子的安排或可被认为指示了一种次序：我们有首先\BibleQuote{在起初已有圣言}，其次\BibleQuote{圣言与天主同在}，第三\BibleQuote{圣言就是天主}，如此便可看出，圣言与天主同在，使得祂成为天主。\par

% structure: p0004 source=h2 target=subsection reason=relative-heading-rank; base=h2

\subsection{二、圣言如何是天主。此问题上应避免的错误。}

接下来我们注意到\Person{若望}在这些句子中使用冠词的方式。他在这方面的写作并非不经心，也非不熟悉希腊语的细微之处。在某些情况下他使用冠词，在某些情况下则省略。他在\OriginalTerm{圣言}{Logos}上加了冠词，但在天主的名称上只偶尔加冠词。当天主的名称指万有的非受造原因时，他使用冠词；而当圣言被称为天主时，他省略冠词。我们在带冠词的天主和无冠词的天主之间观察到的区别，是否也适用于带冠词的圣言和无冠词的圣言？我们必须探究这一点。正如万有之上的那位天主是带冠词的天主，而非无冠词的，同样，\Quote{那圣言}是那居住在每一理性受造物内的理性（圣言）的源头；但存在于每个受造物内的理性，并不像前者那样被称为\Emph{最卓越地}的圣言。如今有许多人真诚关心宗教，在此陷入极大的困惑。他们害怕自己可能宣告两位天主，而恐惧驱使他们陷入错误且邪恶的道理。他们要么否认圣子除了圣父的本性之外还有自己独特的本性，使他们所称为圣子的那一位几乎只是名义上的天主；要么否认圣子的天主性，赋予祂自己独立的存在，使祂的本质领域落在圣父之外，从而彼此可分离。对这样的人，我们必须说：一方面，天主是\OriginalTerm{自己天主}{Autotheos}（自身的天主）；因此救主在向圣父的祈祷中说：\BibleRef{若 17:3}\BibleQuote{使他们认识你，唯一的真天主。}但所有超出那自己天主的，都是因分享祂的天主性而成为天主的，并且不应简单地称为天主（带冠词），而应称为天主（无冠词）。于是，一切受造的首生者，祂是首先与天主同在，并吸引天主性到自身的那一位，是比祂旁边其他诸神更高贵的存在——天主是那些神的天主，正如经上所记：\BibleRef{咏 50:1}\BibleQuote{万神的天主、上主发言，呼唤大地。}正是通过首生者的职分，它们才成了神，因为祂从天主那里慷慨地汲取了使它们成为神的恩赐，并按照祂自己的美善将之传递给他们。因此，真天主就是\Quote{那位天主}，而那些照着祂被塑造的，便是神，如同祂这原型的肖像。然而，这一切肖像的原型肖像又是天主的圣言，祂在起初就与天主同在，并因与天主同在而始终是天主，不是凭自身拥有此性，而是因与圣父同在；而且，若我们思量，除非藉着始终不间断地默观圣父的深渊，祂便不再继续是天主。\par

% structure: p0006 source=h2 target=subsection reason=relative-heading-rank; base=h2

\subsection{三、圣言对人的各种关系。}

现在，可能有人不赞同我们所说的，即圣父是唯一真天主，但在真天主之外也承认其他存有因分有天主而成为神。他们害怕那超越一切受造者的荣耀会被降低到那些被称为神的其他存有层面。我们在祂们之间作了这一区分，即表明天主圣言对于其他所有神祇而言，是祂们神性的分施者。为了消除反对意见，我们尚需补充一点：每个理性受造物内的理性，与那太初与天主同在、且是天主圣言的理性之间的关系，正如天主圣言与天主的关系一样。正如圣父——祂是至真天主和真天主——与祂的肖像以及祂肖像的肖像（人被认为是按照肖像造的，而非天主的肖像）的关系一样，同样，圣言也与每个人内的理性（话语）有这样的关系。每一个都充满泉源的地位：圣父是神性的泉源，圣子是理性的泉源。因此，如同有许多神，但对我们来说只有一个天主圣父，有许多主，但对我们来说只有一位主耶稣基督，同样也有许多\OriginalTerm{逻各斯}{\GreekText{Λόγοι}}，但就我们而言，我们祈求那位太初与天主同在，并与天主同在的\OriginalTerm{天主圣言}{\GreekText{Λόγος}}能与我们同在。因为凡不接纳这太初与天主同在的圣言，或不依附于祂在肉身显现时的样子，或不参与那些分有这圣言者，或曾分有祂却又远离祂的人，他将有份于那被称为与理性最敌对之物。我们从所开始的真理中推演出的内容现在已足够清晰。首先，我们谈论了天主和天主的圣言，以及诸神——要么是分享神性的存有，要么是被称为神而非真神的存有。又谈论了天主的圣言和成了血肉的天主圣言，以及诸逻各斯（\GreekText{λόγοι}），即某种方式分享圣言的存有，以及次等的逻各斯或第三等的逻各斯，它们被认为是在那超越万有的逻各斯之外的逻各斯，但实际并非如此。这些可以被称为非理性的理性；有些存有被称为神却不是神，我们也可以将这些不是神的神与不是理性的理性并列。宇宙的天主是选民的天主，更是选民救主的天主；然后祂是那些真正为神的存有的天主，最后，简言之，祂是活人的天主，而非死人的天主。但天主圣言也许是那些将一切都归功于祂、并视祂为父者的天主。日月星辰，根据古人的记述，与那些不配以众神之神为他们的天主的存有关联。他们之所以持此意见，是由申命记中的一段经文引导，大意如下：\BibleQuote{当你举目望天，看见太阳、月亮和天上的万象时，不要误入歧途去朝拜它们，事奉它们，因为这是主你的天主分给天下万民享用。但主你的天主没有这样赐给你们。}然而，天主将太阳、月亮和天上的万象分给万民，难道祂不是同样赐给了以色列人吗？为使那些不能上升到理性境界的人，能通过可感觉的神祇倾向于思考天主性，并自愿与它们关联，从而被保守而不至于堕入偶像和魔鬼？难道不是有些人以宇宙的天主为他们的天主，而第二类人紧随其后依附于天主的圣子、祂的基督，第三类人则敬拜太阳、月亮和天上的万象——他们固然离弃了天主，但这离弃远优于那些以人手所造之物——金银、人的工艺作品——为神的人的离弃。最后是那些委身于被称为神却不是神的存有的人。同样地，现在有些人相信那太初与天主同在、并与天主同在的天主圣言，如\Person{欧瑟亚}、\Person{依撒意亚}、\Person{耶肋米亚}以及其他宣称主的圣言来到他们那里的人。第二类人只知道耶稣基督，这位被钉在十字架上的，认为成了血肉的圣言就是全部的圣言，并且只按肉体认识基督。这就是被算为信徒的大多数。第三类人投身于某些逻各斯（言论），这些逻各斯分享那被视为超越一切其他理性的逻各斯；这些人是追随\Place{希腊}人中受尊敬且卓越的哲学学派的人。第四类人除此之外，他们信赖腐败而无神的言论，否认如此明显且几乎可见的天主眷顾，并承认人当追求除善之外的另一个目的。有人可能认为我们偏离了主题，但在我看来，我们所得出的关于天主之名的四种事物和关于圣言的四种事物的观点，在此处非常合宜。有带着冠词的天主和不带冠词的天主，然后有分为两个等级的诸神，在更高等级之首的是天主圣言，祂本身又被宇宙的天主所超越。同样地，有带着冠词的圣言和不带冠词的圣言，对应于绝对的天主和一位神；还有两个等级的逻各斯。有些人连结于圣父，是祂的一部分，紧随其后的是那些现在我们的论证更清楚地揭示的人，即那些来到救主面前并完全立足于祂内的人。第三是我们前面提到的人，他们视日月星辰为神，并立足于它们。第四也是最后是那些屈服于无生命和死寂偶像的人。在关于圣言方面，我们也发现类似的情形。有些人以圣言本身为装饰；有些人以仅次于祂、看起来像是原始圣言本身的东西为装饰，即那些只知道耶稣基督并祂钉在十字架上、将圣言视为血肉的人。第三类正如我们稍早所述。何须我提那些自认为在圣言内，却已堕落，不仅远离了善本身，而且远离了善的痕迹和那些有分于善的人呢？\par

% structure: p0008 source=h2 target=subsection reason=relative-heading-rank; base=h2

\subsection{四、圣言是唯一的，而非众多。论忠信与真实的圣言，以及祂的白马。}

\Quote{\Emph{祂在起初就与天主同在。}} 圣史借前面的三个命题使我们认识了三个秩序，现在他将这三个归结为一个，说道：\Quote{这位（圣言）在起初就与天主同在。} 在第一个前提中，我们得知圣言在哪里：祂在起初；然后我们得知祂与谁同在，与天主同在；再然后得知祂是谁，祂是天主。现在，祂藉着\InnerQuote{祂}这个字指出那身为天主的圣言，并将前面的三个命题归结为第四个命题：\Quote{在起初已有圣言}、\Quote{圣言与天主同在} 和 \Quote{圣言就是天主}。于是圣史说：\Quote{这位（圣言）在起初就与天主同在。} \InnerQuote{起初}一词可以理解为世界的起初，这样我们可以从所述之言得知圣言比从起初受造的事物更古老。因为如果\Quote{在起初天主创造了天地}，而\Quote{祂}已在起初，那么圣言显然比那些在起初受造的事物更古老，不仅比穹苍和旱地古老，甚至比诸天和大地更古老。现在有人可能合情合理地问道：为何不说\InnerQuote{在起初有天主的圣言，天主的圣言与天主同在，天主的圣言就是天主}？但提出这种问题的人可以证明是预先假定了有众多逻各斯（\OriginalTerm{逻各斯}{Logoi}），彼此种类不同，一个是天主的圣言，另一个可能是天使的言语，第三个是人的言语，其他逻各斯亦如此。如果圣言如此，那么智慧和正义也是如此。但若有众多事物都同样可被称为\InnerQuote{圣言}就荒谬了；智慧和正义亦然。若我们审视真理的情况，便不得不承认我们不应寻求众多逻各斯、众多智慧或众多正义。任何人都会承认真理只有一个；绝不能在此说有天主的真理、天使的真理和人的真理——事物的本性决定任何事物的真理都是唯一的。既然真理是唯一的，那么显然，真理的预备和展示，即智慧，在理性上也应被视为唯一，因为被视为智慧的，若缺少了唯一的真理，就不能正当地称为智慧。但若真理是唯一的，智慧也是唯一的，那么宣布真理并使真理单纯清晰地呈现给适宜接受者的理性（\OriginalTerm{逻各斯}{Logos}）也将是唯一的。我们这样说，绝非否认真理、智慧和理性来自天主，而是想指出在本段中略去\InnerQuote{天主}二字以及采用\InnerQuote{在起初逻各斯与天主同在}这一表述的用意。同一位\Person{若望}在\BibleBook{若望}{默示录}中给祂加上\InnerQuote{天主的}这个称号，说道：\BibleRef{默 19:11} \Quote{我看见天开了，看，有一匹白马，骑马的那位称为\InnerQuote{忠信}和\InnerQuote{真实}；他凭正义审判和作战。他的眼睛如同火焰，头上戴着许多冠冕，写着一个名字，除他自己外无人认识。他身披染过血的衣服，他的名字称为\InnerQuote{天主的圣言}。天上的军队也骑马跟着他，穿着洁白的细麻衣。从他口中射出一把利剑，用来击打万民；他要以铁杖牧养他们，并要践踏全能者天主烈怒的酒榨。在他衣服和大腿上写着他的名字：\InnerQuote{万王之王，万主之主}。} 在此段经文中，\InnerQuote{圣言}是绝对地（不带冠词）被提及，同时也加上\InnerQuote{天主的圣言}；若前一种情况不成立（即若使用了冠词），我们可能会错误地理解其意义，从而偏离关于圣言的真理。因为若它仅被称为\InnerQuote{圣言}，而没有说\InnerQuote{天主的圣言}，我们就不能清楚地知道圣言是天主的圣言。再者，若它被称为\InnerQuote{天主的圣言}，却没有绝对地被称为\InnerQuote{圣言}，那么我们就可能根据每个存在理性受造物的构造想象众多逻各斯；那么我们就可能假设有众多真正意义上的逻各斯。在\BibleBook{若望}{默示录}关于天主的圣言的描述中，宗徒兼圣史（\BibleBook{若望}{默示录}也使他配称为先知）说他看见天开了，天主的圣言骑着白马。现在我们必须思考，他所说的天开了、白马、天主的圣言骑着白马，以及\InnerQuote{天主的圣言是忠信和真实}、\InnerQuote{他凭正义审判和作战}是什么意思。这一切将大大推进我们对天主的圣言这一主题的研究。我认为天对不虔敬的人和那些带有属土肖像的人是关闭的，而对义人和那些带有属天肖像的人是敞开的。因为对于前者，他们处于低下且仍居住在肉身中，美好的事物是关闭的，因为他们无法理解，也没有能力或意愿看见其美丽，他们向下看而不努力向上看。但对于卓越的人，或那些在天上拥有国度的人，\BibleRef{斐 3:20} 他以\Person{达味}的钥匙开启天上的事物，向他们显示，并藉着骑马使他们明白一切。这些话也有其含义；马是白色的，因为更高知识的本性是清晰、洁白且充满光明的，这更高知识（\OriginalTerm{更高知识}{\GreekText{γνῶσις}}）如此。而骑在白马上的那位被称为\InnerQuote{忠信}者，他更稳固地、也可说更庄严地骑在那些不可废弃的话语上，这些话语比任何马匹更尖锐、更迅疾地奔跑，并在其飞驰过程中胜过每一个模拟圣言的所谓言语和每一个模拟真理的所谓真理。骑在白马上的那位被称为\InnerQuote{忠信者}，不是由于祂所珍视的信德，而是由于祂所激发的信德，因为祂是值得信赖的。按照\Person{梅瑟}的话，\BibleRef{申 32:4} 天主上主是忠信和真实的。祂相对于影、预像和形象也是真实的；因为在敞开的天空中的圣言就是这样，祂在地上不同于在天上；在地上祂成了血肉，并通过影、预像和形象发言。因此，那些被认为是信者的大多数，是圣言之影的门徒，而不是在敞开的天上的真天主圣言的门徒。为此\Person{耶肋米亚}说：\BibleRef{哀 4:20} \Quote{我们鼻口中的气息是基督上主，我们曾说过：在他的荫庇下我们将在万民中生存。} 因此，被称为\InnerQuote{忠信}的天主圣言也被称为\InnerQuote{真实}，祂凭正义审判和作战；因为祂从天主那里领受了以绝对的正义和绝对的判断进行审判，并将应得的份归于每个受造物的能力。因为那些分有部分正义和判断能力的人，没有一个人能使自己的灵魂接受到如此完美的正义和判断的拷贝和印像，以至于在绝对正义和绝对公平上毫无欠缺，就像画家无法将原画的所有特质都传达给复制品一样。我认为这就是\Person{达味}所说\Quote{在祢面前，凡有生命的都不能成义}的原因。他不是说没有一个人，或没有天使，而是说没有有生命的，因为即使任何有生命者分有生命并完全摆脱了死亡，在祢面前（祢就是生命本身）也不能成义。分有生命因而被称为有生命的，也不能成为生命本身；分有正义因而被称为正义的，也不能等同于正义本身。天主的圣言的职责不仅是以正义审判，而且是以正义作战，祂藉着理性和正义向祂的敌人作战，以致不理性与邪恶被摧毁，祂就可以居住在为了自己的得救而成为基督俘虏的人的灵魂内，并成义那灵魂，从她身上驱逐一切敌人。如果我们记得祂是真理的使者，而另有自称为圣言却不是的，以及自称为真理却不是的，而是谎言，那么我们就更能理解圣言所进行的这场战争。然后圣言武装自己对抗谎言，以口中的气息斩杀它，并藉祂来临的显现使之归于无有。\BibleRef{得后 2:8} 我们要思考，宗徒给得撒洛尼人写的这些话是否可以按灵性意义理解。因为那被基督口中的气息所毁灭的，基督就是圣言、真理和智慧，除了谎言还能是什么？那被基督来临的显现所消灭的，当基督被理解为智慧和理性时，除了那自称智慧实则不然的东西还能是什么？那东西实际上是天主像宗徒所描述的那样对待的事物之一：\BibleRef{格前 3:19} \Quote{他以自己的诡计捕捉智者，那些不是以真智慧为智慧的人。} 关于骑白马者，\Person{若望}还加上这句奇妙的话：\Quote{他的眼睛如同火焰。} 因为火焰既明亮又照耀，但同时炽热并销毁物质事物，同样，如我可以说，圣言的眼睛——祂用这眼睛观看，以及凡有分于祂的人——不仅具有把握心智事物的内在品质，也具有消耗并除去那些更物质和粗俗观念的能力，因为任何虚假的东西都逃避真理的直接和轻盈。按非常自然的顺序，在谈到那凭正义审判并依照祂公义的判断作战，以及祂作战、赐予光明之后，作者接着说：\Quote{头上戴着许多冠冕。} 因为如果谎言是单一的、只有一种形式，那真实而忠信的圣言与之争战并因战胜而加冕，那么自然只会赐给祂一顶冠冕以表彰胜利。但实际上，因为自称真理的谎言很多，圣言与它们争战并因此加冕，所以冠冕也很多，环绕这战胜一切者的头部。祂已战胜每一种悖逆的势力，许多冠冕标志着祂的胜利。冠冕之后，祂被说成有一个名字，除了祂自己，没有人知道。因为有些事情只有圣言知道；在祂之后存在的受造物本性比祂低劣，没有一个能看见祂所领悟的一切。也许只有那些有分于圣言的人，才知晓那些对没有分于祂的人隐藏的事情。在\Person{若望}的神视中，天主的圣言骑着白马时并非赤身，祂穿着溅血的衣服，因为成了血肉并因此死亡的圣言，身上带有祂的血倾流在地上的印记——士兵刺穿了祂的肋旁。因为即使我们有一天有幸达到对圣言最高和最圆满的默观，我们也不会完全忘记那苦难，也不会忘记我们的进入是由祂在我们身体内的居留所促成的这一真理。天主的圣言由所有天上的军队跟随；他们跟随圣言作为统帅，并在一切事上效法祂，尤其是在他们也骑上白马这事上。对于明白的人，这奥秘是敞开的。正如在万物的终结时忧伤和疼痛和哀号都消失了，同样，我想，晦暗和疑惑也消失了，天主智慧的一切奥秘都得到精确而清晰的揭示。再看圣言跟随者的白马和他们所穿的洁白纯净的细麻衣。既然麻衣出自土地，那么这些细麻衣不正是代表那些声音所穿的语言（这些声音清楚宣布事物）吗？我们已经比较详细地讨论了\BibleBook{若望}{默示录}中关于天主的圣言的陈述；清楚地认识祂对我们十分重要。\par

% structure: p0010 source=h2 target=subsection reason=relative-heading-rank; base=h2

\subsection{五、祂（这一位）在起初就与天主同在。}

对于不能区分上下文的不同命题的人来说，福音作者似乎是在重复自己。\BibleQuote{祂在起初就与天主同在}似乎并没有给\BibleQuote{圣言与天主同在}增加什么。我们必须更加仔细地观察。在\BibleQuote{圣言与天主同在}这一陈述中，我们并未得知任何关于何时或何地的信息；这一点在第四公理中得到了补充。有四个公理，或者如某些人所称的命题，第四公理是\BibleQuote{祂在起初就与天主同在。}现在，\BibleQuote{圣言与天主同在}与\BibleQuote{祂在起初就与天主同在}并不是一回事；因为在这里，我们不仅被告知祂与天主同在，还得知了祂何时何地如此：\BibleQuote{祂在起初就与天主同在。}\Quote{祂}这个词，用作指示代词，将被认为是指圣言，或者被不太仔细的探究者认为是指天主。以前所注意到的，现在被概括在这个名称\Quote{祂}中，即圣言和天主的概念；随着论证的进行，不同的概念被聚集在一个中；因为天主的概念并不包含在圣言的概念中，圣言的概念也不包含在天主的概念中。也许我们面前的这个命题是对前面三个命题的一个总结。接受\BibleQuote{圣言在起初}这一陈述，我们尚未得知祂与天主同在；接受\BibleQuote{圣言与天主同在}这一陈述，我们尚不清楚祂在起初就与天主同在；接受\BibleQuote{圣言就是天主}这一陈述，既没有表明祂在起初，也没有表明祂与天主同在。\par

当福音作者说：\BibleQuote{祂在起初就与天主同在}，如果我们把代词\Quote{祂}应用于圣言和天主（因为祂是天主），并认为\Quote{在起初}是与它结合，\Quote{与天主}是附加于它，那么三个命题中没有任何东西不被总结和汇集于这一命题中。而且，由于\Quote{在起初}已经被说了两次，我们可以考虑是否可以从中学到两个教训。首先，圣言在起初，好像祂独自存在而不与任何人同在；其次，祂在起初与天主同在。我认为，说祂既在起初，又在起初与天主同在，并没有任何不真实之处，因为祂既不是单独与天主同在（因为祂也在起初），也不是单独在起初而不与天主同在，因为\BibleQuote{祂在起初就与天主同在}。\par

% structure: p0013 source=h2 target=subsection reason=relative-heading-rank; base=h2

\subsection{六、圣言如何是万物的创造者，甚至连圣神也是藉着祂受造的。}

\Quote{\Emph{万有是藉着祂造成的。}}\Quote{藉着祂}从不放在首位，而总是放在第二位，如在\WorkTitle{罗马书}中：\BibleRef{罗 1:1}\Quote{保罗，耶稣基督的仆人，蒙召的宗徒，被选拔为传天主的福音——这福音是天主先前藉着祂的先知在圣经上所预许的，论及自己的圣子，按肉身是生于达味的后裔，按圣德的神，藉着从死者中复活，我们的主耶稣基督，被立为具有大能的天主之子，藉着祂我们领受了恩宠和宗徒职务，为使万民服从信德，以光荣祂的名。}因为天主先前藉着先知预许了自己的福音，先知是祂的仆役，他们所说的关于\Quote{藉着祂}的话。而天主又将恩宠和宗徒职务赐给了保罗和其他人，为使万民服从信德，这是藉着救主耶稣基督赐下的，因为\Quote{藉着祂}是属于祂的。宗徒保禄在\WorkTitle{希伯来书}中说：\Quote{在末后的日子，祂藉着圣子对我们说话，立祂为万有的继承者，也\InnerQuote{藉着祂}创造了诸世代，}向我们显示天主藉着祂的圣子创造了诸世代，当诸世代被创造时，\Quote{藉着祂}属于独生子。因此，若万有都是被造成的，如同此处经文所说，是藉着圣言\OriginalTerm{圣言}{Logos}造成的，那么它们就不是由圣言造成的，而是由一位比祂更强大、更伟大的。这位还能是谁呢？岂不就是父吗？既然如我们所见，万有都是藉着祂造成的，我们就必须查考圣神是否也是藉着祂造成的。在我看来，那些认为圣神是被造的，同时也承认\Quote{万有是藉着祂造成的}的人，必然要假设圣神是藉着圣言造成的，因此圣言比祂更古老。而若有人不愿承认圣神是藉着基督造成的，若他承认这部福音的叙述是真实的，就必须假设圣神是非受造的。除了这两种意见（即承认圣神是由圣言造成的，或认为圣神是非受造的）之外，还有第三种，即宣称圣神在父和子以外没有自己的本质。但若进一步思考，或许有理由认为圣子与父有区别，祂与父是同一的，而圣神与圣子之间显然有区别，如同在\BibleRef{玛 12:32}中：\Quote{凡说话干犯人子的，可得赦免；但谁说亵渎圣神的话，必不得赦免，无论是在今世，还是在来世。}因此，我们认为有三个位格\OriginalTerm{位格}{hypostases}：父、圣子和圣神；同时，我们相信只有父是非受造的。因此，作为更虔诚、更真实的做法，我们承认万有都是由圣言造成的，而圣神是在父藉着基督所造的一切中，最卓越、最先被造的。这或许就是圣神不被称作天主子嗣的原因。只有独生子\OriginalTerm{独生子}{Only-begotten}是本性和从起初就是子，而圣神似乎需要圣子，以赐予祂本质，使祂不仅存在，而且成为智慧、理性的、正义的，以及一切我们认为祂所是的那样。这一切祂都是藉着参与基督的特性而获得的，这我们在上文已经谈论过。我认为，圣神供给那些藉着祂并参与祂而被称为圣人的人，从天主来的神恩\OriginalTerm{神恩}{charisms}的质料；因此，这神恩的质料由天主使之有力，由基督分施，其在人内的实际存在则归功于圣神。我之所以对神恩持此看法，是出于保禄某处所写的：\BibleRef{格前 12:4}\Quote{神恩虽有区别，却是同一的圣神；职分虽有区别，却是同一的主；功效虽有区别，却是同一的天主，在一切人身上行一切事。}说万有是藉着祂造成的，其推论似乎是圣神必然是由圣言被造的，这当然会引起一些困难。有些经文将圣神置于基督之上；例如在\WorkTitle{依撒意亚}中，基督宣称祂不仅被父派遣，也被圣神派遣。祂说：\BibleRef{依 48:16}\Quote{现今，主派遣了我和祂的圣神。}在福音中，祂声明得罪祂自己的罪可得赦免，但亵渎圣神的罪，无论是在今世还是在来世都不得赦免。这是什么原因呢？难道是因为圣神比基督更有价值，所以得罪祂的罪不得赦免吗？更合理的解释是：所有理性受造物都有份于基督，当他们悔改时罪就得赦免；而只有那些堪当的人才会有份于圣神，那些在如此伟大而有力的协助下仍堕入邪恶、击败内住在他们内的圣神之谋划的人，很难得到赦免。当我们发现主在\WorkTitle{依撒意亚}中说到祂被父和祂的圣神派遣时，我们也必须指出，此处圣神并非原本高于救主，而是救主为了执行已定的计划——天主子成为人——而取了较低的地位。若有人因我们说救主成为人时被置于圣神之下而跌倒，我们请他思考\WorkTitle{希伯来书}中的话，保禄在那里显示耶稣因受死之苦而比天使低微一点。他说：\Quote{我们看见那位暂时比天使低微一点的耶稣，因所受的死亡之苦，被尊崇以光荣和荣耀。}此外，无疑必须补充：受造界为了脱离败坏奴役，而人类尤其如此，需要将一种幸福而神圣的权能引入人性中，以纠正世上错谬之事；这一行动，可以说是落在圣神的份上。但圣神无法承担如此重任，便推出了救主，只有祂能承受这样的战斗。因此，作为首要的父派遣了圣子，但圣神也派遣了祂，并指示祂先行，应许到适当的时候将降临到天主子身上，与祂合作以拯救人类。这圣神确实做到了，当祂在圣洗后以鸽子的形体降临到祂身上，留在祂内，并未离去，如同祂对待那些不能持续承受祂荣耀的人那样。因此，若翰在解释他如何认识基督时，不仅提到圣神降临在耶稣身上，还提到祂留在祂身上。因为经上记载若翰说：\BibleRef{若 1:32}\Quote{那派遣我来以水施洗的对我说：你看见圣神降下并停在谁身上，谁就是那要以圣神和火施洗的。}并非只说\Quote{你看见圣神降下}，因为圣神无疑也降在别人身上，而是\Quote{降下并停在祂身上}。我们对这个问题的探讨已经稍长，因为我们急于澄清：若万有是藉着祂造成的，那么圣神也是藉着圣言造成的，并且被视为\Quote{万有}之一，低于其造物主。这个观点非常稳固，不会被一些相反的文字所动摇。若有人相信\WorkTitle{希伯来人福音}，其中救主亲自说：\Quote{我的母亲，圣神刚刚抓住我的头发，带我到大博尔山。}他就要面对如何解释圣神是基督的母亲，而圣神本身又是藉着圣言造成的困难。但这段经文和这个困难都不难解释。因为凡在天上承行父的旨意的\BibleRef{玛 12:50}，就是基督的兄弟、姊妹和母亲；如果基督的兄弟这个名号不仅可以应用于人类，也可以应用于比人类更高贵的存有，那么圣神作为祂的母亲就没有什么荒谬之处，因为凡承行天父旨意的人，都是祂的母亲。\par

关于\Quote{万物是借着他造成的}这句话，还有一个要点需要审视。作为概念，\Quote{圣言}来自于\Quote{生命}，然而我们读到：\Quote{那在圣言内造成的，是生命，这生命是人的光。}既然万物都是借着他造成的，那么，那人的光即生命，以及救主向我们展示的其他概念，是否也是借着他造成的呢？或者，我们必须将\Quote{万物是借着他造成的}理解为不包括那些在他内的事物？后者似乎是更可取的。因为，倘若我们承认那作为人的光的生命是借着他造成的，既然经文说生命\Quote{成了}人的光，那么对于那被认为是先于圣言的智慧，我们该说什么呢？因此，关乎圣言的事物（他的关系或状态）并非由圣言所造，结果是，除了基督所呈现的概念之外，万物都是借天主的圣言造成的，圣父在智慧中造成万物。经上说：\Quote{在智慧中祢造成了它们全体}，不是\Emph{借着}，而是\Emph{在}智慧中。\par

% structure: p0016 source=h2 target=subsection reason=relative-heading-rank; base=h2

\subsection{七、未经圣言所造之物}

然而，让我们看看为何要加上\Quote{没有他，什么也没有造成}这句话。有些人可能认为，在\Quote{万物是借着他造成的}之后再加上\Quote{没有他，什么也没有造成}是多余的。因为如果万物都是借着圣言造成的，那么没有他什么也没有造成。然而，从\Quote{没有圣言什么也没有造成}这个命题，并不能推出万物都是借着圣言造成的。可能虽然没有什么是不借着圣言造成的，但万物并非仅仅借着圣言造成，有些是凭着他（由他）造成的。因此，我们必须弄清楚\Quote{万物}是在什么意义上理解，\Quote{什么也没有}是在什么意义上理解。因为，如果不先对这些术语作明确的界定，就会有人主张，如果万物都是借着圣言造成的，而恶是万物的一部分，那么整个罪的事以及一切邪恶之事也都是借着圣言造成的。但我们必须认为这是错误的。认为受造物是借着圣言造成的，人的英勇行为也是借着他成就的，以及如今在福乐中的圣人们的一切有益行为都是借着他成就的，这并无荒谬之处；但人的罪和不幸则不然。有些人认为，既然恶并非基于事物的本质——因为它起初不存在，最终也将终止——那么，我们所说的那些恶就是\Quote{无}；正如有些希腊人说，类和形式（如一般的动物和人）属于\Quote{无}的范畴，同样，人们认为一切不属于天主的都是\Quote{无}，甚至连它似乎具有的存有也并非借圣言而获得的。我们问，是否有可能从圣经中以令人信服的方式证明这一点。至于\Quote{无}和\Quote{非有}这两个词的意义，它们似乎是同义词，因为\Quote{无}可以被称为\Quote{非有}，\Quote{非有}也可以被称为\Quote{无}。然而，宗徒似乎将那些不是的事物不算作完全不存在的事物，而是算作邪恶的事物。对他而言，\Quote{非有}就是恶；他说，\BibleRef{罗 4:17} \Quote{天主叫那不存在的如同存在的一样。}此外，在《七十贤士译本》的《艾斯德尔传》中，摩尔德开称呼以色列的敌人为\Quote{那些不存在者}，说道：\BibleRef{艾 4:22} \Quote{主，不要把祢的权杖交给那些不存在的人。}我们也可以注意，恶人因他们的邪恶而被说成是不存在的，这源于《出谷纪》中归于天主的名字：\BibleRef{出 3:14} \Quote{因为主对梅瑟说：我是自有者，这是我的名。}善的天主也这样对我们说，我们祈求能成为他集会的一部分。救主赞扬他说：\BibleRef{谷 10:18} \Quote{除了独一天主圣父，没有善的。}因此，善与\Quote{那存在者}是同一的。与善相对的是恶或邪恶，与\Quote{那存在者}相对的是\Quote{那不存在者}，由此推出，恶和邪恶就是\Quote{那不存在者}。这或许导致有些人断言魔鬼不是天主所造的。就他是魔鬼而言，他不是天主的工程，但那作为魔鬼的存在是个受造物，既然除了我们的天主之外没有别的造物主，他就是天主的工程。这好比我们说，杀人犯不是天主的工程，同时我们也可以说，就他是人而言，天主造了他。他作为人的存在是从天主领受的；我们并不断言他从天主领受了他作为杀人犯的存在。因此，所有分享于\Quote{那存在者}（圣人们分享于他）的人，可以恰当地被称为\Quote{存在者}；而那些放弃分享于\Quote{那存在者}，借此剥夺了自己存在的人，就成了\Quote{非存在者}。但我们在开始讨论时说过，\Quote{非存在}和\Quote{无}是同义词，因此那些不是存在者的人就是\Quote{无}，一切恶就是\Quote{无}，因为它是\Quote{非存在}，而既然它们被称为\Quote{非存在}，它们就是未经圣言而存在的，不被算在那些借着他造成的万物之中。这样，我们就已尽己所能，指明了哪些\Quote{万物}是借圣言造成的，以及哪些是未经他而存在的，后者从来不是\Quote{存在}，因此被称为\Quote{无}。\par

% structure: p0018 source=h2 target=subsection reason=relative-heading-rank; base=h2

\subsection{八、赫拉克利翁关于圣言并非创造媒介的观点}

我认为，\Person{赫拉克利翁}（据说他是\Person{瓦伦廷}的朋友）在讨论这句话：\Quote{万物是藉着祂造成的}时，采取了一种粗暴且不合理的做法。他排除了整个世界及其中的一切，按照他的假设，从\Quote{万物}中排除世界上及其内容中最好的东西。因为他说，\OriginalTerm{永世（时代）}{æon}及其中的事物并非由\OriginalTerm{圣言}{Logos}所造；他认为它们存在于圣言之先。他处理\Quote{没有祂，什么也没有造成}这句话时，有些大胆，也不怕\BibleRef{箴 30:6}的警告：\Quote{不可加添祂的话，免得祂责备你，你就显为说谎的}，因为他在\Quote{无}上添上了\Quote{世界上和受造物中的}。由于他对这段经文的解释明显十分牵强，且与证据相悖，因为他认为神圣的东西被排除在万有之外，而他视为纯粹邪恶的东西就是万有本身，所以我们无需浪费时间来反驳这种表面荒谬的观点，因为他在没有圣经依据的情况下，在\Quote{没有祂，什么也没有造成}这句话上添加了\Quote{地上和受造物中的}这些字。在这种毫无内在可能性支撑的提议中，他事实上是要求我们像信任先知或宗徒那样信任他；先知和宗徒有权柄，不需要对世人负责，他们传给了身边的人和后来的人那些属于人类救恩的著作。他对\Quote{万物是藉着祂造成的}这句话也有自己的私下解释，他说是圣言促使\OriginalTerm{德穆革}{demiurge}造世界，但圣言不是从祂或凭祂，而是藉着祂；他以不同于通常用法的角度理解经文的措辞。因为，如果事情的真相如他所想，那么作者应当说万物是藉着德穆革由圣言造成的，而不是藉着圣言由德穆革造成的。我们接受通常理解的\Quote{藉着祂}，并提供了支持我们解释的证据；而他不仅提出了自己的新读法——没有神圣圣经的支持——而且似乎甚至藐视真理，无耻地公开反对它。因为他说：\Quote{不是圣言造了万物，如同在另一位施行者的下面}，他以这种意义理解\Quote{藉着祂}，\Quote{而是另一位造了它们，圣言本身是施行者。}现在不是证明不是德穆革成为圣言的仆人而造了世界、而是圣言成为德穆革的仆人而形成了世界的合适时机。因为，按照先知\Person{达味}，\Quote{天主发言，万物就形成；祂一出命，万物就受造}。因为无始的天主命令了一切受造物的首生者（\BibleRef{哥 1:15}），万物就被造成，不仅是世界及其中的事物，而且还有所有其他事物，无论是宝座、统治区、掌权者还是权势，因为万物都是藉着祂并为了祂而造，祂在万有之先。\par

% structure: p0020 source=h2 target=subsection reason=relative-heading-rank; base=h2

\subsection{九、我们内住的圣言不为我们之罪负责}

关于这话，还有一点：\BibleQuote{没有他，什么也没有受造}。关于恶的问题必须得到充分的讨论；关于它已说的话，诚然似乎不太可能，但我认为不应简单地忽略它。问题是，邪恶是否也是藉着\OriginalTerm{圣言}{Logos}受造的——请注意，现在是将圣言理解为那内在于每个人的理性，而这理性是由那从起初就存在的理性所带来（受造）的。宗徒说：\BibleRef{罗 7:8} \BibleQuote{没有法律，罪是死的}，并加上：\BibleQuote{但诫命一来，罪又复活了}，这样一般地教导关于罪：罪在法律和诫命面前没有能力（但圣言在某种意义上就是法律和诫命），如果没有法律，就没有罪，因为\BibleRef{罗 5:13} \BibleQuote{没有法律，罪就不被归算}。再者，若不是因着圣言，就没有罪，因为\BibleRef{若 15:22} 基督说：\BibleQuote{我若没有来，没有对他们说话，他们就没有罪}。因为那想为自己的罪找借口的人，他的每一个借口都被除去，如果——虽然圣言在他内，指示他当作的事——他却不听从。那么，似乎一切事物，连更坏的事物也不例外，都是由圣言所造，而没有他——这里\Quote{没有}取了其更简单的含义——什么也没有受造。我们也不应责备圣言，如果万物都是藉着他造的，没有他就什么也没有造，正如我们不责备那位向学生指明其责任的老师，当他的教导使学生——若他犯罪——毫无借口或余地声称他是因无知而犯错。当我们考虑到老师与学生是不可分离时，这一点就更加明显。因为正如老师与学生是相关的，彼此相属，同样，圣言临在于具有理性的受造物的本性本身，总是提示他们当作的事，即使我们不理睬他的命令，却沉溺于享乐，任由他最善的劝告从我们身边溜过而不被留意。正如眼睛是为最好的目的赐予我们的仆人，然而我们却用它来观看不当看的事物；也正如我们错误地使用听觉，当我们花时间听歌唱竞赛和其他被禁止的声音时；同样，我们也侮辱我们内的圣言，并以不当的方式使用他，当我们使他协助我们的过犯时。因为他临在于犯罪的人，为了他们的定罪，并且他谴责那不以他超越一切的人。为此我们读到：\BibleRef{若 12:48} \BibleQuote{我所讲的话，要审判你们}。这好像他说：\Quote{我，圣言，常在你们内提高声音的，我自己要审判你们，那时你们将没有任何避难所或借口。}然而，这个解释可能显得有些牵强，因为我们将圣言在一个意义上理解为起初与天主同在的圣言，是天主圣言，而现在又在另一个意义上理解它，所论及的不仅关乎受造的主要工程，如\BibleQuote{万物都是藉着他造的}这句话，而且也关乎所有理性受造物的行为，这最后一点就是那圣言（理性），没有它的临在，我们就没有任何罪被犯下。问题出现了：我们内的圣言是否应与那从起初就有的、与天主同在的、是天主圣言的那一位被宣称为同一存在？宗徒当然似乎并不把在我们内的圣言与那起初与天主同在的圣言看作不同的存在。宗徒说：\BibleRef{罗 10:6} \BibleQuote{你不要心里说：\InnerQuote{谁能上天去呢？那就是把基督带下来；或谁能下到深渊去呢？那就是把基督从死者中带上来。}但圣经说了什么？圣言离你很近，就在你的口里，就在你的心里。}\par

% structure: p0022 source=h2 target=subsection reason=relative-heading-rank; base=h2

\subsection{十、\Quote{那受造的在他内是生命，这生命是人的光。}这包含一个悖论：不从圣言获得生命的，根本就不活。}

希腊人有某些被称为\Quote{悖论}的格言，其中将其贤者的智慧展示到极致，并给出一些证明或看似证明的证据。据说，唯有智者且每一位智者都是司祭，因为唯有智者且每一位智者都拥有关于敬事天主的智识。再者，唯有智者且每一位智者都是自由的，并从神律领受了权柄去做他自己所想做的事，这种权柄他们称为合法的决断权。为何还要多谈这些所谓的悖论呢？关于它们的讨论很多，它们要求将圣经的意义与此传授的教义进行比较，以使我们能够确定宗教教义在何处与它们一致，何处与它们不同。这是我们研究\Quote{在他内的造物就是生命}这句话而获得的启发；因为似乎可以遵循这里的圣经言辞，找出一系列具有悖论性质、甚至比这些希腊句子更悖论的事情。若我们思量在起初就与天主同在的圣言、就是天主圣言，或许我们能够宣称：唯有分有这种存有、且就这种特性而言的人，才可被称为理性的（\OriginalTerm{合乎圣言}{logical}）；由此我们就应证明唯有圣人才是理性的。再者，若我们领悟到生命已来到圣言中，即那说过\Quote{我是生命}的那一位，那么我们就应该说：凡在基督信仰之外的人都不是活着的，凡不向天主而活的人都是死的，他们的生命是向罪而活的生命，因此，恕我直言，那是死亡的生命。然而，请想一想，圣经是否在许多地方教导这一点？如救主所说（\BibleRef{谷 12:26}）：\BibleQuote{你们没有念过在荆棘篇中所说的吗？我是亚巴郎的天主，依撒格的天主，雅各伯的天主。他不是死人的天主，而是活人的天主。}以及\Quote{在你面前，凡有血肉的都不能成义。}但我们又何须谈论天主本身或救主呢？因为先知所说的\BibleQuote{我指着我的永生起誓——上主说}（\BibleRef{户 14:28}）这句话，究竟是出自谁的口，还有争议。\par

% structure: p0024 source=h2 target=subsection reason=relative-heading-rank; base=h2

\subsection{十一、在比较中，无人能相对于天主成义或真可谓生活。}

首先让我们看\Quote{他不是死人的天主，而是活人的天主}这句话。这相当于说：他不是罪人的天主，而是圣人的天主。因为对先祖们是一份极大的恩赐：天主以自己的名号为他们的名号加上，称自己为他们的天主，如保禄所说（\BibleRef{希 11:16}）：\Quote{所以天主不以被称为他们的天主为耻。}因此，他是列祖和所有圣人的天主；恐怕很难找到一处经文说天主是任何恶人的天主。那么，如果他是圣人的天主，并被称作活人的天主，则圣人是活人，活人是圣人；既没有在活人之外的圣人，也没有称某人为活人时没有进一步的含意：除了他有生命之外，他还是圣者。与此相近的是从\Quote{我将在活人之地上主面前中悦他}这句话中得到的教训。他似乎说，主的中悦是在圣人的行列中，或在圣人的处所，而他希望自己在其中。凡未进入圣人的行列或圣人的处所之人，不能中悦天主；而每个人若要到达那处所，必须预先在今世将真正中悦天主的影子与形像放在自己身上。那宣布在上帝面前凡活着的都不能成义的经文表明：与天主及在他内的义相比，即使最成全的圣人也不能成义。我们可从另一处打个比方：没有蜡烛能在太阳面前发光——不是因为蜡烛不能发光，而是当太阳照耀时它不能。同样，每一个\Quote{活着的人}都将成义，只是不在天主面前，当它与下面这些在黑暗权势中的人相比时。对他们，圣人的光将要照耀。或许这里我们找到了那节经文的含义（\BibleRef{玛 5:16}）：\BibleQuote{你们的光当在人前照耀。}他没有说：你们的光在天主面前照耀；若他这样说，他必是下了不可能实现的诫命，如同命令那些有灵魂的光去在太阳面前发光一样。因此，不仅普通的活人群众不能在天主面前成义，而且甚至在活人中那些超出其余者、或更确切地说，所有活人的义在与天主的义比较时，都不能在天主面前成义，如同我召聚所有在夜间照耀大地的光，说它们不能在太阳的光芒面前发光一样。当我们把\Quote{我指着我的永生起誓——上主说}这句话放在我们眼前时，我们便从这些考量上升到更高的层次。生命，按这个词的圆满意义，尤其在我们刚才所谈论的之后，或许只属于天主而非任何他者。这难道不是宗徒在谈论天主生命的至高卓越并被引向关于它的最高表达后，说关于天主的话（这表明他真正理解\Quote{我指着我的永生起誓——上主说}那句话）\BibleQuote{那独享不死不灭者}（\BibleRef{弟前 4:16}）的原因吗？除了天主之外，没有其他活物拥有不变易的生命。我们何需再怀疑？甚至基督也没有分有圣父的不死不灭；因为他\Quote{为每个人尝了死味}。\par

% structure: p0026 source=h2 target=subsection reason=relative-heading-rank; base=h2

\subsection{十二、救主是其所是，对众人皆然？}

我们先据此考察了天主的生命、即基督的生命，以及自成一类的活人，并看到活人在天主面前如何不成义，我们注意到了那句相关的声明：\Quote{惟有祂拥有不死}。现在可以探讨可能与此相关的假设，即任何赋有理智的存在者，其真福并非其本质的一部分，亦非其本性不可分离的部分。因为如果真福和至高的生命是理智存在者不可分离的特征，那么\Quote{天主惟独拥有不死}又如何能真实地说出口呢？因此我们应当注意，救主有些事物不是为祂自己，而是为他人；有些事物既为祂自己也为他人的；我们必须考察是否有某些事物是祂为自己而不为任何其他的。显然，祂为他人是牧人，但不是像人间那些以牧职谋利的人那样的牧人；除非因为祂对人的爱，羊群所受的益处也可视为祂自己的益处。同样，祂为他人是道路和门，并且（人人都会承认）是杖。祂为自己和他人则是智慧，或许也是\OriginalTerm{圣言}{Logos}。可以问：既然祂作为智慧自身拥有一套思辨体系，是否有某些思辨是任何受生的本性（唯独祂自己的本性除外）所不能接受的，并且是祂只为自己的缘故而知道的？我们应对圣神所怀的敬畏也不应阻止我们寻求回答此问题。因为圣神自己领受教导，这从关于护慰者和圣神所说的话中是清楚的，正如经上所载：\BibleRef{若 16:14} \BibleQuote{祂要从我所领受的，并传告给你们。}那么，祂从这些教导中，是否领受了圣子从起初注视圣父而自己知道的一切？这需要进一步考虑。如果救主有些事物是为他人，有些事物可能是为祂自己而不为任何其他的，或只为一位，或只为少数，那么我们要问：就祂是在圣言内来到的生命而言，祂是生命是为祂自己并为他人，还是只为他人？如果只为他人，又是为哪些他人？生命和人的光是同一回事吗？因为经文说：\Quote{凡受造的，在祂内有生命，这生命是人的光。}但人的光只是某些而非所有理性受造物的光；所加上的\InnerQuote{人}这个词表明这一点。但祂是人的光，因此祂也是祂所光照之人的生命。就祂是生命而言，祂可被称为救主，不是为祂自己，而是为成为他人的生命，祂也是他们的光。这生命来到圣言，且一旦来到，就与祂不可分离。但净化灵魂的圣言必须先存在于灵魂中；是在祂以及从祂而来的净化之后，当灵魂中一切死亡或软弱的东西被除去时，纯净的生命才来到每一个使自己成为圣言（被视为天主）的适当居所的人。\par

% structure: p0028 source=h2 target=subsection reason=relative-heading-rank; base=h2

\subsection{十三、圣言内的生命如何在起初之后来临}

此处，我们必须仔细观察，我们有两样事物，它们是一，我们必须界定它们之间的区别。首先，是在\Emph{在起初的圣言}中呈现于我们眼前的；其次，是在\Emph{圣言内的生命}中所隐含的。圣言不是在起初\Emph{受造的}；没有一时一刻起初是没有圣言的，因此经上记载说：\Quote{在起初已有圣言}。至于生命，另一方面，我们所读到的，不是它如同圣言那样存在，而是它受造的；至少如果生命是人的光这一说法成立的话。因为当人尚未存在时，就没有人的光；因为人的光只是就人与光的关系而言才被设想。并且，不要有人以这样的反对来烦扰我们：即我们把这归入了时间的范畴，尽管它是事物本身的次序，使它们成为第一、第二等等，甚至即使没有时间使圣言所安排的第三、第四的事物不存在。因此，既然万物是藉着祂\Emph{受造的}，不是万物藉着祂\Emph{存在}；既然没有祂，\Emph{没有}什么\Emph{受造}，但并非没有祂，什么也不\Emph{存在}；所以，在祂内\Emph{受造的}，而不是在祂内\Emph{存在的}，就是生命。并且，再次地说，不是那在起初\Emph{受造的}是圣言，而是那在起初\Emph{存在的}是圣言。确实，有些抄本有一种并非没有可能性的读法：\Quote{凡受造的，在祂内是生命}。但如果生命与人的光是同一回事，那么没有一个人是在黑暗中而活着的，也没有一个活着的人是在黑暗中；但凡活着的人也都在光中，凡在光中的也都是活着的，如此不仅那活着的人，而是每一个活着的人，都是光明之子；而光明之子就是那个其行为在人前发光的人。\par

% structure: p0030 source=h2 target=subsection reason=relative-heading-rank; base=h2

\subsection{十四、人的本性如何并非从起初就如此固定，以致他们可以从黑暗进入光明}

我们一直在讨论某些相反的事物，并且关于它们的论述可以提示那些未被提及的部分。我们正在谈论生命和人的光，与生命相反的是死亡；与人的光相反的是人的黑暗。因此，显然，处于人的黑暗中的人就是在死亡中，而行死亡之工的人除了在黑暗中，无处可寻。但是，思念天主的人，如果我们考虑思念意味着什么，他就不在死亡中，正如经上所说：\Quote{在死亡中，没有人纪念祢。}人的黑暗和死亡，按其本性是这样的吗？关于这一点，我们有另一段经文：\BibleRef{弗 5:8}\BibleQuote{我们从前是黑暗，但现在在主内是光明}，即使我们现在完全是圣人和属神的人。这样，从前是黑暗的人，像\Person{保禄}一样，成了能够在主内成为光明的人。有些人认为有些本性从一开始就是属神的，例如\Person{保禄}和圣宗徒们；但我几乎看不到如何将这种观点与上述经文协调，经文告诉我们，属神的人从前是黑暗，后来成了光明。因为如果属神的人从前是黑暗，那么属土的人会是什么？但如果黑暗变成了光明显然是真实的，正如经文所示，那么假设所有的黑暗都能变成光明，这有什么不合理呢？难道\Person{保禄}没有说过：\BibleQuote{我们从前在黑暗中，但现在在主内是光明}吗？从而暗示那些他们认为天生失落的人，曾经是黑暗，或者现在仍是黑暗，关于不同本性的假设或许可以接受。但\Person{保禄}明确说他从前是黑暗，但现在在主内是光明，这暗示黑暗可能转变为光明。但谁若察觉到这种向更好或更坏变化的可能性，就不难洞察每个人的黑暗，或那种存在于人的黑暗中的死亡。\par

% structure: p0032 source=h2 target=subsection reason=relative-heading-rank; base=h2

\subsection{十五、\Person{赫拉克利翁}的观点：主只把生命带给属神的人。对此的驳斥。}

\Person{赫拉克利翁}在读到这段经文时采取了一种相当粗暴的方法：\BibleQuote{在祂内所造的是生命。}他没有按原文的\Quote{在祂内}，而是理解为\Quote{对于那些属神的人}，好像他认为圣言和属神的是同一回事，尽管他没有明确说出这一点；然后他接着给出事情的起源，仿佛是一种解释，说：\Quote{祂（圣言）在他们出生时提供了最初的形式，进一步推行并显明了由另一位所播下的，使之成为形式、成为光照、成为自己的轮廓。}他没有注意\Person{保禄}如何谈论属神的人（\BibleRef{格前 2:14}），以及他如何避免说他们是\Emph{人}。\BibleQuote{属血气的人不接受天主圣神的事，因为在他看来是愚妄；但属神的人能判断一切。}我们坚持认为，他没有在\Quote{属神的}一词上加\Quote{人}字，这不是没有意义的。属神的是比人更好的东西，因为人是在灵魂或肉体或两者中接受其形式，而不是在比这些更神圣的东西，即精神中接受；只有当人开始拥有这精神的主导部分时，他才被称为\Quote{属神的}。此外，提出这样的假设，他甚至没有提供一点证明的假象，显示出他甚至在论证上连中等程度的可信度都达不到。关于他，就说这么多吧。\par

% structure: p0034 source=h2 target=subsection reason=relative-heading-rank; base=h2

\subsection{十六、生命也可能是其他人的光。}

让我们提出另一个问题，即生命是否只是人的光，而不是每一个处于福乐中的存有的光。因为如果生命等同于人的光，如果基督的光只为人存在，那么生命也就只为人存在。但这样的观点既愚蠢又不敬，因为其他圣经经文驳斥了这种解释，并宣告当我们稍微进步时，我们将与天使同等。（\BibleRef{玛 22:30}）这个问题应按如下原则解决：当某个谓词被应用于某些人时，不应立即认为它仅适用于他们。因此，当提到人的光时，它不仅仅是人的光；如果意思是那样，就会加一个词来表达；经文就会读作：生命只是人的光。因为人的光也可能成为人以外的其他人的光，正如某些动物和某些植物可以构成人的食物，而这些动物和植物也可能成为其他受造物的食物。这是一个来自日常生活的例子；另外，从受启发的书中引出一个类比是合适的。我们现在面临的问题是，为什么人的光不能也是其他受造物的光，我们已经看到，提到人的光绝不排除这光也可能是人以外的其他存有的光，无论这些存有是低于他还是类似于他。天主有一个名字，祂被称为\Person{亚巴郎}的天主、\Person{依撒格}的天主、\Person{雅各伯}的天主。那么，谁若从\BibleQuote{生命是人的光}这句话推断光只为人存在，他也应该得出结论：\Person{亚巴郎}的天主、\Person{依撒格}的天主、\Person{雅各伯}的天主只是这三位族长的天主。但祂也是\Person{厄里亚}的天主（\BibleRef{列下 2:14}），并且如\Person{友弟德}所说（\BibleRef{友 9:2}），是她的祖先\Person{西默盎}的天主，也是希伯来人的天主。通过类推论证，如果没有什么阻止祂成为其他人的天主，那么也没有什么阻止人的光成为人以外的其他人的光。\par

% structure: p0036 source=h2 target=subsection reason=relative-heading-rank; base=h2

\subsection{十七、更高的权能是人；基督也是他们的光。}

另一人再次引用经文：\BibleQuote{让我们按我们的肖像和模样造人}\BibleRef{创 1:26}并坚持认为凡按天主的肖像和模样造成的都是人。为了支持这一点，举出了无数例子，表明在圣经中\Quote{人}与\Quote{天使}可以互换使用，并且同一主体既可称为天使也可称为人。确实如此，如\Person{亚巴郎}所款待的那三位，以及来到\Place{索多玛}的那两位；在整部圣经中，人物有时被称为人，有时被称为天使。持此观点者会说，既然有些明显是人的人物也被称为天使，例如\Person{匝加利亚}说：\Quote{上主的使者，我与你们同在，这是万军的上主说的。}又如关于\Person{洗者若翰}所写的：\Quote{看，我派遣我的使者在你面前。}那么，天主的使者（天使）乃是因其职分而得名，而非因其本性被称为人。这一观点得到证实：应用于高位神体的名称并非生物种类的名称，而是等级的名称，由天主委派给这个或那个理性受造物。\Quote{宝座}不是生物的种类，\Quote{宰制}也不是，\Quote{率领}也不是，\Quote{掌权}也不是；这些都是受委任者的职务名称；主体本身无非是人，但该主体已成了宝座、宰制、率领或掌权。在\Person{若苏厄}书中，我们读到在\Place{耶里哥}有一个人向\Person{若苏厄}显现，说：\Quote{我是上主万军的统帅，现在我来。}结论是：人的光必须等同于每一个理性受造物的光；因为每一个理性受造物都是人，因为他们是按天主的肖像和模样造的。这光照以三种不同方式被提及：\Quote{人的光}、单纯\Quote{光}以及\Quote{真光}。它是人的光，要么正如我们之前所表明的，因为没有任何事物阻止我们视之为除了人之外的其他受造物的光，要么因为所有理性受造物都被称为人，因为他们都是按天主的肖像造的。\par

% structure: p0038 source=h2 target=subsection reason=relative-heading-rank; base=h2

\subsection{十八、天主如何也是光，但方式不同；以及生命如何在光之前。}

救主在此简称为光。但在这位\Person{若望}的公函中，我们读到天主是光。有人认为这证明了圣子在本质上与圣父并无不同。但另一位学者更仔细地审视此事，并以更健全的判断力说道：那照耀在黑暗中而不被黑暗征服的光，与那里面毫无黑暗的光，并非相同。照耀在黑暗中的光，可以说是降临在黑暗之上，并被黑暗追逐，尽管黑暗试图捕捉它，却没有被抓住。但毫无黑暗的光，既不照耀黑暗，起初也不被黑暗追逐，直到证明它是得胜者，并记载它未被追逐者抓住。第三种称号是\Quote{真光}。但是，既然天主是真理之圣父，就比真理更伟大、更超越；既然祂是智慧之圣父，就比智慧更伟大、更卓越；同样，祂比真光更超越。也许我们可以从\Person{达味}那里更深刻地学到圣父与圣子是两个光，他在第三十五篇圣咏中说：\Quote{在你的光中，我们必得见光。}这照耀在黑暗中的人的光，就是真光，在福音中又被称作世界的光；\Person{耶稣}说：\Quote{我是世界的光。}我们还必须注意，尽管经文大可写成：\Quote{那受造的，在他内是人的光，人的光就是生命}，但作者选择了相反的顺序。他把生命放在人的光之前，即使生命和人的光本是同一回事；在思想那些有分于生命的人时，虽然这生命也是人的光，我们应当首先注意到他们活在我们之前讨论过的那种神圣生命，然后再谈到他们的光照。因为生命必须先存在，才能光照那活着的人；如果先谈到一个尚未被设想为活着的人的光照，然后让生命在光照之后才出现，那就不合秩序了。因为虽然\Quote{生命}和人的\Quote{光}是同一回事，但这两个概念是分开考虑的。这人的光也被\Person{依撒意亚}称为\Quote{外邦人的光}，他说：\BibleQuote{看，我已立你作人民的盟约，外邦人的光。}\BibleRef{依 42:6}\Person{达味}也信赖这光，在第二十六篇圣咏中说：\Quote{上主是我的光明，我的救援，我还畏惧谁呢？}\par

% structure: p0040 source=h2 target=subsection reason=relative-heading-rank; base=h2

\subsection{十九、此处所说的生命是更高的生命，即理性的生命。}

至于那些编造关于\OriginalTerm{永世}{æons}的神话、并将它们排列成\OriginalTerm{对偶}{syzygies}（轭或对），且认为\OriginalTerm{圣言}{Logos}和生命是由\OriginalTerm{理智}{Intellect}和\OriginalTerm{真理}{Truth}所发出的人，提出以下困难或许并非离题。在他们的体系中，生命作为圣言的对偶，怎能从其同侪那里获得起源？因为他说：\Quote{凡受造的，是在祂内}，显然是指前面刚提到的圣言，\Quote{是生命}。他们会告诉我们，生命作为圣言的对偶（如他们所说）是如何在圣言内生成的？以及为何生命而非圣言是人类的光？如果心智健全的人被这些问题所困惑，并发现我们所提出的这一点难以解决，他们很可能会反过来邀请我们讨论：为何被称为人类之光的不是圣言，而是源于圣言的生命？对此问题，我们回答：此处所说的生命并非理性与非理性存在物所共有的那种生命，而是当理性在我们内达至圆满时加给我们的生命；我们分享那生命，源自\OriginalTerm{元始圣言}{first reason (Logos)}。只有当我们远离那仅外表是光而非真光的生命，渴望被真生命所充满时，我们才成为其分享者；当这生命在我们内升起时，它便成为\OriginalTerm{更高知识}{gnosis}之光的基础。对某些人而言，这生命可能只是潜能而非现实的光，即那些不努力探求更高知识之事的人；对另一些人，它则是现实的光。显然，那些遵循保禄训诫\Quote{切求更好的恩赐}的人便是如此；在最大的恩赐中，有一项是所有人都被吩咐寻求的，即智慧之言，随后是知识之言。智慧与知识彼此相邻；至于它们之间的区别，此处不宜探讨。\par

% structure: p0042 source=h2 target=subsection reason=relative-heading-rank; base=h2

\subsection{二十、不同种类的光；和黑暗。}

\Quote{光在黑暗中照耀，黑暗没有胜过光。}我们仍在探讨人类之光，因为这是前一节所提及的；并且我认为，也在探讨黑暗，它被称为光的对手，若定义正确，这黑暗也是人类的。人类之光是一个包含两种特殊事物的总概念；人类之黑暗也是如此。获得人类之光并分享其光辉的人，将行光明之事，并以更高意义认识，被更高知识之光所光照。我们必须认识到对立者及其恶行、以及那些被认为是知识却非真正知识之事的类似情况，因为行此事者具有的不是光而是黑暗的\OriginalTerm{理性}{Logos}。由于圣言知道产生光的事物，依撒意亚说：\Quote{因为祢的诫命是地上的光}，达味在圣咏中说：\Quote{主的诫命是明了的，能开启眼睛。}但除了诫命和规条之外，还有更高知识之光，我们在十二先知之一中读到：\BibleRef{欧 10:12}\Quote{你们当为自己播种正义，收割慈爱的果实，为自己制造知识之光，直到我来给你们降下正义。}除了诫命之外，还有知识之光，因此我们读到：\Quote{为自己制造知识之光}，不是简单的\Quote{制造}，而是哪种制造？——知识之光。因为若一个人为自己点燃的任何光都是知识之光，那么附加的\Quote{为自己制造知识之光}就毫无意义。此外，黑暗由恶行带给人类，我们从若望自己那里得知，他在其书信中说：\Quote{如果我们说我们与祂相通，却在黑暗中行走，我们就是撒谎，不履行真理。}又说：\Quote{那说自己是在光中，却恼恨自己弟兄的，到现在还是在黑暗中。}又说：\Quote{恼恨自己弟兄的，就是在黑暗中，并且在黑暗中行走，不知道往哪里去，因为黑暗弄瞎了他的眼睛。}黑暗中行走意指恶行，恼恨自己的弟兄，难道不就是偏离那真正被称为知识的吗？但那些对神圣事物无知的人也在黑暗中行走，正是因为那无知；如达味所说：\Quote{他们不知道，也不明白，在黑暗中走来走去。}然而，请思考这段经文：\BibleRef{若一 1:5}\Quote{天主是光，在祂内没有一点黑暗。}并看此说的理由是否在于黑暗不是单一的，而是双重的，因为有两种黑暗；或是许多的，因为黑暗被分散地、个别地指向许多恶行和许多虚假教义；所以有许多黑暗，没有一种是在天主内的。救主的话不能对圣者说：\Quote{你们是世界的光}；因为圣者是世界的光（绝对的，而非特殊的），在祂内没有任何黑暗。\par

% structure: p0044 source=h2 target=subsection reason=relative-heading-rank; base=h2

\subsection{二十一、基督并非如天主那样完全脱离黑暗：因为祂承担了我们的罪。}

现在有人会问，既然救主也是毫无罪污，那么说在祂内毫无黑暗这句话怎能被视为祂独有的特性呢？难道不也能同样说祂：\Quote{祂是光，在祂内毫无黑暗}吗？这两种情况之间的区别，我们已在上文中部分阐述过。然而，我们现在要进一步论述，并补充说，如果天主使那不认识罪的基督替我们成了罪，\BibleRef{格后 5:21} 那么就不能说在祂内毫无黑暗。因为如果耶稣带着罪恶肉身的形状\BibleRef{罗 8:3}，并为罪恶，且借着祂取得了罪恶肉身的形状而定了罪，那么就不能绝对直接地说在祂内毫无黑暗。我们还可以补充说：\BibleQuote{祂承担了我们的脆弱，担负了我们的疾病}\BibleRef{玛 8:17}，包括灵魂的脆弱和内心隐密之人的疾病。因了这些祂从我们身上背负的脆弱与疾病，祂宣称自己的灵魂忧闷至极、痛苦非常，\BibleRef{玛 26:38} 并且在匝加利亚书中被说成是穿上了污秽的衣服，\BibleRef{匝 3:4} 当祂要脱去时，那衣服就被称为罪恶。\BibleQuote{看，有话说，我已取走了你的罪恶。}因为祂将那些相信祂之人的罪恶担在自己身上，祂使用了诸多此类表达：\Quote{我的过犯之言远离我的救恩}，以及\Quote{你知道我的愚昧，我的罪恶没有向你隐藏。}且不要让任何人以为我们这样说是因为对天主的基督缺乏虔敬；因为正如唯独圣父具有不死不灭，而我们的主为了对人的爱，为自己承担了祂为我们所受的死亡，同样，\Quote{在祂内毫无黑暗}这句话也唯独适用于圣父，因为基督为了对人的善意，承担了我们的黑暗。祂这样做，是为了借祂的能力摧毁我们的死亡，除去我们灵魂中的黑暗，以应验依撒意亚的话：\BibleQuote{那坐在黑暗中的百姓看见了皓光}\BibleRef{依 9:1}。这光，即是在\OriginalTerm{圣言}{Logos}内受造且也是生命的光，照耀在我们灵魂的黑暗中；它已来到这黑暗的统治者与人类争战、并竭力迫使那些不全然坚定的人屈服于黑暗的地方；为使那些人得以被光照，光已如此远来，并使他们被称为光明之子。且这光在黑暗中照耀，被黑暗追逐，但黑暗没有胜过它。\par

% structure: p0046 source=h2 target=subsection reason=relative-heading-rank; base=h2

\subsection{二十二、黑暗如何未能胜过光}

若有人以为我们添加了未写之事，即黑暗对光的追逐，就请他思考：除非黑暗曾追逐光，否则\BibleQuote{黑暗没有胜过它}\BibleRef{若 1:5}这句话就没有意义。若望是为那些有智慧看出省略的部分并根据上下文加以补足的人而写，所以他写了：\BibleQuote{黑暗没有胜过它}\BibleRef{若 1:5}。若它没有胜过它，则必是先追逐了它；而黑暗确实追逐了光，可以从救主所受的苦以及那些接受祂教导的人，即祂的儿女，显明出来，那时黑暗尽其所能反对光明之子，并意图将光从人那里驱走。但是，如果天主偕同我们，\BibleRef{罗 8:31} 谁能反对我们呢？无论何等心怀恶意，都不能反对我们；他们越自谦，就越成长，并且极其得胜。黑暗没有胜过光，有两种方式：要么它被远远抛在后面，自己如此缓慢，而光在其行程中如此敏锐迅速，以至于它甚至无法继续追赶；要么，若光设下圈套捕捉黑暗，并依其形成的计划等候黑暗，那么黑暗走近光时，就被终结了。无论哪种情形，黑暗都没有胜过光。\par

% structure: p0048 source=h2 target=subsection reason=relative-heading-rank; base=h2

\subsection{二十三、有一种并非邪恶、且最终变为光的神圣黑暗}

关于这个问题，我们必须指出，\Quote{黑暗}一词并非每次提到都带有贬义；圣经有时也在褒义上使用它。异端者未能注意到这一区分，因而对世界的创造者采取了最可耻的教义，实际上背叛了祂，沉溺于虚构和神话。因此，我们必须说明\Quote{黑暗}之名在何时以及如何被用于褒义。在《出谷纪》中，黑暗、云彩和风暴被描述为环绕天主；在第十七篇《圣咏》中，也写道：\Quote{祂以黑暗为祂的密处，以云彩中的暗水为祂四围的帐幕。}的确，如果人们思考关于天主的那些超出人性接受能力的、或许也超出除基督和圣神之外所有受造物接受能力的浩繁思辨和知识，那么他们可能就会明白天主如何被黑暗所环绕，因为那本需用来解释祂在何处以黑暗为藏身处——当祂安排关于祂的事物不为人知、超出知识的把握时——的言谈被隐藏在无知之中。如果有人因这些解释而困惑，他可以通过\Quote{暗语}和\Quote{黑暗的宝藏}（即隐藏、不可见的宝藏，由天主赐予基督）来与它们调和。我认为，\Quote{黑暗的宝藏}隐藏在基督里，与经上所写的\Quote{祂以黑暗为祂的密处}以及（圣人）\Quote{将明白比喻和暗语}并无二致。\BibleRef{箴 1:6}请思考，我们是否在此发现了救主对其门徒所说的\Quote{你们在暗中所听到的，要在光明中宣讲出来}的原因。祂吩咐他们将在秘密中、在少数人能听到之处所托付的、难以知晓且隐晦的奥秘，在得到光照（因此被称为在光明中）后，向每一个已成为光的人宣告。我还可以补充这被称赞的黑暗的一个更奇特的特征，即它趋向光明并追上光明，以致最终，在作为黑暗而不为人知之后，对于那看不见其能力的人来说，它发生了如此变化，以至于他得以认识它，并宣告以前被他视为黑暗的如今已变为光明。\par

% structure: p0050 source=h2 target=subsection reason=relative-heading-rank; base=h2

\subsection{二十四、\Person{洗者若翰}被派遣。从何处？他的灵魂来自更高之处。}

\Quote{曾有一人，由天主派遣，名叫若翰。} \BibleRef{若 1:6} 他被派遣，是从某处被派遣到某处；因此细心的学者会探询若翰是从哪里被派遣，又往哪里去。\Quote{往哪里}在故事表面很清楚：他被派往以色列，并派往那些愿意听他宣讲的人，当时他住在犹太旷野，在约旦河畔施洗。然而，按更深的意义，他是被派遣到世界，世界被理解为这尘世之地，即人类所在之处；细心的学者在探询若翰是从哪里被派遣时，会注意这一点。更仔细地审视这些话，他或许会宣称，正如论及亚当所写的：\BibleQuote{天主就遣发他离开伊甸乐园，去耕种他所由出的土地} \BibleRef{创 3:23}，同样，若翰也是被派遣的，或是从天，或是从乐园，或是从别处到地上的此地。他被派遣，是为给光作证。然而，这种解释有一个不容轻忽的反对意见。依撒意亚书中有：\Quote{我将派遣谁呢？谁肯为我们去呢？}先知回答：\Quote{我在这里，请派遣我！}因此，反对那看似更深的经文解释的人可以说，依撒意亚不是从别处被派遣到这世界，而是在看见了\Quote{主坐在崇高的宝座上}之后，被派往人民那里，去说：\Quote{你们听了，却总不明白}等等；同样，若翰的使命开始并未被叙述，他是按照依撒意亚受使命的类比被派遣，去施洗，\BibleRef{路 1:17}为主预备一个合用的百姓，并为光作证。以上我们说了第一种意义；现在提出一些有助于证实关于若翰的更深意义的解决方式。在同一段中，接着又说：\Quote{这人来，是为作证，为给光作证}。那么，如果他是\Quote{来}，他从哪里来？对于那些难以跟随我们的人，我们指出若翰后来所说关于圣神如鸽子降在救主身上的话：\Quote{那派遣我来以水施洗的，给我说：你看见圣神降下，停在谁身上，谁就是那要以圣神及火施洗的人} \BibleRef{若 1:33}。祂什么时候派遣他并给他这个命令？这个问题的答案大概是：当祂派遣他开始施洗时，那位与他打交道者说出了这话。然而，认为若翰是从另一区域被派遣进入身体、他进入此生唯一目的是为真理作证的观点，有一个更有说服力的论证，可以从他诞生的记述中得出。加俾额尔在向匝加利亚宣布若翰的诞生，并向玛利亚宣布救主降临人间时，说：若翰将\Quote{从母胎中就充满圣神}。我们还有这话：\Quote{看，你问安的声音一入我耳，胎儿就在我腹中欢喜踊跃}。凡在对待圣经时谨慎地避免牵强、随意或任性做法的人，必然要假定若翰的灵魂比他的身体更早存在，且在它被派遣执行给光作证的使命之前，已经独立存在。我们也不可忽视这句经文：\BibleQuote{这位就是那要来的厄里亚} \BibleRef{玛 11:14}。因为，如果接受灵魂的一般教义，即它并非与身体同时被播下，而是先于身体，然后为各种原因被披上血肉，那么\Quote{由天主派遣}这句话就不只适用于若翰一人。甚至那最邪恶的、那罪人、那丧亡之子，保禄也说他是天主派遣的：\BibleQuote{天主就给他们一种错谬的能力，使他们相信谎言，为使一切不信真理而喜爱不义的人，都被定罪} \BibleRef{得后 2:11}。但当下问题或许可以用以下方式解决：正如每个人都是属于天主的人，仅仅因为天主创造了他，但并非每个人都被称为天主的人，只有那些奉献自己于天主的人，如厄里亚和圣经中被称为天主的人，类似地，每个人在一般意义上都可以说是由天主派遣的，但在绝对意义上，没有人可以被这样称呼，除非他进入此生是为了神圣的使命并为人类救恩的服务。我们发现，除了圣人之外，没有人被说成是由天主派遣的。依撒意亚是这样说的，正如我们之前所展示的；耶肋米亚也是如此：\Quote{我派你到谁那里，你就应去} \BibleRef{耶 1:7}；厄则克耳也是如此：\Quote{我派遣你到背叛的人民那里，他们却不相信我} \BibleRef{则 2:3}。然而，这些例子并没有明确说从生命之外的区域进入生命的使命，既然我们探究的是进入生命的使命，它们可能看起来与我们的主题关系不大。但将从中得出的论证转移到我们的问题上，并没有什么荒谬之处。它们告诉我们，只有圣人（我们刚才谈论的就是他们）才被天主说成是派遣的，在这个意义上，它们可以适用于那些被派遣进入此生的人。\par

% structure: p0052 source=h2 target=subsection reason=relative-heading-rank; base=h2

\subsection{二十五、从若瑟的祈祷而来的论证，以表明洗者若翰可能是一位成为人的天使。}

既然我们现在讨论的是关于若翰的事，并且追问他的使命，我正好借此机会陈述我对他的看法。我们读到了关于他的这个预言：\Quote{看，我派遣我的使者（天使）在你面前，他要在你前面预备你的道路。}由此我们询问，是否有一位圣天使被派遣到这项职务中，作为我们救主的前驱。当受造物中的首生者取了人身时，一些天使竟充满对人的爱，成为基督的钦佩者和追随者，并认为以类似人的身体来服务他的仁慈是好的，这并不奇怪。谁不因他还在母腹中时就欢喜跳跃、超越普通人性而感动呢？若那部在希伯来人中流传的伪经作品\WorkTitle{若瑟的祈祷}被认为值得信赖，那么这一教义在其中是明确表达的。它表明，起初的某些灵魂，因他们是天使而具有超越众人的显著区别，比其他灵魂更伟大，它们降到了人性中。因此雅各伯说：\Quote{我，雅各伯，正对你们说话，还有以色列，我是天主的天使，一个统治的灵；亚巴郎和依撒格是在天主的每个创造之先被造的；我是雅各伯，人称之为雅各伯，但我的名字是以色列，被天主称为以色列，一个看见天主的人，因为我是一切天主使其存活的受造物中的首生者。}他接着又说：\Quote{当我从叙利亚的美索不达米亚来的时候，天主的天使乌黎耳出来说：我降到地上，居住在人中间，我被称为雅各伯。他向我发怒，与我争斗，与我摔跤，说他的名字以及那位在诸天使之前者的名字应在我名字之前。我告诉了他他的名字，以及他在天主之子中如何伟大：你不是我第八位乌黎耳吗？我是以色列，是主权能的总领天使，是天主之子中的统帅。难道我不是以色列，是天主面前的首席仆役吗？我以不灭之名呼求了我的天主。}很可能这些话确实是雅各伯说的，因而被记录下来，而我们所听说的\Quote{他在母腹中取代了他的兄长}也含有更深的含义。请思考雅各伯与厄撒乌那个著名的问题是否有解答。我们读到：\BibleRef{罗 9:11}\BibleQuote{孩子还没有出生，也没有行善或作恶，为使天主拣选的计划坚定不移，不是出于人的行为，而是出于召唤人的那一位，就对她说：\InnerQuote{年长的要服事年幼的。}}正如所记载的：\BibleQuote{我爱了雅各伯，而恨了厄撒乌。}那么，我们该说什么呢？难道天主有不义吗？绝对不可！那么，当他们还没有出生，也没有行善或作恶时，为使天主拣选的计划坚定不移，不是出于行为，而是出于召唤人的那一位，若那时就说了这话，那么若不追溯到此生之前的行为，怎能说天主没有不义呢？年长的服事年幼的，且（被天主）恨，是在他做了任何配得奴役或憎恨的事之前。我们引入关于雅各伯的故事并引用一部我们不能轻蔑的作品，多少有些偏离主题；但这确实增强了我们关于若翰的论点，即正如依撒意亚的声音宣告的：\BibleRef{依 40:3}他是一位为了给光作证而取了肉身的天使。以上就是关于被视为人的若翰的话。\par

% structure: p0054 source=h2 target=subsection reason=relative-heading-rank; base=h2

\subsection{二十六、若翰是声音，耶稣是言语。二者的关系。}

如今我们知道声音与言语是不同的。声音可以不具任何意义而产生，其中没有言语；同样，言语可以在没有声音的情况下被心灵所接收，正如我们在内心进行思想漫游时那样。因此，救主，从某一角度看，是言语，而若翰与祂不同；正如救主是言语，若翰是声音。若翰自己邀请我这样看待他，因为对那些问他\Quote{你是谁}的人，他回答说：\Quote{我是荒野中呼喊者的声音：修直上主的道路！使祂的路径笔直！}这或许解释了，为何匝加利亚在诞生那指出天主圣言的声音时失了声，直到那声音——圣言的前驱——诞生时才恢复。声音必须被耳朵所感知，以便心灵随后接收声音所指示的言语。因此，若翰在出生上比基督稍微年长，因为我们的声音在言语之先来到。但若翰也指向基督；因为言语由声音引出。基督受若翰的圣洗，尽管若翰自称需要受基督的圣洗；因为在人身上，言语由声音净化，尽管自然的方式是言语应该净化那指示它的声音。总之，当若翰指出基督时，是人指出天主，是肉身指出无形的救主，是声音指出圣言。\par

% structure: p0056 source=h2 target=subsection reason=relative-heading-rank; base=h2

\subsection{二十七、若翰及其双亲名字的意义。}

名字中蕴涵的力量可用于许多事，此时此刻，我们不妨追问若翰和匝加利亚这两个名字的意义。亲戚们认为给孩子取名不是一件可以轻率对待的事，因此想叫他匝加利亚，并且对依撒伯尔想叫他若翰感到惊讶。匝加利亚随即写道：\Quote{他的名字是若翰}，立刻就从困扰他的缄默中解脱出来。于是，在考察这些名字时，我们发现\Quote{\OriginalTerm{若翰}{Joannes}}是\Quote{\OriginalTerm{若阿}{Joa}}去掉\Quote{\OriginalTerm{乃斯}{nes}}。\WorkTitle{新约}将希伯来名字赋予希腊形式，并视之为希腊词；雅各伯变为雅各布斯，西默盎变为西满，而若翰等同于若阿。匝加利亚被认为是\Quote{记忆}，依撒伯尔则是\Quote{我天主的盟誓}或\Quote{我天主的力量}。于是，若翰从天主的恩宠（=若阿=若翰）来到世界，他的父母是（关于天主的）记忆和我们的天主的盟誓，关乎祖先们。他为此而生，好为主预备一个相称的子民，在现已衰老的盟约之末，也就是安息时期的终结。因此，安息日后的安息不可能出自我们天主的第七日；相反，是我们的救主，按照祂自己安息的样式，使我们相似祂的死亡，从而也相似祂的复活。\par

% structure: p0058 source=h2 target=subsection reason=relative-heading-rank; base=h2

\subsection{二十八、先知们为基督作证，并预言了许多关于祂的事。}

\BibleQuote{他来作证，为给光作证，为使众人藉着他而信。}\BibleRef{若 1:7}有些脱离教会教义的人，自称信基督，却因他们的体系需要另一位神，而否认创造者，因此不能接受祂的降临曾被先知预言。他们试图抹去先知关于基督的见证，说圣子不需要见证，祂自带证据，部分在于祂所宣讲充满能力的有力言辞，部分在于祂所行的奇事，这些足以立时令任何人信服。他们说：若梅瑟因言语和作为而被信服，不需任何见证预先宣告他；若每位先知都被这些人接纳为天主的使者，那么远超过梅瑟和先知的那位，岂不更应无需先知作见证，就完成使命造福人类？他们认为基督被先知预言是多余的，因为先知们的关注点（如这些反对者所说）在于使信基督的人不致接纳祂为新神，因此尽力把他们引向梅瑟和先知在耶稣之前所教导的同一位天主。对此我们必须说：正如有多种原因可引导人相信，因不为一种论证所动者可被另一种打动，同样，天主能为人提供多种机缘，每种都可使他们的心对那统御万有的天主取了人性这一真理敞开。众所皆知，有些人因先知著作而钦佩基督。他们惊异于这么多先知在祂之前的预言，这些预言确立了祂的出生地、成长地、教导能力、行奇事的能力、以及以复活告终的人性苦难。我们还须注意，基督当时所行的大能奇事能使当时的人信服，但随时间流逝其证明力减弱，开始被视为神话。如今将这些奇迹与预言相比，其证据价值远大于当时奇事本身，这使得研究者无法怀疑前者。先知见证不仅宣告默西亚来临；绝非仅教导这一点。它们教导大量神学。关于圣父与圣子的关系、圣子与圣父的关系，从先知所宣告的基督事迹，与宗徒所叙述的圣子荣光，同样可以学到。我们可斗胆举一个类似例子：殉道者因作证基督而受尊荣，绝非因他们为圣子作证而赐予祂任何恩惠。那么，若许多基督的真门徒因这样为祂作证而受尊荣，那么先知们从天主领受的特殊恩赐岂不就是明白关于基督的事并预先宣告祂？他们不仅教导基督降临之后的人如何看待圣子，也教导祂之前世代的人？因为在这时代，不认识圣子的也不认识圣父（\BibleRef{若一 2:23}），同样，在更早时代也应如此理解。因此，\BibleQuote{亚巴郎欢欣踊跃地希望看见我的日子，他看见了，且喜乐。}\BibleRef{若 8:56}所以，那宣告他们不应为基督作证的人，是在剥夺先知歌咏团最大的恩赐；因为若从预言中除去一切与我们的主和师傅的救世工程的联系，那么受圣神默感的预言还留下什么同等重要的职务呢？正如那些经由中保、大司祭和护慰者走向宇宙之神的人，其信仰井然有序；而不经门进到父前的人，其宗教是蹩脚的；同样，古代的人也是如此。他们的宗教因对基督的认识、信仰和期待而得以圣化，并为天主所悦纳。我们观察到天主自称是见证，并劝勉众人也为基督作证，效法祂，向一切需要的人为祂作证。因为祂说（\BibleRef{依 43:10}）：\BibleQuote{你们是我的证人，我也是证人——上主天主说——以及我所拣选的仆人。}凡是为真理作证的人，无论是用言语或行为，或任何方式，都可恰当地称为证人（殉道者）；但弟兄们已有习惯，因敬佩那些为真理和勇德奋斗至死的人，更恰当地将殉道者之名保留给那些流血为虔敬奥迹作证的人。救主称每一位为祂所宣告的真理作证的人为殉道者；例如在升天时，祂对门徒说（\BibleRef{宗 1:8}）：\BibleQuote{你们将在地上，在犹太和撒玛黎雅，并直到地极，为我作证人。}被洁净的癞病人（\BibleRef{玛 8:4}）仍要献上梅瑟所命令的礼物，给那些不信基督的人作证据。同样，殉道者为不信者作证据，所有在世人面前善行发光的圣人也如此。他们度着以基督十字架为乐的生活，为真光作证。\par

% structure: p0060 source=h2 target=subsection reason=relative-heading-rank; base=h2

\subsection{二十九、洗者若翰的六个见证列举。耶稣的\BibleQuote{你们来看看}。第十时辰的意义。}

于是\Person{若翰}来为\OriginalTerm{光}{the light}作证，在他作证时，他呼喊说：\BibleQuote{那在我以后来的，存在在我以前，因为祂在我以前。从祂的满盈中，我们都领受了，而且恩宠上加恩宠。因为法律是藉\Person{梅瑟}传授的，但恩宠和真理却是通过\Person{耶稣基督}而来的。从来没有人见过天主；那\OriginalTerm{独生子天主}{the only-begotten God}，在圣父怀里的，祂给我们详述了。}整段话是洗者向基督作证时所说的。有些人另有看法，认为从\BibleQuote{从祂的满盈中}到\BibleQuote{祂给我们详述了}这些话是作者\Person{若望}宗徒写的。实际情况是，\Person{若翰}的第一个证词，正如我们之前所说，以\BibleQuote{那在我以后来的}开始，以\BibleQuote{祂给我们详述了}结束；他的第二个证词是向从\Place{耶路撒冷}派来的司祭和肋未人所说的，他们是犹太人派来的。他对他们承认，并不否认事实，即他不是基督，也不是\Person{厄里亚}，也不是那位先知，而是\BibleQuote{在旷野里呼喊者的声音：修直主的道路，正如\Person{依撒意亚}先知所说的。}此后，同一位洗者对基督又有另一个证词，仍教导祂更崇高的本性，这证词传遍全世界，进入理性的灵魂。他说：\BibleQuote{有一位站在你们中间，你们却不认识祂；就是那在我以后来的，我连给祂解鞋带也不配。}试想，既然心脏位于整个身体的中部，而\OriginalTerm{统御原则}{the ruling principle}在心脏中，那么\BibleQuote{有一位站在你们中间，你们却不认识祂}这句话是否可以理解为每一个人内在的\OriginalTerm{理性}{the reason}？此后，\Person{若翰}关于基督的第四个证词指向祂的人间苦难。他说：\BibleQuote{请看，\OriginalTerm{天主的羔羊}{the Lamb of God}，除免世罪者！这就是我曾说过的：在我以后来了一个人，祂存在在我以前，因为祂在我以前。我本来不认识祂，但为使祂彰显于以色列，因此我来以水施洗。}第五个证词记载为：\BibleQuote{我看见圣神像鸽子从天降下，停留在祂身上。我本来不认识祂，但那派我来以水施洗的，对我说：你看见圣神降下并停留在谁身上，谁就是那以圣神施洗的。我看见了，也作证：这就是天主子。}第六次，\Person{若翰}向两个门徒为基督作证：\BibleQuote{他看见\Person{耶稣}走过，就说：请看，天主的羔羊。}这个证词之后，那两个门徒听见了，就跟随了\Person{耶稣}；\Person{耶稣}转过身来，看见他们跟着，就对他们说：\BibleQuote{你们找什么？}也许并非没有意义，在六个证词之后，\Person{若翰}停止了作证，\Person{耶稣}在第七次提出了祂的\BibleQuote{你们找什么？}那些因\Person{若翰}的证词而受益的人所说的话非常合适：他们称基督为老师，并表示希望看见天主子居住的地方；因为他们对祂说：\OriginalTerm{辣彼}{Rabbi}，即\Quote{老师}，以我们的语言说，\BibleQuote{你住在哪里？}并且，因为凡寻找的，就必找到；当\Person{若翰}的门徒寻找\Person{耶稣}的住所时，\Person{耶稣}就指给他们看，说：\BibleQuote{你们来，看吧！}藉着\BibleQuote{来}这个词，祂或许劝勉它们进入生活的实践部分；而\BibleQuote{看}则暗示他们，那随正确行为而来的默观，必会赐给那些渴望它的人；在\Person{耶稣}的住所里，他们将获得它。他们问过\Person{耶稣}住在哪里，并跟随了老师，看到了之后，便希望与祂同住，并与天主子一起度过那一天。数字十是神圣的，为数不少的奥秘藉此被指明；因此我们应当明白，提到\OriginalTerm{第十时辰}{the tenth hour}是这些门徒与\Person{耶稣}同住的时候，并非没有意义。在这些门徒中，有一个是\Person{西满伯多禄}的兄弟\Person{安德肋}；他因这一天与\Person{耶稣}相处而受益，找到了自己的兄弟\Person{西满}（或许他之前没有找着他），告诉他说：他找到了\OriginalTerm{默西亚}{the Messiah}，默西亚翻译出来就是基督。经上记载：\BibleQuote{寻找的，就必找到。}现在，他曾寻找\Person{耶稣}的住所，跟随了祂，观看了祂的住所；他在\BibleQuote{第十时辰}与主同住，找到了天主子、\OriginalTerm{圣言}{the Word}和\OriginalTerm{智慧}{Wisdom}，并受祂以君王的方式治理。这就是为什么他说：\BibleQuote{我们找到了默西亚。}这是一件每一个找到天主圣言并受祂的天主性如同君王般治理的人都能说的事。作为果实，他立刻带他的兄弟到基督那里；基督垂顾\Person{西满}，也就是说，藉着注视他，眷顾并启迪了他的统御原则；\Person{西满}因\Person{耶稣}的注视而得以强壮，从而因那坚定与力量的工作获得一个新名字，并被称为\Person{伯多禄}。\par

% structure: p0062 source=h2 target=subsection reason=relative-heading-rank; base=h2

\subsection{三十、若翰如何为基督作证，尤其是为\OriginalTerm{光}{The Light}作证。}

可能有人会问，既然我们面前的经文是\BibleQuote{他来是为作证，为给光作证}，我们为何还要经过所有这些讨论？但有必要给出若翰对光的见证，并表明这些见证发生的次序，而且，为了显示若翰的见证是多么有效，还须说明他后来对所见证的人提供的帮助。但在所有这些见证之前，有一个更早的见证，即当洗者若翰在依撒伯尔腹中因玛利亚的问安而跳跃时。那是对基督的见证，证实了他神圣的受孕和诞生。我还有什么可说的？若翰无处不在是基督的见证人和前驱。他先于基督的诞生，并稍微早于天主之子的死亡而去世，因此不仅为那些在基督诞生时的人作证，也为那些等待通过基督死亡而获得人类即将到来的自由的人作证。如此，在他整个一生中，他略早于基督，处处为主预备了一群归向他的百姓。若翰的见证也先于基督第二次、更神圣的来临，因为我们读到：\BibleRef{玛 11:14} \BibleQuote{如果你们愿意接受，他就是那要来的厄里亚。有耳听的，就让他听吧。} 现在，有一个太初，圣言就在那里——我们从\BibleBook{}{箴言}看到，那太初就是智慧——圣言早已存在，并且在圣言中生命被造，生命就是人的光；这一切既已如此，我问：那位从天主那里差来的名叫若翰的人，他为何来是为作证，尤其是为给光作证？他为何不来为生命、或为圣言、或关于太初、或关于基督显现的许多其他方面作证？请思考这些经文：\BibleQuote{那坐在黑暗中的百姓，看见了皓光} 和 \BibleQuote{光在黑暗中照耀，黑暗却没有胜过它}，并思考那些处于黑暗中的人，即人类，需要光。因为如果人的光在黑暗中照耀，而黑暗没有能力达到它，那么我们必须分享基督的其他方面；目前我们根本没有真正分享到他。因为我们这些仍处于死亡之身中的人，我们的生命与基督一同藏在天主内，我们对生命有什么分享呢？\BibleRef{哥 3:3} \BibleQuote{因为当基督，我们的生命显现时，那时我们也要与他一同在光荣中显现。} 因此，那位来到的人不可能为那仍然与基督一同藏在天主内的生命作见证。他也不是来为圣言作见证，因为我们认识那太初与天主同在、且是天主圣言的圣言；因为圣言在人间成了血肉。即使见证表面上似乎是关于圣言的，实际上也将是关于成了血肉的圣言，而非关于天主的圣言。因此，他不是来为圣言作证的。再者，如何能为智慧作见证呢？对于那些即使似乎知道一些事，却不能理解纯正的真理，反而透过镜子观看、在谜语中观看的人？然而，很可能在基督第二次、更神圣的来临之前，若翰或厄里亚会来，略早于基督——我们的生命——显现时，为生命作见证，然后他们将为圣言作见证，并也献上他们对智慧的见证。有必要探究一下，是否像若翰那样的见证必须先于基督的每个方面。关于\BibleQuote{他来是为作证，为给光作证}这句话，就说这么多。对于后续的话\BibleQuote{为使众人都藉着他而相信}，我们可稍后再思考其含义。\par

\SourceRecord{Translated by Allan Menzies.；Ante-Nicene Fathers；Vol. 9.；Edited by Allan Menzies.；Buffalo, NY: Christian Literature Publishing Co.,；1896.；<http://www.newadvent.org/fathers/101502.htm>.}

% structure: p0001 source=h1 target=section reason=page-title

\section{\WorkTitle{若望福音释义（第四卷）}}

[\Emph{开头散佚三页}]。\par

1. 凡在自己内区分言语、意义以及意义所指之事的人，不会因言语的粗俗而见怪，只要他发现所论之事是可靠的。当我们记起圣人们承认他们的言语和宣讲不是依靠智慧动听的言辞，而是藉圣神和德能的明证时，更是如此……\par

[\Emph{接着，在论及福音文体之粗俗后，他继续说道：}]\par

2. 宗徒们并非不知自己在某些事上令人反感，在某些方面学识有所欠缺，他们承认自己俗言俚语，却不昧于知识，正如\BibleRef{格后 11:6}所载；因为我们应当认为，其他宗徒也会说同样的话，正如保禄一样。至于经文\BibleRef{格后 4:7}\BibleQuote{但我们有这宝贝放在瓦器里，要显明这莫大的能力是出于天主，不是出于我们。}我们这样解释：此处所理解的\Quote{宝贝}，与其他经文一样，是指\OriginalTerm{知识}{gnosis}和隐藏的智慧之宝藏；\Quote{瓦器}则指圣经那看似卑微的文辞，希腊人容易因此而轻视它们，而天主能力的卓越在其中彰明较著。真理的奥秘和所说之事的德能，并未因卑微的文辞而受阻，未能传至地极，也未能使世界上的愚拙人、有时甚至连智慧人都服从基督的圣言。因为我们看看你们的蒙召，\BibleRef{格前 1:26}不是按肉体没有智慧的人，而是按肉体有智慧的不多。但保禄在宣讲福音时，\BibleRef{罗 1:14}有义务不但将圣言传给野蛮人，也传给希腊人；不但传给愚拙人（他们容易同意他），也传给智慧人。因为天主使他胜任\BibleRef{格后 3:6}作新约的仆役，行使圣神和德能的明证，好叫信徒赞同他时，他们的信仰不在于人的智慧，而在于天主的德能。因为，假如圣经像希腊人推崇的那些作品一样，具有优雅和言辞的驾驭力，那么就有可能让人以为，不是真理本身抓住了人，而是言语表面上的连贯和华丽俘虏了听众，并且是诡诈地俘虏了他们。\par

\SourceRecord{Translated by Allan Menzies.；Ante-Nicene Fathers；Vol. 9.；Edited by Allan Menzies.；Buffalo, NY: Christian Literature Publishing Co.,；1896.；<http://www.newadvent.org/fathers/101504.htm>.}

% structure: p0001 source=h1 target=section reason=page-title

\section{\WorkTitle{若望福音释义（第五卷）}}

% structure: p0002 source=h2 target=subsection reason=relative-heading-rank; base=h2

\subsection{前言}

您不满足于当我在您身边时，扮演监工的角色，驱使我劳碌于神学；即使我不在时，您也要求我将大部分时间花在您和我为您所做的工作上。就我而言，我倾向于逃避辛劳，并避免那从天主而来的、威胁那些投身于写作神学著作之人的危险；因此，我要躲藏在圣经中，藉以克制自己不多写书。因为\Person{撒罗满}在\BibleBook{}{训道篇}中说：\BibleQuote{我儿，应留心勿多写书；毫无止境，多习劳倦身体。}至于我们，除非这段经文有某种我们尚未察觉的隐秘含义，我们直接违犯了禁令，我们没有防备自己多写书。\par

\EditorialNote{然后，在指出对福音中仅几句话的讨论已写至四卷之后，他接着说：}\par

% structure: p0005 source=h2 target=subsection reason=relative-heading-rank; base=h2

\subsection{二、圣经如何警告我们不要多写书}

因为，从这句话的字面看，\BibleQuote{我儿，应留心勿多写书，}似乎指示了两件事：第一，我们不应收藏许多书；第二，我们不应撰写许多书。如果第一点不是其含义，那必定是第二点；如果第二点是其含义，第一点则不一定随之而来。无论哪种情况，我们似乎都被告知不应多写书。我本可以藉此直面我们的箴言，将经文作为借口寄给您，并援引连圣人也没有闲暇撰写许多书这一事实来支持这一立场；这样我就可以摆脱我们彼此订立的约定，放弃撰写要寄给您的作品。您方面无疑会感受到我所引用的经文的力量，将来或许会原谅我。但我们必须凭良心对待圣经，不可因看到经文的字面意义就自诩完全理解了它。因此，我不免要提出我认为能够为自己辩护的理由，并指出如果您在我违背约定时定会对我的反驳。首先，神圣历史似乎与所引经文一致，因为没有任何一位圣人撰写过多部著作，或以多书阐述其观点。我将提出这一点：当我着手撰写许多书时，批评者会提醒我，连\Person{梅瑟}这样的人物也仅留下了五部书。\par

% structure: p0007 source=h2 target=subsection reason=relative-heading-rank; base=h2

\subsection{三、宗徒们所写很少}

但那成为\WorkTitle{新约}的仆役，不是凭文字，而是凭圣神的\Person{保禄}，他传报了福音，从\Place{耶路撒冷}直到\Place{依里黎苛}，\BibleRef{罗 15:19}并没有给他教导的所有教会写信，而给他写信的教会他也只写了寥寥几行。\Person{伯多禄}，\Person{基督}的教会建立在他身上，地狱的门不能胜过她，\BibleRef{玛 16:18}只留下一封公认的真迹书信。假设我们允许他留下第二封，因为这是可疑的。我们该如何说那位倚靠在\Person{耶稣}胸膛上的，即\Person{若望}，他留下了一部\BibleBook{若望}{福音}，尽管承认他能写这么多以致世界容不下它们？但他也写了\BibleBook{}{默示录}，被命令缄默，不可写下七种雷声。\BibleRef{默 10:4}但他也留下一封只有寥寥几行的书信。假设还有第二封和第三封，因为并非所有人都宣称这些是真实的；但两者加在一起也不足一百行。\par

\EditorialNote{然后，在列举了先知和宗徒，并指出每位只写了很少、甚至很少之后，他接着说：}\par

% structure: p0010 source=h2 target=subsection reason=relative-heading-rank; base=h2

\subsection{4}

我感到自己因这一切而眩晕，疑惑顺从您是否就是在顺从天主，是否行走在圣人的足迹上，除非是我对您过分的爱和不愿使您痛苦，使我误入歧途并想起了所有这些借口。我们从训道者的话开始，他说：\BibleQuote{我儿，应慎防撰写许多书籍。}对此我引用同一\Person{撒罗满}的\WorkTitle{箴言}中的一句话：\BibleQuote{多言多语难免犯过；但抑制口唇是明智的。}我在此问道：无论讲述何种言论，只要是许多言语，是否就是训道者所说的\Quote{多言多语}？即使一个人所说的许多话是神圣的并与救恩相关？如果是这样，且那使用许多有益言辞的人犯了\Quote{多言多语}之罪，那么\Person{撒罗满}自己也没有逃脱这罪，因为\BibleQuote{他讲了\BibleRef{列上 4:32}三千句箴言和五千首歌，又讲论草木，自黎巴嫩的香柏树直到墙上的牛膝草，也讲论走兽、飞鸟、爬虫和鱼类。}我请问，谁能在最简单的意义上不通过多言多语而进行任何教导？智慧本身不是对那些丧亡的人说：\BibleRef{箴 1:24}\BibleQuote{我伸出了我的言语，你们却不理会。}我们不是也发现\Person{保禄}从早晨一直讲到半夜，\BibleRef{宗 20:7}那时\Person{厄提约}因沉睡而坠落，令听众惊慌，以为他死了？那么，如果\Quote{多言多语难免犯过}这句话是真的，而\Person{撒罗满}在谈论上述主题时并未犯大罪，\Person{保禄}也将讲话延长到半夜，那么问题就出现了：这里所指的多言多语是什么？然后我们可以继续思考什么是许多书籍。如今，整个天主的圣言，即那在起初与天主同在的圣言，并不是多言多语，也不是\Quote{言语}；因为圣言是唯一的，由许多的理论组成，每个理论都是整个圣言的一部分。凡是在这唯一的圣言之外、承诺描述和阐述的任何话语，即使它们关于真理，也都不是圣言或理性，它们只是言语或理由。它们不是一元，远非如此；它们不是那内在一致和统一的东西，由于其内在的分裂和冲突，统一已离它们而去，它们成了数字，也许是无限的数字。因此，我们不得不说，凡宣讲与宗教相异之事的人是在说许多话，而宣讲真理之言的人，即使他遍及整个领域且无所遗漏，总在说那唯一的圣言。圣人们并不犯多言多语之罪，因为他们总是持守与那唯一圣言相关的目标。如此看来，所谴责的多言多语更多地是根据所阐述观点的性质，而非根据所发言语的数量。让我们看看是否也能得出同样的结论：所有神圣的书籍是一本书，而外在的则是训道者所说的\Quote{许多书籍}。这证明必须出自\WorkTitle{圣经}，如果能表明并非只有一本书（按通常意义）是关于基督而写的，那将最令人满意。基督甚至被写在\WorkTitle{梅瑟五书}中；每位先知和\WorkTitle{圣咏集}中都提到他，简言之，正如救主自己所说，在全部\WorkTitle{圣经}中。祂将我们指向它们全部，当祂说：\BibleRef{若 5:39}\BibleQuote{你们查考圣经，因为你们以为在其中得到永生；正是这些为我作证。}既然祂指给我们圣经为祂作证，祂不是指向其中一本而排除另一本，而是指向所有谈论祂的经文，那些在\WorkTitle{圣咏集}中被称为\Quote{书卷}的，说：\BibleQuote{在书卷上已为我记载了。}如果有人按字面理解这句话\Quote{在书卷上已为我记载了}，并将其应用于某段特定经文，让他告诉我们他有何理由偏爱一本书胜过另一本。如果有人以为我们自己也做了类似的事，将这句话应用于\WorkTitle{圣咏集}，我们否认；我们坚持在这种情况下，经文本应是\Quote{在这卷书上为我记载了。}但祂把所有书籍称作一个书卷，这样将关于基督的一切记载总结为一，以教育我们。事实上，这书卷被\Person{若望}看见，\BibleRef{默 5:1}\BibleQuote{内外都写着字，用印封严；没有人能展开阅读和揭开封印，唯有\Person{犹大}支派的狮子，\Person{达味}的根，祂有\Person{达味}的钥匙，\BibleRef{默 3:7}祂开了，没有人能关；祂关了，没有人能开。}因为这里所说的书卷是指全部\WorkTitle{圣经}；它外面写着字，因其明显的意义；里面也写着字，因其深奥和属灵的意义。此外，请注意：我们能否从以下事实找到证明，即神圣著作是一本书，而相反的著作是许多书：有一本生命册，那些不配在其中的人被涂抹，如经上所载：\BibleQuote{愿他们从生命册上被涂抹。}而对于那些受审判的人，有许多书卷，正如\Person{达尼尔}所说：\BibleRef{达 7:10}\BibleQuote{审判者就坐，书卷就展开了。}\Person{梅瑟}也为这神圣书籍的统一性作证，他说：\BibleRef{出 32:32}\BibleQuote{如果你赦免这百姓的罪，就赦免；不然，就把我从你所写的册上涂抹吧。}以同样的方式，我读到\Person{依撒意亚}的段落。他的预言并非独有书卷上的话被密封，既不识字的人因不识字不能读，有学识的人也不能读，因为书卷被封住。这适用于所有著作，因为每部著作都需要那封闭它的圣言来开启。\BibleQuote{祂必关闭，没有人能开；}\BibleRef{依 22:22}而当祂开启时，无人能质疑祂所带来的解释。因此，据说祂开启而无人能关闭。我从\Person{厄则克耳}所说的书卷中推断出类似的教训，那书卷上写着哀歌、歌曲和祸灾。因为整部书充满了丧亡者的祸灾、得救者的歌曲和介于两者之间的哀歌。\Person{若望}在谈到他吞下那书卷时，\BibleRef{默 10:9}\BibleQuote{那书卷内外都写着字}，也指全部\WorkTitle{圣经}，一本书，开始咀嚼时极为甘甜，但在启示给每位认识它的人的良心时却变得苦涩。我为此补充一句被\Person{马尔西翁}的门徒误解的宗徒言论，他们因此藐视福音。宗徒说：\BibleRef{罗 2:16}\BibleQuote{按我的福音，在基督耶稣内}；他没有说复数\Quote{福音}，因此他们争论说，既然宗徒只提到单数\Quote{福音}，那就只有一部福音存在。但他们没有看到，正如所有福音作者所写的同一位基督，福音虽由几人执笔，实质上是一部福音。事实上，福音虽由四人书写，但是一部。从这些考虑，我们明白了什么是一本书，什么是许多书，我此刻关心的不是可能写下的数量，而是我所说的话的效果，以免在这方面失足，说出违背真理的话，即使在我的一部著作中，我便会证明自己违犯了诫命，成了撰写\Quote{许多书籍}的人。然而，我看见异端者在这时代以更高智慧为借口攻击天主的圣教会，并推出多卷著作，在其中讲解福音和宗徒著作，我担心如果我保持沉默，不向我们的成员提出有益和正确的道理，这些教师可能会抓住好奇的灵魂，使他们在缺乏健康食粮的情况下，追求被禁止的、实际上不洁和可怕的食物。因此，在我看来，有必要让一位能真实地代表教会道理、并反驳那些所谓知识的贩卖者的人，站出来反对历史虚构，并以真实崇高的福音信息与之对立，在这信息中，所谓的\WorkTitle{旧约}和所谓的\WorkTitle{新约}中的道理一致，清楚而充分。您自己一度感到缺乏优秀代表来支持更好的事业，您对一种与理性冲突并荒谬的信仰感到不耐烦，然后，由于您对主的爱，您曾致力于写作，然而凭借您所拥有的判断力，您后来停止了。我认为这是对拥有言谈和写作能力者的辩护。但我也是在为自己辩护，因为我如今以我所拥有的勇气致力于诠释工作；因为或许我并不具备那种习惯和性情，即那由天主装备为新的盟约服务者所应有的，不是文字的，而是圣神。\par

\SourceRecord{Translated by Allan Menzies.；Ante-Nicene Fathers；Vol. 9.；Edited by Allan Menzies.；Buffalo, NY: Christian Literature Publishing Co.,；1896.；<http://www.newadvent.org/fathers/101505.htm>.}

% structure: p0001 source=h1 target=section reason=page-title

\section{若望福音释义（第六卷）}

% structure: p0002 source=h2 target=subsection reason=relative-heading-rank; base=h2

\subsection{一、在暴力中断后重新开始工作——这场中断曾将作者逐出\Place{亚历山大}。他感谢天主的拯救并祈求指引，再次着手写作。}

当一座房屋要建得尽可能坚固时，建筑工程应在晴朗平静的天气中进行，这样才不会妨碍结构获得所需的坚固性。于是，它变得如此坚固稳定，以致能经受洪水的冲击和河水的冲刷，以及伴随风暴而来的各种力量——这些力量善于发现建筑中腐朽之处，并显示出哪些部分建造得当。尤其应当如此建造那能容纳真理之思辨的房屋，即那在应许或文字中、几乎如同理性的房屋，这样天主才能将他自由的合作赐予如此崇高工程的策划者，此时灵魂平静，享有那超越一切理解的平安，远离一切纷扰，不被任何波涛击打。在我看来，先知精神的仆人和福音讯息的侍奉者们深知这一点；他们使自己配得从那位永远赐给配得者、且说\BibleQuote{我將平安留給你們，我將我的平安賜給你們；我所賜的，不像世人所賜的。}（\BibleRef{若 14:27}）的主那里，领受那隐秘的平安。试看，关于圣殿的记载中，是否也隐含着类似的教导，涉及\Person{达味}和\Person{撒罗满}？达味曾为主作战，并坚定地对抗许多敌人，包括他自己的和以色列的敌人，他想为天主建造一座圣殿。但天主透过\Person{纳堂}阻止了他，纳堂对他说：\BibleQuote{你不可為我建造殿宇，因為你是個流血的人。}（\BibleRef{编上 22:8}）但另一方面，撒罗满在梦中看见了天主，并在梦中领受了智慧，因为那异象的真实性是为那位说\Quote{看，這裏有一位大於撒羅滿的！}的人保留的。那时正值最深切的平安，以至于人人都可以安坐在自己的葡萄树和无花果树下，撒罗满的名字本身就象征着那时代的平安，因为撒罗满意为平安；因此他得以自由建造那位著名的天主圣殿。在\Person{厄斯德拉}时期，当\Quote{真理勝過葡萄酒、敵對的君王和女人}时，天主的圣殿再次被重建。我这样说是向您，尊敬的\Person{盎博罗削}，致歉。是您神圣的鼓励使我有决心以文字建造福音之塔；因此我坐下来计算费用，看是否够完成（\BibleRef{路 14:28}），以免被旁观者嘲笑，因为我立了根基却不能完工。计算的结果，诚然，我并没有完成所需的资源；但我信赖天主——祂以一切智慧和一切知识充满我们（\BibleRef{格前 1:5}）。如果我们努力遵守祂的属灵法律，我们相信祂必充满我们；祂将供应所需，使我们得以继续建造，甚至到达建筑的栏杆。那栏杆正是防止上圣言房屋的人坠落之物；因为只有那些没有栏杆的房屋才会使人坠落，以致建筑本身要为它们的坠落和死亡负责。尽管\Place{亚历山大}风暴带来阻碍，我们仍进行到第五卷，并说了所赐给我们的话，因为耶稣斥责了风和海。我们走出了风暴，被带出了\Place{埃及}，那解救祂子民的天主解救了我们。随后，当敌人以他的新著作极其恶毒地攻击我们——这些著作直接反对福音——并鼓动埃及所有邪恶之风对抗我们时，我觉得理智要求我宁可坚持战斗，拯救我内里更高的部分，以免恶谋成功地将风暴导向淹没我的灵魂，而不是在不合时宜的时候、在我心灵尚未恢复平静之前完成我的工作。事实上，平时为我记录口授的速记员因没有出现而使我工作停顿。但现在，那些射向我的众多火箭已失去锋芒——因为天主熄灭了它们——我的灵魂也渐渐习惯了那为天上圣言而赐予我的安排，并学会了从必要性中忽略敌人的圈套，仿佛极大的平静降临在我身上，我立即继续这项工作。我祈求天主与我同在，并在我灵魂的门廊里以教师的身份说话，使我已开始的\BibleBook{若望}{福音}释义工程得以完成。愿天主俯听我的祈祷，恩赐全书的整体现在得以汇聚，不再有任何中断阻止我遵循圣经的顺序。请放心，我现在非常乐意地开始这第二次起步，进入第六卷，因为我在亚历山大所写的前文，不知何故没有随身带来。我担心自己若不尽早着手，会忽略这项工作，因此认为最好利用现在的时间，立即开始剩余的部分。我不确定先前所写的部分是否会重见天日，也不愿浪费时间等待。序言已足，现在让我们进入经文。\par

% structure: p0004 source=h2 target=subsection reason=relative-heading-rank; base=h2

\subsection{二、旧约中的先知和圣人们如何知晓基督的事。}

\BibleQuote{这就是若翰所作的见证。} \BibleRef{若 1:19} 这是关于洗者若翰为基督作证的第二次记载。第一次以\BibleQuote{这一位就是我论及曾说：在我以后来的那个人}开始，至\BibleQuote{那在父怀里的独生子，他给我们详述了。}为止。赫拉克翁猜想\BibleQuote{从来没有人见过天主}等语并非出于若翰，而是出于那位门徒。但在此他并不正确。他本人承认\BibleQuote{从他的满盈中我们都领受了，而且恩宠上加恩宠；因为法律是藉梅瑟传授的，恩宠和真理却由耶稣基督而来}等语出自若翰。那么，那领受了基督的满盈，并在原有恩宠之外又领受第二个恩宠，且宣告法律是藉梅瑟赐下，而恩宠和真理是由耶稣基督而来的那个人，岂不显然明白，从他由基督满盈所领受的，如何\BibleQuote{从来没有人见过天主}，以及如何\BibleQuote{那在父怀里的独生子}将关于天主的启示传给他和所有领受其满盈的人？他并非在此首次宣告那在父怀里的那一位，仿佛从前从未有过适合接受他告诉宗徒们之事的任何人。难道他没有教导我们，他先于亚巴郎存在，而亚巴郎曾欢喜并高兴见到他的日子？\BibleQuote{从他的满盈中我们都领受了}及\BibleQuote{恩宠上加恩宠}等语表明，正如我们已阐明的那样，先知们也从基督的满盈中领受了他们的恩赐，并获得了第二个恩宠以替代原有的；因为他们也被圣神引导，从他们在预像中所领受的初步教导，进步到真理的视见。因此，并非所有先知，而是他们中的许多人，\BibleRef{玛 13:17} 渴望见到宗徒们所见的事物。因为如果先知之间有差别，那更完美、更出众的，并不渴望见到宗徒们所见的，而是实际看见了它们；而那些未能完全达到圣言高度的人，则满心渴望宗徒们藉基督所认识的事物。\Quote{看见}一词我们并未取肉体意义，\Quote{听见}我们则理解为属灵的交流；惟有具备耳朵的人才能预备好聆听耶稣的话——这并不经常发生。此外，在耶稣肉身来临之前的圣人们在领会神性奥秘方面享有优于多数信徒的优势，因为天主的圣言在成为血肉之前已是他们的导师，因为祂常在工作，效法祂的父，如祂所说：\BibleQuote{我父到现在一直工作。}关于此点，我们可以引述祂对否认复活道理撒杜塞人所说的话：\BibleQuote{你们没有读过}祂说，\BibleRef{谷 12:20} \BibleQuote{在荆棘篇中天主向梅瑟说的：\InnerQuote{我是亚巴郎的天主，依撒格的天主，雅各伯的天主。}他不是死人的，而是活人的天主。}那么，如果天主不以被称为这些人的天主为耻，且他们被基督算在活人之列，而所有信徒都是亚巴郎的子孙，\BibleRef{罗 4:11} 因为所有外邦人都与有信德的亚巴郎一同蒙福，亚巴郎被天主立为外邦人的父亲，我们还能犹豫，不承认那些活在灵魂中的人结识了活人的学问，并受教于那在晨星之前出生的基督，在祂成为血肉之前？正因如此，他们活着，因为他们有分于那说\BibleQuote{我就是生命}的那一位，并作为如此伟大许诺的继承者，不仅看见了天使，更是在基督内看见了天主。因为他们也许看见了无形天主的肖像，因为看见子就是看见了父，因此他们被记载为认识天主，并相称地听见了天主的话语，从而看见了天主并听见了祂。现在，我认为那些完全且真正是亚巴郎的子孙的人，就是他的行动、属灵理解的智慧以及启示给他的知识的儿子。他所知道的和他所做的，在他子孙身上重现，正如圣经教导那些有耳可听的人：\BibleRef{若 8:39} \BibleQuote{如果你们是亚巴郎的子女，你们就该作亚巴郎所作的事。}如果有一句真箴言\BibleRef{箴 16:23} 说\BibleQuote{智慧人的心指教他的口，又使他的嘴增长学问}，那么我们必须立即摒弃某些关于先知的意见，认为他们不是智慧人，不明白自己口中所出的。我们应当相信关于先知的正面而真实的事：他们是有智慧的人，确实明白自己口中所出的，且他们嘴上带着学问。显然，梅瑟心中明白法律的真义，以及他书中所记载故事的高深解释。若苏厄也明白消灭二十九位君王之后分地的意义，并且比我们更能看见他事迹之实体的预像。显然，依撒意亚看见了那坐在宝座上的那位、两位色辣芬、牠们掩面遮脚、翅膀、祭坛和火剪的奥秘。厄则克耳也明白革鲁宾及其行动、牠们上面的穹苍和坐在宝座上那位的真实意义，还有什么比这一切更高超、更荣美？不必详加例举；我旨在确立的观点已足够清楚，即：在较早世代达至完美的人，对于基督所启示的事，其知识并不少于宗徒们，因为同一位导师与他们同在，就是那启示宗徒们虔敬的难以言传奥秘的导师。我再补充几点，然后交由读者判断，并就本题形成自己的看法。保禄在《罗马书》中说：\BibleRef{罗 16:25} \BibleQuote{愿光荣归于那能坚固你们，合乎我所传的福音，和所宣讲的耶稣基督，并照那从永远以来隐蔽而不宣的奥秘的启示，借着诸先知的经书，且借着我主耶稣基督的显现，现已显明出来。}因为如果那自古隐藏的奥秘现正通过先知著作向宗徒们显明，而先知们是智慧人，明白自己口中所出的，那么先知们早已知道那显明给宗徒们的事。但许多人并未得到启示，如保禄所说：\BibleRef{弗 3:5} \BibleQuote{这奥秘在以往的其他世代，没有告诉过任何人，如同现在一样，藉着圣神启示给祂的圣宗徒和先知们：外邦人在基督耶稣内，借着福音，同为承继人，同为一体，同享恩许。}在此，不同意我们所持观点的人可能提出异议；定义\Quote{启示}一词的涵义变得重要。它有两种可能意义：其一是指所讨论之事被理解；其二，如果论及预言，是指它得以成就。现在，外邦人将同为承继人、同为一体、同享在基督内的恩许这一事实，先知们得知至此程度：他们知道外邦人将成为同为承继人、同为一体、并在基督内同享恩许。此事何时发生、为何原因、指哪些外邦人、以及他们如何虽与盟约远离、与恩许无分，却仍成为一身并共享福乐——这一切先知都已知道，因他们得了启示。但预言之事属于未来，对那些知道它们但未见其实现的人，并不像对亲眼目睹其实现的人那样显明。这正是宗徒们的情形。所以我认为，他们知道这些事并不比父辈和先知们更多；然而，真正可以说：\Quote{在其他世代没有启示的，现在启示给宗徒和先知了，即外邦人是同为承继人、同为一体、同享基督恩许。}因为除了知道这些奥秘，他们还在已成事实中看见了运作的大能。经文\BibleQuote{许多先知和义人愿意看你们所看见的，却没有看见；愿意听你们所听见的，却没有听见}也可同样解释。他们也曾渴望看到天主子降生成人、祂降世以完成拯救众人的苦难计划的奥秘，在实际施行中呈现。这可用另一例子说明。假设有一位宗徒明白了\BibleQuote{不可言说的话，是人不可说的} \BibleRef{格后 12:4}，但并未亲眼看到耶稣对信徒荣耀的肉体显现——尽管这应许是他渴望看见的——而另一个人不仅未留意并看见那位宗徒所注意和看见的，而且对属神希望的认识薄弱得多，却亲身经历了我们救主第二次来临——如同上述平行例子中宗徒所渴望但未看到的。如果我们说两者都看见了宗徒（或宗徒们）渴望看见的事物，却因此不被视为比宗徒更有智慧或更有福，我们不会错失真理。同样，宗徒们也不被视为比父辈、比梅瑟和先知们更有智慧，即比那些因德行而堪当获得显现、神圣启示和奥秘启示的人们更有智慧。\par

% structure: p0006 source=h2 target=subsection reason=relative-heading-rank; base=h2

\subsection{三、\Quote{恩宠和真理是藉着耶稣基督而来的。} 这些话属于洗者若翰，不属于福音作者。洗者若翰藉这些话所作的见证。}

我们在这些讨论上耽搁得相当久了，但这是有原因的。有许多人，假借光荣基督来临的名义，宣称宗徒比先祖或先知更有智慧；这些教师中，有些人发明了一个后来的更大的神，有些人虽不敢走那么远，但按照他们自己的说法，因教义上的困难而取消了天主藉着基督——万物是藉着祂造成的——所赐给先祖和先知的全部恩赐。如果万物都是藉着祂造成的，那么显然，赐给先祖和先知的那些辉煌启示也必是藉着祂造成的，这些启示对他们而言成了宗教神圣奥迹的象征。基督的真战士必须时刻准备为真理而战，并且只要力所能及，绝不让错误的信念潜入。我们因此不可忽视此事。可能有人说，若翰对基督较早的见证见于这话：\Quote{那在我以后来的，成了在我以前的，因祂原先在我以前}，而\Quote{从祂的圆满中我们都领受了，而且恩宠上加恩宠}这句话是出于门徒若望之口。现在，我们必须指出这种解释是牵强的，是对上下文的一种强解；假定洗者若翰的讲话如此突然地、似乎不合时宜地被门徒的话打断，实在是一种大胆的做法。任何稍微能判断上下文的人都很清楚，在\Quote{这就是我曾说过的那一位：那在我以后来的，成了在我以前的，因祂原先在我以前}之后，讲话是连续不断的。洗者若翰提出一个证据，证明耶稣在祂以前存在，因为祂是在祂以前，祂是受造物的首生者；他说：\Quote{从祂的圆满中我们都领受了。}这就是为什么他说：\Quote{祂成了在我以前的，因祂原先在我以前。}这就是我何以知道祂是首要的，在圣父那里有更高的尊荣，因为从祂的圆满中，我和在我以前的先知都领受了更神圣的先知恩宠，代替我们先前因被拣选而从祂手中领受的恩宠。这就是为什么我说：\Quote{祂成了在我以前的，因祂原先在我以前}，因为我们知道我们从祂的圆满中领受了什么；即法律是藉着梅瑟赐下的，但不是由梅瑟赐下的，而恩宠和真理不仅是藉着耶稣基督赐下的，而且是藉着祂得以存在的。因为祂的天主和圣父曾藉着梅瑟赐下法律，又藉着耶稣基督造了恩宠和真理，就是那临到人的恩宠和真理。如果我们对\Quote{恩宠和真理是藉着耶稣基督而来的}这句话给予合理的解释，我们就不会因它与另一句话\Quote{我是道路、真理和生命}可能有的不一致而感到惊慌。如果耶稣说\Quote{我是真理}，那么真理怎能藉着耶稣基督而来呢？因为没有人能藉着自己存在。我们必须认识到，这个真理本身——本质的真理，可以这么说，是在有理性的灵魂中存在的真理的原型，那从它仿佛印出影像给关心真理的人的真理——不是藉着耶稣基督造的，也不是藉着任何别的东西造的，而是由天主造的——正如那在起初与圣父同在的圣言不是藉着任何人造的；正如天主在造化之始所造的道路的智慧不是藉着任何人造的，同样，真理也不是藉着任何人造的。但是，那与人同在的真理是藉着耶稣基督来的，正如在保禄和宗徒们身上的真理是藉着耶稣基督来的。这并不奇怪，既然真理是唯一的，许多真理从那个唯一的真理涌流出来。先知达味无疑知道许多真理，正如他说的：\Quote{上主察究真理}；因为真理之父所察究的不是那唯一的真理，而是许多真理，藉着这些真理，拥有它们的人得以得救。正如唯一的真理和许多真理，同样也有唯一的正义和许多正义。因为本质的正义就是基督，\Quote{祂由天主成了我们的智慧、正义、圣化和救赎}。但是从那正义形成了每个人里面的正义，以致在得救的人中有许多正义，因此经上写道：\Quote{因为上主是正义的，祂喜爱正义。}这是精确抄本以及除七十贤士译本以外的其他译本和希伯来文中的读法。试想，基督被说成是统一体的其他事物，是否也可以同样地被繁殖并用作复数呢？例如，基督是我们的生命，正如救主亲自说的：\BibleRef{若 14:6} \Quote{我是道路、真理和生命。} 宗徒也说：\BibleRef{哥 3:4} \Quote{当基督——我们的生命显现时，你们也要与祂一同在光荣中显现。} 在圣咏中我们又读到：\Quote{祢的仁慈比生命更宝贵}；因为由于基督是在每人里面的生命，所以有许多生命。这也许也是理解\BibleRef{格后 13:3} \Quote{你们既然寻求基督在我内说话的凭证}的钥匙。因为基督在每位圣人身上被发现，所以从唯一的基督产生了许多基督，他们效法祂，并照着祂——天主的肖像——被塑造；因此天主藉着先知说：\Quote{不要触碰我的受傅者。}这样，我们就顺便解释了似乎在我们注释中被忽略的一段话，即：\Quote{恩宠和真理是藉着耶稣基督而来的}；并且我们也表明这些话属于洗者若翰，是他对天主之子见证的一部分。\par

% structure: p0008 source=h2 target=subsection reason=relative-heading-rank; base=h2

\subsection{四、若翰否认自己是厄里亚或\Quote{那位}先知，然而他确实是\Quote{一位}先知。}

现在让我们思考若翰的第二次见证。从耶路撒冷来的犹太人，\BibleRef{若 1:19} 与洗者若翰同族，因为他也属于司祭家族，派遣了司祭和肋未人去问若翰他是谁。当他说\Quote{我不是基督}时，他做了对真理的宣认。这句话并非如人所想的是一种否定；因为为了基督的荣耀而说自己不是基督，这并不是否定。从耶路撒冷派来的司祭和肋未人，首先在那里听说他不是所期待的默西亚，便问起他们所期待的第二个大人物，即厄里亚，问他若翰是否就是厄里亚；他说他不是厄里亚，并用他的\Quote{我不是}做了第二次对真理的宣认。又因为在以色列曾出现过许多先知，并且根据梅瑟的预言，特别有一位将被期待，梅瑟说：\BibleRef{申 18:15} \BibleQuote{主你的天主要从你们弟兄中，为你兴起一位像我一样的先知，你们应听从祂；将来谁若不听从那位先知，必从民间消灭。}因此，他们问了第三个问题，不是问他是不是一位先知，而是问他是不是那位先知。他们并没有把这个名称用于基督，而是认为那位先知是基督之外的另一个人物。但若翰却知道祂所预备道路的那位既是基督，也是预言中的那位先知，所以他回答\Quote{不是}；而如果他们问的是他是不是一位先知，他会回答\Quote{是}；\BibleRef{若 1:25} 因为他并非不知道自己是先知。在所有这些回答中，若翰对基督的第二次见证尚未完成；他还要给问话的人一个他们可以带回去给派遣者的回答，并用依撒意亚预言的话来表明自己，那话说：\BibleQuote{在旷野中呼喊者的声音：你们要预备主的道路。}\par

% structure: p0010 source=h2 target=subsection reason=relative-heading-rank; base=h2

\subsection{五、有两个代表团去见洗者若翰；他们的不同性质。}

这里有一个问题值得探讨：第二次见证是否结束，以及是否还有第三次见证，即针对那些从法利塞人被派来的人所说的话。他们想知道，如果他既不是基督，也不是厄里亚，又不是那位先知，那他为什么施洗。他说：\BibleQuote{我用水施洗；但有一位站在你们中间，你们却不认识，祂就是在我以后来的，我连解祂的鞋带也不配。}这是第三次见证，还是他们应向法利塞人报告的第二次见证的一部分？据我所能从字句推断的，我认为给法利塞人使者的话是第三次见证。然而，值得注意的是，第一次见证宣告了救主的天主性；第二次则排除了那些怀疑若翰是否就是基督的人的疑虑；第三次则宣告一位已经临在人间、虽肉眼看不见、且其来临已不再是未来的那一位。在继续后面的见证——其中他指明基督并为祂作证——之前，让我们逐字考察第二次和第三次见证，并首先注意到有两个代表团去见洗者若翰：一个\BibleQuote{从耶路撒冷}来自犹太人，他们派遣司祭和肋未人去问他\BibleQuote{你是谁？}；第二个由法利塞人派遣，他们对给司祭和肋未人的回答感到疑惑。请注意，第一批使者所说的话与司祭和肋未人的身份相符，表现出温和与学习的意愿。他们说：\BibleQuote{你是谁？}\BibleQuote{那么，你是厄里亚吗？}\BibleQuote{你是那位先知吗？}然后，\BibleQuote{你是谁，好让我们给派遣我们来的人一个答复？你说你自己是什么？}这些人的询问毫无粗暴或傲慢之处；一切都符合真正而谨慎的天主仆人应有的态度，他们对所得到的回答也没有提出任何难题。相反，那些从法利塞人派来的人，用傲慢和冷漠的话攻击洗者若翰：\BibleQuote{那么，你若不是基督，不是厄里亚，也不是那位先知，为什么你施洗呢？}这个使命几乎不是为求了解，像之前司祭和肋未人那样，而是为了阻止洗者若翰施洗，仿佛认为除了基督、厄里亚和那位先知以外，没有人有权施洗。凡愿意理解圣经的学生，必须始终以这种谨慎的方式进行；他必须针对每一段话问：说话者是谁，以及在什么场合说的。只有这样，我们才能辨别话语如何与说话者的性格和谐一致，正如在整部圣经中一样。\par

% structure: p0012 source=h2 target=subsection reason=relative-heading-rank; base=h2

\subsection{六、与洗者若翰关于默西亚的讨论。}

那时，犹太人从耶路撒冷派遣司铎和肋未人到他那里问他说：你是谁？他承认并且不否认，他承认说：我不是基督。\BibleRef{若 1:19} 犹太人应当派遣什么样的使者到若翰那里，又应当从何处派遣？难道不应当是那些被认为在同胞中蒙天主拣选而站立的人，并且应从那从全地中拣选出来的地方，即虽有全境称为善，但却是耶路撒冷，那里有天主的圣殿？这样，他们怀着如此尊荣询问若翰。至于基督，从未记载犹太人对祂做过类似的事；但犹太人对若翰所做的，若翰也对基督做了，他派遣自己的门徒去问祂（\BibleRef{玛 11:3}）：\Quote{那要来的是你吗，还是我们等另一位？} 若翰向那些派来的人承认，并且不否认，随后他宣告：\Quote{我是旷野中呼喊者的声音。} 但基督，因有比洗者若翰更大的见证，以言以行作答，说：\Quote{你们去，把你们所听见、所看见的报告给若翰：瞎子看见，瘸子行走，癞病人洁净，聋子听见，贫穷人得听福音。} 关于这段经文，若蒙天主许可，我将在适当的地方加以阐述。然而在此处，我们完全可以合理地提出这个问题：为何若翰对向他提出的问题作此回答？司铎和肋未人并没有问他：\Quote{你是基督吗？} 而是问：\Quote{你是谁？} 而洗者对这问题的答复本应当是：\Quote{我是旷野中呼喊者的声音。} 对问题\Quote{你是基督吗？} 的恰当回答是：\Quote{我不是基督；} 而对问题\Quote{你是谁？} 的回答则是：\Quote{旷野中呼喊者的声音。} 对此我们可以说，他很可能在司铎和肋未人的问题中看出一种谨慎的尊敬，这引导他们暗示他们心中的想法，即那施洗的可能是基督，但又阻止他们公开说出，因为那可能是冒失的。因此，他非常自然地首先解除他们可能对他产生的任何错误印象，并公开声明事实真相：\Quote{我不是基督。} 他们的第二个问题，以及第三个问题，表明他们对他怀有某种猜测。他们猜想他可能是他们所期待的第二位，即厄里亚，他在他们心中居于仅次于基督的地位；因此，当若翰回答\Quote{我不是基督} 后，他们问：\Quote{那么你是谁？你是厄里亚吗？} 他说：\Quote{我不是。} 他们第三次想知道他是否是那位先知，当他回答\Quote{不是} 后，他们不再有任何名字来称呼他们所期待的人物，便说：\Quote{你到底是谁？好让我们给那派遣我们来的人一个答复。你说你自己是什么？} 他们的意思是：\Quote{你说你不是以色列所期望和盼望的那些人物中的任何一个，那么你是谁，竟如此施洗？我们不知道；因此告诉我们，好让我们向那些派遣我们来的人报告，他们的使命是弄清这一点。} 我们补充一点，这与上下文有关：民众因认为基督来临的时期已近而受感动。从耶稣诞生之前不久直至传道开始的那些年间，这个时期可以说是迫在眉睫。因此，很可能正如经师和法学家们从圣经中推算出时期并期待那将要来临的一位，这想法被特乌达所接受，他自称默西亚并聚集了相当多的群众，随后在户口登记的日子又有著名的加里肋亚人犹达。\BibleRef{宗 5:36} 这样，默西亚的来临被更热切地期待和讨论，犹太人从耶路撒冷派遣司铎和肋未人到若翰那里，问他：\Quote{你是谁？} 并了解他是否自称是基督，这是很自然的。\par

% structure: p0014 source=h2 target=subsection reason=relative-heading-rank; base=h2

\subsection{七、论若翰的诞生，及其与厄里亚被认为的同一性；论\OriginalTerm{轮回}{Transcorporation}的学说。}

\Quote{他们便问他说：\InnerQuote{那么，你是厄里亚吗？}他回答：\InnerQuote{我不是。}\BibleRef{若 1:21}} 在这件事上，无人不记得耶稣论及若翰所说的话：\Quote{倘若你们愿意接受，他就是那要来的厄里亚。\BibleRef{玛 11:14}} 那么，若翰为何对那些问他\Quote{你是厄里亚吗？}的人说\Quote{我不是}？同时，按照玛拉基亚的话，若翰又真是那要来的厄里亚呢？——\Quote{看，在主伟大而可畏的日子来临以前，我必派遣厄里亚先知到你们这里来；他必使父亲的心转向儿子，人的心转向他的近人，免得我来临击打此地。\BibleRef{拉 4:5-6}} 此外，主的天使向在香坛右边出现的匝加利亚所说的话，也与玛拉基亚的预言大致相同：\Quote{你的妻子依撒伯尔要给你生一个儿子，你要给他起名叫若翰。\BibleRef{路 1:13}} 稍后又言：\Quote{他要以厄里亚的精神和能力，在他前面走，使父亲的心转向儿女，使悖逆的人转向义人的明达，并为主准备一个相称的子民。\BibleRef{路 1:17}} 关于第一点，有人可能说若翰并不知道自己是厄里亚。那些在本段中发现支持他们转生教义的人会如此解释，仿佛灵魂披上新的身体，却不完全记得前世的事。这些思想家还会指出，有些犹太人赞同此教义，他们论及救主时，仿佛他是一位从前的先知，不是从坟墓中复活，而是从出生时就是先知。他的母亲玛利亚是众所周知的，木匠若瑟被认为是他的父亲，人们很容易以为他是从死者中复活的古先知之一。同一位人会引用创世纪中的经文：\Quote{我要消灭全部的复活}，从而让那些致力于从圣经中寻找虚妄可能之答案的人，在此教义上陷入极大困境。然而，另一位教会人士，他拒绝转生教义为虚假，也不承认若翰的灵魂曾经是厄里亚，可以引用上述天使的话，指出若翰出生时提及的并非厄里亚的灵魂，而是厄里亚的精神和能力。经上说：\Quote{他要以厄里亚的精神和能力，在他前面走，使父亲的心转向儿女。} 现在，从成千上万的经文可以证明，精神与灵魂是不同的，所谓的\Quote{能力}也与灵魂和精神不同。我目前无法详述这些点；这部著作不能过度膨胀。为确立能力不同于精神，引用以下经文便已足够：\Quote{圣神要临于你，至高者的能力要庇荫你。\BibleRef{路 1:35}} 至于先知的精神，它们由天主赐予他们，并在某种程度上被称为他们的财产（仆人），如\Quote{先知的精神服从先知}（\BibleRef{格前 14:32}），以及\Quote{厄里亚的精神安息在厄里叟身上}（\BibleRef{列下 2:15}）。因此，据说假设若翰\Quote{以厄里亚的精神和能力}使父亲的心转向儿女，并且因这精神，他被称为\Quote{要来的厄里亚}，这并无荒谬之处。为加强这一观点，可以论证：如果宇宙的天主与他的圣人如此认同，以至于被称为亚巴郎的天主、依撒格的天主和雅各伯的天主，那么圣神更能与先知如此认同，以致被称为他们的精神，以致当提到精神时，可能指的是厄里亚的精神或依撒意亚的精神。我们的教会人士继续论述，他说那些以为耶稣是从死者中复起的先知之一的人，可能是被上述教义所误导，部分是因为他们以为他是一位先知；至于他是一位先知这种误解，这些人可能是由于不知道耶稣的假定父亲和实际母亲，而以为他是从坟墓中复活的人。至于创世纪中关于复活的那句经文，教会人士会用一句意义相反的经文回应：\Quote{天主给我另立了一个儿子替代亚伯尔，因为加音杀了他。}（\BibleRef{创 4:25}）这表明复活发生在创世纪中。至于提出的第一个困难，我们的教会人士会这样回应转生信徒的观点：若翰毫无疑问在某种意义上是那要来的厄里亚，正如他已表明的；他之所以用\Quote{我不是}来回答司铎和肋未人的询问，是因为他察觉到他们提问的意图。因为司铎和肋未人向若翰提出的问题，并不是要查明两者是否有同一精神，而是要问若翰是否就是那个被接去的厄里亚，如今按照犹太人的期待出现，却没有出生（因为使者们或许不知道若翰的出生）；对于这样的询问，他自然回答\Quote{我不是}；因为被称为若翰的人并不是那被接去的厄里亚，也没有为这次出现改变身体。我们第一位学者，我们已看到他从本段得出转生观点，他会继续仔细考察经文，并反对他的对手：如果若翰是如司铎匝加利亚这样的人的儿子，并且他在父母年迈时出生，违反一切人的常理，那么耶路撒冷这么多犹太人不太可能如此不了解他，或者所派的司铎和肋未人不知道他出生的细节。路加不是宣称\Quote{凡住在他们四周的人都害怕起来}（显然是指匝加利亚和依撒伯尔四周），并且\Quote{这一切事就传遍了全犹太山地}（\BibleRef{路 1:65}）吗？如果若翰生于匝加利亚是众所周知的事，而耶路撒冷的犹太人却派遣司铎和肋未人去问\Quote{你是厄里亚吗？}那么显然，说这话时他们假设转生教义是真实的，并且这是他们国家的流行教义，而非秘密教导。因此若翰说\Quote{我不是厄里亚}，因为他不知道自己的前世生活。因此，这些思想家的观点并非可轻蔑。然而，我们的教会人士可以回击，并质问：一位先知，蒙圣神光照，被依撒意亚预言，出生前被如此伟大的天使预告，领受了基督的圆满，分享如此大的恩宠，知道真理是藉耶稣基督而来，并教导了关于天主和那位在父怀里的独生子的深刻之事，这样的人竟然说谎，或由于对自己是什么的无知而犹豫不决，这合适吗？因为对于模糊之事，他应当拒绝承认，既不肯定也不否定摆在他面前的命题。如果所讨论的教义确实广泛流传，若翰岂不应该对此犹豫不决，以免他的灵魂实际上曾在厄里亚里面吗？在此，我们的教会人士会援引历史，并吩咐他的对手去询问希伯来秘密教义的专家，看看他们是否真的持有这种信念。因为如果证明他们没有，那么基于该假设的论证便显得毫无根据。然而，我们的教会人士仍可自由求助于先前给出的解决方案，并坚持注意提问的方式。因为正如我所表明的，如果派遣者知道若翰是匝加利亚和依撒伯尔的孩子，而使者们身为司铎家族，更不可能不知道他们的亲戚匝加利亚如此非凡地得到儿子，那么他们的问题\Quote{你是厄里亚吗？}可能有什么含义？难道他们没有读过厄里亚被接升天，并期待他显现吗？既然他们期待厄里亚在末世基督之前来，基督随后而来，或许他们的问题更不是字面意义，而是比喻意义：你是那位在末世基督之前预告圣言的人吗？对此，他非常恰当地回答\Quote{我不是}。然而，对手试图证明司铎们不可能不知道若翰的出生如此非凡，因为\Quote{这一切事在犹太山地已广为流传}；教会人士必须回应这一点。他通过展示类似的错误也广泛存在于救主本身周围来回应；因为\Quote{有人说他是若翰洗者，有的说是厄里亚，有的说是耶肋米亚，或先知中的一位}（\BibleRef{玛 16:13}）。当门徒在斐理伯的凯撒勒雅地区时，主询问他们此事，他们如此告诉主。黑落德也说：\Quote{若翰，我斩了他，他复活了}（\BibleRef{谷 6:16}）；所以他似乎不知道福音中关于基督所说的话：\Quote{这不就是木匠的儿子吗？他的母亲不是叫玛利亚吗？他的弟兄不是雅各伯、若瑟、西满和犹达吗？他的姊妹不是都在我们这里吗？}（\BibleRef{玛 13:55-56}）因此，就救主而言，虽然许多人知道他是玛利亚所生，但其他人却对他有误解；同样，就若翰而言，虽然一些人知道他是匝加利亚所生，但其他人却怀疑所期待的厄里亚是否在他身上显现，这也就不足为奇了。关于若翰是否厄里亚的疑问，并不比关于救主是否若翰的疑问更多。两者之间，关于厄里亚外貌的问题可以从经文的言辞解决，尽管不是通过实际观察，因为经上记载：\Quote{他身穿毛衣，腰束皮带}（\BibleRef{列下 1:8}）。相反，若翰的外貌是众所周知的，且不像耶稣；然而仍有人猜测若翰已从死者中复活，并取了耶稣的名字。至于名字的改变，这使人想起奥秘，我不知道希伯来人如何会说关于丕乃哈斯、厄肋阿匝尔之子的事，他显然延长生命到许多民长时代，正如我们在民长纪中读到（\BibleRef{民 20:28}），并说出我现在提到的事。他们说他是厄里亚，因为（在户籍纪中，\BibleRef{户 25:12}）他因和平的盟约被许诺永生，因为他以神圣的嫉妒而嫉妒，在愤怒的激情中刺穿了米德杨女人和以色列人，并被称为平息了天主的忿怒，正如经上所写：\Quote{丕乃哈斯、厄肋阿匝尔之子、亚郎的孙子，使我的忿怒转离以色列子民，因为他在他们中以我的嫉妒而嫉妒。} 因此，如果那些认为丕乃哈斯和厄里亚是同一人的人（无论他们对此判断是否正确，因为这不是当前问题）也认为若翰和耶稣是同一人，那就不足为奇了。因此，他们怀疑这一点，并想知道若翰和厄里亚是否同一。在另一个时间，这一点肯定需要仔细调查，论证必须仔细权衡灵魂的本质、她的构成原则以及她进入这个尘世身体的方式。我们还须调查每个灵魂生命的分配、她离开此世的方式，以及她是否可能再次进入身体开始第二生命，这是否发生在同一时期、按照相同安排，还是不同；并且她进入同一个身体还是不同的身体，如果相同，主体是否保持不变而性质改变，或者主体和性质都保持不变，以及灵魂是否将永远使用同一身体还是改变它。除了这些问题，还有必要问什么是转生，它与化身有何不同，以及若坚持转生是否必然坚持世界永恒。这些学者的观点也必须考虑，他们认为根据圣经，灵魂与身体一起被播种，这种观点的后果也必须审视。事实上，灵魂的主题是宽广的，难以解开，必须从圣经分散的表述中仔细拣选。因此，需要单独处理。我们关于厄里亚和若翰的问题所作的简要考量现在可以了；我们继续福音接下来的内容。\par

% structure: p0016 source=h2 target=subsection reason=relative-heading-rank; base=h2

\subsection{八、若翰是先知，但不是那位先知。}

\BibleQuote{你是那位先知吗？他回答说：不。}（\BibleRef{若 1:21}）如果法律和先知到若翰为止（\BibleRef{路 16:16}），我们怎能说若翰仅是一位先知呢？他的父亲匝加利亚充满圣神，预言说（\BibleRef{路 1:76}）：\BibleQuote{你，孩子，要称为至高者的先知，因为你要在主前面预备他的道路。}（有人或许可以通过强调\Quote{称为}这个词来避开这段经文：说他将要被称为先知，而不是说他是先知。）更有分量的是，救主对那些认为若翰是先知的人说（\BibleRef{玛 11:9}）：\BibleQuote{你们出去看什么？一位先知吗？是的，我告诉你们，而且他比先知还大。}\Quote{是的，我告诉你们}这些话明确肯定了若翰是先知，而后再没有被否认。那么，如果救主说他不仅是先知，而且是\InnerQuote{比先知还大}，为什么当司祭和肋未人前来问他：\BibleQuote{你是那位先知吗？}他回答\Quote{不}呢？对此我们必须指出，说\Quote{你是那位先知吗？}和\Quote{你是一位先知吗？}不是同一回事。这两种说法的区别我们已经看到，当我们追问\Quote{那位天主}和\Quote{天主}、以及\Quote{那位圣言}和\Quote{圣言}之间的区别时已经注意到了。现在申命纪上写着：\BibleQuote{上主你的天主要从你中间为你兴起一位像我一样的先知，你们要听从他；凡不听从那先知的，必将从民中铲除。}所以，人们期待一位特定的先知，与梅瑟相似，在天主和子民之间做中保，从天主那里接受新盟约，并赐给那些接受他教导的人；而对于其他每一位先知，以色列子民都认出他不是梅瑟所说的那位。因此，正如他们怀疑若翰是不是基督（\BibleRef{路 3:15}）一样，他们也怀疑他能不能是那位先知。毫不奇怪，那些怀疑若翰是不是基督的人，并不明白基督和那位先知是同一位；他们对若翰的怀疑必然意味着他们在这个问题上不清楚。然而，\Quote{那位先知}和\Quote{一位先知}之间的区别，大多数研究者都没有注意到；赫拉克利翁就是这样，他原话这样说：\Quote{因此，若翰承认他不是基督，甚至不是一位先知，也不是厄里亚。}如果他这样解释我们面前的这些话，他就应该查考各处经文，看他若翰说自己不是先知也不是厄里亚时，说的是否真实。然而，他并不留意这些经文，在他其余的注释中，他也没有查考就略过了这些点。随后，他的评论——我们马上要说到——也非常简略，不能证明他仔细研究过。\par

% structure: p0018 source=h2 target=subsection reason=relative-heading-rank; base=h2

\subsection{九、若 1:22}

\BibleQuote{于是他们对他说：\InnerQuote{你是谁？好让我们给那些派遣我们的人一个答复。关于你自己，你说什么？}}这些使者的话大意如下：我们曾猜测你是什么人，来查明是否如此，但现在我们知道你不是那一位。因此，我们还需要听你对自己的陈述，好让我们把你的回答报告给那些派遣我们的人。\par

% structure: p0020 source=h2 target=subsection reason=relative-heading-rank; base=h2

\subsection{十、若翰洗者作为呼声。}

\BibleQuote{他说：我是旷野中呼喊者的声音：修直主的道路，正如依撒意亚先知所说的。}正如那特为天主子、就是\OriginalTerm{圣言}{Logos}的那一位，仍使用圣言（理性）——因为祂在起初就是圣言，并与天主同在，是天主的圣言——同样，那圣言的仆人\Person{若翰}，如果我们按经文所说的来理解，他无非是一个声音，却用他的声音来指向圣言。他于是这样理解依撒意亚先知论及他的预言，说自己是一个声音，不是在旷野中呼喊，而是\Quote{旷野中呼喊者的}，即那位站着呼喊者的声音，\BibleRef{若 7:37}\BibleQuote{谁若口渴，让他到我这里来喝吧。}那一位也曾说：\BibleRef{路 3:4}\BibleQuote{预备主的道路，修直祂的路径。一切山谷将被填满，一切山岳丘陵将被削平；弯曲的将成为正直。}因为正如我们在\WorkTitle{出谷纪}中所读到的，天主对\Person{梅瑟}说：\Quote{看，我已立你作为\Person{法郎}的天主，你的兄弟\Person{亚郎}将是你的先知；}同样我们应当理解——这些情形至少是类似的，即使不完全相同——关于在起初就与天主同在的圣言，以及\Person{若翰}。因为若翰的声音指向那圣言并彰显它。因此，\Person{匝加利亚}对天使说\BibleRef{路 1:18}\BibleQuote{我凭什么知道这事呢？因为我已经老了，我的妻子也上了年纪。}时，落在其上的惩罚是非常恰当的。因为他对于这声音的诞生缺乏信德，他自己被剥夺了声音，正如\Person{加俾额尔}天使对他说：\Quote{看，你要成为哑巴，不能说话，直到这些事成就的日子，因为你没有相信我的话，这些话将在适当的时候应验。}后来当他\Quote{要了一块写字板，写上：他的名字是若翰；众人都惊讶}，他恢复了声音；因为\Quote{他的口立刻开了，他的舌头也松了，就说话赞美天主。}我们上面讨论了如何理解圣言是天主子，并探讨了与此相关的想法；在此处我们也可以看到类似的想法序列。若翰来是为作证；他是从天主那里被派来为光作证的，为使众人借着他而信；那么，我们要明白，他就是那声音，唯独它配得上宣布圣言。如果我们想起我们在解释经文时引用的句子，就会正确理解这一点：\Quote{为使他人都借着他而信}，以及\Quote{这人就是经上所记载的：看，我派遣我的使者在你面前，他要在你前面预备你的道路。}\BibleRef{玛 11:10}说他不是旷野中说话者的声音，而是旷野中呼喊者的声音，这也是贴切的。那呼喊\Quote{预备主的道路}的人也在说它；但他可以不喊而说。但他呼喊并大声宣布，使那些远离说话者的人也能听见，甚至聋子也能明白这消息的伟大，因为它用大声宣布；他因此帮助那些离开天主的人，以及那些丧失听觉敏锐的人。这也是\Quote{耶稣站着呼喊说：谁若口渴，让他到我这里来喝吧。}的原因。因此，\BibleRef{若 1:15}\Quote{若翰为祂作证，呼喊说：}因此天主也命令\Person{依撒意亚}呼喊，\Quote{有声音说：你呼喊。我说：我呼喊什么？}我们在祈祷时使用的肉体声音不必宏大或惊人；即使我们不发出大声或呼喊，天主仍会垂听我们。祂对\Person{梅瑟}说：\BibleRef{出 14:15}\Quote{你为什么向我呼喊？}当时梅瑟根本没有出声喊叫。在\WorkTitle{出谷纪}中没有记载他这样做了；但梅瑟在祈祷中曾用那唯独天主听见的声音向天主大声呼求。因此\Person{达味}也说：\Quote{我用我的声音向主呼求，祂就垂听了我。}在旷野中呼喊的人需要有一个声音，使那被剥夺天主、被真理遗弃的灵魂——还有什么比被天主和一切德行遗弃、仍然行走弯曲且需要教导的灵魂更可怕的旷野呢？——被劝诫去修直主的道路。那道路被人修直，这人远非效仿蛇的弯曲行程；而具有相反性情的人则扭曲自己的道路。因此，对这样的人和所有像他的人有这样的斥责：\Quote{你们为什么歪曲主的正道？}\BibleRef{宗 13:10}\par

% structure: p0022 source=h2 target=subsection reason=relative-heading-rank; base=h2

\subsection{十一、论主的道路：它如何狭窄，以及耶稣如何就是道路。}

主的路以两种方式被修直。首先，在默观的方式上，当思想在真理中变得清晰，没有丝毫虚假的掺杂时；其次，在行为的方式上，在正确默观了当行之事后，所产生之行为与健全的行为理论相协调。为使我们更清楚地理解经文\Quote{修直主的路}，最好将其与\WorkTitle{箴言}中所说的\Quote{不偏右，也不偏左}(\OriginalTerm{不偏右，也不偏左}{Depart not, either to the right hand or to the left})相比较。因为凡偏向任何一方的人，都已放弃保持路径笔直，不再配得关注，因为他已偏离了旅程的笔直，因为\Quote{主是公义的，喜爱公义，祂的面容注视正直}。因此，那成为关注对象并因此获得眷顾益处的人说：\Quote{上主，祢慈颜的光辉曾照耀我们。}那么，让我们像耶肋米亚（见：\BibleRef{耶 4:16}）所劝勉的那样，站在路上，让我们观看并询问上主古老的道路，让我们看清哪条是善道，并行在其中。宗徒们就是这样站着询问上主古老的道路；他们询问圣祖和先知，探究他们的著作，当他们开始理解这些著作时，他们看见了善道，即耶稣基督，祂说：\Quote{我就是道路}，他们便行在其中。因为这是一条善道，引领善人到达善父，那人从心中善库中发出善事，并且是良善忠信的仆人。这条路确实狭窄，因为许多人不能忍受行在其上，而贪爱自己的肉身；但它也被那些以暴力（见：\BibleRef{玛 11:12}）行在其上的人所逼迫，因为它不被称为\Quote{使人受苦}，而是\Quote{受苦的}。因为那条路是活路，感受行走者的品性，当那没有脱下脚上鞋（见：\BibleRef{出 3:5}）、也没有真正意识到他所站立或践踏之地是圣地的人行在其上时，它就被压迫和受苦。它将引领人到达那生命本身，祂说：\Quote{我就是生命。}因为救主，在其内结合了所有德行，具有许多面貌。对于尚未接近终点却在前进的人，祂是道路；对于已脱去一切死物的人，祂是生命。行在这条路上的人被告知不要在路上带任何东西，因为它提供食粮和一切生活所需，仇敌在此无能为力，他不需要手杖，又因它是圣的，他不需要鞋。\par

% structure: p0024 source=h2 target=subsection reason=relative-heading-rank; base=h2

\subsection{十二、赫拉克利昂关于呼声及洗者若翰的看法。}

然而，那些话：\Quote{我是旷野里呼喊者的声音}等等，可以被视为等同于\Quote{我就是那经上写着\InnerQuote{在旷野里之声}的那一位}。这样，\Person{若翰}就是那呼喊的人，他的声音就是在旷野里呼喊的：\BibleQuote{修直主的道路}。\Person{赫拉克里昂}在讨论\Person{若翰}和先知时，有些诽谤性地说：\Quote{圣言就是救主；声音是\Person{若翰}在旷野里解释的那一位；响声是全部先知的行列}。对此我们可以回答，提醒他这段经文，\BibleRef{格前14:8}：\BibleQuote{若号角发出模糊的声音，谁能准备作战呢？}以及那段说即使人懂得奥迹或有先知之恩却没有爱德，他就是发声的锣或响的钹。\BibleRef{格前13:1}如果先知的声音不过是响声，我们的主怎能指点我们去看它，正如祂说：\BibleRef{若5:39}\BibleQuote{你们查考圣经，因为你们认为其中有永生，正是这些为我作证；}以及\BibleRef{若5:46}\BibleQuote{你们若是相信\Person{梅瑟}，也必相信我；}以及\BibleQuote{依撒意亚论你们预言得真好，说：这民族用嘴唇尊敬我。}我不知道是否有人能合理地承认，救主这样说话是赞美模糊的声音，或者我们能从圣经——我们被引到那里如同从号角之声——得到任何准备，为我们对抗敌对势力的战争，假若它们的声音发出模糊的响声。如果先知们没有爱德，并且正因如此他们是发声的铜锣或响的钹，那么主如何差遣我们去听他们的声音——如这些作者所主张——好像我们能从那里得到帮助？他断言，声音当很好适应于言语时，就变成言语，如同有人说女人变成了男人；这个断言没有论证支持。并且，好像他有资格就此提出一个教理并这样前进，他宣称响声可以类似地变成声音，而声音，变成言语，他说，处于门徒的地位，而响声变成声音则处于奴隶的地位。如果他曾费心建立这些观点，我们需要费一些注意力去驳斥它们；但实际如此，仅仅否认就足以反驳。前面有一点我们推迟了讨论，即\Person{若翰}讲话的动机。我们现在可以讨论它。据\Person{赫拉克里昂}说，救主称他为先知和\Person{厄里亚}，但他自己否认他是其中任何一个。当救主，\Person{赫拉克里昂}说，称他为先知和\Person{厄里亚}时，祂不是在说\Person{若翰}本人，而是关于他的环境；但当祂称他比先知和妇人所生的人更大时，祂是在描述\Person{若翰}本人的性格。另一方面，当\Person{若翰}被问及他自己时，他的回答涉及他自己，而非他的环境。我们已尽可能仔细地检查了这一点，将所讨论的每一条术语与\Person{赫拉克里昂}的陈述相比较，以免他表达得不够准确。因为为什么说他是\Person{厄里亚}和他是先知适用于他周围之人，而说他是旷野里呼喊者的声音适用于他自己，完全没有展示。\Person{赫拉克里昂}只给出一个例子，即：他的环境，可以说是他的衣服，和他自己不同，当人家问及他的衣服，他是否是衣服时，他不能回答\Quote{是}。如今，他是那要来的\Person{厄里亚}，是他的衣服，这几乎与我所能看到的\Person{赫拉克里昂}的观点不一致；或许它与我们自己给\BibleQuote{以\Person{厄里亚}的精神和能力}这句话的解释一致；在某种意义上，可以说这\Person{厄里亚}的精神相当于\Person{若翰}的灵魂。他接着试图确定为什么那些由犹太人派去询问\Person{若翰}的是司铎和肋未人，他的回答并不坏：这样的人，既然献身于天主的服务，应当忙于打听这样的事。然而，当他继续说\Quote{因为\Person{若翰}属于肋未支派}时，这个考虑就不那么周到了。我们自己在上面提出了问题，并看到如果被派去的犹太人知道\Person{若翰}的出生，他们就不能问他是不是\Person{厄里亚}。然后，在处理\Quote{你是那位先知吗？}这个问题时，\Person{赫拉克里昂}不认为加冠词有什么特别力量，并说：\Quote{他们问他是否是一位先知，想要知道这个更一般的事实。}再次，不仅\Person{赫拉克里昂}，而且据我所知，所有与我们的观点相异的人，好像他们不能处理微小的歧义并作出适当区分，他们假设\Person{若翰}比\Person{厄里亚}和所有先知都大。经文是：\Quote{妇人所生的，没有比\Person{若翰}更大的}；但这容许两种意义：\Person{若翰}比他们都大，或者，他们中有一些与他相等。因为虽然许多先知与他相等，但就赋予他的恩宠而言，仍然可以说他们中没有比他更大的。他认为这证实了\Person{若翰}更大的观点，即\Quote{他被\Person{依撒意亚}预言了}；因为所有发出预言的人中，没有一个人被天主认为配得上这一殊荣。然而，这是一个冒失的陈述，并暗示对所谓的\WorkTitle{旧约}不尊重，完全无视\Person{厄里亚}本人是预言的对象这一事实。因为\Person{厄里亚}被\Person{玛拉基亚}预言了，他说：\BibleQuote{看，我派遣提市贝特人\Person{厄里亚}到你们这里来，他要使父亲的心转向儿子。}同样，\Person{约史雅}，正如我们在\BibleBook{列王}{纪上}读到，\BibleRef{列上13:2}被来自\Place{犹大}的先知指名预言了；因为他说，\Person{雅洛贝罕}也在祭台旁时：\BibleQuote{上主这样说：看，\Person{达味}要生一个儿子，他的名字是\Person{约史雅}。}还有一些人说，\Person{三松}被\Person{雅各伯}预言了，当他\BibleRef{创49:16}说：\BibleQuote{\Place{丹}将审判自己的百姓，如同\Place{以色列}的一个支派}，因为审判\Place{以色列}的\Person{三松}属于\Place{丹}支派。以上是证据，表明\Person{赫拉克里昂}在尝试解释\Quote{我是旷野里呼喊者的声音}这句话时，所宣称的\Person{若翰}独自是预言的对象这一论断的鲁莽。\par

% structure: p0026 source=h2 target=subsection reason=relative-heading-rank; base=h2

\subsection{十三、若望福音第一章第廿四、廿五节：论若翰的洗礼、厄里亚的洗礼与基督的洗礼。}

那被派遣的是一些法利塞人。他们问他说：\BibleRef{若 1:24} \BibleQuote{若你不是基督，也不是厄里亚，也不是那位先知，那么你为什么施洗呢？} 那些从耶路撒冷派来问若翰这些问题的司祭和肋未人，在得知若翰不是什么人以及他是什么人之后，保持了得体的缄默，仿佛是默许并表明他们接受所说的话，并看出这洗礼正适合那在旷野中呼喊、为主预备道路的声音。但法利塞人，正如他们的名字所表明的，是一群分裂和煽动的人，显示他们不同意京城中的犹太人和服事天主的职员、司祭及肋未人。他们派遣了以斥责为事的使者，并在他们权力范围内禁止他施洗；他们的使者问：\Quote{若你不是基督，也不是厄里亚，也不是那位先知，那么你为什么施洗呢？} 如果我们把记载在各福音书中的内容拼接成一个陈述，我们就应该说，此时他们说话正如这里所记载的，但后来当他们想接受洗礼时，却听到了若翰的训言：\BibleRef{玛 3:7} \BibleQuote{毒蛇的种类！谁指示你们逃避那将要来临的忿怒呢？那么，结出与悔改相称的果实吧！} 这就是洗者若翰在玛窦福音中所说的，当他看见许多法利塞人和撒杜塞人来到他的洗礼时，显然没有悔改的果实，并法利塞式地自夸有亚巴郎为父。为此，若翰斥责他们，他因圣神的沟通而拥有厄里亚的热忱。因为那是一句斥责的话：\BibleQuote{你们不要心里说：我们有亚巴郎为父。} 这也是教师的话，当他论及那些因心硬而被称为无信之石的人，并说藉天主的权能，这些石头可变为亚巴郎的子孙；因为他们出现在先知的眼前，并未因他的神圣注视而退缩。因此他的话：\BibleQuote{我告诉你们，天主能从这些石头给亚巴郎兴起子孙来。} 既然他们来到他的洗礼而没有做出与悔改相称的果实，他就最合宜地对他们说：\BibleQuote{斧头已放在树根上；凡不结好果子的树，就砍下来丢在火里。} 这无异于对他们说：既然你们来到洗礼而没有做出与悔改相称的果实，你们就是一棵不结好果子的树，必须被那锋利刺人的斧头砍倒，这斧头就是那生活的、有力的、比任何双刃剑更锋利的圣言。法利塞人自以为义的评价也由路加在经文中呈现：\BibleRef{路 18:10} \BibleQuote{有两个人上圣殿去祈祷：一个是法利塞人，另一个是税吏。那法利塞人站着，自言自语这样祈祷：天主，我感谢你，我不像别人一样勒索、不义、奸淫，也不像这个税吏。} 这话的结果是，那税吏成了义人回家，而非那法利塞人；并且引出了教训：凡高举自己的，必被贬抑。于是，他们以救主斥责的话所描述的性格来到若翰的洗礼，身为假善人；洗者若翰也注意到他们舌下有蝮蛇的毒汁和毒蛇的毒，因为\BibleQuote{蝮蛇的毒汁在他们唇舌下。} 蛇的比喻恰当地表明他们的性情，并且在他们那个更尖锐的问题中清楚显露：\Quote{若你不是基督，也不是厄里亚，也不是那位先知，那么你为什么施洗呢？} 对此，我很乐意回答，如果情况是基督、厄里亚和那位先知是施洗的，而旷野中呼喊的声音却没有权柄这样做的话：\Quote{朋友们，你们最严厉地责问那位在基督面前被派遣、为祂预备道路的使者。与这一点相关的奥秘全都向你们隐藏；因为耶稣，无论你们愿意与否，就是基督，祂自己没有施洗，而是祂的门徒施洗，祂自己就是那位先知。你们怎么会相信那要来的厄里亚会施洗呢？} 在阿哈布的时代，他没有给祭台上的木柴施洗，尽管它们需要这样的沐浴才能被焚烧，那时主在火中显现。不，他命令司祭为他这样做，而且不止一次；因为他说：\BibleQuote{第二次倒水} 他们就第二次倒水，又说：\BibleQuote{第三次倒水} 他们就第三次倒水。那么，如果那时他自己没有施洗，而是将工作留给别人，他又如何在玛拉基亚所说的时期施洗呢？因此，基督不用水施洗，而是祂的门徒施洗。祂为自己保留用圣神和火施洗。赫拉克利翁接受法利塞人的话，认为它明确暗示施洗的职务属于基督、厄里亚和每一位先知，因为他用了这些话：\Quote{唯有他们的职务才是施洗。} 他被我们刚才所说的反驳，尤其因他把\Quote{先知}一词作一般意义来理解这一点；因为他不能证明有任何先知施过洗。他补充说，不无道理地，法利塞人出于恶意发问，而非出于求知欲。\par

% structure: p0028 source=h2 target=subsection reason=relative-heading-rank; base=h2

\subsection{十四、四福音关于洗者若翰的记述、关于他的预言、他对群众和法利塞人的讲话等的比较。}

我们认为有必要将我们所考虑的这段经文与其他福音书中类似的表达进行比较。我们将逐点进行，直至本卷的结尾，以便那些看似不一致的用语得以显明其和谐，并且每处特有的含义得以解释。在当前这段经文中，我们就这样做。\BibleQuote{旷野中有呼号者的声音：修直主的道路} 是由门徒若望放在洗者若翰口中的。而在\BibleBook{马尔谷}{福音}中，同样的经文被记录在耶稣基督福音的开头，正如\BibleBook{依撒意亚}{先知书}上所记载的：\BibleQuote{耶稣基督福音的开始，正如依撒意亚先知书上所记载的：\InnerQuote{看，我派遣我的使者在你面前，他要在你前面预备你的道路。旷野中有呼号者的声音：预备主的道路，修直祂的路径。}} 现在，若望所添加的 \BibleQuote{修直主的道路} 并不在先知的书中。也许若望试图压缩 \BibleQuote{预备主的道路，修直我们天主的路径}，所以写成了 \BibleQuote{修直主的道路}；而马尔谷则将在不同地方由两位不同先知所说的两段预言组合成一个预言：\BibleQuote{正如依撒意亚先知书上所记载的：\InnerQuote{看，我派遣我的使者在你面前，他要在你前面预备你的道路。旷野中有呼号者的声音：预备主的道路，修直祂的路径。}} \BibleQuote{旷野中有呼号者的声音} 紧接在希则克雅病愈的叙述之后（\BibleRef{依 40:3}），而 \BibleQuote{看，我派遣我的使者在你面前} 则是玛拉基亚所写。若望在此所行的，即缩短他所引用的经文，我们发现在别处马尔谷也是这样做的。因为先知的话是 \BibleQuote{预备主的道路，修直我们天主的路径}，而马尔谷却写成 \BibleQuote{预备主的道路，修直祂的路径}。若望在经文 \BibleQuote{看，我派遣我的使者在你面前，他要在你前面预备你的道路} 中也进行了类似的缩写，他没有加上原文中的 \BibleQuote{在你前面}。现在，我们来到这句话：\BibleQuote{他们是由法利塞人派来的，他们问祂}（\BibleRef{若 1:24}）。我们因对这段经文的考察，已先行将法利塞人的询问（\BibleBook{玛窦}{福音}未提及）置于玛窦所记载的事件之前，即当若翰看见许多法利塞人和撒杜塞人来到他的洗礼前，对他们说 \BibleQuote{毒蛇的种类} 等话时。因为自然的顺序是：他们先询问，然后才来。我们还须注意，在\BibleBook{玛窦}{福音}中记载，耶路撒冷、全犹太和约但河全区域的人都出去到若翰那里，在约但河里受他的洗，承认自己的罪；但听到洗者任何责备或斥责之言的并不是这些人，而是他所看见的那些许多法利塞人和撒杜塞人。正是这些人受到 \BibleQuote{毒蛇的种类} 等话的迎接（\BibleRef{玛 3:7}）。\BibleBook{马尔谷}{福音}则没有记载若翰对来者说过任何斥责的话，因为来者是全犹太地和所有耶路撒冷的人，他们在约但河受他的洗，承认自己的罪。这是因为马尔谷没有提到法利塞人和撒杜塞人来到若翰这里。我们还需提及另一个情况：玛窦和马尔谷都说，在一处是耶路撒冷、全犹太和约但河全区域的人，在另一处是全犹太地和所有耶路撒冷的人，都受洗，承认自己的罪；但当玛窦引入法利塞人和撒杜塞人来到洗礼时，他没有说他们承认自己的罪，这很可能就是他们被称为 \BibleQuote{毒蛇的种类} 的原因。读者啊，不要以为我们在讨论那些由法利塞人派来向若翰提问的人时，引用其他福音书的平行经文是不合适的。因为如果我们已经指出法利塞人的询问（由门徒若望记载）与他们的洗礼（见于\BibleBook{玛窦}{福音}）之间的适当联系，我们就很难不考察所提及的经文，也不记录对这些经文所作的观察。\BibleBook{路加}{福音}同\BibleBook{马尔谷}{福音}一样，记下了 \BibleQuote{旷野中有呼号者的声音} 这段经文，但他处理如下（\BibleRef{路 3:2}）：\BibleQuote{天主的话在旷野中来到匝加利亚的儿子若翰那里。他来到约但河全区域，宣讲悔改的洗礼，为使罪得赦；正如依撒意亚先知书上所记载的：\InnerQuote{旷野中有呼号者的声音：预备主的道路，修直祂的路径。}} 然而，路加加上了预言的连续部分：\BibleQuote{一切山谷都要填满，一切山丘都要削平；弯曲的要修直，崎岖的要成为平坦的道路；凡有血肉的，都要看见天主的救恩。} 他像马尔谷一样写 \BibleQuote{修直祂的道路}，如我们之前所见，缩短了 \BibleQuote{修直我们天主的道路} 这句话。在 \BibleQuote{一切弯曲的要修直} 这句话中，他省略了 \BibleQuote{一切}，并将 \BibleQuote{修直} 从复数变为单数。此外，他不用 \BibleQuote{崎岖之地变为平原}，而写 \BibleQuote{崎岖的道路变为平坦的道路}，并省略了 \BibleQuote{上主的荣耀要显示出来}，而给出接下来的 \BibleQuote{凡有血肉的，都要看见天主的救恩}。这些观察有助于显示福音作者们习惯于缩短先知们的话语。还须注意，\BibleQuote{毒蛇的种类} 等话，玛窦说是对前来受洗的法利塞人和撒杜塞人说的，这些人不同于那些承认自己罪的人，对那些人没有说过这样的话。相反，在路加福音中，这些话是对出来接受若翰洗礼的群众说的，并且受洗的人没有分为两类，就如我们在玛窦福音中所见。但细心观察者会看到，玛窦并没有以称许的方式提到群众，他可能意指洗者的讲话 \BibleQuote{毒蛇的种类} 等也是对群众说的。另一点是，他对法利塞人和撒杜塞人说：\BibleQuote{结出} 单数的 \BibleQuote{与悔改相称的果实}，但对群众则用复数：\BibleQuote{结出与悔改相称的果实}。也许法利塞人被要求结出悔改的特殊果实，即圣子及对祂的信德；而群众连善事的开端都没有，所以要求他们结出悔改的所有果实，因此对他们用复数。此外，对法利塞人说：\BibleQuote{不要心里说：\InnerQuote{我们有亚巴郎为父。}} 因为群众现在有了开端，似乎被引入天主圣言，接近真理；于是他们开始心里说：\BibleQuote{我们有亚巴郎为父。} 相反，法利塞人不是刚开始这样想，而是早已如此。但两类人都看见若翰指着前述的石头，宣称甚至能从这些石头中给亚巴郎兴起子孙，从无意识和死寂中兴起。请注意，根据先知的话（\BibleRef{欧 10:13}），对法利塞人所说的话（\BibleRef{玛 3:10}）：\BibleQuote{你们吃了虚假的果实}，他们有虚假的果实——\BibleQuote{凡不结好果子的树，就砍下来，丢在火里}；而对那些根本不结果子的群众（\BibleRef{路 3:9}）：\BibleQuote{凡不结果子的树，就砍下来。} 因为那根本不结果子的，就没有好果子，所以该被砍下来。但那结果子的，绝非好果子，因此也需要斧头来砍倒它。但是，如果我们更仔细地考察这关于果子的事，我们就会发现，那刚刚开始栽培的树，即使它不白费力气，也不可能结出最初的佳果。农夫满足于那刚被栽培的树最初能结出它可能有的果子；以后，当他按照自己的艺术修剪并训练这树时，他收取的就不再是起初偶然结出的果子，而是佳果。法律本身也支持这种解释，因为它说（\BibleRef{申 19:23}）栽种的人要等三年，让树修剪，不吃它的果子。\BibleQuote{三年之久，} 它说，\BibleQuote{果子对你们算是不洁净的，不可吃；但第四年，所有的果子都是圣的，为赞美上主。} 这就解释了为什么在向群众说话时，\BibleQuote{好} 字被省略了：\BibleQuote{因此，凡不结果子的树，就砍下来，丢在火里。} 那继续结起初那种果子的树，是不结好果子的树，因此被砍下来，丢在火里，因为三年过后，到第四年时，它还不为赞美上主而结好果子。我这样引用其他福音书的经文，似乎有些离题，但我不能认为这毫无用处或与当前主题无关。因为法利塞人在从耶路撒冷来的司祭和肋未人之后，派人到若翰那里，这些来的人问他是谁，并询问：\BibleQuote{你既不是基督，又不是厄里亚，也不是那位先知，你为什么施洗呢？} 他们这样询问后，就立刻来受洗，正如玛窦所记载的，然后他们听到了适合他们欺诈和伪善的话。但对群众所说的话与对他们所说的话非常相似，因此有必要仔细考察这两段话并加以比较。当我们这样进行时，在事情的顺序中出现了各种点，我们必须加以考虑。除了已经说过的，我们还须补充以下内容。在\BibleBook{若望}{福音}中，我们发现有两次派人的记载：一次是犹太人从耶路撒冷派司祭和肋未人；另一次是法利塞人想知道他为什么施洗。我们发现，在询问之后，法利塞人前来受洗。这不就是犹太人（他们派了从耶路撒冷来的第一次代表团）在第二次代表团（即法利塞人）之前就听到了若翰的话，因而比他们先到吗？因为耶路撒冷和全犹太，以及约但河全区域的人，都在约但河里受他的洗，承认自己的罪；或者如马尔谷所说：\BibleQuote{全犹太地和所有耶路撒冷的人都出去到他那里，在约但河里受他的洗，承认自己的罪。} 现在，玛窦没有说法利塞人和撒杜塞人（就是那些听到 \BibleQuote{毒蛇的种类} 等话的人）承认自己的罪；路加也没有说那些受到同样斥责的群众承认自己的罪。可能会有人问：如果全耶路撒冷城、全犹太和约但河全区域的人都在约但河受了若翰的洗，救主怎能说（\BibleRef{玛 11:13}）\BibleQuote{若翰来了，也不吃，也不喝，你们说他有魔鬼}；祂又怎能对那些问祂的人说（\BibleRef{玛 21:23}）：\BibleQuote{你凭什么权柄做这些事？我也要问你们一句话，你们若告诉我，我就告诉你们我凭什么权柄做这些事。若翰的洗礼是从哪里来的？是从天上来的，还是从人来的？他们彼此商议说：\InnerQuote{我们若说从天上来，祂必对我们说：\InnerQuote{你们为什么不信他？}} } 这个难点的解决办法如下。法利塞人，正如我们之前所见，被若翰称为 \BibleQuote{毒蛇的种类} 等，他们来受洗，却不信他，大概是因为害怕群众，并且由于他们惯常对群众的伪善，认为接受洗礼是合适的，以免显得与受洗的人为敌。那么，他们的信念是：他的洗礼是从人来的，而不是从天上来；但为了群众，怕被石头砸死，他们不敢说出自己的想法。因此，救主对法利塞人的话与福音书中关于群众常去若翰洗礼的叙述并无矛盾。法利塞人的厚颜无耻表现在他们说若翰有魔鬼，同样也说耶稣是靠魔王贝耳则步、魔鬼的首领行奇事。\par

% structure: p0030 source=h2 target=subsection reason=relative-heading-rank; base=h2

\subsection{十五、洗者若翰如何回答法利塞人的问题并高举基督的本性。论他所不能解开的鞋带。}

\Person{若翰}\BibleRef{若 1:26}回答他们说：\BibleQuote{我用水施洗，但有一位站在你们中间，是你们所不认识的；祂在我以后来，我就是给祂解鞋带也不配。}\Person{赫拉克利昂}认为若翰对法利塞人所差来的人的回答，并非针对他们所问的，而是针对他所愿意的；他没有注意到，他指责先知缺乏礼貌，因为当他被问一件事时，却回答另一件事；这是在谈话中应避免的过错。反之，我们断言，回答准确地回应了问题。他们问：\Quote{你既然不是基督，为什么施洗？}对此还能有什么别的回答，除了表明他的洗礼在本质上是一件身体的事？他说：\BibleQuote{我用水施洗}；这是对\Quote{你为什么施洗？}的回答。至于他们问题的第二部分\Quote{如果你不是基督}，他通过高举基督的超越本性来回答，说明祂具有如此德能，以致于祂的天主性不可见，却临在于每一个人，并延伸至整个宇宙。这就是\BibleQuote{有一位站在你们中间}这句话所表明的。此外，虽然法利塞人期待基督的来临，但他们没有在祂身上看到若翰所说的那种本性；他们只相信祂是一位完美圣洁的人。因此，若翰责备他们对祂超越性的无知，并在\BibleQuote{有一位站在你们中间}这句话之后加上\BibleQuote{是你们所不认识的}。并且，为了避免有人以为那不可见的、临于每一个人甚至全世界的祂，与那降生成人、显现于世上并与人交谈的祂是不同的一位，他在\BibleQuote{有一位站在你们中间是你们所不认识的}这句话之后又加上\BibleQuote{在我以后来的那位}，即在我以后将要显现的那位。若翰深知，与祂那无与伦比的卓越相比，自己的本性远远不及，尽管有人怀疑他可能就是基督；因此，为了表明他距离基督的伟大有多么遥远，以免任何人对他存有超出所见所闻的想法，他继续说：\BibleQuote{我不配给祂解鞋带。}通过这话，他如同比喻般地传达出，他不配解决和解释关于基督取了人性身体的道理，这道理对那些不理解的人而言（如同鞋带）是系紧而隐藏的——以至于对这如此压缩的来临说出任何相称的话。\par

% structure: p0032 source=h2 target=subsection reason=relative-heading-rank; base=h2

\subsection{十六、比较不同福音书中若翰对耶稣的见证。}

在我们正在研究的经文\BibleQuote{我用水施洗}之时，比较福音作者们与\Person{若翰}平行的表述或许并不失礼。玛窦记载，当洗者\Person{若翰}看到许多法利塞人和撒杜塞人来到他的洗礼处时，在指责他们之后，正如我们已经研究过的，他接着说：\BibleRef{玛 3:11} \BibleQuote{我固然用水洗你们，为使你们悔改；但在我以后要来的那一位比我更强，我连提祂的鞋也不配；祂要以圣神及火洗你们。}这与\BibleBook{若望}{福音}中洗者向法利塞人所差来的人表明自己用水施洗的话语是一致的。马尔谷又说：\BibleRef{谷 1:6} \BibleQuote{若翰宣讲说：那比我更有力量的，要在我以后来，我连俯身解祂的鞋带也不配。我用水洗你们，祂却要以圣神洗你们。}而路加记载：\BibleRef{路 3:16} 当众人期待，并在心中思量若翰是否就是基督时，若翰回答众人说：\BibleQuote{我固然以水洗你们，但比我强的那一位要来，我就是解祂的鞋带也不配；祂要以圣神和火洗你们。}\par

% structure: p0034 source=h2 target=subsection reason=relative-heading-rank; base=h2

\subsection{十七、论\BibleBook{玛窦}{福音}中若翰对耶稣的见证，}

这些就是四部福音的平行经文；让我们尽可能清楚地去探究每段经文的旨意以及它们之间的差异。我们将从《玛窦福音》开始，据传统记载，玛窦比其他宗徒更早向希伯来人——即那些受割礼而信的人——发表了他的福音。他说：\\BibleQuote{我以水洗你们，为使你们悔改}，就好比洁净你们，使你们远离恶行，并召唤你们悔改；因为我前来是为给主预备一个合用的百姓，并藉我的悔改之洗，为那在我之后来临的一位预备道路，祂将比我的能力更有效、更有力地使你们得益。因为祂的圣洗不仅限于身体；祂以圣神充满悔改者，祂那更神圣的火焰消除一切物质性的东西，焚尽所有属地的——不仅从接受祂进入生命之人身上，甚至连那些从拥有祂的人那里听闻此事的人也是如此。那在我之后来的，比我更强有力，我甚至不能承担祂周围那些离祂最远的权能的边缘（它们并非公开显露，以致任何人都能看见），也不能承担那些支持它们的人。我不知道该说些什么。我该说我极大的软弱，连这些在基督内与祂更伟大的事相比算为最小的事都不能承担，还是该说祂那超越整个世界的超然神性？我，这领受了如此恩宠，竟被尊为先知，预言我来到这人间生活的人，经上记载：\\BibleQuote{荒野中有呼号者的声音}，以及\\BibleQuote{看，我派遣我的使者在你面前}；我，那位站在天主面前的加俾额尔曾向年迈的父亲出乎意料地宣布了我诞生的消息；我，因我的名字，匝加利亚恢复了声音并得以预言；我，我的主为我作证，在女人所生者中没有比我再大的，我竟然连祂的鞋也不配提！若不配提祂的鞋，那关于祂的衣裳还能说什么呢？谁能如此伟大，以至于能守护祂的外衣？谁能认为自己能理解那件从上到下浑然一体的无缝长衣所包含的意义？应当注意，虽然四部福音都记载若翰宣称自己来用水施洗，但唯有玛窦加上了\\BibleQuote{为使你们悔改}这句话，教导说圣洗的益处与受洗者的意向有关：对悔改者，圣洗是得救的；对无忏悔而来的人，它反会招致更大的惩罚。在此我们必须注意，正如救主在祂所行的治愈奇迹中所行的奇事——它们象征着任何时候藉天主圣言从疾病或病痛中获得释放的人——虽然这些奇迹是针对身体并带来身体的缓解，但也召唤那些受益者去操练信德；同样，这水的洗涤象征着灵魂洗净自己一切恶行的污秽，就其本身而言，对那将自己交付给呼求天主圣三之神圣能力的人，它是神圣恩赐的起始和源头；因为\\BibleQuote{神恩虽有区别}。这一观点从《宗徒大事录》所记载的叙述中得到确认，该叙述表明，在那些恰当领受圣洗的人中，当水为他们预备了道路之后，圣神是多么明显地降临在领受圣洗的人身上。西满术士看到这一切，大为惊讶，希望从伯多禄那里领受这份恩赐，虽然他渴望的是一件好事，却想藉不义的钱财获得它。我们顺便再指出一点：若翰的圣洗低于耶稣藉祂的门徒所施行的圣洗。在《宗徒大事录》中，那些受若翰圣洗且未曾听说有圣神的人（\\BibleRef{宗 19:2}），被宗徒重新付洗。藉若翰，并没有发生重生；但藉耶稣并通过祂的门徒，重生确实发生了，那被称为重生的洗涤伴随着圣神的更新而实现；因为圣神现在进一步降临，祂来自天主，超越水之上，并且并非在众人受水洗之后都降临。到目前为止，这就是我们对《玛窦福音》记载的考察。\par

% structure: p0036 source=h2 target=subsection reason=relative-heading-rank; base=h2

\subsection{十八、论《马尔谷福音》中的见证。救主的鞋和解开祂的鞋带是什么意思。}

现在让我们思考玛尔谷所记述的。玛尔谷关于若翰宣讲的记述与其他福音一致。他的话是：\Quote{在我以后要来一位比我更强有力的，} 这等同于 \Quote{在我以后来的那位比我更强有力。} 然而，接下来的话有所不同：\Quote{我连俯身解祂的鞋带也不配。} 因为，替人提鞋是一回事——显然，鞋已经从穿着者脚上解下来了——而俯身解他的鞋带又是另一回事。由此可知，既然信徒不能认为两位圣史中有任何错误或曲解，那么洗者若翰必然是在不同时间做了这两次宣告，并且意在表达不同的事物。并非如某些人所猜想的那样，报道指向同一事件，却因对某些事实或言辞记忆不清而出现差异。现在，提起耶稣的鞋子是一件大事，俯身探究祂使命的肉身层面，即发生在较低境界的事物，以便在较低领域沉思祂的肖像，并解开与祂降生成人的奥秘相连的每一重困难，这些困难就好比祂的鞋带，这是一件大事。因为蒙昧的束缚是一，知识的钥匙也是一；就连妇人所生的最伟大者，凭自身也不足以解开或开启此类事物，因为那位最初系紧并锁住的主，也将祂鞋带的松开和祂所关闭之物的开启赐予祂所愿意的人。若关于鞋子的段落具有奥秘意义，我们不应轻视而应加以思考。现在我认为，当天主圣子取肉身与骨骼时的\OriginalTerm{降生成人}{inhumanisation}是祂的鞋子之一，另一只是下降到\OriginalTerm{阴府}{Hades}——无论那阴府为何——以及与圣神一同往监狱的旅程。关于下降到阴府，我们读到第十六篇圣咏：\BibleQuote{你不会把我的灵魂留在阴府。} 至于与圣神一同在监狱中的旅程，我们在伯多禄的公函中读到：\BibleRef{伯前 3:18} \BibleQuote{在肉身内被处死，} 他说，\BibleQuote{却在圣神内生活；祂也曾去给那些在狱中的灵魂宣讲过，他们从前在诺厄时代建造方舟期间，天主的耐心等待时，曾经叛逆。} 那么，谁有能力相称地阐述这两段旅程的意义，谁就能解开耶稣鞋子的鞋带；他俯身于思想中，随耶稣下降阴府，从天上和基督神性的奥秘，下降至祂穿上人（作为祂的鞋子）时所必须完成的与我们同在的降临。现今，那位穿上人的，也穿上了死者，因为 \BibleRef{罗 14:9} \BibleQuote{正是为此，耶稣死了并复活了，为使祂成为死者和生者的主。} 这就是祂穿上生者与死者——即地上的居民和阴府中的居民——的原因，为使祂成为死者和生者的主。那么，谁能俯身解开如此鞋子的鞋带，并且在解开后不让它们掉落，而是以他所领受的第二种能力将它们拿起并提起，把它们的含义存记在心呢？\par

% structure: p0038 source=h2 target=subsection reason=relative-heading-rank; base=h2

\subsection{十九、路加和若望暗示，人可以在不俯身的情况下解开\OriginalTerm{圣言}{Logos}的鞋带。}

然而，我们不可不问，路加和若望如何给出这句话而没有出现\Quote{俯身}一词。祂也许认为，俯身者可以被理解为在我们所陈述的意义上解开。另一方面，也可能有人将目光注视\OriginalTerm{圣言}{Logos}崇高荣耀的高度，从而找到那些鞋子的解开方法——这些鞋子在人们寻找时似乎是被绑着的——以便他也解开那些可与圣言分离的鞋子，并看见圣言脱去了低等事物，就如祂原本那样，天主子。\par

% structure: p0040 source=h2 target=subsection reason=relative-heading-rank; base=h2

\subsection{二十、\Quote{足够}与\Quote{相配}之间的区别}

若望记载说，洗者若翰自称不配，玛尔谷自称不够，这两者并非相同。一个不配的人可能足够，而配得的人可能不够。因为即使恩赐是按益处赐下，而不仅仅按信德的比例，但似乎属于爱人的天主，祂预见到自我意见或自负兴起所必带来的损害，不将足够甚至赐予配得的人。但属于天主的良善，以慷慨施恩征服祂恩惠的对象，预先接纳那注定要成为配得的人，甚至在他成为配得之前就以足够装饰他，使他因足够而后来成为配得；他不是先成为配得，然后抢先于施恩者，在时机未到之前领受祂的恩赐而达到足够。现在，三位（圣史）记载中洗者若翰都说他不够，而在若望中他说自己不配。但可能，先前宣称不够的人后来变得足够，尽管或许他仍不配；或者，当他自认不配且确实不配时，他却达到了配得。除非有人说，人的本性永远不能相称地完成这解开或提起的工作，因此若翰说真话：他从未足够到能解开救主鞋子的鞋带，也不配。无论我们在思想中摄入多少，仍有尚未理解的事物存留；因为我们在耶稣，息辣之子所作的智慧书中读到：\BibleRef{德 18:7} \BibleQuote{人做完时，他才开始；他停止时，他便会怀疑。}\par

% structure: p0042 source=h2 target=subsection reason=relative-heading-rank; base=h2

\subsection{二十一、第四福音只说一只鞋，其他福音说两只。此事的含义。}

至于在三部福音中所提及的鞋，我们有一个需要考量的问题；我们必须将它们与若望门徒所说的单只鞋相比。我们在那里读到：\Quote{我是不配的}给祂解鞋带。也许他因天主的恩宠而获胜，并接受了那凭自己本不配去行的事的恩赐，即解开其中一只鞋的鞋带，就是在他看见了救主在人间居住之后，他为此作证。但他不知道随后要发生的事，即耶稣是否也要到那地方去，就是他在狱中被斩首后所要去的，或者他是否要期待另一位；因此他隐晦地暗示了那后来向我们显明的疑惑，并说：\Quote{我不配解祂的鞋带。}若有人认为这是多余的推测，他可以将关于鞋和鞋带的话合而为一，仿佛若望说：我绝不配解开祂的鞋带，连起初的一只鞋带也不配。或者，以下是将四福音所说的话结合起来的一种方法：如果若望明白耶稣在此世的居住，但对未来有疑惑，那么他说自己不配解祂的鞋带是完全真实的；因为虽然他解开了那只鞋的鞋带，但并未解开两只。另一方面，他关于鞋带的话也完全真实；因为如我们所见，他仍然疑惑耶稣是否就是那要来的那位，或者是否要在另一区域期待另一位。\par

% structure: p0044 source=h2 target=subsection reason=relative-heading-rank; base=h2

\subsection{二十二、圣言如何站在人中间而不被人认识。}

至于那句话：\BibleQuote{有一位站在你们中间，你们却不认识祂}，这引导我们思考天主之子，即圣言，万物藉祂而造成的，因为祂存在于万有本质的整个基础中，祂本身就是上智。因为祂从起初就渗透了一切受造物，使得任何时候所造之物都藉祂而造成，并且对于任何事物总是可以说：\BibleQuote{万物是藉着祂造成的；凡受造的，没有一样不是藉着祂造的}；以及这句话：\BibleQuote{祢以智慧造万有。}如今，如果祂渗透一切受造物，那么祂也在那些发问者中间，他们问：\BibleQuote{祢既不是基督，又不是厄里亚，也不是那位先知，那么祢为什么施洗呢？}在他们中间站着圣言，祂是同一而坚定不移的，由圣父安置在一切处。或者，\BibleQuote{有一位站在你们中间}这句话可以理解为：在你们人中间，因为你们是有理性的存在，站着祂——圣经证明祂是每一身体中的主权原则，因此祂临在于你们心中。因此，那些拥有圣言在他们中间，却不思考祂的本性，也不思想祂从何泉源和原则而来，以及祂如何赐予他们现有的本性的人，这些人虽拥有祂在他们中间，却不认识祂。但是洗者若翰认识祂：因为\BibleQuote{你们却不认识祂}这句话是对法利塞人的责备，表明若翰很认识他们所不认识的圣言。因此，洗者若翰认识祂，看见祂在自己以后来——祂如今站在他们中间，这就是说，在他之后并在他的圣洗教训中，住在那些按照理性（或圣言）接受那净化之礼的人里面。然而，\Quote{以后}这个词在这里的意义不同于耶稣命令我们\Quote{跟随}祂时的意义；因为后一种情况是命令我们跟从祂，踏着祂的足迹，从而到达圣父面前；而在前一种情况，意思是说在若翰的教训之后（因为\BibleQuote{祂来是为使众人藉着祂而信}），圣言与那些预备好自己的人同住，他们被较卑微的话所净化，为的是那成全的圣言。首先，圣父站着，毫无转变或变化；然后，祂的圣言也站着，始终进行祂的救恩工作，即使站在人中间，也不被理解，甚至不被看见。祂也站着施教，邀请所有人从祂丰盛的泉源中饮用，因为\BibleRef{若 7:37}：\BibleQuote{耶稣站着大声说：\InnerQuote{谁若渴，到我这里来喝罢！}}\par

% structure: p0046 source=h2 target=subsection reason=relative-heading-rank; base=h2

\subsection{二十三、赫拉克利昂对洗者若翰此话的看法，以及对耶稣之鞋的解释。}

但是赫拉克利翁声称，\BibleQuote{你们中间站着一位}这句话等同于\Quote{祂已经在这里，祂在世界中，在众人内，并且祂已经显示给你们众人。}借此，他抹去了这句话中也包含的另一层含义，即圣言已渗透整个世界。因为我们不得不对他说：祂何时不是临在的？何时不在世界中？这部福音不是说过：\BibleQuote{祂在世界中，世界是由祂造成的，世界却不认识祂。}正因如此，那些对圣言是\BibleQuote{你们不认识的那位}的人不认识祂：他们从未离开世界，但世界不认识祂。但祂何时停止在人间？当祂在依撒意亚中说：\BibleQuote{上主的神临于我，因为祂给我傅了油}（\BibleRef{依 61:1}），以及\BibleQuote{我显现给那些不寻求我的人}（\BibleRef{依 65:1}）时，祂不就在依撒意亚内吗？他们也要说说，祂难道不在达味内吗？那时达味不是出于自己说：\BibleQuote{但祂立我为王，在熙雍祂的圣山上}，以及圣咏中那些以基督的口吻所说的话。我何必细数这证明的细节？它们实在难以数计，因为我能清楚显示祂始终在人间。而这就足以表明赫拉克利翁对\BibleQuote{你们中间站着一位}的解释是错误的，他说这句话等同于\Quote{祂已经在这里，祂在世界中，在众人内。}我们倾向于同意他说\BibleQuote{在我以后来的那一位}这句话表明若翰是基督的前驱，因为他事实上是一种仆人在主人前奔跑。然而，\BibleQuote{我连给祂解鞋带也不配}这句话，当他说\Quote{在这些话中，若翰承认他连为基督做最卑微的服务也不配}时，却得到了过于简单的解释。在此解释之后，他不无道理地补充说：\Quote{我不配让祂因我的缘故从祂的伟大中降下，并取肉体作为祂的鞋履，关于这鞋履，我无法给出任何解释或描述，也无法解开其安排。}赫拉克利翁以其鞋履指代世界，显示出某种心胸开阔，但紧接着他近乎不虔，宣称所有这些话都应理解为若翰在此所意谓的那个人物。因为他认为正是世界的\OriginalTerm{造物主}{\GreekText{δημιουργός}}借这些话承认自己比基督低等；这是不虔之极。因为那派遣祂的圣父，即亚巴郎、依撒格和雅各伯的天主——正如耶稣亲自作证的、活人的天主——祂且比天地更大，因为祂是二者的创造者，祂也是唯一的美善，比祂所派遣的那位更大。即使如我们所说，赫拉克利翁的想法是高超的，即整个世界是耶稣的鞋履，但我认为我们不应同意他。因为这与\BibleQuote{天是我的宝座，地是我的脚凳}这一说法如何协调？耶稣接受这句话是论及圣父的。祂说：\BibleQuote{不要指着天发誓，因为那是天主的宝座；也不要指着地，因为那是祂的脚凳}（\BibleRef{玛 5:34}）。如果他认为整个世界是耶稣的鞋履，他又如何接受经文：\BibleQuote{我岂不充满天地？}（\BibleRef{耶 23:24}）——上主说。还有必要探究：既然圣言和智慧渗透了整个世界，且圣父在圣子内，那么以上这些话究竟是应如上理解，还是应理解为：那首先被整个受造界环绕的那一位（除了圣子在他内之外），也赐给了救主——作为仅次于他的第二位，且是天主圣言——渗透整个受造界？对那些能够注意到大天穹不间断运行、如何从东到西携带如此众多星辰的人来说，最需要探究的是：那存在于整个世界的巨大且是何本质的力量是什么。因为断言那力量是圣父和圣子之外的什么，或许是不符合虔敬的。\par

% structure: p0048 source=h2 target=subsection reason=relative-heading-rank; base=h2

\subsection{二十四、若翰施洗之地的名称，不是许多抄本中的贝塔尼亚，而是贝塔巴拉。此事的证据。类似地，在猪群的故事中，应读作\Quote{革尔格萨}而非\Quote{革拉萨}。应注意圣经中的专有名词，它们常被错误书写，且对解释至关重要。}

\BibleQuote{这些事发生于伯达巴辣，在约但河对岸，若翰施洗的地方。}\BibleRef{若 1:28} 我们知道，几乎所有抄本中都可见到另一种读法：\Quote{这些事发生于伯达尼。}此外，这似乎也是更早的读法；在赫辣克伦那里，我们看到\Quote{伯达尼。}然而，我们确信不应读作\Quote{伯达尼}，而应读作\Quote{伯达巴辣。}我们曾亲临那些地方，探询耶稣及其门徒和先知的足迹。如今，正如同一位福音作者告诉我们的，伯达尼是拉匝禄、玛尔大和玛利亚的村镇；它离耶路撒冷有十五斯塔狄，而约但河距离它约一百八十斯塔狄。约但河附近也没有其他同名的地点；但人们指出，约但河岸有一处名为伯达巴辣，据说若翰曾在那里施洗。这名字的语源也与那位为主预备合用的百姓的施洗者相符；因为它的意思是预备之所，而伯达尼的意思是顺服之所。那作为使者被派在基督面前、为祂预备道路的人，难道不正是应当在预备之所施洗吗？至于玛利亚——她选择了更好的一份，这一份不会从她手中夺去，\BibleRef{路 10:41}以及为接待耶稣而操劳的玛尔大，和她们那位被称为救主朋友的兄弟，还有什么比伯达尼，即顺服之所，更适合他们的家呢？由此我们看出，凡渴望完全理解圣经的人，不可忽略仔细查考其中的专有名称。在专有名称方面，希腊抄本常常出错，在福音书中，我们可能会被它们的权威误导。关于那群被魔鬼驱赶下陡坡并淹死在海中的猪的记载，据说发生在革辣撒人的地区。然而，革辣撒是阿拉伯的一座城，附近既无海也无湖。福音作者们不会做出如此明显且可证明的错误陈述，因为他们都是详细了解犹太地一切事务的人。但我们在少数抄本中找到了\Quote{进入加达辣人的地区。}根据这一读法，必须说明加达辣是犹太地的一座城，附近有著名的温泉，但那里没有带着陡岸的湖，也没有海。而革尔盖撒——革尔盖撒人这个名称便来源于此——是靠近现在称为提庇黎雅的海的一座古城，湖边缘有一处陡峭的悬崖，据说猪就是从这里被恶魔赶下去的。革尔盖撒的意思是驱逐者之居，其中包含着对那地居民对待救主态度的预言性指涉，他们\Quote{恳求祂离开他们的境内。}在法律和先知书的多处经文中，也有类似的专名不准确之处，我们曾不辞辛劳地从希伯来人那里得知，将自己的抄本与他们的抄本（有阿基拉、特奥多提翁和辛玛库斯的译本作为印证，这些译本从未遭受窜改）进行比较。我们再加几个例子，以鼓励学者更加注意这些细节。肋未的一个儿子，即长子，在多数抄本中被称为革松，而不是革尔雄。他的名字与梅瑟的长子同名，\BibleRef{出 2:22}两者都恰当，因为这两个孩子都因在埃及寄居而生于异地。犹大的第二个儿子，\BibleRef{创 38:4}在我们这里名叫阿南，但在希伯来人那里是敖难，意为\Quote{他们的劳苦。}此外，在户纪中以色列子民的起程记载中，我们看到\Quote{他们从索科特出发，在布堂安营；}但希伯来文，在布堂的位置读作阿依曼。我又何必再列举这样的例子呢？凡愿意的人都可以亲自查考专有名称，了解实际情况。圣经中的地名尤其值得怀疑，尤其是在许多地名同时出现在一个目录中时，例如若苏厄书中分地的记载，以及编年纪上从开头一直到关于丹的段落，厄斯德拉书也是如此。名称不可忽略，因为从中可以收集到有助于解释所在段落的提示。然而，我们不能在此脱离我们正当的主题，来探讨名字的哲学。\par

% structure: p0050 source=h2 target=subsection reason=relative-heading-rank; base=h2

\subsection{二十五、约但河意为\Quote{他们的下降}。其属灵意义与应用。}

让我们现在查看面前的福音话语。\Quote{约但}意思是\Quote{他们的下降}。名字\Quote{耶勒得}在词源上与它相近，如果我可以这样说的话；它也产生\Quote{下降}的意思；因为\Person{耶勒得}是\Person{玛拉肋耳}所生，正如在\WorkTitle{哈诺客书}上所记载的——如果有人愿意把那本书视为神圣——在那些日子，天主的儿子下降到人的女儿那里。在这个下降下，有些人认为有灵魂降入身体的隐喻，把\Quote{人的女儿}这句话当作地上帐幕的比喻表达。若是如此，那么到哪条河才能得到净化呢？\Quote{他们的下降}是指哪条河呢？一条河流，不是以自己的下降，而是以\Quote{他们的}下降，即人的下降，除了我们的救主之外还能是什么呢？祂将那些从\Person{梅瑟}领受产业的人与那些通过\Person{若苏厄}获得自己份额的人区分开来。祂的波浪，在下降的河流中流淌，使天主的城欢乐，正如我们在\WorkTitle{圣咏}中所见，不是那可见的\Place{耶路撒冷}——因为它旁边没有河流——而是无可指责的天主的教会，建立在宗徒和先知的基础上，\Person{基督耶稣}我主为房角石。因此，在\Place{约但}之下，我们必须理解那成为血肉并住在我们中间的圣言，\Person{耶稣}，祂将自己所取的人性作为产业赐给我们，因为那就是头块角石，它被提升到天主圣子的神性里，因这提升而被洗净，然后接受到自己内纯净无伪的圣神的鸽子，与之结合，再也不能从中飞离。因为经上记载：\BibleQuote{你看見聖神降下，並停留在誰身上，誰就是那要以聖神施洗的人。}所以，那领受圣神停留在\Person{耶稣}本身上的，也就能以那常驻的圣神为前来的人施洗。但是\Person{若翰}在约但河彼岸施洗，在靠近\Place{犹太}外境的地区，在\Place{伯大巴辣}，作为那一位的前驱，祂来不是召义人，而是召罪人，并教导说健康的人不需要医生，而是有病的人才需要。因为这洗涤是为赦免罪过而赐予的。\par

% structure: p0052 source=h2 target=subsection reason=relative-heading-rank; base=h2

\subsection{二十六、以色列在\Person{若苏厄}带领下渡过\Place{约但}河的故事是基督之事的预象，且是为教训我们而记载的。}

现在，很可能有一位不熟悉救主各方面的人，会因上文对约旦的解释而跌倒；因为若翰说：\Quote{我用水施洗，但在我以后来的那一位比我更强，祂要用圣神给你们施洗。}对此我们回答说，作为可饮之物的天主的圣言，对某些人是水，对另一些人是酒，能悦乐人心，对另一些人又是血，因为经上说：\BibleRef{若 6:53} \BibleQuote{你们若不吃我的肉，并喝我的血，在你们内便没有生命。}而且，作为食物，祂被不同地理解为生活的食粮或肉体，同样，这同一位也是水的圣洗、圣神和火的圣洗，对某些人而言，也是血的圣洗。祂所说的最后一种圣洗，正如有些人认为的，是指祂的话：\BibleRef{路 12:50} \BibleQuote{我应受的圣洗，如何受其约束，直到完成呢？}这与门徒若望在其书信中论及圣神、水和血为一相符合。\BibleRef{若一 5:8}并且祂又宣称自己是道路和门，但显然，对那些祂是道路的人，祂并非门；对那些祂是门的人，祂也不再是道路。所以，所有那些在天主圣言开端受初学的人，前来听从在旷野里呼喊者的声音：\Quote{预备上主的道路}，这声音在约旦对岸、预备之所响起，让他们自己预备好，以便他们能借着圣神的光照，领悟那向他们启示的神性言语。现在，按我们主题的要求，汇集一切与约旦相关的，让我们来看这\Quote{河流}。天主藉梅瑟带领百姓穿过红海，使水在他们左右成为墙壁；又藉若苏厄带领他们穿过约旦。保禄处理这段圣经，他的争战不是按肉体的，因为他知道法律是属神的，应当以属神的意义理解。他向我们表明他明白关于过红海之事的说法；因为他在致格林多人第一封书信中说：\BibleQuote{弟兄们，我不愿意你们不知道，我们的祖先都曾在云下，都从海中走过，都曾在云中和海中受洗归于梅瑟，都吃过同样的神粮，都喝过同样的神饮；因为他们所饮的，是来自伴随他们的神石，那神石就是基督。}在这段经文的启发下，我们也祈祷，能从天主那里得蒙理解若苏厄穿过约旦的神性意义。关于这事，保禄也可能说：\Quote{弟兄们，我不愿意你们不知道，我们所有的祖先都穿过约旦，都曾在圣神和河里受洗归于耶稣。}而继承梅瑟的若苏厄，是耶稣基督的预像，祂承继了法律时期的安排，并以福音的宣讲取代之。即使保禄所说的人是在云中和海中受洗，他们的圣洗仍带有某些严厉和咸味。他们仍恐惧仇敌，向主和梅瑟呼号说：\BibleRef{出 14:11} \BibleQuote{难道在埃及没有坟墓，你就带我们来死在旷野吗？你为什么这样待我们，将我们从埃及领出来？}但若苏厄的圣洗，发生在颇为甘甜可饮的水中，在许多方面优于先前的圣洗，此时宗教已变得更加清晰，并呈现出合宜的秩序。因为上主我们天主的约柜由司祭和肋未人列队抬着，百姓跟随天主的仆役，也接受了圣洁的法律。因为若苏厄对百姓说：\BibleRef{苏 3:5} \BibleQuote{你们自洁，明天上主要在你们中间行奇事。}他命令司祭带着约柜走在百姓前面，其中清楚地显示了圣父关于圣子的救恩计划的奥秘，这奥秘被那赋予圣子这职位者所高举：\BibleRef{斐 2:9} \BibleQuote{致使天上、地上和地下的一切，一听到耶稣的名字，无不屈膝叩拜，并且一切唇舌无不承认耶稣基督是主，以光荣天主圣父。}这由我们在称为\WorkTitle{若苏厄书}的书中所指明的：\Quote{在那一天，我要开始在你面前，在以色列子民前，显扬你。}并且我们听到我们的主耶稣对以色列子民说：\BibleRef{苏 3:9} \BibleQuote{你们前来，听上主你们天主的话。你们由此知道，生活的天主在你们中间。}因为当我们受洗归于耶稣时，我们就知道生活的天主在我们内。并且，在先前的情况下，他们在埃及举行逾越节，然后开始旅程；但随若苏厄，在正月十日渡过约旦后，他们在\Place{基耳加耳}扎营；因为在若苏厄的圣洗之后，必须先备办一只羊，才能发出宴会的邀请。然后，以色列子民，由于那些从埃及出来者的子女未受过割损，被若苏厄用很锋利的石头行了割损；上主声明，在若苏厄圣洗的日子，当若苏厄洁净了以色列子民时，祂除去了埃及的耻辱。因为经上记载：\BibleRef{苏 5:9} \BibleQuote{上主对若苏厄说：今天我由你们身上割去了埃及的耻辱。}然后，以色列子民在月十四日举行逾越节，比在埃及更加喜乐，因为他们享用了圣地出产的无酵饼和比玛纳更好的新鲜食物。因为当他们领受预许之地时，天主没有以较逊色的食物款待他们，在像若苏厄这样的领袖领导时，他们也没有得到次等的食粮。对于默想真正圣地和天上耶路撒冷的人，这一点将是明显的。因此，在同一福音中记载：\Quote{你们的祖先在旷野中吃过玛纳，却死了；谁吃这食粮，将永远活着。}因为玛纳，虽由天主所赐，却是旅行中的食粮，供给仍在受管教者的食粮，很适合那些在师保和管家之下的人。而若苏厄从他们在地里收割的谷物所得的新食粮，即预许之地的出产——别人劳苦，他的门徒收割——那是更为充满生命的食粮，分给那些为成全而能承受祖先产业的人。因此，仍在受管教于那食粮的人，可能会因它而死亡；但已达到那随后之食粮的人，吃它便永远活着。凡此种种，我认为，都并非不适当地被加入我们对约旦河圣洗的研究中，即若翰在\Place{伯大巴喇}施行的圣洗。\par

% structure: p0054 source=h2 target=subsection reason=relative-heading-rank; base=h2

\subsection{二十七、厄里亚与厄里叟过约但河。}

另一件我们不可忽视的事是：当厄里亚要被旋风接升，仿佛升天的时候，\BibleRef{列下 2:8}他拿起自己的外衣，卷在一起击打水面，水便左右分开，他们两人——即他和厄里叟——都走了过去。他在约但河中的洗礼使他更适合被接升，因为正如我们之前指出的，保禄将如此非凡的过水称为洗礼。而通过同一条约但河，厄里叟经由厄里亚获得了他所渴望的恩赐，说：\Quote{愿你的灵双倍地临于我。}使他能领受厄里亚之灵恩赐的，或许是因为他两次经过约但河：一次是与厄里亚一起，第二次是在他接过厄里亚的外衣后，击打水面说：\Quote{厄里亚的天主在哪里？祂在哪里？他击打水面，水就左右分开。}\par

% structure: p0056 source=h2 target=subsection reason=relative-heading-rank; base=h2

\subsection{二十八、叙利亚人纳阿曼与约但河。没有其他河流具有同样的治愈力。}

若有人因为我们之前关于约但河得出的结论——即它是为我们而下降的圣言之类型——而对\Quote{他击打水}这一说法提出反对，我们反驳说，在宗徒看来，磐石明明就是基督，且这磐石被杖击打两次，好让百姓喝那跟随他们的属灵磐石。在这新难题中，\Quote{击打}指的是那些喜欢提出与结论相矛盾之事的人，他们甚至在尚未了解手头的问题是什么之前就提出异议。愿天主从这样的人中解救我们，因为一方面，祂在我们口渴时给我们水喝，另一方面，祂在浩瀚无垠的深渊中为我们预备了一条可通行的路——即通过分开祂的圣言，因为借着分辨（分开）的理性，大多数事物对我们变得明了。然而，为了让我们对这条如此甘美、如此满溢恩宠的约但河获得正确的解释，比较叙利亚人纳阿曼的癞病得洁净，以及关于以色列敌人宗教河流的记载，或许是有益的。\BibleRef{列下 5:9}记载纳阿曼带着车马来到，站在厄里叟的家门口。厄里叟派了一个使者到他那里说：\Quote{你去在约但河中洗七次，你的肉就会复原，你必得洁净。}纳阿曼就发怒；他没有看到我们的约但河能洁净那些因癞病而不洁的人，使他们脱离那种不洁，并恢复健康；这是约但河所做的，而非先知；先知的职责是指向治愈的媒介。纳阿曼不明白约但河的伟大奥秘，便说：\Quote{看啊，我以为他必定会出来见我，呼求上主他的天主的名，在患处摇手，治好这癞病人。}因为把手放在癞病上\BibleRef{玛 8:2}并洁净它，唯独属于我们的主耶稣的作为；当那癞病人怀着信心恳求他说：\Quote{主，你若愿意，就能使我洁净。}祂不仅说：\Quote{我愿意，你洁净了吧。}而且在言语之外还摸了摸他，他就洁净了。纳阿曼仍在错误中，没有看到其他河流在治愈疾病方面远逊于约但河；他夸耀大马士革的河流，阿巴纳和帕尔帕尔，说：\Quote{阿巴纳和帕尔帕尔，大马士革的河流，不是比以色列所有的水更好吗？我难道不能在那里洗而得洁净吗？}因为除了独一的圣父天主之外，没有良善的；同样，在河流中，除了约但河之外，没有良善的，也没有能洁净癞病的，除非那人怀着信心在耶稣里洗涤自己的灵魂。我想，这就是为什么以色列人记载说，他们坐在巴比伦的河边，一想起熙雍就哭了；那些因自己的邪恶而被掳的人，在尝过神圣的约但河之后，再尝别的河水，便怀着渴望想起自己救恩的河流。因此，论到巴比伦的河流说：\Quote{我们曾坐在那里，}显然是因为他们无法站立，\Quote{并且哭了。}耶肋米亚斥责那些想喝埃及的水，而离弃那从天降下、并因其降下而得名——即约但河——的水的人。他说：\Quote{你在埃及的路上做什么？要喝基红的水，要喝大河的水？}或如希伯来文所说，\Quote{要喝息红的水？}我们现在要讲的就是这些水。\par

% structure: p0058 source=h2 target=subsection reason=relative-heading-rank; base=h2

\subsection{二十九、埃及的河及其龙，与约但河对比。}

然而，圣神在默感圣经中所说的，主要不是指肉眼可见的河流，这可以从厄则克耳论埃及王法郎的预言中看出：\Quote{看哪，我与你为敌，法郎埃及王，你这卧在河中的大龙，你说：\InnerQuote{我的河是我的，是我造的。}我必把钩子放进你的腮，使河中的鱼贴在你的鳞上，将你从你的河中拉上来，连同河中的一切鱼；我必急速将你和河中的一切鱼丢下；你必倒在你的地面上，你不得聚集，也不得被妆饰。}因为有什么真实的肉身之龙曾被报道出现在埃及的物质河流中？但请思想：埃及的河岂不是那与我们为敌之龙的住处吗？这龙甚至不能杀死婴儿梅瑟。然而，龙在埃及的河中，天主却在使天主的城欢乐的河中；因为圣父在圣子内。因此，那些来在祂内洗濯的人，摆脱了埃及的羞辱，更适合被恢复。他们从那最污秽的癞病中得洁净，领受双倍的属灵恩赐，并预备好领受圣神，因为属灵的鸽子不降在任何其他水流上。这样，我们以更相称于神圣主题的方式，思考了约但河及其中的净化，以及在约但河中被洗的耶稣和预备的住所。那么，让我们从河里汲取我们所需要的帮助吧。\par

% structure: p0060 source=h2 target=subsection reason=relative-heading-rank; base=h2

\subsection{三十、论玛利亚探访山地中的依撒伯尔时若翰从耶稣所学到的事。}

\BibleQuote{第二天，若翰见耶稣向他走来。} \BibleRef{若 1:29} 耶稣的母亲先前一怀孕，就与若翰的母亲同住，那时若翰的母亲也怀了孕；那先在者便将自己形像精确地传达给那成形者，使他符合自己的光荣。由于这外在的相似，那些不能分辨形像本身与按照形像之事物的人，便认为若翰是基督 \BibleRef{路 3:14}，而耶稣被认为是若翰从死者中复活 \BibleRef{玛 14:2}。因此现在，在若翰为他所作并我们已考察过的那些见证之后，耶稣亲自被洗者若翰看见向他走来。值得注意的是，在先前，当玛利亚问安的声音传到依撒伯尔耳中时，胎儿若翰在他母亲腹中跳跃，而他的母亲好像从地上领受了圣神。因为经上记载 \BibleRef{路 1:41} \BibleQuote{依撒伯尔一听到玛利亚问安，胎儿就在她腹中跳跃；依撒伯尔就充满了圣神，且大声呼喊说，} 等等。同样，这次，若翰看见耶稣向他走来，就说：\BibleQuote{请看，天主的羔羊，除免世罪者。} 因为关于重大的事，人首先通过听闻受教，然后亲眼看见。若翰所穿戴的形状，是由仍在他母亲腹中成形的主帮助形成的，当主接近依撒伯尔时，这对任何理解我们证明的人都是清楚的，即若翰是声音，耶稣是圣言，因为当依撒伯尔因玛利亚问安而充满圣神时，她内有一个大声，正如经文字句本身所表明的；因为经上说：\BibleQuote{她大声呼喊说。} 显然，依撒伯尔这样做，\BibleQuote{她说话。} 因为玛利亚问安的声音传入依撒伯尔耳中，使若翰充满了自己；因此若翰跳跃，他的母亲成了她儿子的口和一位女先知，大声呼喊说：\BibleQuote{你在妇女中应受赞美，你胎中的孩子也应受赞美。} 现在我们清楚看见玛利亚匆忙赶到山地，进入匝加利亚的家，并向依撒伯尔问候，这一切是为了将她从所怀的那位那里获得的一些能力传递给还在母亲腹中的若翰，并且若翰也将一些临到他的先知恩宠传递给他的母亲。最恰当的是，这些事发生在山地，因为低微之人不能容纳任何伟大的事。因此，在若翰的这些见证之后——第一次，当他呼喊并谈论祂的神性；第二次，对从耶路撒冷由犹太人派来的司祭和肋未人；第三次，作为对来自法利塞人的更尖锐问题的回答——耶稣被见证人看见向他走来，而此时他仍在进步和变得更好。这种进步和改善在\BibleQuote{次日}这句话中象征性地表明。因为耶稣随之而来的光照中，以及在先前事件之后的那一天，不仅被知道是站在那些不认识祂的人中间，而且现在被清楚看见向从前宣讲过祂的人走来。第一天，见证发生；第二天，耶稣来到若翰那里；第三天，若翰同他的两个门徒站在那里，看见耶稣走过，就说：\BibleQuote{请看，天主的羔羊，} 从而催促那里的人跟随天主子。第四天，祂也意欲前往加里肋亚，那出来寻找迷失者的祂找到了斐理伯，并对他说：\BibleQuote{跟随我。} 在那一天之后，即第四天之后，也就是我们所列举开始后的第六天，在加里肋亚的加纳有婚宴，我们将在论及那段落时加以考虑。也请注意这点：身为更大的玛利亚来到较小的依撒伯尔那里，天主子来到洗者若翰那里；这应鼓励我们毫不迟疑地帮助那些地位较低的人，并为自己培养谦逊的品性。\par

% structure: p0062 source=h2 target=subsection reason=relative-heading-rank; base=h2

\subsection{三十一、论若翰与耶稣在受洗时的交谈，仅玛窦记载。}

门徒若望没有告诉我们救主从何处来到洗者若翰那里，但我们从玛窦得知，他写道：\BibleQuote{那时，耶稣从加里肋亚来到约但河若翰那里，为受他的洗。} 玛尔谷加上了加里肋亚的地点；他说：\BibleQuote{在那些日子里，耶稣从加里肋亚的纳匝肋来，在约但河里由若翰受了洗。} 路加没有提及耶稣从哪里来，但他告诉了我们从其他福音中不知道的事：即在受洗后，当祂上来时，天为祂开了，圣神以形体如同鸽子降在祂身上。此外，只有玛窦告诉我们若翰阻止主，对救主说：\BibleQuote{我需要受祢的洗，而祢却来就我吗？} 其他福音作者在玛窦之后没有添加这些，以免重复同样的事。主回答的话：\BibleQuote{你暂且容许吧，因为我们应当这样履行全义。} 也只有玛窦记载了。\par

% structure: p0064 source=h2 target=subsection reason=relative-heading-rank; base=h2

\subsection{三十二、若翰称耶稣为\Quote{羔羊}。为何他特别提到这种动物？论一般祭献的预像。}

他说：\BibleQuote{请看，天主的羔羊，除免世罪者。} \BibleRef{若 1:29} 有五种动物被带到祭台前，三种行走的，两种飞翔的；似乎值得一问的是，为何若翰称救主为羔羊，而不称这些其他受造物之一；以及，为何当每种行走的动物都有三种时，他对羊类使用\BibleQuote{羔羊}一词。五种动物如下：公牛、绵羊、山羊、斑鸠、鸽子。至于行走的动物，有三种——公牛、牛、牛犊；公绵羊、绵羊、羔羊；公山羊、山羊、山羊羔。至于飞翔的动物，我们只听说鸽子要两只雏鸽；斑鸠只要一对。因此，凡想准确理解祭献的属灵理据之人，必须探究这些事物是哪些天上事物的模型与阴影，以及每样牺牲的祭献是为了何种目的，他必须特别收集与羔羊相关的要点。祭献的原则必须参照某些天上奥迹来理解，这从宗徒的话中可以看出，他在某处说：\BibleQuote{他们供奉的事，不过是天上事物的模型与阴影。} 又说：\BibleQuote{既然天上的模型必须用这些祭品来洁净，那天上的本物，就必须用比这些更美的祭献了。} 要探讨这一切的细节，并说明在耶稣基督内所实现的属神法律之祭献与它们的关系（这是人性难以理解的真理），除非是成全的人，\BibleRef{希 5:14} 他因使用而感官习练了分辨善恶，并能以爱真理的心志说，\BibleRef{格前 2:6} \BibleQuote{我们在成全的人中讲论智慧。} 关于这些事以及类似的事，我们可以引用\BibleRef{出 29:38}的话说：\BibleQuote{今世的首领中没有一个知道。}\par

% structure: p0066 source=h2 target=subsection reason=relative-heading-rank; base=h2

\subsection{三十三、在清晨和傍晚的祭献中献上一只羔羊。其意义。}

我们发现羔羊是作为日祭（每日祭献）而献上的。正如经上所记：\BibleRef{出 29:38} \BibleQuote{这是你当在祭台上献的：两只一岁的羔羊，每天不断，作为日常祭献。早晨献一只羔羊，傍晚献另一只羔羊；还要献十分之一细面，与捣成的油四分之一辛混合，作为素祭；又献酒四分之一辛，作为奠祭，献给第一只羔羊。傍晚你献第二只羔羊，按照早晨的素祭和奠祭，作为馨香的祭献，献给上主，作为你们世世代代在会幕门口、在上主面前的常燔祭，在那里我要与你们相会，对你们说话。我要在那里为以色列子民指定你，并要在我的荣耀中被尊为圣，且要以我的荣耀使会幕成为圣。} 但是，在理性人的心灵世界中，还有什么其他的日常祭献，除了那渐趋成熟的圣言——即象征着被称为羔羊、在灵魂获得光照时献上的那一位？这就是早晨的日常祭献；当心灵与神圣事物的交往结束时，它又再次被献上。因为心灵不能恒久维持与更高事物的交流，因为灵魂被注定要与属于尘世且沉重的身体结合在一起。\par

% structure: p0068 source=h2 target=subsection reason=relative-heading-rank; base=h2

\subsection{三十四、圣人在其思想生活中的清晨和傍晚祭献。}

如果有人问，圣人在清晨与傍晚之间该做什么，就让他跟随礼仪中的做法，从中推断他所要求的原则。在礼仪中，司祭们从日常祭献开始，在傍晚的日常祭献之前，他们献上法律所规定的其他祭献，例如赎罪祭、误犯之祭、求拯救的祈祷、嫉妒祭、安息日祭、月朔祭等等，不可胜数。同样，我们从那代表基督之典型的言论开始献祭，然后可以继续论述许多其他至为有益的事物。当我们结束于基督之事时，仿佛到了傍晚和夜晚，即我们抵达祂显现的肉身层面。\par

% structure: p0070 source=h2 target=subsection reason=relative-heading-rank; base=h2

\subsection{三十五、耶稣在其人性方面是羔羊。}

若我们进一步探讨耶稣被若翰指出的意义，当他说道：\BibleQuote{请看，天主的羔羊，除免世罪者。} 我们可以立足于天主圣子降生成人、居于人间的身体性来临的经世中，在这种情况下，我们将羔羊理解为不是别的，正是那个人。因为那人\BibleQuote{被牵去宰杀，如同羊羔；又像羔羊在剪毛者前缄默不语，} \BibleRef{依 53:7} 又说：\BibleQuote{我好像一只驯良的羔羊被牵去宰杀。} \BibleRef{耶 11:19} 因此，在\WorkTitle{默示录}中，我们看到一只羔羊，仿佛被宰杀后站立着。根据某些隐秘的理由，这被宰杀的羔羊已成为全世界的净化，为了祂因圣父对世人的慈爱而接受了死亡，用自己的血将我们从那因罪而被贩卖、掌控我们的人手中赎买回来。带领这只羔羊去宰杀的是天主在人间，即伟大的大司祭，正如祂的话所显示的：\BibleRef{若 10:18} \BibleQuote{没有人夺去我的生命，是我自己舍弃的。我有权舍弃，也有权再取回来。}\par

% structure: p0072 source=h2 target=subsection reason=relative-heading-rank; base=h2

\subsection{三十六、殉道者的死亡被视为祭献，以及它如何有益于他人。}

与这祭献相似的是法律之祭献所象征的其他祭献；在我看来，还有一种祭献也属于同一性质，即高贵的殉道者流血，门徒\Person{若望}看见他们站在天上的祭台旁，这并非毫无意义。\Quote{谁明智，\BibleRef{欧 14:10} 就会明白这些事；谁聪敏，就会知道它们？} 稍微思考那些为所献之人洁净的祭献的理据，是一件更高的思辨之事。\Person{依弗大}献其女的祭献应受关注；正是藉此誓愿他征服了阿孟子民，而牺牲者认可了他的誓愿，因为当她的父亲说：\Quote{\BibleRef{民 11:35} 我已向主开口论到你，} 她回答说：\Quote{你既已向主开口论到我，就照你所许的去做吧。} 这故事暗示，那神明必然非常残忍，竟为世人的救恩而接受这样的祭献；我们需要有相当的胸襟和解决对天主眷顾所提疑难的能力，才能解释这些事，并看出它们是奥秘，超越我们的人性。然后我们要说：\Quote{\BibleRef{智 17:1} 天主的判断伟大，难以描述；为此未受教化的灵魂迷失了。} 在外邦人中，也记载许多人在瘟疫爆发时，将自己献为公共利益的牺牲。忠信的\Person{克莱孟}依据这些叙述假定确有其事，\Person{保禄}为他作证说：\Quote{\BibleRef{斐 4:3} 也请与克莱孟以及其余的同工，他们的名字在生命册上。} 如果有人因存心挑剔向少数人启示的奥秘，而觉得这些叙述中有任何不合宜之处，那么同样的困难也归于委托给殉道者的职分，因为天主的旨意是，我们宁愿忍受因宣认祂为天主而带来的一切可怕辱骂，也不愿因用言语顺从真理仇敌的意愿而暂时逃脱这些（人所视为恶的）痛苦。因此，我们确信邪恶势力因圣殉道者的死亡而受挫败；仿佛他们的忍耐、他们至死不移的宣认、以及他们对虔敬的热忱，钝化了邪恶势力向受苦者进攻的锋芒；他们的力量因此削弱耗尽，许多曾被他们征服的人抬起头来，从邪恶势力先前压迫和伤害他们的重负中获得解放。甚至殉道者本身也不再受苦，即使那些先前加害他人的势力并未耗尽；因为献上如此祭献的人战胜了敌对势力，我可以用一个适合本主题的比喻来说明。谁杀死一只毒兽，或用咒语使它沉睡，或以任何方式消除其毒液，他就造福了许多人，否则这些人就会因那毒兽而受害，若非它被杀死、或被咒语催眠、或被除去毒液。此外，如果曾受咬伤的人得知此事，并治愈了疾病、目睹那伤害他的动物的死亡，或践踏它、或触摸已死的它、或尝其一部分，那么，这先前受苦的人就会将痊愈和益处归于杀死毒兽的人。我们必须以类似的方式理解至圣殉道者的死亡如何运作，许多人因它而受益，其影响是我们无法描述的。\par

% structure: p0074 source=h2 target=subsection reason=relative-heading-rank; base=h2

\subsection{三十七、论基督死亡的效果、祂死后的凯旋、以及祂的死除去人的罪恶。}

我们已对殉道者这一主题以及因瘟疫而死者之记述流连许久，因为这让我们看见那被牵到宰杀之地、在剪毛人面前默不作声如羔羊者的卓越。因为，若在希腊人的这些故事中有任何意义，若我们关于殉道者所说的话有根据——宗徒们也为同样的理由成了世上的污秽、万物中的渣滓，\BibleRef{格前 4:13}——那么，对天主的羔羊，那正是为此而被祭献，为能除掉不只是一少部分人而是整个世界的罪，且祂也因此而受苦，又当如何说，该说何等伟大之事？若有人犯罪，我们读到，\BibleRef{若一 2:1} \BibleQuote{在父那里我们有一位中保，就是那义者耶稣基督；祂为我们的罪作了赎罪祭，不仅为我们的，也为全世界的。} 因为祂是众人的救主，\BibleRef{弟前 4:10} 特别是信徒的。祂\BibleRef{哥 2:14} 用祂自己的血涂抹了那反对我们的债卷，并把它撤去，以致连一个痕迹，连我们被涂抹之罪的痕迹，都找不到，且把它钉在十字架上；祂解除了率领者和掌权者的武装，公然羞辱他们，仗着十字架向他们夸胜。我们受教导要因在世上的苦难而喜乐，知道我们喜乐的理由是：世界已被征服，并已明确地臣服于它的征服者。于是，这位救主，以谦卑自下羞辱了那毁谤者，与可见的太阳一起，在祂那比喻性地称为月亮的荣光教会之前，世世代代常存。祂以苦难摧毁了仇敌，祂是战斗中的勇士、大能的主，在祂大能的作为之后，需要一种唯有祂圣父才能赐予祂的净化；这就是为什么祂禁止玛利亚触摸祂，说：\BibleRef{若 20:17} \BibleQuote{不要摸我，因为我还没有升到圣父那里去；但你去告诉我的门徒们，我要升到我的圣父和你们的圣父那里，到我的天主和你们的天主那里。} 当祂满载胜利和战利品，带着从死者中复活的躯体来临——因为我们在\Quote{我还没有升到圣父那里去}和\Quote{我要升到我的圣父那里去}这些话中还能看出别的什么意思呢？——那时，某些权势说：\Quote{那从厄东来、从波责辣来、穿着红衣服、这美丽者是谁？} \BibleRef{依 63:1} 然后护送祂的对那些在天上门上的说：\Quote{众城门哪，抬起你们的头来！永久的门户，你们要被举起！荣耀的君王将要进来！} 但他们又看见祂的右手似乎染着血，全身布满祂英勇的痕迹，问道：\Quote{你的衣服为何是红色的？你的衣服为何像踹酒醡的呢？} 祂回答说：\Quote{我踹了它们。} 为此，祂需要把\BibleRef{创 49:2} \BibleQuote{祂的衣服在酒中洗净，祂的袍子在葡萄汁中洗净。} 因为当祂担当了我们的软弱，背负了我们的疾病，承当了全世界的罪，并且赐福给这么多人之后，或许祂接受了那比人类所能想像的任何洗礼都更大的洗礼，我想当祂说：\BibleRef{路 12:50} \BibleQuote{我有当受的洗，我是何等迫切直到它完成！} 就是指此而言。我在此大胆探究，并挑战大多数作者提出的观点。他们说那最大的、无法想像有更大者的洗礼就是祂的苦难。但若如此，为什么祂在受难之后对玛利亚说\Quote{不要摸我}？祂本该在她触摸时将自己献给她，因为祂已通过苦难接受了完美的洗礼。但若如我们先前所说，在祂对仇敌行了所有英勇事迹之后，祂需要洗净\Quote{祂的衣服在酒中，祂的袍子在葡萄汁中}，那么祂正往真葡萄树的栽培者——圣父那里去，以便在那里洗净，并升到高处后，掳掠了俘虏，降临带来多样的恩赐——如火焰的舌头，分开分给宗徒们，以及圣天使们，要与他们在每项行动中同在并拯救他们。因为这些安排之前，他们尚未被洁净，天使不能与他们同住，因为天使或许也不愿与那些没有预备自己、未被耶稣洁净的人在一起。因为唯独出于耶稣的良善，祂与税吏和罪人一同吃喝，容许那忏悔的罪妇用眼泪洗祂的脚，甚至为不敬虔的人降卑至死，祂不以自己与天主同等为强夺的，反倒虚己，取了奴仆的形像。在成就这一切时，祂所成就的更是那为罪人舍弃祂的圣父的旨意，而非祂自己的旨意。因为圣父是良善的，而救主是祂良善的肖像；祂在一切事上向世界行善，因为天主在基督里使世界与祂自己和好，这世界从前因邪恶而与祂为敌，祂按次序和顺序成就祂的善行，并不立刻将所有仇敌放在祂脚下。因为圣父对祂，我们众人的主说：\Quote{你坐在我的右边，等我使你的仇敌作你的脚凳。} 这要一直持续到最后的仇敌——死亡被祂消灭。若我们思考这臣服于基督的意思，并主要从这句话找到解释：\BibleRef{格前 15:26} \BibleQuote{当一切都被服在祂以下时，那时子自己也要服在那叫万有服祂的以下，} 那么我们就会看到这一概念如何与万有的天主的良善相符，因为那正是除去世人罪孽的天主羔羊的观念。然而，并不是所有人的罪都被天主的羔羊除去，不是那些不忧伤、不受苦直到罪被除去之人的罪。因为荆棘不仅刺入，而且深深扎根在每个被邪恶灌醉、失去清醒的人的手中，正如箴言所说：\Quote{荆棘长在醉汉的手中}，而它们给那在灵魂实质中容许这样生长的人带来何等痛苦，甚至难以言说。谁若容许邪恶如此深植于他的灵魂，以致成为布满荆棘的土地，他必被天主那活泼而有力、比双刃剑更锋利、比任何火更灼热的话语所砍伐。对于这样的灵魂，必须降下那寻见荆棘的火，以其神圣的德能停留在荆棘所在之处，却不烧毁禾场或站立的庄稼。但那除去世界之罪、并以自己的死亡开始这样做的羔羊，有几种方式，其中一些大多数人都能清楚理解，另一些则对大多数人隐藏，只有那些配得上神圣智慧的人才知道。为何我们要数算我们人类相信的所有方式？那是每个活在肉身中的人都能亲自看见的事。在我们相信的这些方式中，罪也被除去：通过苦难、邪灵、危险的疾病和沉重的病痛。谁知道其后还有什么？我们所说的一切并非多余——我们不能忽视那与\Quote{请看，天主的羔羊，除去世人罪孽的}这些话如此明显相连的思想，因此必须较为仔细地关注我们主题的这一部分。这使我们看到，天主以祂的忿怒定罪一些人，以祂的怒气管教他们，因为祂对人的爱如此之大，以致祂不会让任何人不受定罪和管教；因此我们应尽力而为，以免因最严厉的试炼而遭受这样的定罪和管教。\par

% structure: p0076 source=h2 target=subsection reason=relative-heading-rank; base=h2

\subsection{三十八、其所除之罪的世界，被称为教会。不同意此意见的理由。}

读者最好思考一下上文从各方面对圣经中\Quote{世界}一词的涵义所作的说明和举例；我认为有必要重复这一点。我知道有位学者认为，\Quote{世界}仅指教会，因为教会是世界的装饰，且被称为世界的光。\Quote{你们}，他说，\BibleRef{玛 5:14}，\BibleQuote{是世上的光。}如今，世界的装饰是教会，基督是她的装饰，祂是世界最初的光。我们必须思考，基督是否也被称为祂门徒所同属的那个世界的光。当基督是世界的光时，或许意指祂是教会的光；但当祂的门徒是世界的光时，或许他们就是其他呼求主之名者的光，即教会之外的其他呼求者，正如保禄在\BibleBook{格林多}{前书}开头论及这一点时所写的：\BibleQuote{致天主的教会，以及所有在各地呼求我们主耶稣基督之名的人。}若有人认为教会被称为世界的光，意指其余的人种包括不信者，这若从预言性和神学角度理解可能是正确的；但若指现在而言，我们提醒他：一物的光照亮该物，并请他指出其余的人种如何因教会在世上的临在而被光照。若持此见解者不能证明这一点，就让他们思考我们的解释是否不成立：光就是教会，世界就是那些呼求圣名的人。\BibleBook{玛窦}{福音}上述经文之后的语句，会指引细心的探究者得出正确的解释。经上说：\BibleQuote{你们是地上的盐}；其余的人被设想为地，信友是他们的盐；正因为信友信德坚固，大地才得以保存。因为如果盐失了味，不能再腌渍和保存大地，结局就来临了；因为很明显，如果罪恶增多，许多人的爱情冷淡，\BibleRef{玛 24:12} 正如救主亲自对那些将目睹祂来临的人表示怀疑，说：\BibleRef{路 18:8} \BibleQuote{人子来的时候，能在世上找到信德吗？}那时世代的结局就会来到。假设教会被称为世界，因为救主的光照耀着它——我们必须在经文\BibleQuote{请看天主的羔羊，除免世罪者}的语境中问，此处的\Quote{世界}是否应从理智上理解为教会，而除罪只限于教会。既然如此，我们该如何理解同一位门徒关于救主作为罪之赎罪祭所言的话？我们读到：\BibleQuote{若有人犯罪，我们在圣父那里有一位中保，就是那义者耶稣基督；祂为我们作了赎罪祭，不仅为我们的罪，也为整个世界的罪。}保禄的言论在我看来也有同样的效果，他说：\BibleRef{弟前 4:10} \BibleQuote{祂是众人的救主，尤其是信徒的救主。}再者，赫拉克利翁在处理我们这段经文时断言，没有任何证据或引证证人，\Quote{天主的羔羊}这话是若翰以先知身份说的，而\Quote{除免世罪者}这话是若翰作为高于先知者说的。他认为前者用于祂的身体，后者用于那在身体内的祂，因为羔羊是羊类中不完善的成员；同样，身体比起住在其中的那一位也是如此。若他本意是赋予身体完美性，他就会提到一只将要被祭献的公绵羊。在上文仔细讨论之后，我认为没有必要就此段落重复论述，或反驳赫拉克利翁的轻率言论。只需指出一点：既然这世界几乎不能容纳那位自我空虚者，就需要一只羔羊而非一只公绵羊，好使它的罪被除免。\par

\SourceRecord{Translated by Allan Menzies.；Ante-Nicene Fathers；Vol. 9.；Edited by Allan Menzies.；Buffalo, NY: Christian Literature Publishing Co.,；1896.；<http://www.newadvent.org/fathers/101506.htm>.}

% structure: p0001 source=h1 target=section reason=page-title

\section{\WorkTitle{若望福音释义（卷十）}}

% structure: p0002 source=h2 target=subsection reason=relative-heading-rank; base=h2

\subsection{一、耶稣来到葛法翁。四圣史关于此事的陈述。}

\BibleQuote{此后（\BibleRef{若 2:12}）祂和祂的母亲、兄弟及门徒下到葛法翁，在那里住了不多几天。犹太人的逾越节近了，耶稣就上耶路撒冷去。祂在圣殿里看见有卖牛、羊、鸽子的，和坐在那里兑换银钱的人，便用绳索做了一条鞭子，把众人连羊带牛都从圣殿赶出去，倒出兑换银钱之人的钱币，推翻他们的桌子，并对卖鸽子的人说：把这些东西拿出去，不要将我父的圣殿当作商场。祂的门徒就想起经上记载的话：对祢殿的热忱将吞噬我。犹太人于是问祂说：祢行这些事，给我们显什么记号呢？耶稣回答说：你们拆毁这座圣殿，三天之内我要把它重建起来。犹太人说：这圣殿建造了四十六年，祢三天之内就能重建起来吗？但耶稣所说的是以祂的身体为圣殿。所以当祂从死者中复活以后，祂的门徒就想起祂说过这话，便相信了圣经和耶稣所说的话。耶稣在耶路撒冷过逾越节庆日时，许多人看见祂所行的神迹，就信了祂的名。但耶稣却不信任他们，因为祂认识所有的人，也不需要谁对人有见证，因为祂自己知道人心里有什么。}\par

在名为\WorkTitle{户籍纪}的书中记载的数字，是按照某种与各事物比例相适应的原则在圣经中占有一席之地的。因此，我们应当探讨，梅瑟的\WorkTitle{户籍纪}是否——如果我们能探究出来——以某种特殊方式教导了我们关于此事的原理。我在第十卷书的开头向你们作此说明，因为我观察到，在圣经的许多章节中，数字十享有特殊的恩宠，你们可以仔细考虑，是否有理由希望本卷将为你们从天主那里带来某种特别的益处。为使此事成真，我们将尽力完全投身于天主，祂乐于赐予祂最精选的恩赐。本书开头的话是：\BibleQuote{此后祂和祂的母亲、兄弟及门徒下到葛法翁，在那里住了不多几天。}其他三位圣史说，主在与魔鬼的较量后，去了加里肋亚。玛窦和路加表示祂先是到了纳匝肋，然后离开那里，来到葛法翁居住。玛窦和玛尔谷也陈述了祂离开的一个原因，即祂听说若翰被监禁。他们的话如下：玛窦说：\BibleQuote{于是魔鬼离开了祂，看，有天使前来伺候祂。但当祂听说若翰被交付以后，就退往加里肋亚，离开纳匝肋，来住在海边葛法翁，在则步隆和纳斐塔里的境内，这应验了依撒意亚先知所说的话：则步隆地，纳斐塔里地；}在引用依撒意亚之后，接着说：\BibleQuote{从那时起，耶稣开始宣讲说：你们悔改吧！因为天国临近了。}玛尔谷这样说：\BibleQuote{祂在旷野里四十天四十夜，受撒殚的试探，与野兽在一起，并有天使伺候祂。但若翰被交付以后，耶稣来到加里肋亚，宣讲天主的福音，说：时期已满，天主的国临近了，你们悔改，信从福音吧！}然后，在记述了安德肋、伯多禄、雅各伯和若望之后，玛尔谷写道：\BibleQuote{祂进入了葛法翁，一到安息日，祂就在会堂里教训人。}路加说：\BibleQuote{魔鬼退去，暂时离开了祂。耶稣因圣神的德能回到加里肋亚，祂的名声传遍了整个周围地区。祂在他们的会堂里教训人，受到众人的光荣。祂来到了祂长大的纳匝肋，照祂的习惯，在安息日那天进了会堂。}然后路加记述了祂在纳匝肋所说的话，以及会堂里的人们如何对祂发怒，把祂赶出城，带到他们城建造的山崖上，要把祂推下去，但主却从他们中间走过，离开了那里；接着，他写道：\BibleQuote{祂下到加里肋亚的葛法翁城，在安息日教训他们。}\par

% structure: p0005 source=h2 target=subsection reason=relative-heading-rank; base=h2

\subsection{二、若望与前三福音在此段叙事中的分歧：按字面读，叙事不能协调；必须作属灵的诠释。}

这些事情的真理必定在于心灵所见。如果福音之间的差异得不到解决，我们就必须放弃对福音的信任——认为它们是真实的、由神圣之灵所写，或是值得信赖的记录，因为这两种特征都被认为属于这些著作。接受四福音、却不认为它们表面的差异应当以寓意法（即神秘解释）来解决的人，必须澄清上文提出的关于四十天试探的难题，这段时期在若望的叙述中无论如何都找不到位置；他们还必须告诉我们，主何时首次来到葛法翁。如果是在祂受洗后的六天之后（第六天是加里肋亚加纳的婚宴），那么显然试探从未发生，祂也从未到过纳匝肋，若翰也尚未被解送。而在葛法翁——祂在那里住了不多几日——之后，犹太人的逾越节近了，祂便上耶路撒冷去，在那里把羊和牛从圣殿里赶出去，倒出钱庄的零钱。在耶路撒冷，似乎尼苛德摩——首领和法利塞人——在夜间首次来到祂那里，并听到了我们在福音中可以读到的话。\Quote{这些事以后，\BibleRef{若 3:23}耶稣和祂的门徒来到犹太地，在那里同他们住下并施洗，同时若翰也在厄农靠近撒冷施洗，因为那里水多，众人都来受洗；而若翰尚未被下狱。}此时，若翰的门徒与犹太人关于洁净的话题起了争论，他们来到若翰那里，论及救主说：\Quote{看，祂施洗，众人都到祂那里去了。}他们曾从施洗者口中听到话，这些话的确切措辞最好从圣经本身获取。现在，如果我们问基督何时首次在葛法翁，我们的回应者若跟随玛窦和其他两位作者的话，会说是在试探之后，当时\Quote{祂离开纳匝肋，来住在海边葛法翁。}但他们如何证明两段陈述都是真实的呢？一方面，玛窦和马尔谷说，祂是因听见若翰被解送才退往加里肋亚；另一方面，若望在那里记载，在祂住在葛法翁之后、在祂上耶路撒冷以及从那里往犹太地去之后许多其他事情发生之时，若翰尚未被下狱，仍在厄农靠近撒冷施洗。还有许多其他的点，细心的福音研读者会发现它们的叙述并不一致；我们将尽己所能，随其出现，一一摆放在读者面前。研读者因思虑这些事而困惑，要么会放弃找出所有福音都是真实的尝试，不敢断定所有关于主的信息都不可靠，而随意选择其中一本作为引导；要么会接受四福音，并认为它们的真理不应在外在的、物质的文字中寻求。\par

% structure: p0007 source=h2 target=subsection reason=relative-heading-rank; base=h2

\subsection{三、我们该如何看待不同福音之间的差异？}

然而，我们必须设法了解福音作者在这类事情上的意图，我们转向这个问题。假设有几个人，藉着圣神看见了天主，知道了祂向祂的圣人所说的话，以及祂按拣选之时为了他们的进步而赐予他们的临在。有这样几个人，他们在不同的地方，从上头领受的恩惠在形状和性质上各不相同。让他们各自分别报告他们在圣神内关于天主、祂的话以及祂向圣人显现的事，这样其中一人报告天主在某地对某一义人的显现、话语和作为，另一人报告关于主的其他神谕和伟大工程，第三人报告与前两人不同的事。再假设有第四人，就某一特定事项做了与前三人同类的事。假设这四人对于圣神启示给他们所有人的某件事彼此一致，并且除了那件事之外还简短报告了其他事；那么他们的叙述大概会如此发生：天主在某时某地显现在某人面前，如此这般地对待了他；仿佛祂以某种形状显现，并牵手领他到某地，然后如此这般地对待他。第二人会报告，就在上述事件发生的同时，天主在某一个被提及的城镇，显现在一个有名有姓的第二人面前，地点与前面叙述相距甚远，并且他会报告在同一时间对这第二人所说的一套不同的话。让我们假设第三人和第四人也是这样。而且，正如我们所说，这些见证人——他们报告关于天主的真事以及祂赐予某些人的恩惠——要在他们报告的某些叙述中彼此一致。那么，凡把这些人的著作当作历史、或当作通过历史图像对真实事物的再现，并认为天主在空间上局限于一定范围、无法在同一时间向不同地方的几个人呈现几个不同的异象、或同时发表几篇演说的人，都会认为我们的四位作者不可能都讲真话。对他而言，由于天主在空间上有限，祂不可能在同一规定时间对一人说一件事而对另一人说另一件事，也不可能同时做一件事和相反的事，直白地说，就是不可能同时坐着又站着——假设一位作者将祂描述为届时站着，在某地向某人发表某篇讲话，而另一位作者却描述祂坐着。\par

% structure: p0009 source=h2 target=subsection reason=relative-heading-rank; base=h2

\subsection{四、圣经包含许多矛盾，以及许多并非字面真实的陈述，但必须以属神和奥秘的方式阅读。}

在我假定的情况下，即史家们想要通过图像教导我们他们在心中所见之物，那么，如果这四位史家有智慧，他们的含义便会显得毫无分歧；我们必须明白，四位福音作者的情况也是如此。他们为了自己的目的，充分利用了耶稣藉发挥其奇妙非常的能力所行的事；他们也同样运用祂的言论，并在某些地方，用显然暗示感官事物的语言，把纯粹以理智方式向他们显明的事物附接到他们的写作中。即使他们有时随意处理那些从历史眼光看发生得不同的事，并改变它们以服务于他们心目中的奥秘目的，我也不谴责他们；以至于把在某地发生的事说成仿佛发生在另一地，或在某时发生的事说成仿佛发生在另一时，并在以某种方式说出的话中引入他们自己的一些改动。他们立意要在可能的地方既从物质上也从精神上讲出真理，在不可能的地方，他们则意图以精神优先于物质。精神真理常常，正如有人所说，保存在物质的虚假中。例如，我们可判断关于雅各伯和厄撒乌的故事。\EditorialNote{创世纪二十七章}雅各伯对依撒格说：\Quote{我是厄撒乌你的长子。}从精神上讲，他说的确是真理，因为他已分享了长子的权利，而这些权利在他兄弟身上正在消逝；他披上山羊皮，取用了厄撒乌的外貌，除了赞美天主的声音外，他就是厄撒乌，好让厄撒乌后来也能找到接受祝福的地方。因为如果雅各伯没有作为厄撒乌而受祝福，厄撒乌或许也不能接受属于他自己的祝福。耶稣也是如此，按照对祂的观念，祂有许多方面，很可能福音作者们领受了不同的看法；但在他们各自所写不同的事上，他们仍然彼此一致。关于我们的主，有在言辞上相互矛盾的说法，即祂是达味的后裔，又不是达味的后裔。\Quote{祂是达味的后裔}这一说法是真实的，正如宗徒所说：\BibleRef{罗 1:3}\BibleQuote{按肉身说，祂生于达味的后裔。}如果我们将此应用于祂的肉身部分；但若我们将祂生于达味的后裔理解为祂的神性能力，那么同一个说法便不真实；因为祂以大能被宣示为天主子。也或许因为这个原因，神圣的预言时而称祂为仆人，时而称祂为儿子。他们因祂所取仆人的形状，并因祂是达味的后裔，而称祂为仆人；但他们按照祂作为首生者的身份，称祂为天主子。因此，称祂为人是真实的，称祂不是人也是真实的；称祂为人，因为祂能死；称祂不是人，因为祂比人更神圣。我想，\OriginalTerm{马尔西翁}{Marcion}曲解了正统的言词，他拒绝祂由玛利亚诞生，并宣称就祂的神性而言，祂并非由玛利亚所生，因而大胆地从福音中删除了有这层意思的段落。类似的下场似乎也落到了那些取消祂的人性而只接受祂的天主性的人身上；以及那些与此相反，取消祂的天主性而承认祂仅是一个圣洁的人、是众人中最正义的人身上。那些持守\OriginalTerm{幻象论}{Dokesis}教义的人，不记得祂谦抑自己，以至于死（\BibleRef{斐 2:8}），且服从至死，甚至死于十字架，只想象祂没有受苦，超脱一切此类偶然事件，他们尽其所能地剥夺我们那个比所有人更正义的人，结果只剩下一个不能拯救他们的形象，因为正如死亡藉一人而来，生命的成义也藉一人而来。我们从\OriginalTerm{圣言}{Logos}所领受的如此恩惠，若非祂取了人，若非祂保持自始与天主圣父同在的模样，若非祂接纳了人——那最先的、比众人更珍贵、比众人更纯洁、且堪当接纳祂的人——我们便无法领受。但在那人之后，我们也将能够接纳祂，接纳如此伟大、如祂本性的祂——如果我们在自己的灵魂中为祂准备一个与祂相称的地方。关于福音中表面上的矛盾，以及我愿它们以精神诠释的方式被处理，我已说了这么多。\par

% structure: p0011 source=h2 target=subsection reason=relative-heading-rank; base=h2

\subsection{五、保禄也作出关于自身的矛盾陈述，并在不同时间以相反方式行事。}

关于同一段落，也可以采用保禄这样的例子：他在一处（见：\BibleRef{罗 7:14}）说自己属血肉，被卖于罪恶之下，因此不能判断什么事；而在另一处，他却是属神的人，能判断一切，而自己却不受任何人判断。属血肉的人所说的话是：\Quote{我所愿意的，我并不做；我所恨的，我倒去做。}那被提到第三层天并听到了不可言说的话语的人，与那说\Quote{我要因这样的人夸耀，但为我自己，我不夸耀}的保禄是不同的。如果他对犹太人，就成为犹太人（见：\BibleRef{格前 9:20}），为要赢得犹太人；对法律之下的人，就成为法律之下的人，为要赢得法律之下的人；对没有法律的人，就成为没有法律的人（其实我不是没有天主的法律，而是在基督的法律之下），为要赢得没有法律的人；对软弱的人，就成为软弱的人，为要赢得软弱的人；那么显然，这些说法必须逐一审查：他成为犹太人，有时在法律之下，有时没有法律，有时软弱。例如，当他以许可的方式说话（见：\BibleRef{格前 7:6}），而非以命令的方式时，我们便可认出他是软弱的；因为他说：\Quote{谁软弱，我不软弱呢？}（见：\BibleRef{格后 11:29}）当他剃头献祭，或给弟茂德行割礼时（见：\BibleRef{宗 16:3}），他是犹太人；但当他对雅典人说（见：\BibleRef{宗 17:23}）：\Quote{我见到一座祭坛，上面写着\InnerQuote{给未识之神}。那么，你们所敬拜而不认识的，我向你们宣告}，以及\Quote{正如你们自己的一些诗人也说过：\InnerQuote{我们也是祂的子孙}}时，他就为那些没有法律的人成了没有法律的人，恳求那些最不虔诚的人归向宗教，并将诗人的话转用于自己的目的，说：\Quote{我们从爱开始；我们是祂的种族。}或许还能找到一些例子，表明他对那些不是犹太人却仍在法律之下的人，他也在法律之下。\par

% structure: p0013 source=h2 target=subsection reason=relative-heading-rank; base=h2

\subsection{六、关于伯多禄蒙召和洗者若翰被囚的不同记述。\Quote{葛法翁}的含义。}

这些例子或许有助于解释不仅关于救主，也关于门徒的陈述，因为这里也存在一些说法上的差异。因为西满被自己的兄弟安得肋找到并被告知\Quote{你将被称为刻法}（见：\BibleRef{若 1:41}），与他被耶稣在加里肋亚海边行走时看见，并与他的兄弟一起被召唤说\Quote{来跟随我，我要使你们成为得人的渔夫}之间，或许在思想上有不同。那位更深入根源、讲述圣言降生成人、因而没有记载那起初与天主同在的圣言的人性出生的作者，不应当告诉我们西满是在海边被发现并从那被召，而是他被那位曾在第十时辰与耶稣一起的兄弟找到，并且因着他这样被找到而得到了\OriginalTerm{刻法}{Cephas}的名字，这样安排是合适的。如果他在加里肋亚海边被耶稣看见，那几乎不可能在后来才被称呼说：\Quote{你是伯多禄，在这磐石上我要建立我的教会。}在若望福音中，法利塞人又知道耶稣与门徒一起施洗，将这加于祂其他的伟大活动之上；但三部福音中的耶稣根本不施洗。洗者若翰，与同名的福音作者一起，很长时间没有被下狱。相反，在玛窦福音中，他几乎在耶稣受试探的同时被下狱，这也成为耶稣退到加里肋亚以免被囚的缘由。但在若望福音中，根本没有关于若翰被囚的事。谁如此智慧且有能力，能了解四位福音作者所记载的有关耶稣的一切，既理解每件事本身，又对祂在各处的逗留、言语和作为有连贯的看法呢？至于我们目前面对的这段经文，按事件顺序记载，在第六天，救主在加里肋亚加纳的婚宴事毕后，便与祂的母亲、兄弟和门徒下到葛法翁，意为\Quote{安慰的田地}。因为在宴饮和饮酒之后，救主理应与祂的母亲和门徒来到这安慰的田地，安慰那些祂正在训练为门徒的人，以及那由圣神怀了祂的灵魂，用那将要立于饱满田地中的果实来安慰。\par

% structure: p0015 source=h2 target=subsection reason=relative-heading-rank; base=h2

\subsection{七、为什么祂的兄弟未被邀请赴婚宴；以及祂为什么在葛法翁只住了不多几日。}

但我们必须追问，为何祂的弟兄未受邀赴婚宴：他们原本不在那里，因为经上并未说他们在；然而他们与祂、祂的母亲及祂的门徒一同下到\Place{葛法翁}。此外，我们必须考察，为何在此场合，他们并非\Quote{\Emph{进入}}\Place{葛法翁}，亦非\Quote{\Emph{上到}}那里，而是\Quote{下到}那里。请思考，我们是否应当将此处祂的弟兄理解为那些与祂一同下到（世间）的众力量，它们未受邀赴婚宴，正如上文所作的解释，因为它们在比那些被称为基督门徒者更低更卑微的地方，且以另一种方式获得援助。因为如果祂的母亲蒙召，那么就有些结出果实的人，而主甚至与圣言的仆役和门徒一同下到这些人那里，以帮助他们，祂的母亲也与祂同在。那些被称为\Place{葛法翁}者看来不能容许\Person{耶稣}及那些与祂同下者与他们长久停留：因此他们与他们同住了不多几天。因为低处的安慰之田并不接纳诸多教义的光照，只堪接受少许。为清楚看出那些接待\Person{耶稣}时间长短不同者之间的区别，我们可将此处\BibleQuote{他们住在那里不多几天}与\BibleBook{玛窦}{福音}中所记载的、\Person{基督}由死者中复活后对奉命教导万民的门徒所说的话作一比较：\BibleQuote{看，我同你们天天在一起，直到今世的终结。}对那些将要知道人类本性尚在此世时所能知道的一切的人，祂强调地说：\BibleQuote{我与你们同在}；并且，正如每个新日子在默观之田上上升，将更多的日子带到蒙福者眼前，因此祂说：\BibleQuote{直到今世的终结的所有日子。}至于在\Place{葛法翁}的那些人，他们向他们下到那里，如同到更贫乏者那里，不但是\Person{耶稣}，连祂的母亲、祂的弟兄及祂的门徒\BibleQuote{在那里住了不多几天}。\par

% structure: p0017 source=h2 target=subsection reason=relative-heading-rank; base=h2

\subsection{八、基督如何与信徒同在直到今世的终结，以及在那一圆满之后，祂是否仍与他们同在。}

有些人很可能且并非无理地询问，当此世的所有日子结束时，是否将再无人可说\BibleQuote{看，我与你们同在}——对那些接纳祂直到今世圆满的人而言，因为\BibleQuote{直到}一词似乎指明了一个时间界限。对此我们必须说，\BibleQuote{我与你们同在}这句话与\BibleQuote{我在你们内}不同。我们或许更正确地认为，救主并非在门徒内，而是与他们同在，只要他们思想上尚未抵达今世的圆满。但及至他们看到，就其努力而言，那被钉死于他们的世界的圆满已近在咫尺，那时\Person{耶稣}将不再与他们同在，而是在他们内，他们要说：\BibleQuote{我生活已不是我生活，而是基督在我内生活}〔\BibleRef{迦 2:20}〕，以及\BibleQuote{你们既然寻求基督在我内说话的证验}〔\BibleRef{格后 13:3}〕。我们这样说，也为我们自己保留了通常的解释，即把\BibleQuote{常常}理解为直到今世圆满的时间，并不追问超出人类本性在此所能达到的。可以坚持那解释，同时恰当地对待那个\BibleQuote{我}。与那些被派遣去教导万民的门徒同在直至圆满的那一位，可能是那位空虚自己、取了奴仆形体的那一位，而此后在状态上可能是另一位；此后祂可能是如祂空虚自己之前那样的，直到祂的父使祂的一切仇敌成为祂的脚凳；此后，当圣子把王国交付于天主父时，可能是圣父对他们说：\BibleQuote{看，我与你们同在}。但究竟是\BibleQuote{所有日子}直到那时，还是仅仅是\BibleQuote{所有日子}，还是不是\BibleQuote{所有日子}而是\BibleQuote{每一天}，任何喜欢的人都可以思量。我们的计划不容许我们目前如此深入离题。\par

% structure: p0019 source=h2 target=subsection reason=relative-heading-rank; base=h2

\subsection{九、\Person{赫拉克利昂}说，经上未记载\Person{耶稣}在\Place{葛法翁}行了任何事。但在其他福音中，祂在那里行了诸多大事。}

但\Person{赫拉克利翁}处理\BibleQuote{此后祂下到葛法翁}这话时，声称这表示另一项行动的开始，并且\Quote{下到}一词并非无意义。\Quote{葛法翁}，他说，\Quote{意指世界这些最遥远的区域，这些物质的领域，祂下到那里；因为那地方不合适，他说，祂在那里没有被记载做过或说过任何事。}如今，若主在其他福音中也未被记载在\Place{葛法翁}做过或说过任何事，我们或许会犹豫是否应该接受此观点。但事实远非如此。\Person{玛窦}说，我们的主离开\Place{纳匝肋}，来到\Place{葛法翁}住在海边，并且从那时起祂开始宣讲说：\BibleQuote{你们悔改吧，因为天国临近了。}\Person{马尔谷}在其叙述中，从魔鬼诱惑开始，记载\Person{若翰}被下狱后，\Person{耶稣}来到\Place{加里肋亚}，宣讲天主的福音，并召叫四位渔夫成为宗徒之后，\BibleQuote{他们进入葛法翁；安息日祂随即在会堂里施教，众人对祂的教导感到惊讶。}\Person{马尔谷}也记载了\Person{耶稣}在\Place{葛法翁}所行的一件事迹，因为他接着说：\BibleQuote{在他们的会堂里，有一个人附着邪魔，他喊叫说：\InnerQuote{啊！纳匝肋人耶稣，我们与祢有什么相干？祢来毁灭我们吗？我知道祢是谁，祢是天主子。}\Person{耶稣}斥责他说：\InnerQuote{不要作声，从他身上出来！}邪魔撕裂那人，大声喊叫，就从他身上出来了。众人都惊讶不已。}在\Place{葛法翁}，\Person{西满}的岳母也热症痊愈了。\Person{马尔谷}还补充说，到了晚上，凡患病和附魔的都被治好。\Person{路加}关于\Place{葛法翁}的记载与\Person{马尔谷}很相似。他说：\BibleQuote{祂来到加里肋亚的葛法翁城，在安息日教导他们；他们对祂的教导感到惊讶，因为祂的话具有权威。在会堂里有一个附着不洁恶魔的人，他大声喊叫说：\InnerQuote{啊！纳匝肋人耶稣，我们与祢有什么相干？祢来毁灭我们吗？我知道祢是谁，祢是天主的圣者。}\Person{耶稣}斥责他说：\InnerQuote{不要作声，从他身上出来！}恶魔把那人摔倒在中间，就从他身上出来了，没有伤害他。}然后\Person{路加}记载主如何从会堂起身，进入\Person{西满}的家，斥责他岳母的热症，治愈了她的疾病；在这件事之后，\BibleQuote{太阳西沉时}，他说\BibleQuote{所有患有各种疾病的人都带到祂跟前，祂按手在每人身上，治好了他们。也有许多魔鬼从许多人身上出来，喊着说：\InnerQuote{祢是天主子。}祂斥责它们，不许它们说话，因为它们知道祂是基督。}我们列举所有这些关于救主在\Place{葛法翁}的言行记载，是为了驳斥\Person{赫拉克利翁}对这段经文的解释：\Quote{因此祂在那里没有被记载做过或说过任何事。}他要么必须赋予葛法翁两个含义，并说明理由；要么如果做不到，就必须放弃说救主到任何地方都无目的的说法。至于我们，如果遇到其他福音比较也无法显示\Person{耶稣}到某个地方没有结果的情况，我们将寻求天主的帮助，阐明祂的到来并非徒然。\par

% structure: p0021 source=h2 target=subsection reason=relative-heading-rank; base=h2

\subsection{十、\Place{葛法翁}的意义。}

\Person{玛窦}则补充说，当主进入\Place{葛法翁}时，百夫长来到祂跟前，说：\BibleQuote{我的仆人在家里瘫痪，痛苦不堪。}又向主多说了些关于他的事，然后得到回答：\BibleQuote{去吧，照你所信的，给你成就。}\Person{玛窦}随后讲述了\Person{伯多禄}岳母的故事，与其他两位福音作者一致。我认为，对于那些渴望聆听\Person{基督}的人来说，从四福音中收集所有关于\Place{葛法翁}的记载、言论、事迹，以及主到访该地多少次，有时是\Quote{下到}那里，有时是\Quote{进入}那里，以及祂从何处来，这是一项值得称赞且合宜的工作。若我们将所有这些要点比较，便不会在\Place{葛法翁}的含义上出错。一方面，病人在那里得到医治，其他奇迹也行在那里；另一方面，宣讲\BibleQuote{你们悔改吧，因为天国临近了}从那里开始。这似乎是一个标志，正如我们在进入这个主题时所表明的，那是一个更需要安慰的地方，或许由\Person{耶稣}使之如此，祂在那里以教导和作为安慰了人，那是一个安慰之地。因为我们知道，地名的含义与\Person{耶稣}所关联的事物相符；例如\Place{革尔格撒}，那地方的居民求祂离开他们的边界，其意思是\Quote{驱逐者之居所}。我们也注意到关于\Place{葛法翁}的这一点：不仅仅宣讲\BibleQuote{你们悔改吧，因为天国临近了}从那里开始，而且根据三位福音作者，\Person{耶稣}在那里行了祂最初的奇迹。然而，三位作者中没有一个在记录在\Place{葛法翁}所行的最初奇迹时，加上\Person{若望}门徒对\Person{耶稣}第一件奇事的注释：\BibleQuote{这是耶稣在加里肋亚的加纳所行的初次神迹。}因为在\Place{葛法翁}所行的并非神迹的开始，因为天主子的首要神迹是喜乐，在人类经验的光照下，也是最代表祂的。因为天主的圣言并非在治愈病人上显示其美，而是在祂提供清醇之饮，使那些健康且能赴宴的人快乐。\par

% structure: p0023 source=h2 target=subsection reason=relative-heading-rank; base=h2

\subsection{十一、为何逾越节被称为\Quote{犹太人的}逾越节。其设立：以及\Quote{主的庆节}与非如此称谓的庆节之间的区别。}

\BibleQuote{犹太人的逾越节临近了。} \BibleRef{若 2:13} 当我探究最明智的\Person{若望}在（此处）的准确性时，我自问：加上\Quote{犹太人}一词有何含义？逾越节本是哪个民族的节庆？难道说\Quote{逾越节临近了}还不够吗？然而，可能的情形是：人的逾越节是一回事，当人们没有按照圣经的本意来举行时；而天主的逾越节是另一回事，真正的逾越节，在心神和真理中朝拜天主的人以心神和真理来举行。那么，文本中所指示的区别就是天主的逾越节与所谓犹太人的逾越节之间的区别。我们应当注意逾越节的法律，观察主在圣经中首次提及它时说了什么。\BibleRef{出 12:1} \BibleQuote{主在埃及地对\Person{梅瑟}和\Person{亚郎}说：\InnerQuote{这个月为你们是正月，为你们是一年中的第一个月。你们要告诉以色列子民全会众说：\InnerQuote{本月初十，每人要取一只羊，按你们家族的户数取。}}}然后，在一些指示之后，其中没有再出现\Quote{逾越节}一词，他补充说：\BibleQuote{你们要这样吃：腰束带，脚穿鞋，手拿杖，急速地吃。这是主的逾越节。}他没有说：\Quote{这是你们的逾越节。}稍后，他又以同样的方式命名这个节庆：\BibleQuote{当你们的子孙问你们说：这是什么礼仪？你们要回答他们说：这是祭献，是主的逾越节，他怎样保护了以色列子民的房屋。}又稍后：\BibleQuote{主对\Person{梅瑟}和\Person{亚郎}说：\InnerQuote{这是逾越节的法律。任何外人都不可吃。}}再稍后：\BibleQuote{若有外方人来到你处，要守主的逾越节，他所有的男子都应受割损。}请注意，在法律中，我们从未发现\Quote{你们的逾越节}的说法；但在所引用的所有段落中，这一短语有一次没有附加任何修饰语，而我们三次看到\Quote{主的逾越节}。为了确认主的逾越节与犹太人的逾越节之间确实存在这样的区别，我们可以考虑\Person{依撒意亚}谈论此事的方式：\BibleRef{依 1:13} \BibleQuote{你们的月朔和安息日，以及你们的盛大节日，我不能容忍；你们的斋戒和节日，你们的月朔和庆节，我的灵魂憎恨。}主不把这些罪人的礼仪称为他自己的（若有这样的礼仪，他的灵魂就憎恨它们）；既不是月朔，也不是安息日，也不是盛大节日，也不是斋戒，也不是节庆。在\WorkTitle{出谷纪}关于安息日的立法中，我们读到：\BibleQuote{\Person{梅瑟}对他们说：\InnerQuote{这是主所说的话：\InnerQuote{安息日是归于主的圣安息日。}}}稍后：\BibleQuote{\Person{梅瑟}说：\InnerQuote{你们吃吧，因为今天是归于主的安息日。}}在\WorkTitle{户籍纪}中，在每个节期所献的祭品之前，仿佛所有节期都服从于连续不断的日常祭品的法律，我们读到：\BibleQuote{主对\Person{梅瑟}说：\InnerQuote{你要通告以色列子民，对他们说：\InnerQuote{我的祭品，我的供物，我的馨香火祭，你们要留心在规定的节期献给我。}你们要对他们说：\InnerQuote{这是你们应献给主的火祭。}}}圣经中所记载的节期，他称之为自己的，而不是领受法律之人的节期；他谈到\Emph{他的}祭品，\Emph{他的}供物。类似的表达方式也见于\WorkTitle{出谷纪}中关于子民的记载；当天主说子民是他自己的子民时，是在他们不犯罪的时候；但在关于牛犊的段落中，他弃绝了他们，称他们为\Person{梅瑟}的子民。\BibleRef{出 8:21} 一方面，\BibleQuote{你要对法郎说：\InnerQuote{主这样说：让我的人民去，好叫他们在旷野中事奉我。但如果你不让我的人民去，看，我要打发狗蝇来攻击你、你的臣仆、你的百姓和你的房屋；埃及人的房屋必充满狗蝇，甚至他们所在的地，我也要降下狗蝇。在那一天，我要区别我的百姓所在的歌笙地，那里没有狗蝇，好叫你认识我是主，是全地的主。我要在我的百姓和你的百姓之间施行区别。}}另一方面，他对\Person{梅瑟}说：\BibleRef{出 32:7} \BibleQuote{下去，赶快下去，因为你的百姓，就是你从埃及地领出来的，已经犯了罪。}因此，正如子民在不犯罪时是天主的子民，但犯罪时就不再被称为他的子民，同样，节期当被主的灵魂憎恨时，就被称为罪人的节期；但当关于它们的法律被颁布时，它们就被称为主的节期。在这些节期中，逾越节是其中之一，在当前的段落中，它被称为不是主的逾越节，而是犹太人的逾越节。在另一处，\BibleRef{肋 23:2} 我们读到：\BibleQuote{以下是主的节期，你们应称为圣召集。}因此，从主自己的口中，我们看到在此点上我们的陈述无可辩驳。无疑，有人会问关于宗徒的话，他在致格林多人书中写道：\BibleRef{格前 5:7} \BibleQuote{因为我们逾越节的羔羊基督，已被祭杀作了牺牲。}他没有说：\Quote{主的逾越节已被祭杀，即基督。}对此，我们必须说：要么宗徒简单地称逾越节为我们的逾越节，因为它是为我们而祭杀的；要么每一个真正属于主的祭献（而逾越节是其中之一）都等待其圆满实现，不在这个世代，也不在地上，而是在未来的世代和天上，当天国显现时。关于那些节期，十二位先知之一说：\BibleRef{欧 9:5} \BibleQuote{在集会的日子和主的节期之日，你们要做什么？}但\Person{保禄}在\WorkTitle{希伯来书}中说：\BibleQuote{但是你们却来到了熙雍山和永生天主的城，天上的耶路撒冷，以及千万天使，和那些登记在天上的首生者的集会与教会。}在\WorkTitle{哥罗森书}中：\BibleQuote{不要让任何人在饮食上，或在节期、月朔或安息日的事上，判断你们；这些事是将来事物的阴影。}\par

% structure: p0025 source=h2 target=subsection reason=relative-heading-rank; base=h2

\subsection{十二、论天上节庆，地上节庆为其预像。}

如今，在那些天上事物中——犹太人在世上所拥有的只是其阴影——那些首先在真法律之下受监护人和管家训练、直到时期圆满的人，将在那里庆祝节庆；即，在上方，当我们能够将天主子圆满的尺度接纳入自己之内时，这乃是那隐藏在奥秘中的智慧的工作，它使这些事明朗；并且它也向我们展示，有关肉食的法律如何是那些将在那里滋养并坚强我们灵魂之事的象征。但若有人想要凭想象探究如此浩瀚思想的汪洋，即便他愿意说明地方性的敬拜如何仍是天上事物的模式与阴影，且祭献与羊群充满意义，却妄想比宗徒走得更远，那便是徒劳的；宗徒确实力求将我们的心思提升到超越对法律的俗世见解之上，但他并未向我们详细展示这些事如何实现。即使我们从未世的角度审视节庆，逾越节是其中之一，我们仍须追问：为何我们的逾越节如今已被祭献，即基督，并且不仅如此，将来还要被祭献？\par

% structure: p0027 source=h2 target=subsection reason=relative-heading-rank; base=h2

\subsection{十三、逾越节的属灵意义。}

对于现在所考虑的教义，可以补充几点，尽管要全面论述法律中一切奥秘性的表述，特别是与节期有关的，尤其是与逾越节有关的，需要专门撰写多卷著作。犹太人的逾越节包括一只被祭献的羊，各人按父家取一只羊；逾越节还伴随着成千上万公绵羊和公山羊的宰杀，数量与各家族数目相称。但我们的逾越节为我们而被祭献，即基督。犹太节期的另一个特点是无酵饼：他们家中所有的酵都要清除；但\BibleQuote{所以我们守这节，不可用旧酵，也不可用奸恶、邪恶的酵，只用真诚和真理的无酵饼。}（\BibleRef{格前 5:8}）在我们提到的两种逾越节和无酵节之外，是否还有其他逾越节和无酵节，这是一个我们必须更仔细考察的问题，因为这些是天上事物的模型和影子，不仅是饮食、月朔、安息日，连节期也是未来事物的影子。首先，当宗徒说\Quote{我们的逾越节基督已被祭杀}，对此可能会产生这样的疑问：如果犹太人的羊是基督祭献的预像，那么应该只献一只羊，而不是许多只，因为基督是一位；如果献了许多只羊，那是为了遵循预像，就好像有许多基督被祭杀。但暂不讨论这一点，我们可以问：那作为牺牲的羊如何包含基督的形像——羊是由遵守法律的人祭杀的，而基督是由违法者处死的；并且，在基督身上如何应用\BibleRef{出 12:8}的命令：\BibleQuote{他们当夜要吃肉，用火烤了，与无酵饼和苦菜同吃。不可吃生的，不可用水煮，只可吃用火烤的：头与腿与内脏；一点不可留到早晨；骨头一根也不可折断。但所剩下的留到早晨，要用火烧了。}（\BibleRef{出 12:8}）\Quote{骨头一根也不可折断}这句话，若望似乎在他的福音中运用了，并写道：\BibleQuote{于是兵丁来了，先打断了第一个人的腿，又打断了与耶稣同钉的另一个人的腿；但来到耶稣那里，见他已经死了，就没有打断他的腿，惟有一个兵丁用长枪刺透他的肋旁，随即有血和水流出来。看见这事的那人就作见证，他的见证是真的，并且他知道自己所说的是真的，叫你们也可以信。这些事成了，为要应验经上的话：\InnerQuote{他的骨头一根也不可折断。}}除了这个之外，宗徒的话中还有无数其他要点需要探究，无论是关于逾越节还是无酵饼，但正如我们上面所说，这些必须在专门的长篇著作中处理。目前，我们只能根据眼前经文作一概要，并力求对主要问题作出简短解答。我们记得这话：\Quote{看，天主的羔羊，除免世罪者！}因为关于逾越节说：\BibleQuote{要从绵羊或山羊中取一只}（\BibleRef{出 12:5}）。圣史在此与保禄一致，两者都牵涉到我们上面提到的困难。但另一方面，我们必须说：如果圣言成了血肉，主说：\BibleQuote{你们若不吃人子的肉，不喝他的血，在你们内便没有生命。谁吃我的肉，并喝我的血，必得永生，在末日我且要叫他复活。因为我的肉，是真实的食品；我的血，是真实的饮料。谁吃我的肉，并喝我的血，便住在我内，我也住在他内。}（\BibleRef{若 6:53}）——那么，所论的肉就是那除免世罪的羔羊的肉；而这血，就是必须涂在门上两边门柱和门楣上、在我们吃逾越节的房屋里的血。我们必须在这世界的夜间吃这羔羊的肉，这肉要用火烤，与无酵饼同吃；因为天主的圣言不单单是血肉。实际上，祂亲自说：\BibleQuote{我是生命之粮}（\BibleRef{若 6:48}），以及\Quote{这是从天上降下来的生命之粮，叫人吃了就不死。我是从天上降下来的生命之粮，人若吃了这粮，必永远活着。}（\BibleRef{若 6:48}）然而，我们不可忽视，在宽松的用语中，任何食物都称为饼，正如我们在梅瑟的申命纪中读到：\Quote{四十天没吃饼，没喝水}，意思是他没有吃任何干湿食物。我因若望的话而有此观察：\Quote{我所要赐给的饼，就是我的肉，是为世界的生命而赐的。}再者，当我们悔改自己的罪，怀着依照天主的忧苦而忧伤——这是一种为我们的救恩而运作、无需后悔的悔改——或者当我们因试炼而转向那些被认为是真理的沉思，并从中获得滋养时，我们就吃羔羊的肉，与苦菜和无酵饼同吃。然而，我们不可吃羔羊的肉是生的，就像那些作字句奴隶的人一样，如同无理性的动物，以及那些对真正讲理性的人发怒，因为他们渴望理解属神的事；实际上，他们具有野兽的本性。但我们必须努力将圣经的生硬转化为熟食，不让所写的变得松弛、湿润和薄弱，就像那些耳朵发痒、转耳不听真理的人（\BibleRef{弟后 4:3}）所做的那样；他们的方法导致生活行为的松弛和软弱。但让我们有热切的心神，持守天主赐予我们的如火之言，就如耶肋米亚从对他说话的那一位所领受的：\BibleQuote{看，我已使我的话在你口中成为火。}（\BibleRef{耶 5:14}）并且让我们看见羔羊的肉被烤熟，好使那些分享的人，如基督在我们内说话，能说：\BibleQuote{当祂在路上给我们开启圣经时，我们的心不是火热的吗？}（\BibleRef{路 24:32}）此外，如果我们的职责是探究诸如用火烤羔羊肉这一点，我们不可忘记耶肋米亚因天主的话所受的苦的相似之处，如他所说：\Quote{它就像一团炽热的火，在我骨中焚烧，我毫无力气，不能忍受。}但是，在这进食中，我们必须从头部开始，也就是说，从关于天上事物的首要和最本质的教义开始，我们必须以脚部结束，即学问的最后分支，探究事物的最终本性，或关于更物质的事物，或关于地下的事物，或关于邪恶的灵和不洁的魔鬼。因为这些事物的记述可能并不明显，就像它们自身一样，而是隐藏在圣经的奥秘中，所以可以借喻地称之为羔羊的脚。我们也不能不处理内脏，即内在且向我们隐藏的部分；我们必须将全部圣经视为一个身体来对待，我们不可撕裂或破坏其整体结构和谐中存在的牢固而紧密的联系，正如那些尽可能撕裂贯穿所有圣经的圣神合一的人所做的那样。但是，上述关于羔羊的预言只在我们这黑暗生命的夜间作我们的滋养；此生命之后的东西，就如白天的黎明，我们为何要留到那时才吃那现在对我们有用的食物呢？但是，当黑夜过去，紧接着的白昼临近时，我们将有饼可吃，这与那更古老、更低级状态的有酵饼毫无关系，而是无酵饼，那将满足我们的需要，直到那无酵饼之后的东西赐给我们，即玛纳，那是天使的食物，而非人的食物。这样，我们中的每个人都可以在我们祖先的每一座房屋中献祭自己的羔羊；一个人破坏法律，根本不献羔羊，而另一个人可以完全遵守诫命，献上祭品，正确烹饪，而不折断一根骨头。因此，简而言之，这就是按照宗徒们所持的看法，以及按照福音中的羔羊，对为我们而牺牲的逾越节，即基督，所作的解释。因为我们不应认为历史的事是历史的事的预像，物质的事是物质的事的预像，而是物质的事是属神的事的预像，历史的事是理智的事的预像。现在，我们的论述无需上升到那第三个逾越节，即要与万千天使在最完美、最蒙福的出谷中庆祝的；关于这些事，我们已经说得比本段经文所需的更多了。\par

% structure: p0029 source=h2 target=subsection reason=relative-heading-rank; base=h2

\subsection{十四、前三个福音仅在传教生涯结束时提及逾越节；若望福音则在开端。对此的评论。\Person{赫拉克利翁}论逾越节。}

然而，我们不可不查考关于犹太人的逾越节临近的说法，当时主与祂的母亲、兄弟及门徒在\Place{葛法翁}。在\BibleBook{玛窦}{福音}中，主被魔鬼所弃、天使前来服事祂之后，当祂听说\Person{若翰}被交付时，便退往\Place{加里肋亚}，离开\Place{纳匝肋}，来到\Place{葛法翁}居住。然后祂开始宣讲，选召四位渔夫为宗徒，在\Place{加里肋亚}全境的会堂施教，治好了那些被带到祂面前的人。随后祂上山，讲论真福八端及其后续；结束那教训后，祂下山，第二次进入\Place{葛法翁}。\EditorialNote{\BibleRef{玛 8}} 然后祂上船，渡到对岸，到\Place{革辣撒}人的地区。他们恳求祂离开他们的境内，祂就上船，渡过去，来到自己的城。然后祂行了某些治愈，走遍各城各村，在他们的会堂施教；此后，发生了福音中大部分事件，然后玛窦才提到逾越节时间的临近。在其他福音作者那里，在\Place{葛法翁}停留之后，也过了许久才提及逾越节；这可能会证实那些在上文所持关于\Place{葛法翁}的看法的人的意见。那次停留，在犹太人的逾越节附近，因着那临近而更加明亮，既因其本身更好，更因在犹太人的逾越节期间，圣殿里有贩卖牛、羊和鸽子的人。这更加强调了那逾越节不是主的逾越节，而是犹太人的逾越节；圣父的家，在那些不圣化它的人眼中，成了商场，而主的逾越节对于那些持低级、物质看法的人，成了犹太人的逾越节。关于逾越节时间（大约在春分时）的查考，以及任何与之相关的其他查考，今后将有更合适的时机进行。至于\Person{赫拉克利翁}，他说：\Quote{这是伟大的节日；因为它是救主受难的预象；不仅羔羊被宰杀，吃它带来松弛，宰杀指向救主在此世的受难，而吃它则指向婚宴中的安息。} 我们已给出他的话，为可见他是何等缺乏谨慎、何等松散地进行，甚至在这一主题上也何等缺乏构建技巧；因之，他实在不值得重视。\par

% structure: p0031 source=h2 target=subsection reason=relative-heading-rank; base=h2

\subsection{十五、福音叙述中与洁净圣殿相关的不一致之处。}

\BibleQuote{耶稣便上耶路撒冷去。\BibleRef{若 2:13} 在殿院里，他发现了卖牛、羊、鸽子的，和坐在钱庄上兑换银钱的人；他就用绳索做了一条鞭子，把羊和牛都从殿院里赶出去，倾倒了兑换银钱之人的小钱，推翻了他们的桌子；并对卖鸽子的人说：\InnerQuote{把这些东西从这里拿出去，不要使我父的殿宇成为商场。}那时他的门徒就想起经上记载的：\InnerQuote{我对你殿宇所怀的热忱，把我耗尽}的话。}应当注意，\Person{若望}把\Person{耶稣}与在殿里卖牛、羊、鸽子的人的这件事当作他的第二件工作；而其他福音史家却在接近末尾、与受难史有关的记载中叙述了一件类似的事。\Person{玛窦}这样说：\BibleRef{玛 21:10} \BibleQuote{耶稣进耶路撒冷时，全城哄动，说：\InnerQuote{这是谁？}群众说：\InnerQuote{这是加里肋亚纳匝肋的先知耶稣。}耶稣进了圣殿，把一切在圣殿内买卖的人都赶出去，推倒了兑换银钱之人的桌子和卖鸽子之人的凳子，对他们说：\InnerQuote{经上记载：\InnerQuote{我的殿宇将称为祈祷之所，你们竟把它变成贼窝。}} }\Person{马尔谷}这样说：\BibleQuote{他们来到耶路撒冷。耶稣一进了圣殿，就开始赶出在圣殿内买卖的人，推倒了兑换银钱之人的桌子和卖鸽子之人的凳子。他不许任何人拿着器皿穿过圣殿。他教训他们说：\InnerQuote{经上不是记载：我的殿宇将称为万民祈祷之所吗？你们竟把它变成贼窝。} }\Person{路加}这样说：\BibleRef{路 19:41} \BibleQuote{耶稣临近的时候，望见京城，便哀哭她说：\InnerQuote{恨不得在这一天，你也知道有关你平安的事；但这事如今在你眼前是隐藏的。的确，日子将临于你，你的仇敌要在你四周筑起壁垒，包围你，四面窘困你；又要荡平你，及在你内的子女；在你内决不留一块石头在另一块石头上，因为你没有认识眷顾你的时期。}耶稣进了圣殿，开始把买卖人赶出去，对他们说：\InnerQuote{经上记载：\InnerQuote{我的殿宇应是祈祷之所，你们竟把它做成了贼窝。}} }还应进一步指出，三位福音史家所记载的主上耶路撒冷时在圣殿中所行的事，被\Person{若望}以一种非常相似的方式叙述为发生在此事很久以后，在另一次与这次不同的耶路撒冷之行之后。我们必须考虑这些叙述，首先是\Person{玛窦}的叙述，我们读到：\BibleRef{玛 21:1} \BibleQuote{耶稣临近耶路撒冷，来到\Place{贝特法革}，在\Place{橄榄山}那里，就打发两个门徒，对他们说：\InnerQuote{你们往对面的村庄去，立时会看见一匹拴着的母驴，和跟它在一起的驴驹。解开，给我牵来！如果有人对你们说什么，你们就说：\InnerQuote{主要用它们。}那人就会立刻放它们来。}这事发生，是为应验先知所说的：\InnerQuote{你们应向\Place{熙雍}女子说：\InnerQuote{看，你的君王来到你这里，温和地骑着一匹驴，一匹驴驹，即驮重驴的驴驹上。}}门徒就去，照耶稣所吩咐的作了。他们牵了母驴和驴驹来，把外衣搭在它们身上，扶耶稣坐在上面。很多群众把自己的外衣铺在路上，还有些人从树上砍下树枝铺在路上。前呼后拥的群众喊说：\InnerQuote{贺三纳于达味之子！因主名而来的，当受赞美！贺三纳于至高之天！} }此后是：\BibleQuote{一进了耶路撒冷，全城都惊动了。}然后我们有\Person{马尔谷}的记载：\BibleRef{谷 11:1} \BibleQuote{当他们临近耶路撒冷，来到\Place{贝特法革}和\Place{伯达尼}，在\Place{橄榄山}那里，耶稣就打发两个门徒，对他们说：\InnerQuote{你们往对面的村庄去，一进去，就看见一匹拴着的驴驹，从来没有人骑过；解开它，牵来。如果有人对你们说：\InnerQuote{你们为什么做这事？}你们就说：\InnerQuote{主要用它。}那人就立刻把它牵回来。}他们去了，便见一匹驴驹拴在门外大路上，就解开它。站在那里的人对他们说：\InnerQuote{你们为什么解开驴驹？}他们就照耶稣所吩咐的答复了，那些人便让他们牵去了。他们把驴驹牵到耶稣跟前，把外衣搭在上面，耶稣就骑了上去。另有人砍下田野的树枝铺在路上。前呼后拥的人喊说：\InnerQuote{贺三纳！因主名而来的，当受赞美！那要来的我们祖先达味之国，当受赞美！贺三纳于至高之天！}耶稣进了耶路撒冷，到了圣殿，周围看了各样事物；时候已晚，便同十二门徒出来，往伯达尼去了。第二天，他们从伯达尼出来，耶稣饿了。}然后，在无花果树枯干的事件以后，\BibleQuote{他们来到耶路撒冷。耶稣进了圣殿，开始赶出那些买卖的人。} \Person{路加}这样记载：\BibleRef{路 19:29} \BibleQuote{当耶稣临近\Place{贝特法革}和\Place{伯达尼}，在名叫\Place{橄榄山}的那里时，就打发两个门徒，说：\InnerQuote{你们往对面的村庄去，一进去，就看见一匹拴着的驴驹，从来没有人骑过；解开它，牵来。如果有人问你们：\InnerQuote{为什么解开它？}你们就说：\InnerQuote{主要用它。}}门徒去了，果然如耶稣所说。当他们解开驴驹时，主人问他们：\InnerQuote{为什么解开驴驹？}他们回答：\InnerQuote{主要用它。}他们就把驴驹牵到耶稣跟前，把外衣搭在上面，扶耶稣骑上。前进时，人们把外衣铺在路上。当他临近橄榄山的下坡时，众门徒因所见过的一切奇能，都欢欣踊跃，大声赞美天主说：\InnerQuote{因主名而来的君王，当受赞美！和平在天上，光荣于至高天！}人群中有几个\Person{法利塞人}对他说：\InnerQuote{师傅，责斥你的门徒吧！}他回答说：\InnerQuote{我告诉你们：这些人若不作声，石头就会喊叫！}当他临近时，望见京城，便哀哭起来，等等，正如我们前面所引用的。}相反，\Person{若望}在几乎相同的记述之后，即\Quote{耶稣上了耶路撒冷，在殿里找到了那些卖牛和羊的人}之后，又第二次记述主上耶路撒冷，然后讲述逾越节前六天在\Place{伯达尼}的晚餐，\Person{玛尔大}伺候，\Person{拉匝禄}同席。第二天，\BibleRef{若 12:12} \BibleQuote{有许多上节来的人，听说耶稣来到耶路撒冷，便拿着棕榈枝出去迎接他，喊说：\InnerQuote{贺三纳！因主名而来的以色列君王，当受赞美！}耶稣找到一匹驴驹，就骑上去，正如经上所记载的：\InnerQuote{熙雍女子，不要害怕！看，你的君王骑着驴驹来了。}}我从福音书中摘录了很长的段落，但我认为有必要这样做，以展示我们这部福音书在这一部分中的不一致。三部福音书将这些我们认为是与\Person{若望}所叙述相同的事件，与主的一次耶路撒冷之行联系起来。而\Person{若望}却将它们与两次相隔甚远的访问联系起来，其间主还到过其他地方旅行。我认为，对于那些在解释中只承认历史事实的人来说，是不可能表明这些不一致的陈述是彼此和谐的。如果有人认为我们没有给出正确的解释，就让他对我们的这一断言写一篇有理有据的反驳吧。\par

% structure: p0033 source=h2 target=subsection reason=relative-heading-rank; base=h2

\subsection{十六、洁净圣殿的灵意化解释。若按字面理解，它呈现出一些非常困难和不大可能的特点。}

然而，我们将按所赐予我们的力量，阐释促使我们在此承认和谐的理由；同时，我们恳求那赐予每一位恳切探求者恩赐的主，我们叩门，愿藉更高知识的钥匙，圣经中隐藏的事向我们开启。首先，让我们专注于若望的话，开始：\Quote{耶稣便上耶路撒冷去。}\BibleRef{若 2:13}耶路撒冷，如主自己在\BibleBook{玛窦}{福音}中所教导的，\BibleRef{玛 5:35}\Quote{是大君王的城市。}它并非位于低洼或低下之处，而是建在高山上，周围有群山环绕，归于同一处所，上主的支派上去，为以色列作证。但那城也称为耶路撒冷，凡地上的无人能上去，也无人能进入；但每个灵魂，由于本性具有某种崇高和洞察精神之物的敏锐，便是那城的公民。甚至耶路撒冷的居民也可能犯罪（因为甚至最敏锐的心智也可能犯罪），如果他们犯罪后不迅速回转，失去心智的能力，不仅可能居住在犹大那些异乡城中，甚至可能被登记为那些城的公民。\newline{}耶稣在帮助了加里肋亚的加纳那些人之后，上耶路撒冷去，然后下到葛法翁，为要在耶路撒冷行所记载的事。祂在圣殿里——确实，被称为救主圣父的殿，即在教会或公教会纯正道理的宣讲中——发现有些人使祂父的殿成为商场。耶稣在圣殿里总是发现这类人。因为在那被称为教会，即永生天主的家、真理的柱石和基础的地方，\BibleRef{弟前 3:15}何时没有兑换银钱的人坐着，需要耶稣用细绳编成的鞭子责打，以及需将他们的金钱倒出、推翻桌子的小贩？何时没有倾向买卖的人，却需要被挽住犁和牛，好让他们手扶着犁不向后看，配得天主的国？何时没有那些偏好不义之财的人，胜过给予他们真正装饰材料的羊？总有许多人轻视真诚、纯洁、无苦涩无胆汁的事物，为可怜的利益出卖那些比喻称为鸽子之人的照顾。因此，当救主在圣殿——祂父的殿——中发现那些卖牛、羊、鸽子的人，以及兑换银钱的人坐着，祂便用祂所编的细绳鞭子将他们连他们生意的牛羊一同赶出，倒出他们的钱币，因为它们不值得保存，如此微小。祂也推翻那些贪爱钱财之人灵魂中的桌子，甚至对那些卖鸽子的人说：\Quote{把这些东西拿出去}，使他们不再在天主的殿中买卖。\newline{}但我相信，在这些话中，祂也指示了更深奥的真理，我们可以将这些事件视为一个象征：那圣殿的敬拜不再由司铎以物质祭献的方式继续进行，时候将到，法律不能再被遵守，尽管按肉身的犹太人多么渴望。因为当耶稣赶出牛和羊，并命令拿走鸽子时，这是因为牛、羊、鸽子将不再按犹太人的惯例在那里被祭献。可能那些带着物质事物印记而非天主印记的钱币被倒出，作为一种预表；因为那看似可敬的法律，及其叫人死的文字，现在耶稣已来到，并用祂的鞭子责打人民，这法律将被解散和倒出，圣职（\OriginalTerm{主教职}{episcopate}）将转给那些信了的外邦人，天主的国将从犹太人\BibleRef{玛 21:43}夺去，赐给那结果子的民族。\newline{}但也可能如此：自然的圣殿是灵魂，精通理性，因其内在的理性而高于身体；耶稣从较低、较卑微的葛法翁上到这圣殿；在耶稣的训诲应用于它之前，其中发现有属地的、无理性的、危险的倾向，以及有名无实的美善之物，这些都被耶稣用由论证和责备的教义编成的言语赶走，为使祂父的殿不再成为商场，而是为了它自身和他人的救恩，接受按照天上和属神法律所行的天主事奉。牛象征属地的事物，因它是耕地的。羊象征无理性、兽性的事物，因为它比大多数无理性生物更奴性。鸽子象征空虚和不稳定的思想。小钱象征那些被认为善实则不然的事物。\newline{}若有人反对这种解释，说经文提到的只是洁净动物，我们必须说，否则经文就显得不可信。事件必然按故事的可能性叙述。不可能叙述一群不洁净的动物进入圣殿，也不可能在圣殿中出售除了用于祭献之外的任何动物。福音作者利用了犹太节期时商人的已知习俗；他们确实将这类动物带到外院；这种习俗，加上祂所知的真实事件，成为祂的材料。然而，任何人若愿意，可以查考：这样做是否与耶稣在这世界的身份相符，因为祂被看作木匠的儿子，竟敢如此行动，将一群商人从圣殿赶出？他们上到节期，为大量人民出售羊——数以万计——他们按父家要祭献这些羊。为富裕的犹太人，他们卖牛；还有鸽子给那些许愿这类动物的人；许多人无疑是为节日盛宴而购买的。耶稣倒出兑换银钱者自己的钱，推翻他们的桌子，难道不是不当之举吗？谁被一个被视为卑微之人的细绳鞭子打中，被赶出圣殿，会不攻击祂、呼喊、亲手报仇？尤其有如此众多群众在场，他们都会感到同被耶稣侮辱。此外，天主子亲手拿起细绳，编成鞭子用以驱赶出圣殿，这难道不表明大胆、鲁莽甚至某种不法吗？\newline{}对于想要辩护这些事并视之为真实历史事件的作者，还有一个避难所，即诉诸耶稣的神性：祂能按祂的意愿熄灭敌人的怒火，藉恩宠胜过千万人，驱散暴乱之人的计谋；因为\Quote{主消散列国的谋略，使人民的计谋归于无有，但主的谋略永远立定。}因此，这段经文中的事件，若确实发生，在它所显示的权能上，不亚于基督所行的任何最奇妙作为，并以它的神性要求旁观者的信德。有人可证明这比在加里肋亚的加纳变水为酒更伟大；因为后者只是无生命的物质被改变，而这里却是千万人的灵魂和意志。然而，要注意，在婚宴上，耶稣的母亲在那里，耶稣和祂的门徒被邀请；但据说只有耶稣下到葛法翁。不过，祂的门徒后来出现，与祂同在；他们记起\Quote{我为祢的殿所发的焦急将耗尽我。}或许耶稣在每位门徒内，当祂上耶路撒冷时，因此并未说\Quote{耶稣和祂的门徒}上耶路撒冷，而是说祂下到葛法翁，\Quote{祂和祂的母亲、祂的兄弟及祂的门徒。}\par

% structure: p0035 source=h2 target=subsection reason=relative-heading-rank; base=h2

\subsection{十七、\Person{玛窦}所叙述的进\Place{耶路撒冷}事件。为按字面理解者带来的困难。}

我们现在要考虑其他福音书关于从圣殿中驱逐那些使之成为商场之人的记述。首先看看我们在\BibleBook{玛窦}{福音}中读到的。关于主进入\Place{耶路撒冷}，他说：\BibleQuote{全城都惊动了，说：这个人是谁？}但在这之前，他记载了驴和驴驹的故事，它们是按主的命令被牵来，并由他从\Place{贝特法革}派往对面村庄的两个门徒所找到。这两个门徒解开拴着的驴，并得到命令：若有人对他们说什么，就回答说：\BibleQuote{主要用它们，并且会立刻把它们送来。}通过这些事件，玛窦宣告那预言应验了，即\BibleQuote{看，你的君王来到你这里，温和地骑着一匹驴，一匹驴驹，就是驴的崽子}，我们在\Person{匝加利亚}先知书中找到的。\BibleRef{匝 9:9}于是，门徒们就去照\Person{耶稣}所吩咐的做了，他们牵来驴和驴驹，把他们的外衣搭在上面，他说，主就骑上它们，显然是指驴和驴驹。然后\BibleQuote{群众多半把外衣铺在路上，另有些人从树上砍下树枝铺在路上。前呼后拥的群众喊道：贺三纳归于达味之子！奉主名而来的当受赞美！贺三纳于至高之天！}因此，当他进\Place{耶路撒冷}时，全城都骚动起来，说：这是谁？\BibleQuote{群众说}（显然是那些前呼后拥的人），对那些问他是谁的人说：\BibleQuote{这是\Place{加里肋亚}\Place{纳匝肋}的先知\Person{耶稣}。}\Person{耶稣}进了圣殿，把一切在圣殿里买卖的人都赶出去，推倒了兑换银钱者的桌子和卖鸽子者的凳子，对他们说：\BibleQuote{经上记载：我的殿宇应称为祈祷之所，你们竟使它成为贼窝了。}我们要问那些认为\BibleBook{玛窦}{福音}写福音时心中只有历史事实的人，为什么有必要派两个门徒去\Place{贝特法革}对面的村庄，找到拴着的驴和它旁边的驴驹，解开它们并牵来？他骑在驴和驴驹上进城，这事为何值得记载？当\Person{匝加利亚}说\BibleRef{匝 9:9} \BibleQuote{熙雍女子，你应欢乐！耶路撒冷女子，你应欢呼！看，你的君王来到你面前，是正义的，带来救恩，温和地骑在驴和驴驹上}，他是如何预言基督的呢？如果这个预言仅仅预言了福音作者所描述的物质事件，那些拘泥于字句的人又如何能坚持认为预言的下文也是如此呢？即\BibleQuote{他必从\Place{厄弗辣因}毁灭战车，从\Place{耶路撒冷}毁灭战马，战士的弓必被折断，并向外邦人宣布和平；他必统治那海到海，从河到地极}等等。还应注意到，玛窦没有使用先知的原话，因为他把\BibleQuote{熙雍女子，你应欢乐！耶路撒冷女子，你应欢呼！}改为\BibleQuote{你们要对熙雍女子说}。他删减了预言，省略了\BibleQuote{正义的，带来救恩}，然后保留了\BibleQuote{温和地骑着}，如原文一样，但把\BibleQuote{骑在驴和驴驹上}改为\BibleQuote{骑在驴和驴驹，就是驴的崽子}。犹太人查考这预言应用于\Person{耶稣}的事迹上，以一种我们不能忽视的方式追问我们：\Person{耶稣}如何毁灭了\Place{厄弗辣因}的战车和\Place{耶路撒冷}的战马？他如何毁灭了敌人的弓，并做了经文中提到的其他事？关于预言的讨论就到此为止。然而，我们的字面解释者，如果觉得在驴和驴驹上没有什么堪当天主子显现的事，可能会指出路途遥远来解释。但是，首先，十五斯塔狄不是什么大距离，不能为这事提供合理的解释；其次，他们不得不告诉我们，为什么这么短的路程需要两头役畜？经上说：\BibleQuote{他骑在上面}，又说：\BibleQuote{若有人对你们说什么，你们就说：主要用它们，他立刻就会把它们送来。}在我看来，宣称像他这样本性的人承认需要解开一头驴并和一头驴驹一起来，并不配得上圣子神性的伟大；因为天主之子所需要的一切，都应当是伟大且配得上他的美善的。然后，那一大群人把外衣铺在路上，而\Person{耶稣}允许他们这样做，并没有责备他们，这从另一段经文中的话可以明显看出：\BibleRef{路 19:40} \BibleQuote{这些人若不出声，石头就要喊叫了}。我不知道，如果这些事除了表面意思外别无他意，那么作者对此感到喜悦是否不表明了他具有一定的愚蠢程度。从树上砍下树枝铺在驴经过的路上，对于作为人群中心的那一位来说，与其说是精心设计的迎接，不如说是一种妨碍。我们在被\Person{耶稣}逐出圣殿的那些人身上遇到的困难，在这里以更大的程度出现。在\BibleBook{若望}{福音}中，他赶出那些买的人，但玛窦说他赶出在圣殿里卖的人和买的人。买的人自然比卖的人多。我们必须考虑，赶出圣殿里的买家和卖家，是否与一个被认为是木匠之子的人的名声不相称，除非如我们之前所说，他是以神能制服了他们。在其他福音作者中，他对他们说的话也比若望的更严厉。因为若望说\Person{耶稣}对他们说：\BibleQuote{不要使我父的殿宇成为商场}，而在其他福音中，他们因使祈祷之所成为贼窝而受责备。他父的殿宇不能容许变成贼窝，虽然因罪人的行为，它被变成了商场。它不仅是祈祷之所，事实上是天主的殿，由于人的疏忽，它窝藏了贼，并且不仅变成了他们的屋，还变成了他们的巢穴——这不是任何建筑技巧或理性所能造成的。\par

% structure: p0037 source=h2 target=subsection reason=relative-heading-rank; base=h2

\subsection{十八、驴和驴驹代表旧约与新约。故事各细节的属灵含义。若望记述与其他福音作者记述之间的差异。}

现在，要洞察这些事情的真正真理，是那真正理智的一部分，这理智赐给了那些能说\BibleRef{格前 2:16}\BibleQuote{我们却有基督的心，为能看透天主白白赐给我们的一切}的人；无疑，这超出了我们的能力。因为我们灵魂中的主宰原则并非免于动荡，我们的眼睛也不像基督那美丽的新娘应有的那样，新郎论及她说，\BibleRef{歌 1:15}\BibleQuote{你的眼睛有如鸽子}，这或许是以谜语的方式意指那寓居于灵性中的观察力，因为圣神曾像鸽子一样降临到我们的主身上，也降临到每个人里面的那\Quote{主}身上。然而，按我们现在的状况，我们不会迟延，而是要摸索那已向我们说出的生命之言，并努力抓住其中的力量，这力量会流向以信德触碰它们的人。如今耶稣是天主圣言，他进入那被称为耶路撒冷的灵魂，骑在门徒从束缚中释放的驴上。也就是说，骑在旧约的简单文字上，由那两位释放它的门徒所解释：第一位是将所写的应用于灵魂的服务，并指出其中与灵魂相关的寓意；第二位是通过隐藏在阴影中的事物，揭示未来美好而真实的事物。但他也骑在小驴驹上，即新约；因为在这两者中，我们都找到真理之言，它净化我们，并驱散我们心中所有倾向于买卖的思想。但他并非独自来到耶路撒冷，即灵魂，也不是只有少数同伴；因为在使我们成全的天主圣言之前，有许多事物必须进入我们内，在他之后也有许多事物必须跟随，这一切都在歌颂、荣耀他，并将他们的装饰和衣服铺在他之下，好使他骑乘的牲畜不接触地面，当他——那从天上降下的那一位——坐在它们上面时。为了使他的坐骑——圣经的旧约与新约之话——能更高于地面，必须从树上砍下树枝，使它们踏在合理的解释上。而那些在前在后跟随他的群众，也可能表示天使的服役，其中一些在我们灵魂中为他预备道路，并帮助装饰他们，另一些则在他临在于我们之后而来，这一点我们常提及，故不必再引证。也许不无理由，我将那些引导圣言本身到灵魂的周围声音比作驴；因为它是负重的牲畜，而许多重担，沉重的负荷，从文本中出现，尤其是旧约，这可以清楚看到，如果观察犹太人在此方面所行的事。但驴驹非同驴那样是负重的牲畜。因为虽然驴的每一条绳索对没有内在的圣神扶持和减轻力量的人是沉重的，但新的话语比旧的话语更轻。我知道有些人解释那被拴住的驴是指受割礼而来的信徒，他们被那些真正在灵性上受教导于圣言的人从许多束缚中释放；而驴驹他们则视为外邦人，他们在接受耶稣的话语之前，不受任何控制，不屈服于任何轭，过着放纵、贪图享乐的生活。我所说的这些作者并未说明那些在前和在后的是谁；但若说在前的是像梅瑟和先知，在后的是圣宗徒，也不算荒谬。如今我们要探询的是，这一切所进入的是怎样的耶路撒冷，以及那被天主子赶出许多买卖人的房屋是什么。也许上面的耶路撒冷，是主将要升上去的，他像驾车者一样驱赶受割礼者和外邦信徒，而先知和宗徒在他前后（或者是服侍他的天使，因为那些在前在后的也可能指他们），也许那座城在他升上去之前，包含着所谓的\BibleRef{弗 6:12}\BibleQuote{天界里邪恶的鬼神}，或是客纳罕人、赫特人、阿摩黎人以及天主人民的其他仇敌，总之，是外方人。因为在那个区域，也可能应验了那预言，即\BibleRef{依 1:7}\BibleQuote{你们的国土已荒凉，你们的城市被火焚毁，你们面前的田地，外邦人正在瓜分}。因为这些就是那污秽并使之成为贼窝的人，即他们自己，这天上的父家，圣洁的耶路撒冷，祈祷之所；他们有伪币，给所有来的人便士和小钱，廉价无价值的货币。这些人就是那些与灵魂相争，从他们那里夺走最宝贵的，抢走他们更好的部分，还给他们毫无价值之物的。但门徒去找到拴着的驴并解开它，因为由于法律加在它上面的遮蔽，它不能承载耶稣。\BibleRef{格后 3:14}驴驹也和它一起被找到，两者都迷失直到耶稣来到；我指的是受割礼的和后来相信的外邦人。但关于耶稣坐在它们上面升入耶路撒冷后，它们又如何被送回去，说来有些危险；因为这其中有某种神秘的成分，与圣人们改变为天使有关。在那改变之后，在下一个时代，它们将被送回去，如同服役的灵，\BibleRef{希 1:11}，它们被派遣服役，为那些将要承受救恩的人。但如果驴和驴驹是旧约和新约圣经，天主圣言骑在其上，那么很容易看出，在圣言显现于它们之后，它们就被送回去，并在圣言进入耶路撒冷（即那些赶出所有买卖思想的人中）之后不再等候。我也认为，驴被找到拴着的地方——以及驴驹——是一个村庄，一个没有名称的村庄，这不是没有意义的。因为与天上的大世界相比，整个地球是一个村庄，驴和驴驹被拴在那里，它被简称为\Quote{那村庄}，没有附加任何别的名称。玛窦说门徒是从贝特法革被派去取驴和驴驹的；贝特法革是一个司铎之地，其名字意为\Quote{颚骨之屋}。关于玛窦的经文，我们已经按我们的能力说了这些，并保留到以后有机会时——当我们被允许单独讨论玛窦福音时——再对他的话做更完整、更准确的讨论。马尔谷和路加说，两位门徒按他们师傅的指示，找到一只拴着的小驴，从未有人骑过，他们解开它，带到主那里。马尔谷补充说，他们发现驴驹拴在门外，在路旁。但谁是\Quote{外面}的？是那些外邦人，他们\BibleRef{弗 2:12}在盟约上是外人，是天主恩许的陌生者；他们在路上，不在屋顶或房屋下安息，被自己的罪恶捆绑，将藉由上述的双重知识，即耶稣的朋友们，被解开。并且那捆绑驴驹的绳索，以及违背健全法律并被其谴责的罪恶——因为法律是生命之门——就此而言，我说，它们不是在门内，而是因为在门外，也许在门内不会有任何这样的罪恶之绑。但马尔谷说，有一些人站在被拴的驴驹旁边；我想是那些拴它的人；如路加所记，驴驹的主人问门徒说：\Quote{你们为什么解开驴驹？}因为这些奴役并捆绑罪人的主人是不合法的主人，当真正的主人释放驴驹时，他们不敢正视他。因此，当门徒说\Quote{主需要它}时，这些邪恶的主人无言以对。门徒随后将驴驹赤身带到耶稣跟前，把自己的衣服搭在上面，好使主能舒适地坐在门徒的衣服上。根据玛窦的记述，进一步所说的将不会造成任何困难；即\BibleRef{谷 11:15}\BibleQuote{他们来到耶路撒冷，耶稣进入圣殿，开始赶出在圣殿里买卖的人}，以及\BibleRef{路 19:41}\BibleQuote{耶稣临近的时候，望见京城，便哀哭她；进入圣殿，开始赶出卖东西的人}。因为在一些自己内有圣殿的人中，他赶出所有在圣殿里买卖的人；但在其他不完全服从天主圣言的人中，他只是开始赶出卖和买的人。此外还有第三类，在其中他开始只赶出卖的人，而不是买的人。反之，在若望那里，他们全部被用细绳编成的鞭子赶出，连同羊和牛。应仔细考虑，是否可能通过寓意解经法满意地解决所描述之事的变化以及其中的差异。每位福音书作者都归给圣言不同的行动方式，这些方式在不同气质的灵魂中产生不同但相似的效果。我们注意到关于耶稣前往耶路撒冷行程的差异——即我们手中这本福音书的记载与其他三本完全不同，如我们解释过他们的话——无法用其他方式调和。若望给出了与其他三本相似但不相同的说法；他不用从树上砍下的树枝或从田里带来的秸秆铺在路上，而是说他们拿了棕榈树枝。他说有许多人来到庆节，他们出去迎接他，喊着说：\Quote{奉主名来的当受赞美！}以及\Quote{以色列王当受赞美！}他还说是耶稣自己找到了基督所骑的小驴，而\Quote{小驴}这个词无疑传达了某种额外的含义，因为这小动物所提供的恩惠不是出于人，也不是通过人，而是通过耶稣基督。此外，若望并不比其他人更准确地复述先知的话；他给我们的不是原文，而是\Quote{熙雍女子，不要害怕！看，你的君王骑着驴驹来了}（代替了\InnerQuote{骑着}和\InnerQuote{驴和驴驹}）。\Quote{熙雍女子，不要害怕}这句话根本不在先知书中。但由于先知的话被所有人都这样应用，让我们看看是否必然：熙雍女子应大大喜乐，而比她更大的耶路撒冷女子，不仅应大大喜乐，还应宣告她的君王来到她那里，正义的，带来救恩，谦逊的，骑着驴和驴驹。那么，谁接待他，就不再害怕那些用异端似是而非的言论武装的人，那些厄弗辣因的战车，主说将毁灭它们，\BibleRef{匝 9:10}也不怕那为安全而虚妄的马，即那习惯了感官事物的疯狂欲望，它伤害许多希望居于耶路撒冷并专注健全圣言的人。因那骑着驴和驴驹者消灭了每一支敌对的箭而喜乐也是合适的，因为敌人的火箭不再能胜过那已将耶稣接到自己圣殿中的人。当救主来到耶路撒冷时，也有大量外邦人带着平安\BibleRef{匝 9:9}前来，那时他统治众水，以击碎水中龙的头，我们要踏上海浪和地上所有河流的出口。然而，马尔谷在写关于驴驹时，记载主说：\Quote{从来没有人骑过}；在我看来，他似乎暗示那些后来相信的人，在耶稣来到他们之前从未服从过圣言。因为或许在人类中，从未有人骑过这驴驹，但在敌对圣言的心灵或权势中，有些曾骑过它，因为在依撒意亚先知书中，敌对势力的财富据说由驴和骆驼驮着。\BibleRef{依 30:6}他写道：\Quote{在困苦和灾难中，有狮子、幼狮，以及飞蛇的后裔，他们用驴和骆驼驮载自己的财富。}对于那些只注重字面意义的人，问题再次出现：按照他们的观点，\Quote{从来没有人骑过}这句话是否并非毫无意义？因为除了人，谁还骑驴驹呢？以上是我们的一些看法。\par

% structure: p0039 source=h2 target=subsection reason=relative-heading-rank; base=h2

\subsection{十九、赫拉克利翁关于洁净圣殿的各种看法。}

我们来看看\Person{赫拉克利翁}对此作何解释。他说，上耶路撒冷意味着主从物质事物上升到属灵的地方，这地方是耶路撒冷的肖像。他认为，经文说的是\Quote{祂在殿里发现}，而不是\Quote{在至圣所里}，因为主不应被理解为只是在那召叫中作工具，这召叫发生在没有圣神的地方。他认为殿是至圣所，除了大司祭外无人能进入，我相信他说，属灵的人去那里；而殿的外院，连肋未人也能进入，是那些得救的\OriginalTerm{属魂之人}{psychical ones}的象征，但他们在\OriginalTerm{满全}{Pleroma}之外。然后，那些在殿里贩卖牛、羊、鸽子的人和坐着兑换银钱的人，他认为象征那些不将任何事归于恩宠，却把外人进入圣殿看作买卖和获利之事的人，以及那些为敬拜天主而献祭，却着眼于自己的利益和贪财的人。至于耶稣用细绳做成、而不是从别人那里得来的鞭子，他按自己的方式解释说，这鞭子是圣神之德能与能力的形象，藉着祂的气息驱赶恶人。他还声称，鞭子、细麻布、手巾以及诸如此类的东西，都象征圣神的德能与能力。然后他假设了没有记载的事，譬如鞭子拴在一块木头上，而这木头他视为十字架的预像；在这木头上，赌徒、商人和一切邪恶都被钉死而除掉了。在探究耶稣的行动、讨论鞭子由两种物质构成时，他以一种异常的方式幻想；他说，祂没有用死皮革做它。祂希望使教会不再作贼窝，而成为祂父的殿宇。我们在此必须就赫拉克利翁文本中所提到的神性，说些最重要的事。如果耶稣称耶路撒冷的圣殿为祂父的殿，而那殿是为尊荣创造天地的那位而建造的，为什么我们不立即被告知，祂就是那位创造天地的独一者的儿子，即天主子呢？对于这耶稣之父的殿，作为祈祷的殿，基督的宗徒们，正如我们在他们的\Quote{\BibleBook{}{宗徒大事录}}中所见，被天使吩咐\BibleRef{宗 5:20}去站在那里，宣讲这生命的一切话。但他们来到祈祷的殿，经过美门，在那里祈祷；这事他们不会做，除非他们知道祂与那些奉献了那殿的人所敬拜的天主是同一的。因此，那些顺从天主而不顺从人的人，伯多禄和宗徒们，也说：\BibleQuote{我们祖先的天主\BibleRef{宗 5:29}，就是你们所钉死在十字架上的耶稣，天主已使祂复活了；}因为他们知道，除了祖先的天主，没有别的天主使耶稣从死者中复活；耶稣也称颂祂为亚巴郎、依撒格和雅各伯的天主，他们不是死人的，而是活人的天主。此外，如果这殿不是与基督的天主同一的天主的殿，门徒们又怎能记得第六十九篇圣咏中的话：\BibleQuote{为祢殿宇所发的热忱把我耗尽；}因为在先知书中是这样写的，而不是\BibleQuote{已耗尽了我。}如今基督主要是为我们每人内在的天主殿宇而热忱；祂不希望它成为买卖的殿，也不希望祈祷的殿成为贼窝；因为祂是忌邪的天主之子。我们应当对圣经的此类表述给予宽泛的解释；它们谈论属人的事物，但以隐喻的方式，表明天主愿意在所有人的灵魂中——但尤其是那些决意接受我们最神圣信仰的讯息的人的灵魂中——没有任何异己之物与祂的旨意相混杂。但我们必须记住，第六十九篇圣咏包含这样的话：\BibleQuote{为祢殿宇所发的热忱把我耗尽，}再往后一点：\BibleQuote{他们拿苦胆给我作食物，我渴了，他们拿醋给我喝。}这两处经文都记载在福音书中，那篇圣咏是以基督的口吻说的，并且没有任何地方显示人称的改变。这显明赫拉克利翁非常缺乏观察力，他认为\BibleQuote{为祢殿宇所发的热忱把我耗尽}这句话是以那些被救主赶出并毁灭的权势的口吻说的；他未能看到圣咏中预言的连贯性。因为如果这些话被理解为出自那被驱逐和毁灭的权势之口，那么他必然也要把\BibleQuote{他们拿醋给我喝}——这是同一篇圣咏的一部分——看作是那些权势所说的话。可能误导他的是，他无法理解\BibleQuote{把我耗尽}怎能由基督说出，因为他不理解\OriginalTerm{拟人化情感的陈述}{anthropopathic statements}如何应用于天主和基督。\par

% structure: p0041 source=h2 target=subsection reason=relative-heading-rank; base=h2

\subsection{二十、基督说祂要重建的圣殿是教会。枯骨将如何复生。}

\BibleQuote{犹太人遂回答祂说：\Quote{祢显什么神迹给我们看，因为祢行这些事？}\BibleRef{若 2:18}耶稣回答说：\Quote{你们拆毁这座圣殿，三天之内我要把它重建起来。}}那些属肉体的、倾向物质事物的人，在我看来就是所指的犹太人。在耶稣赶出那些使天主的殿成为商场的人之后，他们因为耶稣这样处理这些事而向祂发怒，并要求一个神迹，一个能证明他们所不接受的圣言有权做这些事的神迹。救主将祂关于圣殿的话与之前关于祂自己身体的话连接起来，并且针对问题\Quote{祢既然行这些事，显什么神迹给我们看？}回答说：\BibleQuote{你们拆毁这座圣殿，三天之内我要把它重建起来。}祂本可以显示上千个其他神迹，但针对问题\Quote{祢既然行这些事}，祂不能回答别的事；祂恰当地给出了关于圣殿的神迹的答案，而不是与圣殿无关的神迹。现在，这两件事——圣殿和耶稣的身体——在我看来，至少在一种解释中，都是教会的预像，并表明教会是由活石\BibleRef{伯前 2:5}建造的，是圣洁司祭职的属灵殿宇，建立在宗徒和先知的基础上\BibleRef{弗 2:20}，耶稣基督自己为房角石；因此它被称为圣殿。此外，从经文\BibleRef{格前 12:27}\BibleQuote{你们是基督的身体，各自互为肢体}，我们看到，即使圣殿的石头和谐组合似乎被瓦解和分散——正如第二十二篇圣咏所记载的，基督的一切骨头，在迫害和苦难中，因那些在迫害中敌视教会合一之人的阴谋而如此——然而圣殿将被重建，身体将在那威胁它的邪恶之日以及随之而来的终结之日后的第三天复活。因为第三天将出现在新天新地上，当这些骨头——以色列全家——在主的大日复活\BibleRef{则 37:11}，死亡已被征服。因此，救主从十字架的苦难中复活，包含了基督整个身体复活的奥秘。但正如耶稣那个物质的身体为基督牺牲、被埋葬、然后复活，同样，基督所有圣人的整个身体也与祂一同被钉死，现在不再活着；因为每个人，像保禄一样，只夸耀我们主耶稣基督的十字架\BibleRef{迦 6:14}，借着那十字架，他已被钉死于世界，世界也钉死于他。因此，它不仅与基督同钉十字架、被钉死于世界，也与基督同埋葬，因为我们与基督同埋葬了，保禄说\BibleRef{罗 6:4}。然后他说，好像享受了复活的某种保证：\Quote{我们与祂一同复活了}，因为他行走在某种新生活中，虽然尚未在那所盼望的有福而完美的复活中复活。那么，他现在要么被钉死，然后被埋葬；要么现在被埋葬并被从十字架上取下，并且既然被埋葬，将来要复活。但对于我们大多数人来说，复活的奥秘是伟大的，难以默观；圣经中许多其他地方也提到，尤其在厄则克耳以下经文中宣布：\BibleQuote{上主的手临于我，上主的神引我出去，把我放在平原中，那里布满了骨头。祂领我在它们周围走了一圈，看，平原上骨头很多，而且很干枯。祂对我说：\Quote{人子，这些骨头能复活吗？}我答说：\Quote{吾主上主，祢知道。}祂对我说：\Quote{你向这些骨头讲预言，对他们说：干枯的骨头，听上主的话罢！}}再稍后：\BibleQuote{上主对我说：\Quote{人子，这些骨头就是以色列全家。他们说：我们的骨头干枯了，我们的希望失去了，我们已完了。}}因为这些骨头是什么，它们被称呼说：\BibleQuote{听上主的话罢！}好像它们听到了上主的话？它们属于以色列家，或属于基督的身体，关于这身体，主说：\BibleQuote{我的一切骨头都散失了}，虽然祂身体的骨头并没有散失，连一根也没有折断。但是，当真正的、更完美的基督身体的复活本身发生时，那些现在成为基督肢体的人——因为他们那时将成为干枯的骨头——将聚拢在一起，骨与骨相合，接合与接合相合（因为那些缺乏\OriginalTerm{和谐}{\GreekText{ἁρμονία}}的人没有一个会达到完美的成人），达到基督圆满身体的尺度\BibleRef{弗 4:13}。然后许多肢体将成为一身体，它们虽然多，却成为一身体的肢体。但分辨脚、手、眼、听觉和嗅觉，这些在一种意义上充实了头，在另一种意义上充实了脚和其余的肢体，以及较软弱、较卑微的肢体，较尊贵和较不尊贵的，只属于天主。天主将协调身体，到那时，而不是现在，祂将给那缺少的肢体更丰盛的尊荣，使身体没有任何分裂，反而肢体彼此关心，若一个肢体幸福，所有肢体都分享它的幸福，若一个肢体受光荣，所有肢体都一同欢乐。\par

% structure: p0043 source=h2 target=subsection reason=relative-heading-rank; base=h2

\subsection{二十一、子由圣父复活。耶稣受审的指控基于眼前的事件。}

我所说的与目前我们所关注的段落并不相悖，该段涉及圣殿以及那些被逐出圣殿的人，关于这座圣殿，救主说：\Quote{我为你的殿心里焦急，如同火烧；}以及关于那些要求显神迹的犹太人，以及救主对他们的回答，其中祂将关于圣殿的言论与关于祂自己身体的言论结合起来，说：\Quote{你们拆毁这座殿，我三日内要再建立起来。}因为这殿就是基督的身体，凡不合理性和带有商场气味的一切，都必须从其中赶出，使它不再成为商场。这殿必须被那些图谋反对天主圣言的人拆毁，在它被拆毁之后，正如我们前面所讨论的，要在第三日再建立起来；那时，门徒们也要记起祂——\OriginalTerm{圣言}{Logos}——在天主圣殿被毁之前所说的话，并且相信，不仅他们的知识，而且他们的信德也在那时得以成全，且是因耶稣所说的话。每一个具有这种性质的人，耶稣使他洁净，\BibleRef{若 15:3}除去不合理的事物和带有商场气味的事物，使它们因住在他内的圣言的热心而被拆毁。但它们被拆毁是为了被耶稣再建起来，不是第三日——如果我们注意面前的精确措辞——而是\Quote{在三日内}。因为圣殿的再建发生在它被拆毁后的第一日和第二日，它的复活是在所有这三日内完成的。因此，复活既是已发生的，也是将发生的，如果我们确实与基督一同埋葬，也与祂一同复活。既然\Quote{我们与祂一同复活}这句话并不涵盖整个复活，\Quote{在基督内众人都要复活，\BibleRef{格前 15:22}但各人要按自己的次序：首先是初果的基督，然后是在祂来临时候属于基督的人，再后才是终结。}在复活中，有人于第一日就在天主的乐土中，\BibleRef{路 23:43}这属于复活；当耶稣显现并说\Quote{不要摸我，因为我还没有升到我的父那里去}，\BibleRef{若 20:17}这也属于复活；但复活的圆满是在祂去到父那里的时候。现在有一些人在圣父和圣子这个问题上陷入混乱，我们必须对他们说几句话。他们引用经文，\BibleRef{格前 15:15}\Quote{我们甚至被发现有作为天主假见证的，因为我们为天主作证说，祂曾使基督复活，而实际上祂并没有使祂复活，}以及其他类似的经文，这些经文表明那使死者复活的是另一个人，而非那被复活的；以及经文：\Quote{你们拆毁这座殿，我三日内要再建立起来。}好像由此得出结论：圣子在数目上与圣父并无不同，二者并非只是在实体上是一，而且在主体上也是一，并且圣父和圣子只是在某些方面被说成不同，但在祂们的\OriginalTerm{位格}{hypostases}上并无不同。针对这些观点，我们首先要引证主要经文，这些经文证明圣子与圣父不同，并且圣子必然是子因父而生，父必然是父因子为父。然后我们可以很合适地引用基督的声明：祂不能做什么，除非祂看见父正在做的事和所说的话，\BibleRef{若 5:19}因为凡父所做的，子也同样做，并且祂曾复活死者，即身体，这是父准许祂的，父必须被称为从死者中复活基督的主导者。但赫拉克利翁说：\Quote{在三日内}，而非\Quote{在第三日}，他没有考察这一点（尽管他注意到了\Quote{在三日}这些字眼），认为复活是在三日内完成的。但他也称第三日为神性的日子，他们认为其中暗示了教会的复活。由此看来，第一日应被称为\OriginalTerm{属地的}{earthly}之日，第二日为\OriginalTerm{魂性的}{psychical}之日，因为教会的复活并没有发生在它们身上。现在，假见证人的陈述，记载在\BibleBook{玛窦}{福音}和\BibleBook{马尔谷}{福音}的结尾处，以及他们对我们的主耶稣基督的控告，似乎都与祂的这句话有关：\Quote{你们拆毁这座殿，我三日内要再建立起来。}因为祂是指祂身体的殿说的，但他们以为祂的话是指石头的殿，所以他们在控告祂时说：\Quote{这个人说：我能拆毁天主的殿，三日内再建起来。}或者，如马尔谷所记：\Quote{我们听见祂说：我要拆毁这座用人手所造的殿，三日内我要另建一座不是人手所造的殿。}这里大司祭站起来对祂说：\Quote{你什么都不回答吗？这些人作证反对你的是什么事？}但耶稣默不作声。或者，如马尔谷所说：\Quote{大司祭站起来，在中间问耶稣说：你什么都不回答吗？这些人作证反对你的是什么事？}祂仍然默不作声，一言不发。这些语句必定与眼前的这段经文有关。\par

% structure: p0045 source=h2 target=subsection reason=relative-heading-rank; base=h2

\subsection{二十二、撒罗满圣殿并非四十六年建成。关于厄斯德拉的圣殿，我们无法得知建造时间。数字四十六的意义。}

犹太人因此说：\Quote{这座圣殿建造了四十六年（\BibleRef{若 2:20}），你要在三天内把它重建起来吗？}如果我们紧守历史，便无法说明犹太人怎么会说圣殿已建造了四十六年。因为在\BibleBook{列王纪}{上}中记载（\BibleRef{列上 5:18}），他们用了三年预备石头和木材，在第四年第二个月（\BibleRef{列上 6:1}），当撒罗满作以色列王时，王下令，他们便运来建造殿基的大块宝石和未凿的石头。撒罗满的儿子们和希兰的儿子们凿好石头，在第四年安放，在尼散月和第二个月立了主殿的根基；第十年巴耳月，即第八个月，整座殿宇按全部尺寸和全部计划完成。这样，比较完工之时与建造时期，建造时间不足十一年。那么，犹太人如何说圣殿建造了四十六年呢？有人确实可以勉强解释，无论如何也要凑出四十六年，从达味计划建造圣殿时对纳堂先知说（\BibleRef{撒下 7:2}）\Quote{看，我住在香柏木的宫中，而天主的约柜却在帐幕内}算起，尽管他因是血战之人而被阻止建造，但他似乎热心收集建筑材料。在\BibleBook{编年纪}{上}中（\BibleRef{编上 29:1}），达味王确实对全会众说：\Quote{我的儿子撒罗满，主所选拣的，还年轻娇嫩，这工程实在浩大，因为他不是为人建造，而是为主天主建造。我竭尽全力为我天主的殿预备了金子、银子、铜、铁、木头、苏姆石、填缝石、多种宝石、各种贵重木材、大量帕罗斯大理石。此外，因我爱慕我天主的殿，我所有的金银，看，我都献给了我主的殿，直到满盈；从这些预备中，我为圣人的殿准备了三千塔冷通苏斐尔的金子，七千塔冷通印记的银子，以便工匠的手用以覆盖天主的殿宇。}达味在赫贝龙作王七年，在耶路撒冷作王三十三年（\BibleRef{列上 2:11}）；因此，若能证明圣殿预备和达味收集必需材料的工作始于他作王的第五年，那么关于四十六年的说法便可勉强成立。但有人会说，这里所说的圣殿不是撒罗满所建的，因为它在充军时期已被摧毁，而是厄斯德拉时代所建的圣殿（\BibleRef{厄上 6:1}），对于这圣殿，四十六年可以完全准确。但在玛加伯时代，关于人民和圣殿的情况非常不稳定，我不知道圣殿是否真的在那若干年内建成。\Person{赫拉克利昂}不注重历史，却说撒罗满用四十六年预备圣殿，正是救主的预像。他将数字六与物质（即预像）相联，将数字四十（他说这是不接纳组合的四元组）与灵感及灵感中的种子相联。可以考虑，四十是否可归因于在建造圣殿的关键点上所安排的世界的四种元素，而六则归因于人在第六日被创造。\par

% structure: p0047 source=h2 target=subsection reason=relative-heading-rank; base=h2

\subsection{二十三、基督所说的圣殿即是教会。将有关建造撒罗满圣殿的论述及其中所述的数字应用于教会。}

\Quote{但祂所说的是祂身体的圣殿。\BibleRef{若 2:21} 所以，当祂从死者中复活以后，祂的门徒就想起祂说过这话，便相信了圣经和耶稣所说的这句话。}这指的是圣子的身体就是祂的圣殿这一陈述。可能会有人问：这是否应按字面意义理解，还是我们应当尽力将有关圣殿的每项记载与我们关于耶稣身体的观点联系起来——即祂由童贞女所取的身体，还是被称作教会的基督的身体，正如宗徒所说，我们众人都是祂身体的肢体（\BibleRef{格前 12:27}）。一方面，有人可能认为，无论从何种意义上理解身体，要将有关圣殿的一切都说成与身体正确关联是毫无希望的；他们可能采用一种更简单的解释，说身体（无论哪种意义）之所以被称为圣殿，是因为正如圣殿有天主的荣耀住在其中，同样，祂——天主的肖像和荣耀，一切受造物的首生者——就其身体或教会而言，可以正当地被称为包含这肖像的圣殿。就我们而言，我们认为要详尽解释《列王纪上》中关于圣殿的每一项细节是艰巨的任务，远超出我们的语言能力；我们暂时将此事搁置，因为超出本著作的篇幅。我们也坚信，在这样超越人性的事情上，阐明默感圣经的意义必定是天主智慧之工，是那隐藏于奥秘中的智慧，这世界的掌权者没有一个认识它。我们也很清楚，我们需要那卓越的智慧之神帮助，才能理解这些事，正如圣职的仆人所应当理解的那样；在此，我们将尽可能简短地描述我们对这一主题的看法。身体就是教会，我们从伯多禄那里学到（\BibleRef{伯前 2:5}），教会是天主的房屋，由活石建成，是一座神圣司铎职的属灵房屋。因此，建造这房屋的达味之子，是基督的预像。祂在战争结束（\BibleRef{列上 5:3}）和深刻和平来临之后才建造这房屋；祂为天主的荣耀在地上耶路撒冷建造圣殿，使敬拜不再像会幕那样在可移动的建筑中进行。我们要在教会中寻找关于圣殿的每项陈述的真实含义。如果基督的所有敌人都成了祂的脚凳（\BibleRef{格前 15:25}），最后仇敌死亡也被毁灭，那么就会有最完全的和平。基督将是撒罗满，意思是\Quote{平安的}\BibleRef{编上 22:9}，那说\Quote{我与仇恨和平的人曾和睦相处。}的预言将在祂身上应验。然后每一块活石，根据其在此生的工作，将成为那圣殿的一块石头：有的在根基上，是宗徒或先知，承托着放在他们上面的人；有的在根基之后，由宗徒们支持，并连同宗徒们一同帮助驮负那些更有需要的人。有的将是内殿的石头，那里有约柜、革鲁宾和赎罪盖；有的将是外墙上的石头，还有的甚至是在肋未人和司铎的外墙之外，是全燔祭祭台的石头。这些事物的管理和服务将委托给神圣的权能，即天主的使者，它们分别是宰制者、宝座者、掌管者或权能者；还有服从这些者的其他权能，由三千六百名（\BibleRef{列上 5:15}）监督者（他们奉派管理撒罗满的工程）、七万搬运重物者、八万在山中凿石者（他们参加工程，准备石头和木材）所预表。值得注意的是，那些记载为搬运重物者与\OriginalTerm{七之数}{Hebdomad}有关。采石者和凿石者，他们使石头适合于圣殿，与\OriginalTerm{八之数}{Ogdoad}有某种联系。而监督者，数目为六百，与完全数六自乘有关。石头的准备，即开采并使之适合建筑，持续三年；在我看来，这完全指向永恒间隔的时间，与\OriginalTerm{三之数}{Triad}类似。这将在和平成就之时发生，那是在处理与出离埃及有关的事件之后经过一定年数，也就是三百四十年，以及在天主与亚巴郎立约之后四百三十年于埃及所发生的事之后。这样，从亚巴郎到圣殿开始建造，有两个\OriginalTerm{安息数}{sabbatic numbers}，即七百和七十；那时，我们的君王基督也将命令那七万搬运重物者，不为圣殿的根基取随便的石头，而要取大的、宝贵的、未凿过的石头，使它们可被凿琢，但不由随便的工人，而由撒罗满的子孙来凿；因为我们在《列王纪上》中看到如此记载。然后，由于深切的和平，提洛王希兰也合作建造圣殿，并将自己的儿子交给撒罗满的儿子们，与他们一同为圣所凿出大而宝贵的石头；这些石头在第四年被放在上主殿宇的根基中。但\OriginalTerm{八之数}{Ogdoad}的年数后，这殿宇在奠基后的第八年第八月完成。\par

% structure: p0049 source=h2 target=subsection reason=relative-heading-rank; base=h2

\subsection{二十四、撒罗满圣殿建造的叙述含有严重困难，应以属灵方式解读。}

然而，为那些认为这些话语仅仅指叙事本身的人着想，此刻或许不无裨益地引入一些他们几乎无法反驳的思考，以表明这些话语应被视为圣神的话语，并当在其中寻求圣神的心意。国王的儿子们真的花时间凿琢巨大而宝贵的石头，从事如此不合王室出生的手艺吗？挑夫、石匠和监工的人数，以及准备和标记石头的时期长度，这些都如实记载了吗？圣殿也在和平中预备妥当，要不用锤子、斧头或任何铁器为天主建造，以免在天主的殿中发出喧嚣。再者，我要问那些受文字束缚的人：怎么可能有八万石匠，而天主的殿是用坚硬的白色石头建造，建造过程中却听不见锤子、斧头或任何铁器的声响？难道不是在殿外某处无声无息地凿琢活石，然后被运到建筑中等待它们的位置吗？天主的殿有一种上升的构造，不是有角，而是有直线的弯曲。因为经上记载：\BibleRef{列上 6:8} \Quote{殿的旁屋有盘旋的楼梯，通到中层，从中层通到第三层。}因为天主殿中的楼梯必须是螺旋形的，这样在上升时模仿圆，这是最完美的形状。但为使这殿稳固，里面建了五层旁屋，\BibleRef{列上 6:10} 尽可能美观，高一肘，使人仰望时能看出我们如何从可感觉的事物上升到所谓的神圣知觉，从而被引导去知觉那些唯有心智才能看见的事物。但更幸福之石的位置似乎是那称为\OriginalTerm{达彼尔}{Dabir}的地方，那里有上主盟约的约柜，以及，可以这样说，天主的笔迹，即祂亲自用手指写成的版。整个殿都贴上了金子；我们读到：\BibleRef{列上 6:21} \Quote{他用金子贴了殿的里面，直到整个殿都完工。}但在达彼尔有两个\OriginalTerm{革鲁宾}{cherubim}，这个词希伯来圣经的希腊译者未能满意地译出。有些人未能恰当处理语言，译作殿；但它比殿更神圣。殿中一切都用金制成，作为记号，表明完全成全的心智准确地评估理智所感知的事物。但并非所有人都能接近并认识它们；因此设立了院子的帷幔，因为对大多数司祭和肋未人来说，殿最内部的事物并未启示。\par

% structure: p0051 source=h2 target=subsection reason=relative-heading-rank; base=h2

\subsection{二十五、撒罗满圣殿的进一步灵意化。}

值得探究：一方面，撒罗满王据说建造了圣殿；另一方面，撒罗满打发人请来的工匠长，\BibleRef{列上 7:13} \Quote{是提洛人希兰，他是寡妇的儿子，属纳斐塔里支派，他的父亲是提洛人，作铜匠的，满有智慧、聪明、技巧，能作各样铜工；他被带到撒罗满王那里，作他一切的工。}在此我问：撒罗满是否可以视为万有首生者，\BibleRef{哥 1:15} 而希兰视为祂所取的人——从人的束缚而来（因为\Quote{提洛人}一词意为\Quote{束缚者}）——这人从自然而生，满有一切技艺、智慧和聪明，被带来与万有首生者合作，建造圣殿。在这殿中也有窗户，\BibleRef{列上 6:4} 倾斜而隐蔽，以使神圣光明的光照进入，为得救恩——我何必详述？——使基督的身体，即教会，具有属灵殿宇和天主圣殿的样式。如我先前所说，我们需要那隐藏在奥秘中的智慧，只有能说\Quote{但我们有基督的心意}的人才能领悟——我们需要那智慧，按着使此事得以记载的那位的旨意，对所说的一切详细进行属灵的解释。详述这些与我们现在的话题不符。以上所说的足以让我们明白祂如何\Quote{说关于祂身体的殿}。\par

% structure: p0053 source=h2 target=subsection reason=relative-heading-rank; base=h2

\subsection{二十六、先知书对耶路撒冷的应许指向教会，并将要应验。}

在这之后，合理的疑问是：关于圣殿所叙述的事，是否曾发生在或将要发生在属灵的房屋上？无论我们如何取舍，论证似乎都会陷入困境。如果我们说，类似圣殿所记述的事件可能发生在属灵的房屋上，或已经发生其中，那么听众将很难被说服，承认这些美好的事物会发生改变，一则因为他们不愿意，二则因为认为此类事物容许改变是矛盾的。另一方面，若我们试图维持一次赐予圣徒的美好事物之不变性，则我们不能将历史中的记载应用在他们身上，而且我们似乎会像那些异端人士所做的一样，他们未能保持圣经叙述从始至终的一致性。如果我们不采取老妇或犹太人的观点来理解先知书（尤其依撒意亚）中记载的许诺，即，如果我们期待这些许诺在尘世耶路撒冷应验，那么，既然有关重建圣殿和从掳掠中复兴百姓的某些非凡之事被描述为发生在掳掠和圣殿毁灭之后，我们必须说：我们现在就是圣殿，就是曾被掳掠的百姓，但还要再上到犹太和耶路撒冷，用耶路撒冷的宝石被建造。但我不知道，在漫长时代的循环中，同样性质的事是否可能再次发生，但情况会更糟。我们所提及的依撒意亚的预言如下：\BibleRef{依 54:11} \BibleQuote{看啊，我预备红玉作你的石头，蓝宝石作你的根基；我要以碧玉作你的城垛，以水晶作你的城门，以精选的宝石作你的外墙；你的众子都要受主的教导，你的儿女必大享平安，你必在正义中被建造。} 再稍往后，对同一位耶路撒冷说：\BibleRef{依 60:13} \BibleQuote{黎巴嫩的光荣要来到你那里，有柏木、松木和香柏木，连同那荣耀我圣所的人。那谦卑你、侮辱你的儿子们要战战兢兢地来到你面前；你必被称为主的城，圣以色列的熙雍，因为你曾是荒凉被恨恶的，没有帮助你的。我要使你成为永远的喜乐，万代的欢乐。你要吸吮外邦人的奶，要吃君王的财富，你便知道我是拯救你的主，是拣选你的以色列的天主。我要以金代替铜，以银代替铁，以铜代替木，以铁代替石。我要立你的官长得享平安，你的监督得行正义。你地上不再听见强暴，你境内不再有荒凉和灾祸，你的城墙要称为救恩，你的城门要称为雕刻。白昼太阳不再作你的光，夜间月亮也不再发光照耀你，但基督要作你永远的光，你的天主是你的荣耀。你的太阳不再下落，你的月亮也不亏缺，因为你的主要作你永远的光，你悲哀的日子已经结束了。} 这些预言明显是指向将要来临的世代，且是对被掳的以色列子民说的，那被派遣者就是到他们那里去的那位，祂曾说：\BibleQuote{我奉差遣，不过是到以色列家迷失的羊那里去。} \BibleRef{玛 15:24} 这些事情，虽然他们是被掳的，却要在自己的土地上领受；那时，归化者也要通过基督来到他们那里，投靠他们，如经上所说：\BibleRef{依 54:15} \BibleQuote{看啊，归化者要因我来到你那里，并要逃到你那里得庇护。} 若这一切要发生在被掳者身上，那么显然，他们必须围绕他们的圣殿，必须再上到那里被建造，成了最宝贵的石头。因为我们从若望在其默示录中发现，\BibleRef{默 3:12} 那给得胜者的承诺：他要成为天主圣殿中的柱子，不再出去。我说这一切，是为了使我们至少对圣殿、天主的家、教会和耶路撒冷的事有一个粗略的认识，这些我们不能在此系统地论述。然而，那些在读先知书时不回避寻求其属灵意义之辛劳的人，必须以最大的细节探究这些事，必须考虑到每一种可能性。到此为止是关于\Quote{他身体的圣殿}。\par

% structure: p0055 source=h2 target=subsection reason=relative-heading-rank; base=h2

\subsection{二十七、论门徒后来对耶稣言语的信德}

\Quote{当祂从死者中复活后，\BibleRef{若 2:22}祂的门徒就想起祂说过这话，便相信了圣经和耶稣所说的话。} 这告诉我们，耶稣从死者中复活后，祂的门徒看见祂关于圣殿的话更高级地应用于祂的苦难和复活；他们想起\Quote{三天内我要把它建立起来}是指着复活说的；\Quote{他们便相信了圣经和耶稣所说的话。} 我们没有被告诉说他们先前相信了圣经或耶稣所说的话。因为完整意义上的信德，是在圣洗中宣认的人用全心接受的行为。至于更高层次的意义，既然我们已经说过主的整个身体从死者中复活，现在我们注意到，门徒因圣经的应验而想起他们生前没有完全明白的事；其意义如今摆在他们眼前，完全显明给他们，他们知道了那是天上之事的模型和影子。然后他们相信了圣经，从前他们不信，并且相信了耶稣的话，按说话者的意思，他们在复活前并未相信。因为一个人若没有看到圣经中圣神的心意，而天主愿意我们相信那心意胜过字面意义，又怎能说他完全相信圣经呢？从这个角度我们必须说，没有一个随从肉性的人相信法律的属神之事，他们连这些事的最初概念都没有。但是他们说：那些没有看见而相信的人，比那些看见了而相信的人更为有福。为此他们引用了\BibleBook{若望}{福音}末尾对\Person{多默}说的话：\Quote{那些没有看见而相信的人，是有福的。} 但这里并没有说那些没有看见而相信的人比那些看见了而相信的人更有福。按照他们的看法，宗徒之后的人比宗徒更有福；这再愚蠢不过了。要成为有福的人，必须在心中看见他所相信的事物，并且能够与宗徒一起听到对他所说的话：\Quote{你们的眼睛有福，因为看见了；你们的耳朵有福，因为听见了。}\BibleRef{玛 13:16} 和\Quote{许多先知和义人渴望看你们所看见的，却没有看见；听你们所听见的，却没有听见。} 但一个人如果只领受了较低级的真福，即：\BibleRef{若 20:29}\Quote{那些没有看见而相信的人，是有福的}，也可以满足。然而，耶稣称为有福的眼睛所看见的事物，远比那些未能达到如此异象的眼睛更有福；\Person{西默盎}满足于将天主的救恩抱在怀中，看见后他说：\BibleRef{路 2:29}\Quote{主啊！现在可照祢的话，让祢的仆人平安离去，因为我亲眼看见了祢的救恩。} 因此，我们必须像\Person{撒罗满}所说的那样，睁开我们的眼睛，以便用食粮饱足；他说：\Quote{睁开你的眼睛，用食粮饱足。} 我以上关于经文\Quote{他们相信了圣经和耶稣对他们所说的话}所说的，在讨论信德的题目之后，可能引导我们理解：我们信德的圆满，将在耶稣整个身体（即祂的圣教会）从死者中伟大复活时赐予我们。因为关于知识所说的\Quote{现在我知道的只是局部}\BibleRef{格前 13:12}，我认为也可以同样应用于其他一切美善；其中之一就是信德。我们可以说\Quote{现在我信得只是局部}，\Quote{但当那圆满的来到时，这局部的信德就要被废除。} 正如知识那样，信德也是如此：借由看见的方式远胜于（如果可以这样说）借由镜子和谜语的方式。\par

% structure: p0057 source=h2 target=subsection reason=relative-heading-rank; base=h2

\subsection{二十八、相信耶稣之名与相信耶稣本身的区别}

\Quote{当\Person{耶稣}在\Place{耶路撒冷}过逾越节，在庆节期间，许多人看见他所行的记号，便信了他的名。但是\Person{耶稣}，他自己却不信任他们，因为他认识众人，并且他不需要任何人给人类作证，因为他自己知道人心里有什么。}有人可能会问，\Person{耶稣}为何不信任那些我们听说已经相信了的人。对此我们必须说，\Person{耶稣}所不信任的并不是那些信他人，而是那些信他名的人；因为信他的名与信他本人是两回事。谁因他的信德而不受审判，他得以免于审判，不是因为信他的名，而是因为信他本人；因为主说：\BibleRef{若 3:18}\BibleQuote{信我的，不受审判}，而不是\BibleQuote{信我名的，不受审判}；后者相信，因此他尚不值得被定罪，但他次于那信他本人的。因此，\Person{耶稣}不把自己托付给那信他名的人。所以，我们更应该依附他本人，而不是他的名，免得我们以他的名行了奇事后，却听见他对那些仅夸耀他的名的人所说的话。\BibleRef{玛 7:21} 让我们与宗徒\Person{保禄}一同\BibleRef{斐 4:13}喜乐地说：\BibleQuote{我依靠那加强我的基督\Person{耶稣}，能应付一切。}我们还要注意，在前一段中\BibleRef{若 2:13}，福音作者称那逾越节是犹太人的逾越节，而在这里他没有说\Person{耶稣}是在犹太人的逾越节，而是在\Place{耶路撒冷}的逾越节；在前一段中，当逾越节被称为犹太人的时，它没有被称作庆节；但这里记载\Person{耶稣}在庆节期间；他在\Place{耶路撒冷}过逾越节时正在庆节中，许多人相信了他的名，尽管只是他的名。我们当然应当注意，\Quote{许多人}被认为相信，不是信他本人，而是信他的名。而那些信他本人的人，是走那正直而狭窄的路\BibleRef{玛 7:14}，那条路通往生命，找到它的人很少。然而，那些信他名的人中，许多人将会与\Person{亚巴郎}、\Person{依撒格}和\Person{雅各伯}一同坐席在天国里，那是在天父的家中，那里有许多住处。并且应当注意，那些信他名的人，并不像\Person{安德肋}、\Person{伯多禄}、\Person{纳塔乃耳}和\Person{斐理伯}那样相信。他们是相信\Person{若翰}的见证，当他说\Quote{看，天主的羔羊}，或因\Person{安德肋}找到了基督，或因\Person{耶稣}对\Person{斐理伯}说\Quote{跟随我}，或因\Person{斐理伯}说\Quote{我们找到了\Person{梅瑟}和先知们所写的那位，就是\Place{纳匝肋}人\Person{若瑟}的儿子\Person{耶稣}。}而我们现在所说的那些人，\Quote{因看见他所行的记号，便信了他的名}。既然他们相信记号，而不是信他本人，而是信他的名，\Person{耶稣}\Quote{就不信任他们，因为他认识众人，不需要谁给人类作证，因为他知道每个人心里有什么}。\par

% structure: p0059 source=h2 target=subsection reason=relative-heading-rank; base=h2

\subsection{二十九、论\Person{耶稣}需要关于什么受造物的证言}

这些话：\Quote{他不需要别人为人类作证}，很可以说明天主子能自己看清关于每个人的真理，不需要任何他人所提供的这种证言。然而，\Quote{他不需要别人为人类作证}并不等于\Quote{他不需要关于任何受造物的证言}。如果我们把\Quote{人}这个词理解为每一个依照天主肖像的受造物，或每一个理性的受造物，那么他就无需任何人为他作证关于任何理性的受造物，因为他自己借着父所赐的能力，认识他们全部。但如果\Quote{人}这个词仅限于有死的、有灵性的、理性的受造物，那么一方面可以说，他需要关于高于人的受造物的证言，而他对人的认识是足够的，却未延伸到那些其他受造物；另一方面，然而，那自我谦卑的他不需要别人为人类作证，但他对于高于人的受造物，却有这种需要。\par

% structure: p0061 source=h2 target=subsection reason=relative-heading-rank; base=h2

\subsection{三十、\Person{耶稣}如何知道人内住的美善或邪恶的力量}

也许还会问，那些信他的人看见了他所行的什么记号，因为并没有记载他在\Place{耶路撒冷}行了什么记号，尽管可能行了一些没有记载的。然而，有人可以考虑他所行的是否可称为记号，当他用细绳做成鞭子，把众人连同羊和牛都赶出圣殿，倒出兑换银钱者的钱，推翻桌子时。至于那些认为他只是关于人类不需要证言的人，必须说明，福音作者把两件事归给他：他认识一切受造物，并且不需要任何人为人类作证。如果他认识一切受造物，那么他不仅认识人类，也认识高于人类的一切受造物，即一切没有我们这种身体的受造物；并且他知道人心里有什么，因为他比那些借先知之口责备和判断，并照明那些圣神所感动他们如此对待的人心中隐秘的事的人更大。\Quote{他知道人心里有什么}这句话也可以理解为指那些在人内运作的、或善或恶的力量。因为如果有人给魔鬼留了余地，撒殚就会进入他里面；\Person{犹大}就是这样留了余地，因此魔鬼把背叛\Person{耶稣}的意思放在他心里，\BibleQuote{在吃饼块后}，因此，\BibleQuote{魔鬼就进入了他里面}\BibleRef{若 13:2}。但如果有人给天主留了余地，他就成为有福的；因为那从天主获得帮助的人是有福的，他心里的上升是从天主而来的。认识万物的你啊，天主子，你知道人心里有什么。现在，我们的第十卷书已经足够庞大，我们将在此暂停我们的主题。\par

\SourceRecord{Translated by Allan Menzies.；Ante-Nicene Fathers；Vol. 9.；Edited by Allan Menzies.；Buffalo, NY: Christian Literature Publishing Co.,；1896.；<http://www.newadvent.org/fathers/101510.htm>.}
